% This file is a chapter fragment, in the same style as example.tex.
% It is intended to be included by the main blueprint document.
%
% Status convention:
%   \lean{...} records a proposed declaration name.  It is not a claim that
%   the name already exists.  Exact names verified in the current Mathlib
%   documentation are recorded in implementation notes immediately below the
%   relevant block.  For statements already realized in the WeilConjectures
%   Lean development (the nagata-01 ... nagata-09 PR stack, 2026-08), the
%   \lean{...} names have been synced to the actual declarations, and
%   "Formalization route" remarks record where the Lean proof deviates from
%   the route written here.  A \nagataGap paragraph is an intentional, unresolved
%   formalization obligation; it must be closed before a completion marker is
%   added.

\newcommand{\nagataGap}[1]{%
  \par\noindent\textbf{Formalization gap.} #1\par}

% Logical-audit convention.  The blueprint package already defines \uses;
% this wrapper preserves the dependency-graph annotation and prints an
% explicit reference for every listed edge at the point where the annotation
% occurs.  The annotation is placed immediately before the sentence or proof
% paragraph that first uses the dependency.
\let\nagataOriginalUses\uses

\ExplSyntaxOn
\cs_new_protected:Npn \nagataRefs:n #1
  {\clist_map_inline:nn {#1} {\ref{##1}\,}}
\renewcommand{\uses}[1]{%
  \nagataOriginalUses{#1}%
  \par\noindent
  {\footnotesize\emph{uses: }
    \nagataRefs:n {#1}.}%
  \par
}
\ExplSyntaxOff

\chapter{Nagata compactification}

This chapter gives a formalization blueprint for Nagata's compactification
theorem for schemes.  The key result is
\begin{enumerate}
  \item the theorem labelled \texttt{thm:nagata-compactification}, asserting
  that every separated
  finite type morphism $X\to S$ with $S$ quasi-compact and quasi-separated
  factors as a quasi-compact open immersion followed by a proper morphism.
\end{enumerate}

The proof follows the scheme-theoretic proof in the Stacks Project.  The
intermediate statements are deliberately finer than the printed proof: every
change of ambient scheme, every use of a scheme-theoretic image, every
admissible blow-up, and every valuative lifting step is given its own labelled
interface.  This is the level at which the objects and equalities will have to
be supplied to Lean.

Throughout this chapter, ``finite type'' includes quasi-compactness, and a
compactification always has a quasi-compact open immersion as its first map.
Fiber products and all morphisms are taken in the category of schemes.
The current Mathlib interfaces used without a blueprint definition include
rings, ideals, modules, localizations, schemes, fibre products, open and closed
immersions, proper morphisms, finite-presentation morphisms, directed systems,
scheme gluing via \texttt{AlgebraicGeometry.Scheme.GlueData.glued}, and
categorical limits.  Important verified declarations are recorded in the
implementation notes: \texttt{Scheme.Hom.image}, \texttt{Scheme.Hom.imageι},
\texttt{Scheme.Hom.toImage}, \texttt{Scheme.Hom.toImage\_imageι}, the affine
\texttt{Proj} charts (\texttt{AlgebraicGeometry.Proj.basicOpenIsoSpec} and
\texttt{AlgebraicGeometry.Proj.affineOpenCover}), Rees-algebra operations such
as \texttt{reesAlgebra} and \texttt{reesAlgebra.fg}, and properness stability
and target-locality from \texttt{AlgebraicGeometry.IsProper}.  Relative Proj,
geometric blow-ups,
flattening by blow-up, and the Noetherian approximation package are not
assumed to exist: each is defined below only as a blueprint interface and is
marked with an explicit gap until its construction is supplied.

Every other named construction used below has a labelled definition.  A
statement is written in ordinary mathematical language; a \verb|\lean{...}| name is
an implementation note, not part of the statement.  Canonical comparison maps
between chosen closed subschemes are written as isomorphisms unless a chosen
equality has explicitly been fixed.

\section{Compactifications and scheme-theoretic images}


\begin{definition}[Compactification] 
  \label{def:nagata-compactification}
  \lean{AlgebraicGeometry.Compactification}

  Let $f\colon X\to S$ be a morphism of schemes.  A \emph{compactification}
  of $f$ consists of a scheme $\overline X$, morphisms
  \[
    X\xrightarrow{j}\overline X\xrightarrow{p}S,
  \]
  and an equality $j;p=f$, such that $j$ is a quasi-compact open immersion and
  $p$ is proper.
\end{definition}

\begin{definition}[Compactifiability]
  \label{def:nagata-compactifiable}
  \uses{def:nagata-compactification}
  \lean{AlgebraicGeometry.Compactifiable}

  A morphism $f\colon X\to S$ is \emph{compactifiable} if the type of
  compactifications of $f$ is nonempty.
\end{definition}

\begin{lemma}[Mathlib scheme-theoretic image interface]
  \label{lemma:nagata-scheme-image-ideal-qc}
  \lean{AlgebraicGeometry.Scheme.Hom.image}

  For a morphism $g\colon T\to Y$, Mathlib defines its scheme-theoretic image
  as the closed subscheme cut out by the ideal sheaf $\operatorname{ker}(g)$.
  It also supplies the closed immersion
  $g_{\mathrm{im}}\colon\operatorname{SIm}(g)\hookrightarrow Y$ and the
  factor map $g^{\mathrm{im}}\colon T\to\operatorname{SIm}(g)$.
  When $g$ is quasi-compact, the stalk maps of $g^{\mathrm{im}}$ are
  injective and the support of $\operatorname{ker}(g)$ is the closure of the
  topological range of $g$.

  \emph{Implementation note.}  The verified declarations are
  \texttt{AlgebraicGeometry.Scheme.Hom.image}, \texttt{Scheme.Hom.imageι},
  \texttt{Scheme.Hom.toImage}, \texttt{Scheme.Hom.toImage_imageι},
  \texttt{Scheme.Hom.toImage\_app\_injective} (under quasi-compactness), and
  \texttt{AlgebraicGeometry.Scheme.IdealSheafData.support\_map} together with
  \texttt{Scheme.IdealSheafData.map\_bot} (under quasi-compactness).  No
  invented \texttt{Scheme.Hom.support\_ker} or
  \texttt{schemeTheoreticImageIdeal\_isQuasiCoherent} declaration is assumed.
\end{lemma}

\begin{proof}
  This block records the existing Mathlib construction.  The closed
  subscheme is \texttt{Scheme.Hom.image g}, its inclusion is the declaration
  \texttt{Scheme.Hom.imageι g}, and the factor map is
  \texttt{Scheme.Hom.toImage g}; the displayed composite is the theorem
  \texttt{Scheme.Hom.toImage_imageι}.  The additional quasi-compact
  assertions are exactly the two declarations named in the implementation
  note, so no direct-image calculation is silently imported here.
\end{proof}

\begin{definition}[Scheme-theoretic image]
  \label{def:nagata-scheme-theoretic-image}
  \uses{lemma:nagata-scheme-image-ideal-qc}
  \lean{AlgebraicGeometry.Nagata.schemeTheoreticImage}

  Let $g\colon T\to Y$ be a morphism.  Define
  $\operatorname{SIm}(g)$ to mean the Mathlib scheme-theoretic image
  \texttt{Scheme.Hom.image g}, with inclusion $g_{\mathrm{im}}$ and factor map
  $g^{\mathrm{im}}$ supplied by
  \ref{lemma:nagata-scheme-image-ideal-qc}.  Whenever a later argument needs
  quasi-compactness, it must state that hypothesis explicitly.
\end{definition}

\begin{definition}[Scheme-theoretically dense open immersion]
  \label{def:nagata-scheme-theoretically-dense-open}
  \lean{AlgebraicGeometry.Compactification.IsSchemeTheoreticallyDenseOpen}

  For a quasi-compact open immersion $j\colon U\hookrightarrow X$, say that
  $U$ is scheme-theoretically dense in $X$ if the sheaf morphism
  \[
    \mathcal O_X\longrightarrow j_*\mathcal O_U
  \]
  is injective.  This predicate is independent of whether the ambient scheme
  is presented as part of a compactification.
\end{definition}

\begin{definition}[Scheme-theoretically dense compactification]
  \label{def:nagata-dense-compactification}
  \uses{def:nagata-compactification,
    def:nagata-scheme-theoretically-dense-open}
  \lean{AlgebraicGeometry.IsSchemeTheoreticallyDominant}

  A compactification $X\xrightarrow{j}\overline X\to S$ is
  \emph{scheme-theoretically dense} when its first map has the predicate
  introduced in \Cref{def:nagata-scheme-theoretically-dense-open}; explicitly,
  this means that
  \[
    \mathcal O_{\overline X}\longrightarrow j_*\mathcal O_X
  \]
  is injective.
\end{definition}

\begin{lemma}[Characterization of scheme-theoretic density]
  \label{lemma:nagata-dense-compactification-characterization}
  \uses{def:nagata-dense-compactification,
    def:nagata-scheme-theoretic-image,lemma:nagata-scheme-image-ideal-qc}
  \lean{AlgebraicGeometry.isSchemeTheoreticallyDominant_iff_isIso_imageι}

  For an open immersion $j$ (no quasi-compactness is needed), the density
  condition in
  \Cref{def:nagata-dense-compactification} is equivalent to the assertion that
  the image inclusion
  $j_{\mathrm{im}}\colon\operatorname{SIm}(j)\to\overline X$ is an
  isomorphism.  Under the resulting identification
  $\operatorname{SIm}(j)\cong\overline X$, the factor map
  $j^{\mathrm{im}}$ identifies with the original open immersion $j$; in
  general, $j^{\mathrm{im}}$ is not an isomorphism.
\end{lemma}

\begin{proof}
  \uses{lemma:nagata-scheme-image-ideal-qc}
  The kernel of $\mathcal O_{\overline X}\to j_*\mathcal O_X$ is the ideal
  defining $\operatorname{SIm}(j)$ by
  \Cref{def:nagata-scheme-theoretic-image}.  It is zero exactly when the
  closed immersion $\operatorname{SIm}(j)\hookrightarrow\overline X$ is an
  isomorphism, by
  \texttt{Scheme.isIso\_subschemeι\_iff\_eq\_bot}.  Finally,
  \texttt{Scheme.Hom.toImage\_imageι} identifies the composite of the factor
  map with this inclusion with $j$.  This transports $j^{\mathrm{im}}$ to
  $j$ when the inclusion is inverted; it does not make the factor map an
  isomorphism.
\end{proof}

\begin{definition}[Canonical factor through a scheme-theoretic image]
  \label{def:nagata-scheme-image-factor}
  \uses{def:nagata-scheme-theoretic-image}
  \lean{AlgebraicGeometry.Scheme.Hom.toImage}

  Define
  \[
    g^{\mathrm{im}}\colon T\longrightarrow\operatorname{SIm}(g)
  \]
  to be \texttt{Scheme.Hom.toImage g}.  Its factorization equality is the
  verified theorem \texttt{Scheme.Hom.toImage\_imageι}.
\end{definition}

\begin{definition}[Comparison map for a closed factorization]
  \label{def:nagata-scheme-image-comparison}
  \uses{def:nagata-scheme-theoretic-image}
  \lean{AlgebraicGeometry.IsClosedImmersion.lift}

  Suppose $g$ is supplied with a factorization
  $T\xrightarrow{h}Z\xrightarrow{z}Y$, where $z$ is a closed immersion.
  Define
  \[
    c_{g,Z}\colon\operatorname{SIm}(g)\longrightarrow Z
  \]
  from the containment of the ideal defining $Z$ in the kernel ideal defining
  $\operatorname{SIm}(g)$.
  \end{definition}

\begin{lemma}[Universal property of the scheme-theoretic image]
  \label{lemma:nagata-scheme-image-universal}
  \uses{def:nagata-scheme-image-factor,def:nagata-scheme-image-comparison}
  \lean{AlgebraicGeometry.Nagata.schemeImageUniversal}

  The composite of $g^{\mathrm{im}}$ with
  $\operatorname{SIm}(g)\hookrightarrow Y$ is $g$.  For every closed
  factorization through $Z$, the constructed map $c_{g,Z}$ is a closed
  immersion compatible with the maps to $Y$, and it is the unique compatible
  morphism $\operatorname{SIm}(g)\to Z$.
\end{lemma}

\begin{proof}
  The factorization equality is \texttt{Scheme.Hom.toImage_imageι}.
  Let $\mathcal I_Z$ be the ideal of $Z$.  From
  $z\circ h=g$ every local section of $\mathcal I_Z$ maps to zero under $g$,
  so $\mathcal I_Z\subseteq\operatorname{ker}(g)$.  The induced quotient map
  therefore gives a closed immersion
  $c_{g,Z}\colon\operatorname{SIm}(g)\to Z$; its composite with $z$ is the
  image inclusion by construction.  If $c'$ is another compatible map, then
  composing both $c'$ and $c_{g,Z}$ with the closed immersion $z$ gives the
  same image inclusion
  \[
    z\circ c'=g_{\mathrm{im}}=z\circ c_{g,Z}.
  \]
  A closed immersion is a monomorphism, so cancellation of $z$ gives
  $c'=c_{g,Z}$.
  \end{proof}

\begin{lemma}[Scheme-theoretic image of a closed immersion]
  \label{lemma:nagata-scheme-image-of-closed-immersion}
  \lean{AlgebraicGeometry.IsClosedImmersion.instIsIsoSchemeToImage}

  If $i\colon Z\hookrightarrow Y$ is a closed immersion, then the factor map
  $i^{\mathrm{im}}\colon Z\to\operatorname{SIm}(i)$ is an isomorphism.  Thus
  the scheme-theoretic image of a closed immersion is canonically its source,
  compatibly with the maps to $Y$.
\end{lemma}

\begin{proof}
  The current Mathlib instance
  \texttt{AlgebraicGeometry.IsClosedImmersion.instIsIsoSchemeToImage}
  supplies the isomorphism for every morphism carrying the
  \texttt{IsClosedImmersion} property.  Its compatibility with the inclusion
  is the factorization equality
  \texttt{Scheme.Hom.toImage\_imageι}.
\end{proof}

\begin{lemma}[Localization of the affine image kernel]
  \label{lemma:nagata-affine-image-kernel-localization}
  \uses{def:nagata-scheme-theoretic-image}
  \lean{AlgebraicGeometry.Nagata.affineImageKernelLocalization}

  Let $A$ be a ring, let $T\to\operatorname{Spec}(A)$ be quasi-compact and
  quasi-separated, and let $s\in A$.  Put
  \[
    I=\ker\bigl(A\longrightarrow\Gamma(T,\mathcal O_T)\bigr).
  \]
  If $T_s=T\times_{\operatorname{Spec}(A)}\operatorname{Spec}(A_s)$, then
  \[
    \ker\bigl(A_s\longrightarrow\Gamma(T_s,\mathcal O_{T_s})\bigr)=IA_s.
  \]
  This equality is functorial under further localization.
\end{lemma}

\begin{proof}
  Choose a finite affine cover $(T_i)$ of $T$.  Quasi-separatedness makes
  each $T_i\cap T_j$ quasi-compact, so choose a finite affine refinement of
  every such intersection.  The sheaf equalizer for this finite covering
  computes $\Gamma(T,\mathcal O_T)$ from finite products of the coordinate
  rings of the selected affine opens and their finite affine refinements.
  Localizing this equalizer at $s$ is exact, commutes with its finite
  products, and replaces every selected affine by its inverse image in
  $T_s$.  It therefore computes $\Gamma(T_s,\mathcal O_{T_s})$.  Taking the
  kernel of the structural map gives $IA_s$.  Repeating the same construction
  after a second localization gives the stated functoriality.
\end{proof}

\begin{lemma}[Flat base change of the affine image kernel]
  \label{lemma:nagata-affine-image-kernel-flat}
  \uses{lemma:nagata-affine-image-kernel-localization}
  \lean{AlgebraicGeometry.Nagata.affineImageKernelFlatBaseChange}

  Let $A\to B$ be flat, let $T\to\operatorname{Spec}(A)$ be quasi-compact
  and quasi-separated, and put
  \[
    I=\ker\bigl(A\longrightarrow\Gamma(T,\mathcal O_T)\bigr).
  \]
  For $T_B=T\times_{\operatorname{Spec}(A)}\operatorname{Spec}(B)$ one has
  \[
    \ker\bigl(B\longrightarrow\Gamma(T_B,\mathcal O_{T_B})\bigr)=IB.
  \]
  The equality is natural in the flat algebra $B$ and commutes with
  localization on $B$.
\end{lemma}

\begin{proof}
  Use the finite affine cover and finite affine refinements from
  \Cref{lemma:nagata-affine-image-kernel-localization}.  Tensor its finite
  sheaf equalizer with $B$.  Flatness makes tensoring exact, and tensoring
  commutes with the finite products in this diagram.  The resulting affine
  coordinate rings are precisely those of the base-changed cover of $T_B$,
  so the resulting equalizer computes $\Gamma(T_B,\mathcal O_{T_B})$.
  Exactness identifies the kernel of the base-changed structural map with
  $I\otimes_A B$, whose image in $B$ is $IB$.  Naturality follows from the
  naturality of tensor product, and compatibility with localization is
  \Cref{lemma:nagata-affine-image-kernel-localization} applied over $B$.
\end{proof}

\begin{lemma}[Restriction of a scheme-theoretic image to an open target]
  \label{lemma:nagata-scheme-image-restrict-open}
  \uses{def:nagata-scheme-theoretic-image,
    lemma:nagata-affine-image-kernel-localization}
  \lean{AlgebraicGeometry.Nagata.schemeImageRestrictOpen}

  Let $g\colon T\to Y$ be quasi-compact and quasi-separated, let
  $W\subseteq Y$ be open, let $T_W=T\times_YW$, and let
  $g_W\colon T_W\to W$ be the base change of $g$.  Then
  \[
    \operatorname{SIm}(g_W)
      \cong\operatorname{SIm}(g)\times_YW
  \]
  as closed subschemes of $W$, compatibly with their inclusions into $W$.
\end{lemma}

\begin{proof}
  On an affine open $V=\operatorname{Spec}(A)\subseteq Y$, the image ideal is
  the kernel of
  $A\to\Gamma(T\times_YV,\mathcal O)$.  On every principal open
  $D(s)\subseteq V\cap W$,
  \Cref{lemma:nagata-affine-image-kernel-localization} identifies the image
  kernel of the restricted morphism with the localization of this ideal.
  The declaration
  \texttt{AlgebraicGeometry.Scheme.IdealSheafData.ideal\_comap\_of\_isOpenImmersion}
  identifies that localized ideal with the comap of the original image ideal,
  and \texttt{Scheme.IdealSheafData.comapIso} identifies the closed
  subscheme it defines with the fibre product.  Functoriality under further
  localization makes these affine isomorphisms agree on overlaps.  They glue
  to the displayed canonical isomorphism, and the construction from equal
  ideal sheaves makes it compatible with the two inclusions into $W$.
\end{proof}

\begin{lemma}[Flat base change of a scheme-theoretic image,
  {\cite[\href{https://stacks.math.columbia.edu/tag/081I}{Tag 081I}]{stacks-project}}]
  \label{lemma:nagata-scheme-image-flat-base-change}
  \uses{def:nagata-scheme-theoretic-image,
    lemma:nagata-affine-image-kernel-flat}
  \lean{AlgebraicGeometry.Nagata.schemeImageFlatBaseChange}

  Let $g\colon T\to Y$ be quasi-compact and quasi-separated and let
  $q\colon Y'\to Y$ be flat.  For the base change
  $g'\colon T\times_YY'\to Y'$ one has a canonical isomorphism
  \[
    \operatorname{SIm}(g')
      \cong\operatorname{SIm}(g)\times_YY'
  \]
  of closed subschemes of $Y'$.
\end{lemma}

\begin{proof}
  Work over an affine open $V=\operatorname{Spec}(A)\subseteq Y$ and an
  affine open $V'=\operatorname{Spec}(B)\subseteq q^{-1}(V)$.  The ring map
  $A\to B$ is flat.  Apply
  \Cref{lemma:nagata-affine-image-kernel-flat} to
  $T\times_YV\to\operatorname{Spec}(A)$.  It identifies the ideal of the
  image after base change to $V'$ with the extension to $B$ of the original
  image ideal.  The latter is exactly the ideal defining
  $\operatorname{SIm}(g)\times_YV'$ under the closed-subscheme pullback
  isomorphism.  Naturality and localization compatibility in
  \Cref{lemma:nagata-affine-image-kernel-flat} show that these affine
  identifications agree on all overlaps.  Gluing gives the displayed
  canonical isomorphism of closed subschemes of $Y'$.  This is the
  finite-affine-cover proof of Stacks Project, Tag~081I
  (Lemma~29.26.16), with $f=q$ and $g$ the morphism in the present statement.
\end{proof}

\begin{remark}[Formalization route]
  The formal proof (\texttt{Nagata.schemeImageFlatBaseChange}) does not run
  the affine kernel computations of
  \Cref{lemma:nagata-affine-image-kernel-localization} and
  \Cref{lemma:nagata-affine-image-kernel-flat}: it factors the base-changed
  morphism through $\operatorname{SIm}(g)\times_YY'$ and applies Mathlib's
  \texttt{IsSchemeTheoreticallyDominant.of\_isPullback}, which encapsulates
  those computations.  Only quasi-compactness of $g$ is needed, and
  \Cref{lemma:nagata-scheme-image-restrict-open} is obtained in Lean as the
  special case of flat base change along an open immersion, reversing the
  order of these two lemmas.
\end{remark}

\begin{definition}[Scheme-theoretic closure]
  \label{def:nagata-scheme-theoretic-closure}
  \uses{def:nagata-scheme-theoretic-image}
  \lean{Scheme.Hom.schemeTheoreticClosure}

  Let $a\colon T\to Y$ be a quasi-compact and quasi-separated immersion.
  Define the scheme-theoretic closure of $T$ in $Y$ to be the closed
  subscheme $\operatorname{SIm}(a)\hookrightarrow Y$.
\end{definition}

\begin{lemma}[Minimality of the scheme-theoretic closure]
  \label{lemma:nagata-scheme-theoretic-closure-minimal}
  \uses{def:nagata-scheme-theoretic-closure}
  \lean{Scheme.Hom.schemeTheoreticClosure_minimal}

  The scheme-theoretic closure is the smallest closed subscheme of $Y$
  through which $a$ factors: it factors through every such closed subscheme,
  and the comparison is unique over $Y$.
\end{lemma}

\begin{proof}
  \uses{lemma:nagata-scheme-image-universal}
  This is \Cref{lemma:nagata-scheme-image-universal} applied to the immersion
  $a$, together with
  \Cref{def:nagata-scheme-theoretic-closure}.
\end{proof}

\begin{lemma}[{\cite[\href{https://stacks.math.columbia.edu/tag/01RG}{Tag 01RG}]{stacks-project}}]
  \label{lemma:nagata-immersion-open-in-scheme-image}
  \uses{def:nagata-scheme-image-factor}
  \lean{AlgebraicGeometry.IsImmersion.instIsOpenImmersionToImageOfQuasiCompact}

  Let $g\colon T\to Y$ be a quasi-compact immersion.  The constructed map
  $g^{\mathrm{im}}\colon T\to\operatorname{SIm}(g)$ is a quasi-compact open
  immersion, and its image is topologically and scheme-theoretically dense.
\end{lemma}

\begin{proof}
  Apply the verified instance
  \texttt{AlgebraicGeometry.IsImmersion.instIsOpenImmersionToImageOfQuasiCompact}
  to $g$.  It gives that
  $g^{\mathrm{im}}=\texttt{Scheme.Hom.toImage g}$ is an open immersion.
  The verified instance
  \texttt{AlgebraicGeometry.Scheme.instQuasiCompactToImage} gives
  quasi-compactness of the same factor map.
  The verified instance
  \texttt{AlgebraicGeometry.Scheme.instIsDominantToImageOfQuasiCompact} says
  that its topological range is
  dense.  For every affine open $V$ of $Y$, the theorem
  \texttt{Scheme.Hom.toImage\_app\_injective} gives injectivity on the
  corresponding affine open of $\operatorname{SIm}(g)$ obtained by pulling
  $V$ back along the image inclusion.  These opens form an affine basis of
  $\operatorname{SIm}(g)$, so the sheaf morphism
  $\mathcal O_{\operatorname{SIm}(g)}\to
  (g^{\mathrm{im}})_*\mathcal O_T$ is injective on a basis and hence globally.
  Thus the factor map is scheme-theoretically dense, including at boundary
  points.
\end{proof}

\begin{definition}[Compactification constructed from an immersion]
  \label{def:nagata-immersion-compactification}
  \uses{def:nagata-compactification,
    def:nagata-scheme-image-factor}
  \lean{AlgebraicGeometry.Compactification.ofImmersionConstruction}

  Let $g\colon X\to Y$ be a quasi-compact immersion over $S$, with
  $Y\to S$ proper.  Define $\mathsf C(g)$ to be the factorization
  \[
    X\xrightarrow{g^{\mathrm{im}}}
    \operatorname{SIm}(g)\hookrightarrow Y\longrightarrow S.
  \]
\end{definition}

\begin{lemma}[The immersion construction is a compactification]
  \label{lemma:nagata-immersion-into-proper}
  \uses{def:nagata-immersion-compactification}
  \lean{AlgebraicGeometry.Compactification.ofImmersion}

  The constructed factorization $\mathsf C(g)$ is a compactification of
  $X\to S$.
\end{lemma}

\begin{proof}
  \uses{lemma:nagata-immersion-open-in-scheme-image}
  The first map is a quasi-compact open immersion by
  \Cref{lemma:nagata-immersion-open-in-scheme-image}.  The second map is the
  composite of a closed immersion with the proper map $Y\to S$, hence is
  proper.  If the given over-$S$ structure is
  $g;(Y\to S)=f$, the factorization equality gives the required structural
  equality explicitly:
  \[
    g^{\mathrm{im}};g_{\mathrm{im}};(Y\to S)
      =g;(Y\to S)=f.
  \]
\end{proof}

\begin{definition}[Base-change construction for a compactification]
  \label{def:nagata-compactification-base-change}
  \uses{def:nagata-compactification}
  \lean{AlgebraicGeometry.Compactification.baseChangeConstruction}

  Given a compactification $C=(X\to\overline X\to S)$ and $S'\to S$, define
  $C_{S'}$ to be
  \[
    X\times_SS'\longrightarrow\overline X\times_SS'\longrightarrow S'.
  \]
\end{definition}

\begin{lemma}[Base change of a compactification]
  \label{lemma:nagata-compactification-base-change}
  \uses{def:nagata-compactification-base-change}
  \lean{AlgebraicGeometry.Compactification.baseChange}

  The constructed factorization $C_{S'}$ is a compactification of
  $X\times_SS'\to S'$.
\end{lemma}

\begin{proof}
  Open immersions, quasi-compact morphisms, and proper morphisms are stable
  under arbitrary base change.  The factorization equality follows from the
  universal property of the two fibre products.
\end{proof}

\begin{definition}[Dense-refinement construction]
  \label{def:nagata-dense-refinement}
  \uses{def:nagata-compactification,
    def:nagata-scheme-image-factor}
  \lean{AlgebraicGeometry.Compactification.denseRefinement}

  For a compactification $C=(X\xrightarrow{j}\overline X\to S)$, define
  $\mathsf{Dense}(C)$ by replacing $\overline X$ with
  $\operatorname{SIm}(j)$ and using $j^{\mathrm{im}}$ as its first map.
\end{definition}

\begin{proposition}[Properties of the dense-refinement construction]
  \label{prop:nagata-dense-refinement}
  \uses{def:nagata-dense-compactification,def:nagata-dense-refinement}
  \lean{AlgebraicGeometry.Compactification.isSchemeTheoreticallyDominant_denseRefinement_i}

  The object $\mathsf{Dense}(C)$ is a scheme-theoretically dense
  compactification of the same morphism as $C$.
\end{proposition}

\begin{proof}
  \uses{lemma:nagata-immersion-open-in-scheme-image,
    def:nagata-scheme-image-factor}
  The induced map from $X$ is a quasi-compact open immersion and is
  scheme-theoretically dense by
  \Cref{lemma:nagata-immersion-open-in-scheme-image}.  The new ambient scheme
  is closed in $\overline X$, so it remains proper over $S$.  Finally, if
  $p\colon\overline X\to S$ is the original structural map, then
  \[
    j^{\mathrm{im}};j_{\mathrm{im}};p
      =j;p=f
  \]
  by \texttt{Scheme.Hom.toImage\_imageι}; hence the refined factorization has
  the same composite as $C$.
\end{proof}

\begin{definition}[Closed-subscheme compactification construction]
  \label{def:nagata-closed-subscheme-compactification}
  \uses{def:nagata-compactification,
    def:nagata-scheme-theoretic-image,
    def:nagata-scheme-image-factor}
  \lean{AlgebraicGeometry.Compactification.closedSubschemeConstruction}

  Let $i\colon X\hookrightarrow X'$ be closed over $S$ and fix a
  compactification $C'=(X'\to\overline X'\to S)$.  Define
  $\mathsf{Closed}(i,C')$ by taking the scheme-theoretic image
  $\overline X=\operatorname{SIm}(X\to\overline X')$ and the factorization
  \[
    X\longrightarrow\overline X\hookrightarrow\overline X'\longrightarrow S.
  \]
\end{definition}

\begin{proposition}[The closed-subscheme construction is a compactification]
  \label{prop:nagata-closed-subscheme-compactifiable}
  \uses{def:nagata-closed-subscheme-compactification}
  \lean{AlgebraicGeometry.Compactifiable.ofClosedImmersion}

  The constructed factorization $\mathsf{Closed}(i,C')$ is a compactification
  of $X\to S$.
\end{proposition}

\begin{proof}
  \uses{lemma:nagata-scheme-image-restrict-open,
    lemma:nagata-scheme-image-of-closed-immersion}
  Its ambient scheme $\overline X$ is closed in the proper scheme
  $\overline X'$, hence is proper over $S$.  By
  \Cref{lemma:nagata-scheme-image-restrict-open}, its pullback to the open
  $X'\subseteq\overline X'$ is the scheme-theoretic image of
  $i\colon X\hookrightarrow X'$.  By
  \Cref{lemma:nagata-scheme-image-of-closed-immersion}, the factor map for
  $i$ is an isomorphism, and compatibility with the image inclusion identifies
  the pullback image with $X$.  Hence
  \[
    X\cong\overline X\times_{\overline X'}X',
  \]
  so $X\to\overline X$ is identified with the base change of the
  quasi-compact open immersion $X'\to\overline X'$.
\end{proof}

\section{Separated gluing}

\begin{lemma}[The diagonal on a mixed open]
  \label{lemma:nagata-diagonal-mixed-open}
  \lean{Scheme.diagonal_pullback_openPair}

  Let $X=U\cup V$ be an open cover over $S$.  The pullback of
  $\Delta_{X/S}$ to $U\times_SV$ is canonically the map
  \[
    U\cap V\longrightarrow U\times_SV.
  \]
  The analogous assertion holds over $V\times_SU$.
\end{lemma}

\begin{proof}
  Both schemes represent triples consisting of a point of $U$, a point of
  $V$, and an equality of their images in $X$.  The canonical isomorphism is
  obtained from the universal property of the pullback.
\end{proof}

\begin{lemma}[{\cite[\href{https://stacks.math.columbia.edu/tag/0F3U}{Tag 0F3U}]{stacks-project}}]
  \label{lemma:nagata-separated-two-open-cover}
  \lean{Scheme.IsSeparated.of_two_open_cover}

  Let $X=U\cup V$ be an open cover over $S$.  Assume that $U\to S$ and
  $V\to S$ are separated and that
  \[
    U\cap V\longrightarrow U\times_SV
  \]
  is a closed immersion.  Then $X\to S$ is separated.
\end{lemma}

\begin{proof}
  \uses{lemma:nagata-diagonal-mixed-open}
  The target of $\Delta_{X/S}$ is covered by
  \[
    U\times_SU,\quad U\times_SV,\quad
    V\times_SU,\quad V\times_SV.
  \]
  On the two unmixed pieces the diagonal is closed by separatedness of $U/S$
  and $V/S$.  On the mixed pieces use
  \Cref{lemma:nagata-diagonal-mixed-open}.  Closed immersions are local on the
  target, so $\Delta_{X/S}$ is closed.
\end{proof}

\section{Construction of blow-ups from Rees algebras}
\label{sec:nagata-blowup-foundations}

This section expands the blow-up package down to the algebraic and gluing
lemmas needed by Lean.  The later Nagata argument only uses the public
geometric interfaces, but those interfaces should not be marked complete until
the constructions in this section are available.

\subsection{The Proj interface required by the construction}

Mathlib supplies the affine geometric Proj construction and its standard
charts.  The relative Proj needed for a quasi-coherent graded algebra on an
arbitrary scheme is a separate construction and is recorded explicitly below;
it is not silently identified with an unavailable Mathlib primitive.

\begin{definition}[Degree-zero homogeneous localization]
  \label{def:nagata-homogeneous-localization-degree-zero}
  \lean{GradedAlgebra.HomogeneousLocalization.Away}

  Let $R=\bigoplus_{n\geq0}R_n$ be an $\mathbb N$-graded commutative ring and
  let $f\in R_d$ be homogeneous with $d>0$.  Write $R_{(f),0}$ for the
  degree-zero part of the homogeneous localization supplied by Mathlib's
  \texttt{HomogeneousLocalization.Away}.  Its elements are represented by
  \[
    \frac{x}{f^m},\qquad x\in R_{md},\quad m\geq0,
  \]
  modulo the ordinary localization relation.  This block fixes the notation
  used by the chart statements; it does not introduce a second localization
  implementation.
\end{definition}

\begin{lemma}[Coherence of the standard Proj chart rings]
  \label{lemma:nagata-proj-chart-ring-cocycle}
  \uses{def:nagata-homogeneous-localization-degree-zero}
  \lean{AlgebraicGeometry.Proj.basicOpenIsoSpec_cocycle}

  For positive-degree homogeneous $f,g,h$, localization induces canonical
  equivalences between the restrictions of $(R_{(f)})_0$ and
  $(R_{(g)})_0$ to the overlap indexed by $fg$.  These equivalences are the
  identity for $f=g$, are inverse under exchange of $f$ and $g$, and satisfy
  the cocycle equality on the triple localization indexed by $fgh$.
\end{lemma}

\begin{proof}
  Use the canonical maps between the homogeneous localizations at $f$, $g$,
  and $fg$.  On a representative $x/f^m$ each map sends the fraction to its
  image in the localization at $fg$; the localization relation proves
  independence of representatives.  The identity, inverse, and triple
  cocycle equations follow because both sides are the same map into the
  localization at $fgh$.  These are the compatibilities used by the verified
  declarations \texttt{Proj.basicOpenIsoSpec} and
  \texttt{Proj.affineOpenCover}.
\end{proof}

\begin{definition}[Affine Proj]
  \label{def:nagata-affine-proj}
  \uses{def:nagata-homogeneous-localization-degree-zero}
  \lean{AlgebraicGeometry.Proj}

  For an $\mathbb N$-graded commutative ring $R$, use the Mathlib scheme
  $\operatorname{Proj}(R)$, with its structural morphism to the degree-zero
  spectrum.  On a homogeneous element $f$ of positive degree, its standard
  open is identified with $\operatorname{Spec}(R_{(f),0})$ by
  \texttt{Proj.basicOpenIsoSpec}; a homogeneous family generating the
  irrelevant ideal gives the cover through
  \texttt{AlgebraicGeometry.Proj.iSup\_basicOpen\_eq\_top}.
\end{definition}

\begin{proposition}[Standard affine opens of Proj]
  \label{prop:nagata-proj-standard-open}
  \uses{def:nagata-affine-proj,
    def:nagata-homogeneous-localization-degree-zero}
  \lean{AlgebraicGeometry.Proj.basicOpenIsoSpec,\,
    AlgebraicGeometry.Proj.iSup_basicOpen_eq_top}

  Each positive-degree homogeneous $f$ determines an open subscheme
  $D_+(f)\subseteq\operatorname{Proj}(R)$ and a canonical isomorphism
  \[
    D_+(f)\cong\operatorname{Spec}(R_{(f),0}).
  \]
  If homogeneous elements $(f_i)$ generate $R_+$, the $D_+(f_i)$ cover
  $\operatorname{Proj}(R)$.
\end{proposition}

\begin{proof}
  \uses{lemma:nagata-proj-chart-ring-cocycle}
  The standard-open isomorphism is \texttt{Proj.basicOpenIsoSpec}.
  The cover assertion is
  \texttt{AlgebraicGeometry.Proj.iSup\_basicOpen\_eq\_top}; its hypothesis
  that the chosen homogeneous family generates the irrelevant ideal is the
  assumption in the statement.
  Thus no independently glued affine scheme is being used.
\end{proof}

\begin{definition}[Relative Proj]
  \label{def:nagata-relative-proj}
  \uses{def:nagata-affine-proj}
  \lean{AlgebraicGeometry.RelativeProj}

  Let $X$ be a scheme and
  $\mathcal A=\bigoplus_{n\geq0}\mathcal A_n$ a quasi-coherent graded
  $\mathcal O_X$-algebra.  On every affine open $V\subseteq X$, form
  \[
    \operatorname{Proj}(\Gamma(V,\mathcal A)).
  \]
  Quasi-coherence identifies the graded algebra on a basic affine restriction
  with coefficientwise localization.  The intended relative Proj is the scheme
  obtained by gluing these affine Proj schemes along the induced standard-chart
  equivalences, with structural map to $X$ glued from the degree-zero maps.
  The construction includes the chosen affine cover, all overlap
  isomorphisms, their cocycle equalities, and the resulting universal
  restriction isomorphisms.
\end{definition}

\begin{lemma}[Existence and affine compatibility of relative Proj]
  \label{lemma:nagata-relative-proj-exists}
  \uses{def:nagata-relative-proj,def:nagata-affine-proj,
    lemma:nagata-proj-chart-ring-cocycle}
  \lean{AlgebraicGeometry.RelativeProj.exists}

  The relative Proj construction in
  \Cref{def:nagata-relative-proj} exists as a scheme, is independent up to a
  unique isomorphism compatible with the chosen affine-chart
  identifications, and restricts over every affine open $V\subseteq X$ to
  $\operatorname{Proj}(\Gamma(V,\mathcal A))$.  These restriction
  isomorphisms commute with further restriction and with the structural maps.
\end{lemma}

\begin{proof}
  \uses{lemma:nagata-proj-chart-ring-cocycle}
  Choose the index set of all affine opens $V\subseteq X$.  For each $V$ put
  $P_V=\operatorname{Proj}(\Gamma(V,\mathcal A))$.  If $W\subseteq V$ is a
  principal affine open, quasi-coherence gives the coefficientwise graded
  localization
  \[
    \Gamma(V,\mathcal A)[s^{-1}]\;\xrightarrow{\sim}\;
    \Gamma(W,\mathcal A),
  \]
  and the affine Proj construction gives an isomorphism from $P_W$ onto the
  inverse image of $W$ in $P_V$.  To cover an arbitrary overlap $V\cap V'$,
  index the gluing datum by triples $(V,V',W)$ with $W$ affine and
  $W\subseteq V\cap V'$, together with a principal refinement inside $W$.
  The associated isomorphism is the map on each standard chart induced by the
  displayed localization.  Every overlap is covered by such indices, and the
  maps are isomorphisms because localization is transitive.

  The identity and inverse axioms are the cases $V=V'$ and exchange of the two
  affine indices.  On a triple overlap all routes send a fraction
  $x/s^m$ to its image in the localization at the product of the three
  denominators; hence the cocycle axiom is exactly
  \Cref{lemma:nagata-proj-chart-ring-cocycle}.  The datum therefore satisfies
  the scheme-gluing axioms, so the Mathlib
  construction \texttt{AlgebraicGeometry.Scheme.GlueData.glued} produces a
  scheme $P$ and a morphism $P\to X$ (the degree-zero maps agree on every
  overlap).  Its restriction to $V$ is the specified $P_V$ by the restriction
  field of the same gluing construction.  A second affine cover produces the
  same overlap datum after passing to the common refinement by affine opens;
  the uniqueness field of \texttt{Scheme.GlueData.glued} gives a unique
  isomorphism over $X$ that
  respects all the local identifications.  (No uniqueness among arbitrary
  automorphisms over $X$ is asserted.)  The same argument on a third
  refinement proves compatibility with further restrictions.
\end{proof}

\begin{definition}[Relative projective space]
  \label{def:nagata-relative-projective-space}
  \uses{def:nagata-relative-proj,lemma:nagata-relative-proj-exists}

  For $r\in\mathbb N$, define the relative projective-space interface by
  \[
    \mathbf P^r_X
      =\operatorname{Proj}_X
        \bigl(\operatorname{Sym}_{\mathcal O_X}(\mathcal O_X^{r+1})\bigr),
  \]
  with the symmetric-algebra generators placed in degree one.  Its standard
  opens are required to be the relative affine schemes of the degree-zero
  homogeneous localizations at those generators.  The existence and
  compatibility of this object are supplied by
  \Cref{lemma:nagata-relative-proj-exists}.

  This node is not formalized: it presupposes relative Proj, which is not
  formalized either.  Everything that needs only the scheme
  $\mathbf P^r_X\to X$, and not its standard charts, uses
  \Cref{def:nagata-projective-space-base-change} instead.
\end{definition}

\begin{definition}[Projective space over a base scheme, by base change]
  \label{def:nagata-projective-space-base-change}
  \uses{def:nagata-affine-proj}
  \lean{AlgebraicGeometry.ProjectiveSpace}

  Let $n$ be a finite index set and $S$ a scheme.  Put
  \[
    \mathbf P^n_{\mathbf Z}
      =\operatorname{Proj}\bigl(\mathbf Z[T_i:i\in n]\bigr),
    \qquad
    \mathbf P(n;S)
      =S\times_{\operatorname{Spec}\mathbf Z}\mathbf P^n_{\mathbf Z},
  \]
  where the polynomial ring carries its total-degree grading, and take for the
  structural morphism $\mathbf P(n;S)\to S$ the first projection.  Here $n$
  indexes the \emph{homogeneous} coordinates, so the classical
  $\mathbf P^r_S$ is $\mathbf P(\{0,\dots,r\};S)$.

  This is the same scheme over $S$ as
  \Cref{def:nagata-relative-projective-space}, but obtained without any
  gluing, so it does not come with the standard-chart interface required
  there; only the properties of $\mathbf P(n;S)$ as a scheme over $S$ are
  available from it.  In particular it gives no way to construct morphisms
  \emph{into} $\mathbf P(n;S)$, which is what
  \Cref{lemma:nagata-global-blowup-relative-projective-embedding} needs; the
  comparison of the two nodes is left to whoever builds relative Proj.
\end{definition}

\begin{lemma}[Projective space is proper over its base]
  \label{lemma:nagata-projective-space-base-change-proper}
  \uses{def:nagata-projective-space-base-change,def:nagata-affine-proj}
  \lean{AlgebraicGeometry.ProjectiveSpace.instIsProper}

  For every finite index set $n$ and every scheme $S$, the structural
  morphism $\mathbf P(n;S)\to S$ is proper.
\end{lemma}

\begin{proof}
  The polynomial ring $\mathbf Z[T_i:i\in n]$ is of finite type over the
  degree-zero part of its total-degree grading, which is $\mathbf Z$, because
  $n$ is finite.  Hence the verified properness instance for
  $\operatorname{Proj}$ of a graded ring of finite type over its degree-zero
  part (\Cref{def:nagata-affine-proj}) makes
  $\mathbf P^n_{\mathbf Z}\to\operatorname{Spec}\mathbf Z$ proper.  Since
  $\operatorname{Spec}\mathbf Z$ is the terminal scheme, the same holds for
  $\mathbf P^n_{\mathbf Z}\to\ast$, and properness is stable under base
  change, so the first projection $\mathbf P(n;S)\to S$ is proper.
\end{proof}

\begin{definition}[Projective morphism]
  \label{def:nagata-projective-morphism}
  \uses{def:nagata-relative-projective-space}
  \lean{AlgebraicGeometry.IsProjective}

  A morphism $f\colon Y\to X$ is projective if there exist $r\in\mathbb N$
  and a closed immersion $Y\hookrightarrow\mathbf P^r_X$ over $X$ whose
  composite with the structural morphism is $f$.
\end{definition}

\begin{lemma}[Projective morphisms are proper]
  \label{lemma:nagata-projective-is-proper}
  \uses{def:nagata-projective-morphism}
  \lean{AlgebraicGeometry.IsProjective.isProper}

  Every projective morphism of schemes is proper.
\end{lemma}

\begin{proof}
  Let $(V_i)$ be an affine open cover of $X$.  By the relative-Proj
  construction and its restriction isomorphisms,
  $\mathbf P^r_X\times_XV_i$ is the affine Proj of the polynomial graded
  algebra $\Gamma(V_i,\mathcal O_X)[T_0,\ldots,T_r]$.  The verified affine
  Proj instance
  \texttt{Proj.instIsProperToSpecZeroOfFiniteTypeSubtypeMemOfNatNat} makes
  each of these structural maps proper.  Applying the verified target-locality
  instance \texttt{AlgebraicGeometry.IsProper.instIsZariskiLocalAtTarget}
  gives that $\mathbf P^r_X\to X$ is proper.  A closed immersion is proper, and the
  composite of proper morphisms is proper by
  \texttt{AlgebraicGeometry.IsProper.comp}; applying this to
  $Y\hookrightarrow\mathbf P^r_X\to X$ proves the claim.
\end{proof}

\begin{lemma}[Properness is local on the target]
  \label{lemma:nagata-projective-target-local}
  \lean{AlgebraicGeometry.IsProper.instIsZariskiLocalAtTarget}

  Let $f\colon Y\to X$ be a morphism.  If $X$ has an open cover
  $(V_i)$ such that every $f^{-1}(V_i)\to V_i$ is proper, then $f$ is
  proper.
\end{lemma}

\begin{proof}
  Properness is a Zariski-local property on the target.  Apply the verified
  target-locality instance
  \texttt{AlgebraicGeometry.IsProper.instIsZariskiLocalAtTarget} to the open
  cover $(V_i)$; its local hypotheses are exactly the properness hypotheses in
  the statement, and it returns that $f$ is proper.

  \emph{Warning.}  The analogous assertion with ``projective'' in place of
  ``proper'' is false for the finite-projective-space definition used here;
  see Stacks Project Tag 08J0.  No projectivity conclusion is used below.
\end{proof}

\begin{proposition}[Restriction, base change, and functoriality of Proj]
  \label{prop:nagata-relative-proj-functoriality}
  \uses{def:nagata-relative-proj,lemma:nagata-relative-proj-exists,
    lemma:nagata-proj-chart-ring-cocycle}
  \lean{AlgebraicGeometry.RelativeProj.functoriality}

  Relative Proj has the following properties.
  \begin{enumerate}
    \item Its restriction to an affine open $V$ is canonically
      $\operatorname{Proj}(\Gamma(V,\mathcal A))$, compatibly on iterated
      restrictions.
    \item A degree-preserving equivalence of graded algebras inducing the
      chosen degree-zero base identification induces an isomorphism of Proj.
    \item A surjective degree-preserving morphism of graded algebras induces
      a closed immersion in the opposite direction on Proj.
    \item If $g\colon X'\to X$ is flat, then coefficientwise pullback gives a
      canonical base-change isomorphism
      \[
        \operatorname{Proj}_{X'}(g^*\mathcal A)
          \cong X'\times_X\operatorname{Proj}_X(\mathcal A).
       \]
    \end{enumerate}
    The same base-change construction is required for an arbitrary morphism
    $X'\to X$ when the universal blow-up lift is formed; that unrestricted
    comparison is recorded as a separate conditional interface below.
\end{proposition}

\begin{proof}
  For (1), use the restriction isomorphism in
  \Cref{lemma:nagata-relative-proj-exists}; transitivity of coefficientwise
  localization gives the iterated-restriction equality.  For (2), on an affine
  chart the graded equivalence carries each homogeneous denominator to a
  homogeneous denominator and induces an isomorphism of the corresponding
  degree-zero localizations, hence an isomorphism of standard affine opens.
  These maps agree on overlaps by the localization cocycle and glue uniquely.
  For (3), a surjection has surjective maps on every homogeneous localization;
  the induced maps on degree-zero rings give closed immersions of all standard
  affine opens, which glue to the asserted closed immersion.  Finally, for
  (4), on $V=\operatorname{Spec}(A)\subseteq X$, put
  $C=\Gamma(V,\mathcal A)$, and on
  $V'=\operatorname{Spec}(B)\subseteq X'$ over $V$, the chart at a
  homogeneous $f$ is identified on both sides using
  \[
    (C_{(f),0})\otimes_A B
      \xrightarrow{\sim}(B\otimes_A C)_{(f\otimes1),0},
    \qquad (x/f^m)\otimes b\longmapsto (x\otimes b)/(f\otimes1)^m .
  \]
  To justify this without a hidden cancellation argument, identify the
  degree-zero homogeneous localization with the filtered colimit
  $\varinjlim_m C_{md}$ whose transition map is multiplication by $f$.
  Tensor product commutes with this filtered colimit, giving the displayed
  isomorphism (in fact without using flatness).  On pure fractions its inverse
  sends $(x\otimes b)/(f\otimes1)^m$ to $(x/f^m)\otimes b$.  It is compatible
  with further localization, so the chart isomorphisms glue, and their
  structural maps identify the glued schemes over $X$.  This proves all four
  assertions.
\end{proof}

\subsection{The graded Rees algebra}

\begin{definition}[Rees algebra]
  \label{def:nagata-rees-algebra}
  \lean{reesAlgebra}

  Let $A$ be a ring and $I\subseteq A$ an ideal.  The Rees algebra is
  \[
    \mathcal R_A(I)=\bigoplus_{n\geq 0}I^nt^n\subseteq A[t].
  \]
  Its underlying ring is the subring of $A[t]$ consisting of polynomials whose
  coefficient of $t^n$ belongs to $I^n$.

  \emph{Implementation note.}  The current Mathlib object is the subalgebra
  \texttt{reesAlgebra I} of the polynomial ring.  The declarations
  \texttt{mem\_reesAlgebra\_iff}, \texttt{reesAlgebra.monomial\_mem},
  \texttt{adjoin\_monomial\_eq\_reesAlgebra}, and
  \texttt{reesAlgebra.fg} provide the membership and finite-generation facts
  used below.  The grading and the pieces in the next blocks are blueprint
  data and must be defined explicitly.
\end{definition}

\begin{definition}[Homogeneous piece of the Rees algebra]
  \label{def:nagata-rees-piece}
  \uses{def:nagata-rees-algebra}
  \lean{ReesAlgebra.piece}

  For $n\in\mathbb N$, let $\mathcal R_A(I)_n$ be the $A$-submodule of Rees
  elements whose underlying polynomial is supported in degree $n$.
\end{definition}

\begin{lemma}[Membership in one Rees piece]
  \label{lemma:nagata-rees-piece-membership}
  \uses{def:nagata-rees-piece}
  \lean{ReesAlgebra.mem_piece_iff}

  An element $F\in\mathcal R_A(I)$ lies in $\mathcal R_A(I)_n$ if and only if
  there is $a\in I^n$ such that the underlying polynomial of $F$ is
  $at^n$.
\end{lemma}

\begin{proof}
  If $F$ is supported in degree $n$, coefficient extensionality gives
  $F=\operatorname{monomial}(n,F_n)$.  Membership in the Rees algebra gives
  $F_n\in I^n$.  Conversely, if the underlying polynomial is $at^n$ with $a\in I^n$, its
  support is the singleton degree $n$ and the coefficient lies in the required
  piece, so the defining membership condition for
  $\mathcal R_A(I)_n$ holds.
\end{proof}

\begin{lemma}[Multiplication of Rees pieces]
  \label{lemma:nagata-rees-piece-mul}
  \uses{def:nagata-rees-piece}
  \lean{ReesAlgebra.piece_mul_piece}

  Multiplication maps
  \[
    \mathcal R_A(I)_m\cdot\mathcal R_A(I)_n
      \subseteq\mathcal R_A(I)_{m+n}.
  \]
\end{lemma}

\begin{proof}
  \uses{lemma:nagata-rees-piece-membership}
  Write the two elements as $at^m$ and $bt^n$ with $a\in I^m$ and
  $b\in I^n$.  Their product is $(ab)t^{m+n}$ and
  $ab\in I^mI^n=I^{m+n}$.
\end{proof}

\begin{lemma}[Finite support of the homogeneous decomposition]
  \label{lemma:nagata-rees-component-finite-support}
  \uses{def:nagata-rees-piece}
  \lean{ReesAlgebra.component_finiteSupport}

  For $F\in\mathcal R_A(I)$, only finitely many homogeneous components
  $F_nt^n$ are nonzero.
\end{lemma}

\begin{proof}
  The components are indexed by the support of the underlying polynomial,
  which is finite.
\end{proof}

\begin{lemma}[Reconstruction from homogeneous components]
  \label{lemma:nagata-rees-sum-components}
  \uses{def:nagata-rees-piece}
  \lean{ReesAlgebra.sum_components}

  Every $F\in\mathcal R_A(I)$ is the finite sum of the monomials
  $F_nt^n$ regarded as elements of $\mathcal R_A(I)_n$.
\end{lemma}

\begin{proof}
  \uses{lemma:nagata-rees-component-finite-support,
    lemma:nagata-rees-piece-membership}
  Sum over the finite polynomial support.  Equality follows by comparing the
  coefficient in every degree.
\end{proof}

\begin{lemma}[Uniqueness of the homogeneous decomposition]
  \label{lemma:nagata-rees-components-unique}
  \uses{def:nagata-rees-piece}
  \lean{ReesAlgebra.components_unique}

  Let $E$ be a finite set of pairwise distinct degrees and choose one element
  $F_n\in\mathcal R_A(I)_n$ for each $n\in E$.  If
  $\sum_{n\in E}F_n=0$, then every $F_n$ is zero.
\end{lemma}

\begin{proof}
  \uses{lemma:nagata-rees-piece-membership}
  Apply the coefficient projection in a fixed degree $n\in E$.  Every summand
  of degree different from $n$ has zero $n$th component, so the projection of
  the displayed sum is $F_n=0$.  This holds for every $n\in E$.
\end{proof}

\begin{proposition}[Graded-ring structure on the Rees algebra]
  \label{prop:nagata-rees-graded}
  \uses{def:nagata-rees-piece}
  \lean{ReesAlgebra.instGradedRingPiece}

  The family $n\mapsto\mathcal R_A(I)_n$ defines an
  $\mathbb N$-graded-ring structure on $\mathcal R_A(I)$.
\end{proposition}

\begin{proof}
  \uses{lemma:nagata-rees-piece-mul,lemma:nagata-rees-sum-components,
    lemma:nagata-rees-components-unique}
  Use \Cref{lemma:nagata-rees-piece-mul} for graded multiplication,
  \Cref{lemma:nagata-rees-sum-components} for existence of the direct-sum
  decomposition, and \Cref{lemma:nagata-rees-components-unique} for
  uniqueness.
\end{proof}

\begin{lemma}[The degree-zero piece]
  \label{lemma:nagata-rees-degree-zero}
  \uses{def:nagata-rees-piece}
  \lean{ReesAlgebra.pieceZeroEquiv}

  Taking the constant coefficient gives a ring equivalence
  \[
    \mathcal R_A(I)_0\cong A.
  \]
\end{lemma}

\begin{proof}
  \uses{lemma:nagata-rees-piece-membership}
  Since $I^0=A$, a degree-zero Rees element is exactly a constant polynomial.
  The inverse sends $a$ to the constant monomial $at^0$.  The two composites
  are identities by coefficient extensionality.
\end{proof}

\begin{lemma}[Degree-one element associated to an ideal element]
  \label{lemma:nagata-rees-degree-one-element}
  \uses{def:nagata-rees-piece}
  \lean{ReesAlgebra.degreeOneElement}

  Every $a\in I$ determines a homogeneous element $at$ of degree one.  Its
  underlying polynomial and its coefficient in degree one are recorded by
  simp lemmas.
\end{lemma}

\begin{proof}
  \uses{lemma:nagata-rees-piece-membership}
  Since $a\in I=I^1$, the monomial $at$ belongs to the degree-one piece by
  \Cref{lemma:nagata-rees-piece-membership}.  The underlying-polynomial and
  coefficient formulas follow from the corresponding formulas for a
  polynomial monomial.
\end{proof}

\begin{proposition}[Finite generation over degree zero]
  \label{prop:nagata-rees-finite-generation}
  \uses{def:nagata-rees-algebra,def:nagata-rees-piece}
  \lean{AlgebraicGeometry.algebraFiniteType_reesGradingZero}

  If $I=(f_0,\ldots,f_r)$, then $\mathcal R_A(I)$ is generated as an algebra
  over its degree-zero piece by $f_0t,\ldots,f_rt$.  In particular it is of
  finite type over $\mathcal R_A(I)_0$.
\end{proposition}

\begin{proof}
  \uses{lemma:nagata-rees-degree-zero,lemma:nagata-rees-degree-one-element,
    lemma:nagata-rees-sum-components}
  A homogeneous coefficient in $I^n$ is a linear combination of degree-$n$
  monomials in the $f_i$.  Hence every homogeneous Rees element belongs to the
  algebra generated by the $f_it$, and
  \Cref{lemma:nagata-rees-sum-components} treats an arbitrary element.
  Transport finite generation along
  \Cref{lemma:nagata-rees-degree-zero}.
\end{proof}

\subsection{Affine blow-ups and their standard charts}

\begin{definition}[Affine blow-up]
  \label{def:nagata-affine-blowup}
  \uses{def:nagata-rees-algebra,def:nagata-affine-proj}
  \lean{AlgebraicGeometry.affineBlowup}

  For an ideal $I\subseteq A$, define
  \[
    \operatorname{Bl}_I(\operatorname{Spec}A)
      =\operatorname{Proj}(\mathcal R_A(I)).
  \]
  Its structure morphism is induced by the canonical degree-zero map
  $A\to\mathcal R_A(I)_0$.
\end{definition}

\begin{definition}[Affine blow-up chart algebra]
  \label{def:nagata-blowup-chart-ring}
  \lean{ReesAlgebra.blowupChartRing}

  Let $a\in I$.  Define $A[I/a]$ to be the $A$-subalgebra of $A_a$
  generated by the fractions $x/a$ with $x\in I$.
\end{definition}

\begin{lemma}[Normal form in a chart algebra]
  \label{lemma:nagata-chart-ring-normal-form}
  \uses{def:nagata-blowup-chart-ring}
  \lean{ReesAlgebra.mem_blowupChartRing_iff}

  Every element of $A[I/a]$ is represented in $A_a$ by $x/a^n$ for some
  $n\geq 0$ and $x\in I^n$.
\end{lemma}

\begin{proof}
  The generators have the required form with $n=1$.  The set of such
  fractions contains the image of $A$, and is closed under addition and
  multiplication after replacing two denominators by a common power of $a$.
  It therefore contains the generated subalgebra.  Conversely,
  $x\in I^n$ is a sum of products of $n$ elements of $I$, so $x/a^n$ belongs
  to the generated algebra.
\end{proof}

\begin{lemma}[Equality criterion for chart fractions]
  \label{lemma:nagata-chart-ring-fraction-equality}
  \lean{ReesAlgebra.chartFraction_eq_iff}

  Two representatives $x/a^m$ and $y/a^n$ define the same element of
  $A[I/a]$ exactly when they are equal in the localization $A_a$; equivalently,
  some power of $a$ annihilates $a^nx-a^my$ after aligning denominators.
\end{lemma}

\begin{proof}
  \uses{lemma:nagata-chart-ring-normal-form}
  This is the equality relation in localization away from the powers of $a$.
  The statement must use localization equivalence, not cancellation by $a$,
  because $a$ may be a zero divisor in $A$.
\end{proof}

\begin{definition}[Homogeneous-localization chart map]
  \label{def:nagata-homogeneous-away-to-chart}
  \uses{def:nagata-rees-algebra,
    def:nagata-blowup-chart-ring}
  \lean{ReesAlgebra.homogeneousAwayToChart}

  For $a\in I$, define $\theta_a$ on a degree-zero homogeneous fraction by
  \[
    \theta_a\!\left(\frac{xt^n}{(at)^n}\right)=\frac{x}{a^n},
    \qquad x\in I^n.
  \]
\end{definition}

\begin{lemma}[Well-definedness of the homogeneous-localization chart map]
  \label{lemma:nagata-homogeneous-away-to-chart}
  \uses{def:nagata-homogeneous-away-to-chart}
  \lean{ReesAlgebra.homogeneousAwayToChart_wellDefined}

  The assignment $\theta_a$ is independent of the chosen homogeneous
  representative and defines a ring homomorphism from the degree-zero
  homogeneous localization to $A[I/a]$.
\end{lemma}

\begin{proof}
  \uses{lemma:nagata-chart-ring-fraction-equality}
  Multiplication and addition are compatible with the displayed formula.
  Well-definedness follows from
  \Cref{lemma:nagata-chart-ring-fraction-equality}, applied to the equality
  relation in homogeneous localization after forgetting the grading.
\end{proof}

\begin{lemma}[Surjectivity of the chart map]
  \label{lemma:nagata-homogeneous-away-to-chart-surjective}
  \lean{ReesAlgebra.homogeneousAwayToChart_surjective}

  The constructed map $\theta_a$ is surjective.
\end{lemma}

\begin{proof}
  \uses{lemma:nagata-homogeneous-away-to-chart,lemma:nagata-chart-ring-normal-form}
  A normal-form element $x/a^n$ is the image of
  $(xt^n)/(at)^n$.
\end{proof}

\begin{lemma}[Injectivity of the chart map]
  \label{lemma:nagata-homogeneous-away-to-chart-injective}
  \lean{ReesAlgebra.homogeneousAwayToChart_injective}

  The constructed map $\theta_a$ is injective.
\end{lemma}

\begin{proof}
  \uses{lemma:nagata-homogeneous-away-to-chart,
    lemma:nagata-chart-ring-fraction-equality}
  If $x/a^n$ is zero in $A_a$, then $a^N x=0$ for some $N$.  Multiplying the
  homogeneous numerator and denominator by $(at)^N$ shows that the original
  homogeneous fraction is zero.
\end{proof}

\begin{definition}[Canonical affine blow-up chart morphism]
  \label{def:nagata-affine-blowup-chart-map}
  \uses{def:nagata-affine-blowup,
    def:nagata-homogeneous-away-to-chart,
    prop:nagata-proj-standard-open}
  \lean{AlgebraicGeometry.affineBlowup_basicOpenToSpec}

  Define
  \[
    \chi_a\colon \operatorname{Spec}(A[I/a])\longrightarrow D_+(at)
  \]
  to be the affine morphism induced contravariantly by $\theta_a$ after the
  standard identification of $D_+(at)$ with the spectrum of the degree-zero
  homogeneous localization.
\end{definition}

\begin{proposition}[Affine chart identification]
  \label{prop:nagata-affine-blowup-chart}
  \uses{def:nagata-affine-blowup-chart-map}
  \lean{AlgebraicGeometry.affineBlowup_basicOpenIsoSpec}

  The constructed morphism $\chi_a$ is an isomorphism over
  $\operatorname{Spec}(A)$.
\end{proposition}

\begin{proof}
  \uses{lemma:nagata-homogeneous-away-to-chart-surjective,
    lemma:nagata-homogeneous-away-to-chart-injective}
  The ring map $\theta_a$ is bijective by the two cited lemmas.  The induced
  affine morphism is therefore an isomorphism.  Compatibility with the
  structure morphism is checked on the image of every $r\in A$.
\end{proof}

\begin{lemma}[Degree-one generators cover Proj]
  \label{lemma:nagata-rees-generators-cover-proj}
  \lean{AlgebraicGeometry.Proj.iSup_basicOpen_eq_top}

  If $I=(f_0,\ldots,f_r)$, then the standard opens
  $D_+(f_it)$ cover $\operatorname{Proj}(\mathcal R_A(I))$.
\end{lemma}

\begin{proof}
  \uses{prop:nagata-rees-finite-generation}
  The elements $f_it$ generate the irrelevant ideal by
  \Cref{prop:nagata-rees-finite-generation}.  Apply
  \texttt{AlgebraicGeometry.Proj.iSup\_basicOpen\_eq\_top} to obtain the
  asserted cover; equivalently, a homogeneous prime cannot contain all
  generators.
\end{proof}

\begin{proposition}[Finite affine chart cover]
  \label{prop:nagata-affine-blowup-chart-cover}
  \lean{AlgebraicGeometry.affineBlowup_affineOpenCover}

  If $I=(f_0,\ldots,f_r)$, the affine opens
  \[
    \operatorname{Spec}(A[I/f_i])
  \]
  form a finite open cover of
  $\operatorname{Bl}_I(\operatorname{Spec}A)$.
\end{proposition}

\begin{proof}
  \uses{prop:nagata-affine-blowup-chart,lemma:nagata-rees-generators-cover-proj}
  The opens $D_+(f_it)$ cover by
  \Cref{lemma:nagata-rees-generators-cover-proj}.  Replace each one by the
  affine spectrum supplied by \Cref{prop:nagata-affine-blowup-chart}.  The
  compatibility of the open embeddings with the structure map is part of that
  chart isomorphism.
\end{proof}

\begin{lemma}[Blow-up of the unit ideal]
  \label{lemma:nagata-affine-blowup-top}
  \uses{def:nagata-affine-blowup}
  \lean{AlgebraicGeometry.affineBlowup_top_iso}

  The blow-up of $\operatorname{Spec}(A)$ in the unit ideal is canonically
  isomorphic to $\operatorname{Spec}(A)$ over $\operatorname{Spec}(A)$.
\end{lemma}

\begin{proof}
  \uses{prop:nagata-affine-blowup-chart}
  The Rees algebra of the unit ideal is $A[t]$.  The single standard open
  $D_+(t)$ covers Proj, and its chart algebra $A[A/1]$ is $A$.
\end{proof}

\begin{lemma}[Blow-up of the zero ideal]
  \label{lemma:nagata-affine-blowup-bot}
  \uses{def:nagata-affine-blowup}
  \lean{AlgebraicGeometry.affineBlowup_bot_isEmpty}

  The blow-up in the zero ideal is empty: the positive-degree part of its Rees
  algebra vanishes, so the irrelevant ideal is zero and Proj has no points.
\end{lemma}

\begin{proof}
  If $I=0$, then $I^n=0$ for every $n>0$.  Hence the irrelevant ideal of
  $\mathcal R_A(0)$ is zero.  A point of Proj would be a homogeneous prime not
  containing the irrelevant ideal, which is impossible because every prime
  contains zero.
\end{proof}

\subsection{Localization compatibility and coherence}

\begin{lemma}[Localization commutes with powers of an ideal]
  \label{lemma:nagata-localization-ideal-pow}
  \lean{Ideal.map_localization_pow}

  Let $s\in A$.  For every $n$ there is a canonical equality
  \[
    (I^n)A_s=(IA_s)^n.
  \]
\end{lemma}

\begin{proof}
  Extension of ideals along a ring homomorphism preserves products and the
  unit ideal.  Induct on $n$.
\end{proof}

\begin{lemma}[Localization of each Rees piece]
  \label{lemma:nagata-localization-rees-piece}
  \lean{ReesAlgebra.localization_piece_equiv}

  Flatness of $A\to A_s$ gives an $A_s$-linear equivalence
  \[
    A_s\otimes_A\mathcal R_A(I)_n
      \cong\mathcal R_{A_s}(IA_s)_n
  \]
  for every $n$.
\end{lemma}

\begin{proof}
  \uses{lemma:nagata-localization-ideal-pow,lemma:nagata-rees-piece-membership}
  Identify the two modules with
  $A_s\otimes_AI^n$ and $(IA_s)^n$, then use exactness of localization and
  \Cref{lemma:nagata-localization-ideal-pow}.
\end{proof}

\begin{definition}[Coefficientwise Rees-localization map]
  \label{def:nagata-rees-localization-map}
  \uses{def:nagata-rees-algebra}
  \lean{ReesAlgebra.localizationHom}

  Define
  \[
    \rho_s\colon A_s\otimes_A\mathcal R_A(I)
      \longrightarrow\mathcal R_{A_s}(IA_s)
  \]
  degreewise by sending a coefficient $x\in I^n$ to $x/1\in(IA_s)^n$.
\end{definition}

\begin{proposition}[Graded localization of the Rees algebra]
  \label{prop:nagata-rees-localization}
  \uses{def:nagata-rees-localization-map}
  \lean{ReesAlgebra.localizationEquiv}

  The map $\rho_s$ is a graded ring equivalence.
\end{proposition}

\begin{proof}
  \uses{lemma:nagata-localization-rees-piece,
    lemma:nagata-localization-ideal-pow}
  For each $n$, let $\rho_{s,n}$ be the equivalence of
  \Cref{lemma:nagata-localization-rees-piece}.  Taking their direct sum is
  legitimate because every Rees element has finite homogeneous support.  On
  homogeneous tensors of degrees $m,n$, both ways around the multiplication
  square have coefficient $(xy)/1$ in degree $m+n$; the equality of localized
  ideal powers from \Cref{lemma:nagata-localization-ideal-pow} proves
  multiplicativity.  The direct sum of the inverse maps is constructed in
  the same way, and the two composites are identities in each degree, hence
  on finite sums.  This gives the claimed graded ring equivalence.
\end{proof}

\begin{lemma}[Compatibility with degree zero]
  \label{lemma:nagata-rees-localization-degree-zero}
  \lean{ReesAlgebra.localizationEquiv_pieceZero}

  On degree zero, the constructed coefficientwise localization map is the standard ring
  equivalence $A_s\otimes_AA\cong A_s$.
\end{lemma}

\begin{proof}
  \uses{prop:nagata-rees-localization,lemma:nagata-rees-degree-zero}
  Evaluate both maps on a constant monomial and use ring-homomorphism
  extensionality.
\end{proof}

\begin{lemma}[Compatibility with degree-one generators]
  \label{lemma:nagata-rees-localization-degree-one}
  \lean{ReesAlgebra.localizationEquiv_degreeOneElement}

  The element $at$ maps to the degree-one element $(a/1)t$ of the localized
  Rees algebra.
\end{lemma}

\begin{proof}
  \uses{prop:nagata-rees-localization,lemma:nagata-rees-degree-one-element}
  Unfold the degreewise map in
  \Cref{prop:nagata-rees-localization}.  In degree one it is extension of the
  coefficient map $I\to IA_s$, so it sends $a$ to $a/1$ and therefore sends
  $at$ to $(a/1)t$.
\end{proof}

\begin{lemma}[Localization at one]
  \label{lemma:nagata-rees-localization-one}
  \uses{lemma:nagata-rees-sum-components}
  \lean{ReesAlgebra.localizationEquiv_one}

  After the canonical identifications $A_1\cong A$, the Rees localization
  equivalence at $1$ is the identity.
\end{lemma}

\begin{proof}
  \uses{prop:nagata-rees-localization}
  Both the localization equivalence and the identity send each homogeneous
  monomial $at^n$ to itself under $A_1\cong A$.  Use
  \Cref{lemma:nagata-rees-sum-components} and ring-homomorphism extensionality.
\end{proof}

\begin{lemma}[Iterated localization coherence]
  \label{lemma:nagata-rees-localization-mul}
  \lean{ReesAlgebra.localizationEquiv_mul}

  Localizing first at $s$ and then at $t$ gives the same graded equivalence as
  localizing at $st$, after the canonical equivalence
  $(A_s)_t\cong A_{st}$.
\end{lemma}

\begin{proof}
  \uses{prop:nagata-rees-localization}
  Both maps send a homogeneous coefficient $x\in I^n$ to $x/1$ and commute
  with the inclusions of the localized base ring.  Graded ring-homomorphism
  extensionality gives equality.
\end{proof}

\begin{lemma}[Naturality under a ring isomorphism]
  \label{lemma:nagata-rees-localization-natural}
  \lean{ReesAlgebra.localizationEquiv_natural}

  The Rees localization equivalence commutes with transport of $I$ and $s$
  along a ring isomorphism.
\end{lemma}

\begin{proof}
  \uses{prop:nagata-rees-localization}
  It suffices to compare the maps on homogeneous monomials.  Both routes send
  $at^n$ to the monomial with coefficient obtained by transporting $a$ and
  then localizing it.  Naturality of localization of rings gives equality of
  those coefficients.
\end{proof}

\begin{definition}[Affine blow-up restriction morphism]
  \label{def:nagata-affine-blowup-restriction-map}
  \uses{def:nagata-affine-blowup,
    def:nagata-rees-localization-map,
    prop:nagata-relative-proj-functoriality}
  \lean{AlgebraicGeometry.affineBlowup_localizationMap}

  Define
  \[
    r_s\colon
    \operatorname{Bl}_{IA_s}(\operatorname{Spec}A_s)
      \longrightarrow
    \operatorname{Bl}_I(\operatorname{Spec}A)
      \times_{\operatorname{Spec}A}\operatorname{Spec}A_s
  \]
  by applying the relative-Proj morphism of
  \Cref{prop:nagata-relative-proj-functoriality} to $\rho_s$ and the
  canonical degree-zero identifications.
\end{definition}

\begin{proposition}[Restriction of an affine blow-up]
  \label{prop:nagata-affine-blowup-restrict}
  \uses{def:nagata-affine-blowup-restriction-map}
  \lean{AlgebraicGeometry.affineBlowup_localizationIso}

  The constructed morphism $r_s$ is an isomorphism over
  $\operatorname{Spec}(A_s)$.
\end{proposition}

\begin{proof}
  \uses{prop:nagata-rees-localization,
    lemma:nagata-rees-localization-degree-zero,
    prop:nagata-relative-proj-functoriality}
  By \Cref{prop:nagata-rees-localization}, the coefficientwise localization
  of the Rees algebra is the Rees algebra of $IA_s$.  On the standard open
  indexed by a homogeneous generator, the induced map is the isomorphism of
  degree-zero chart rings from the localization equivalence.  These chart
  isomorphisms agree on overlaps by the localization cocycle, and therefore
  glue to an inverse to $r_s$.  The degree-zero piece is $A_s$ on both sides,
  so the glued inverse and $r_s$ are morphisms over $\operatorname{Spec}(A_s)$.
\end{proof}

\begin{lemma}[Identity and cocycle laws for affine restriction]
  \label{lemma:nagata-affine-blowup-restrict-cocycle}
  \lean{AlgebraicGeometry.affineBlowup_localizationIso_cocycle}

  The restriction isomorphism is the identity for $s=1$, and its two possible
  composites on an iterated basic open $D(st)$ agree literally after the
  chosen canonical associativity isomorphisms.
\end{lemma}

\begin{proof}
  \uses{lemma:nagata-rees-localization-one,lemma:nagata-rees-localization-mul,
    prop:nagata-affine-blowup-restrict}
  On each standard chart, the map for $s=1$ is induced by the identity of the
  Rees algebra, hence is the identity.  For $s,t$, both composites send a
  homogeneous fraction to its image in the localization at $st$; the two
  coefficient maps are equal by
  \Cref{lemma:nagata-rees-localization-one,lemma:nagata-rees-localization-mul}.
  The chart maps therefore agree on a cover, proving the asserted identity
  and cocycle equalities globally.
\end{proof}

\subsection{Global blow-up by gluing}

The following gluing datum is formed from the affine restriction isomorphisms;
all its coherence maps are checked on basic affine refinements.

\begin{definition}[Local blow-up model on an affine open]
  \label{def:nagata-global-blowup-local-model}
  \uses{def:nagata-affine-blowup}
  \lean{Scheme.IdealSheafData.localBlowup}

  Let $\mathcal I\subseteq\mathcal O_X$ be quasi-coherent and let
  $V=\operatorname{Spec}(A)\subseteq X$ be affine.  Put
  $I_V=\Gamma(V,\mathcal I)\subseteq A$ and define the local model
  \[
    \mathsf B_V(\mathcal I)=\operatorname{Bl}_{I_V}(V)\longrightarrow V.
  \]
\end{definition}

\begin{definition}[Basic-overlap comparison map]
  \label{def:nagata-global-blowup-overlap-basic-map}
  \uses{def:nagata-global-blowup-local-model,
    def:nagata-affine-blowup-restriction-map,
    prop:nagata-affine-blowup-restrict}
  \lean{Scheme.IdealSheafData.localBlowup_overlapBasicMap}

  If a basic affine open $D(s)$ lies in $V_i\cap V_j$, define
  \[
    \omega_{ij,s}\colon
      \mathsf B_{V_i}(\mathcal I)|_{D(s)}
      \longrightarrow
      \mathsf B_{V_j}(\mathcal I)|_{D(s)}
  \]
  as $r_{j,s}\circ r_{i,s}^{-1}$: first use the inverse of the restriction
  isomorphism $r_{i,s}$ to move from the first local model to the blow-up of
  the localized ideal, and then use $r_{j,s}$ to reach the second local model.
  Both restriction maps have the blow-up of the same localized ideal as their
  common source.
\end{definition}

\begin{lemma}[The basic-overlap comparison is an isomorphism]
  \label{lemma:nagata-global-blowup-overlap-basic}
  \uses{def:nagata-global-blowup-overlap-basic-map}
  \lean{Scheme.IdealSheafData.localBlowup_overlapBasicIso}

  The constructed map $\omega_{ij,s}$ is an isomorphism over $D(s)$.
\end{lemma}

\begin{proof}
  \uses{prop:nagata-affine-blowup-restrict}
  Quasi-coherence identifies both restricted ideals with the same localized
  ideal.  Thus $r_{i,s}$ and $r_{j,s}$ are isomorphisms by
  \Cref{prop:nagata-affine-blowup-restrict}; their composite with the inverse
  of $r_{i,s}$ is an isomorphism, with inverse $r_{i,s}\circ r_{j,s}^{-1}$.
\end{proof}

\begin{lemma}[Independence of the basic refinement]
  \label{lemma:nagata-global-blowup-overlap-independent}
  \lean{Scheme.IdealSheafData.localBlowup_overlapBasicIso_independent}

  On the intersection of two chosen basic refinements, the two overlap
  constructed maps $\omega_{ij,s}$ agree.
\end{lemma}

\begin{proof}
  \uses{lemma:nagata-affine-blowup-restrict-cocycle,
    lemma:nagata-global-blowup-overlap-basic}
  Refine once more to a basic open contained in both.  There the two maps are
  the same composite of restriction isomorphisms, by
  \Cref{lemma:nagata-affine-blowup-restrict-cocycle}.  Since these refinements
  cover the source overlap, equality of morphisms follows by locality on the
  source.
\end{proof}

\begin{lemma}[Identity and symmetry of overlap isomorphisms]
  \label{lemma:nagata-global-blowup-overlap-id-symm}
  \lean{Scheme.IdealSheafData.localBlowup_overlapIso_refl_symm}

  The overlap isomorphism of a chart with itself is the identity, and the
  overlap isomorphism from chart $i$ to chart $j$ is inverse to the one from
  $j$ to $i$.
\end{lemma}

\begin{proof}
  \uses{lemma:nagata-global-blowup-overlap-independent}
  On a basic refinement, $\omega_{ii,s}=r_{i,s}r_{i,s}^{-1}=1$ and
  $\omega_{ji,s}\omega_{ij,s}=r_{j,s}r_{i,s}^{-1}r_{i,s}r_{j,s}^{-1}=1$.
  The equalities descend from the basic cover.
\end{proof}

\begin{lemma}[Triple-overlap cocycle]
  \label{lemma:nagata-global-blowup-overlap-cocycle}
  \lean{Scheme.IdealSheafData.localBlowup_overlapIso_cocycle}

  On every triple overlap, the composite of the $(i,j)$ and $(j,k)$ overlap
  isomorphisms equals the $(i,k)$ overlap isomorphism.
\end{lemma}

\begin{proof}
  \uses{lemma:nagata-affine-blowup-restrict-cocycle,
    lemma:nagata-global-blowup-overlap-independent}
  Cover the triple overlap by basic affine opens.  On each such open,
  $\omega_{jk,s}\omega_{ij,s}=r_{k,s}r_{j,s}^{-1}r_{j,s}r_{i,s}^{-1}
  =r_{k,s}r_{i,s}^{-1}=\omega_{ik,s}$; this is precisely the affine
  restriction cocycle.  Equality on the cover proves the claim.
\end{proof}

\begin{definition}[Raw global blow-up gluing datum]
  \label{def:nagata-global-blowup-glue-data}
  \uses{def:nagata-global-blowup-local-model,
    def:nagata-global-blowup-overlap-basic-map,
    lemma:nagata-global-blowup-overlap-independent}
  \lean{Scheme.IdealSheafData.blowupRawGlueData}

  Fix an affine cover $(V_i)_i$ of $X$.  Define
  $\mathsf G(X,\mathcal I,(V_i))$ to consist of the local schemes
  $\mathsf B_{V_i}(\mathcal I)$, their inverse images of $V_i\cap V_j$, and
  the overlap maps obtained by gluing the maps $\omega_{ij,s}$ on a chosen
  basic affine cover of every intersection.  The equality on intersections
  required to glue these local maps is
  \Cref{lemma:nagata-global-blowup-overlap-independent}.
\end{definition}

\begin{lemma}[Validity of the global blow-up gluing datum]
  \label{lemma:nagata-global-blowup-glue-data-valid}
  \uses{def:nagata-global-blowup-glue-data}
  \lean{Scheme.IdealSheafData.blowupRawGlueData_isGlueData}

  The raw datum $\mathsf G(X,\mathcal I,(V_i))$ has isomorphism-valued
  overlap maps satisfying the identity, symmetry, and triple-overlap cocycle
  axioms; hence it is a scheme gluing datum.
\end{lemma}

\begin{proof}
  \uses{lemma:nagata-global-blowup-overlap-basic,
    lemma:nagata-global-blowup-overlap-id-symm,
    lemma:nagata-global-blowup-overlap-cocycle}
  The overlap maps are isomorphisms by
  \Cref{lemma:nagata-global-blowup-overlap-basic}; independence of refinements,
  identity, symmetry, and the triple cocycle are given by the cited coherence
  lemmas.  Hence the hypotheses of
  \texttt{AlgebraicGeometry.Scheme.GlueData.glued} are satisfied.
\end{proof}

\begin{definition}[Global blow-up]
  \label{def:nagata-global-blowup}
  \uses{def:nagata-global-blowup-glue-data,
    lemma:nagata-global-blowup-glue-data-valid}
  \lean{Scheme.IdealSheafData.blowup}

  Let $\operatorname{Bl}_{\mathcal I}(X)$ be the scheme produced by applying
  \texttt{AlgebraicGeometry.Scheme.GlueData.glued} to the valid datum of
  \Cref{lemma:nagata-global-blowup-glue-data-valid}.
\end{definition}

\begin{lemma}[Compatibility of the local structure maps]
  \label{lemma:nagata-global-blowup-local-maps-agree}
  \uses{def:nagata-global-blowup-overlap-basic-map}
  \lean{Scheme.IdealSheafData.localBlowupMap_compatible}

  On every overlap, the two local structure maps to $X$ agree after applying
  the constructed overlap map $\omega_{ij}$.
\end{lemma}

\begin{proof}
  \uses{prop:nagata-affine-blowup-restrict}
  On a basic refinement, both composites are the structural map from the
  blow-up of the localized ideal to $D(s)\subseteq X$; the degree-zero part of
  the restriction isomorphism makes them equal.  Equality on the basic cover
  gives equality on the overlap.
\end{proof}

\begin{definition}[Global blow-up morphism]
  \label{def:nagata-global-blowup-map}
  \uses{def:nagata-global-blowup,
    def:nagata-global-blowup-local-model,
    lemma:nagata-global-blowup-local-maps-agree}
  \lean{Scheme.IdealSheafData.blowupMap}

  Define
  \[
    b_{X,\mathcal I}\colon\operatorname{Bl}_{\mathcal I}(X)\longrightarrow X
  \]
  by gluing the local structure maps
  $\mathsf B_{V_i}(\mathcal I)\to V_i\to X$.  Their overlap equality is
  \Cref{lemma:nagata-global-blowup-local-maps-agree}, so the morphism-gluing
  interface applies.
\end{definition}

\begin{definition}[Affine-restriction comparison for the global blow-up]
  \label{def:nagata-global-blowup-affine-restriction-map}
  \uses{def:nagata-global-blowup-map,
    def:nagata-global-blowup-local-model}
  \lean{Scheme.IdealSheafData.preimageAffineMap}

  If $V=\operatorname{Spec}(A)\subseteq X$ and $I_V$ represents
  $\mathcal I|_V$, define
  \[
    q_V\colon b_{X,\mathcal I}^{-1}(V)\longrightarrow
      \operatorname{Bl}_{I_V}(\operatorname{Spec}A)
  \]
  by gluing the affine restriction maps on the basic intersections of $V$
  with the chosen affine cover of $X$.
\end{definition}

\begin{proposition}[Affine restriction of the global blow-up]
  \label{prop:nagata-global-blowup-affine-restrict}
  \uses{def:nagata-global-blowup-affine-restriction-map}
  \lean{Scheme.IdealSheafData.preimageAffineIso}

  The constructed comparison $q_V$ is an isomorphism over $V$.
\end{proposition}

\begin{proof}
  \uses{prop:nagata-affine-blowup-restrict,
    lemma:nagata-global-blowup-overlap-cocycle,
    lemma:nagata-global-blowup-local-maps-agree}
  Cover $V$ by basic opens contained in the chosen affine charts.  On each,
  the comparison is an isomorphism by
  \Cref{prop:nagata-affine-blowup-restrict}; the overlap cocycle and local-map
  compatibility show that these maps and their inverses agree.  The gluing
  theorem therefore produces mutually inverse global morphisms over $V$.
\end{proof}

\begin{definition}[Comparison between two affine-cover constructions]
  \label{def:nagata-global-blowup-cover-comparison}
  \uses{def:nagata-global-blowup,
    def:nagata-global-blowup-affine-restriction-map,
    lemma:nagata-global-blowup-overlap-cocycle}
  \lean{Scheme.IdealSheafData.blowupCoverComparison}

  Given two affine covers of $X$, choose their intersection refinement.  On
  every refined affine chart use the two affine-restriction isomorphisms to
  define a comparison, and glue these comparisons after transporting them to
  the original charts.  The overlap equalities are those of
  \Cref{lemma:nagata-global-blowup-overlap-cocycle}.
  The refinement maps are independent of the selected basic chart by the
  overlap cocycle, so the glued comparison is well-defined as a morphism of
  the two gluing data.
\end{definition}

\begin{proposition}[Independence of the affine cover]
  \label{prop:nagata-global-blowup-cover-independent}
  \uses{def:nagata-global-blowup-cover-comparison}
  \lean{Scheme.IdealSheafData.blowupIso_of_cover}

  The constructed comparison is the unique isomorphism over $X$ restricting
  to the identity comparison on every common affine refinement.
\end{proposition}

\begin{proof}
  \uses{prop:nagata-global-blowup-affine-restrict,
    lemma:nagata-global-blowup-overlap-cocycle}
  By \Cref{prop:nagata-global-blowup-affine-restrict}, the comparison is an
  isomorphism on every member of the common refinement.  The overlap cocycle
  makes these local isomorphisms compatible, so they glue; the glued inverse
  is obtained from the local inverses.  Uniqueness follows because morphisms
  of schemes are determined on an affine open cover.
\end{proof}

\subsection{Complement, properness, and principalization}

\begin{lemma}[The center ideal is the unit ideal off its support]
  \label{lemma:nagata-ideal-unit-off-support}
  \lean{Scheme.IdealSheafData.restrict_compl_support_eq_top}

  If $Z=V(\mathcal I)$ and $U=X\setminus Z$, then
  $\mathcal I|_U=\mathcal O_U$.
\end{lemma}

\begin{proof}
  At a point outside $V(\mathcal I)$, the stalk $\mathcal I_x$ is not
  contained in the maximal ideal of $\mathcal O_{X,x}$, so it contains a unit
  and equals $\mathcal O_{X,x}$.  Equality of quasi-coherent ideals can be
  checked on stalks.
\end{proof}

\begin{lemma}[Affine complement-of-center isomorphism]
  \label{lemma:nagata-affine-blowup-complement}
  \uses{def:nagata-affine-blowup}
  \lean{AlgebraicGeometry.affineBlowup_compl_zeroLocus_iso}

  For $Z=V(I)\subseteq\operatorname{Spec}(A)$, the affine blow-up morphism
  restricts to an isomorphism over $\operatorname{Spec}(A)\setminus Z$.
\end{lemma}

\begin{proof}
  \uses{lemma:nagata-affine-blowup-top,prop:nagata-affine-blowup-restrict,
    lemma:nagata-affine-blowup-restrict-cocycle}
  The complement is covered by the basic opens $D(a)$, $a\in I$.  On $D(a)$
  the localized ideal is the unit ideal, so
  \Cref{prop:nagata-affine-blowup-restrict} and
  \Cref{lemma:nagata-affine-blowup-top} identify the blow-up with $D(a)$.
  The identifications agree on overlaps by
  \Cref{lemma:nagata-affine-blowup-restrict-cocycle}, and hence glue.
\end{proof}

\begin{theorem}[Complement of the center]
  \label{thm:nagata-global-blowup-complement}
  \lean{Scheme.IdealSheafData.blowup_compl_support_iso}

  Let $U=X\setminus V(\mathcal I)$.  The restriction
  \[
    b_{X,\mathcal I}^{-1}(U)\longrightarrow U
  \]
  is an isomorphism.
\end{theorem}

\begin{proof}
  \uses{prop:nagata-global-blowup-affine-restrict,
    lemma:nagata-affine-blowup-complement}
  Cover $X$ by affine opens.  Under
  \Cref{prop:nagata-global-blowup-affine-restrict}, the restriction over each
  affine is the affine complement isomorphism of
  \Cref{lemma:nagata-affine-blowup-complement}.  These local isomorphisms glue,
  and isomorphisms are local on source and target.
\end{proof}

\begin{definition}[$U$-admissible blow-up, {\cite[\href{https://stacks.math.columbia.edu/tag/080K}{Tag 080K}]{stacks-project}}]
  \label{def:nagata-admissible-blowup}
  \uses{def:nagata-global-blowup}
  \lean{AlgebraicGeometry.AdmissibleBlowup}

  Let $U\hookrightarrow X$ be a quasi-compact open immersion.  A
  \emph{$U$-admissible blow-up datum} consists of a finite type
  quasi-coherent ideal $\mathcal I\subseteq\mathcal O_X$ with
  $\mathcal I|_U=\mathcal O_U$, a morphism
  $\pi\colon X'\to X$, and a specified isomorphism
  $X'\cong\operatorname{Bl}_{\mathcal I}(X)$ over $X$.  A morphism equipped
  with such data is called a $U$-admissible blow-up.  The definition records
  the ideal and the identification with its blow-up; the equivalent
  finite-presentation closed-center formulation is stated separately once
  the required ideal/closed-immersion comparison is available.
\end{definition}

\begin{lemma}[Density is preserved by an admissible blow-up]
  \label{lemma:nagata-admissible-blowup-preserves-density}
  \uses{def:nagata-global-blowup-map,
    def:nagata-admissible-blowup,
    def:nagata-scheme-theoretically-dense-open}
  \lean{AlgebraicGeometry.IsAdmissible.isSchemeTheoreticallyDominant_preimage_ι,
    AlgebraicGeometry.Scheme.IdealSheafData.isSchemeTheoreticallyDominant_blowupMap_preimage_compl_support}

  Let $U\hookrightarrow X$ be a quasi-compact scheme-theoretically dense open
  immersion, and let $\mathcal I$ be a finite type quasi-coherent ideal with
  $\mathcal I|_U=\mathcal O_U$.  Then the inverse image of $U$ is
  scheme-theoretically dense in $\operatorname{Bl}_{\mathcal I}(X)$.
  Consequently the same assertion holds after any finite composite of
  $U$-admissible blow-ups.
\end{lemma}

\begin{proof}
  \uses{thm:nagata-blowup-principalizes-center,
    def:nagata-invertible-ideal,
    thm:nagata-global-blowup-complement,
    def:nagata-admissible-blowup}
  First, the complement of the support of an invertible ideal is
  scheme-theoretically dense: on an affine open where the ideal is generated
  by a single non-zero-divisor $a$, the restriction map to the complement of
  the support is the localization away from $a$, which is injective.
  Applying this on $\operatorname{Bl}_{\mathcal I}(X)$ to the extended ideal
  $\mathcal I\mathcal O_{\operatorname{Bl}_{\mathcal I}(X)}$, which is
  invertible by \Cref{thm:nagata-blowup-principalizes-center}, shows that the
  inverse image of the complement $W$ of the center is scheme-theoretically
  dense in the blow-up.  Density now descends along the partial isomorphism:
  the blow-up map $\pi$ is an isomorphism over $W$ by
  \Cref{thm:nagata-global-blowup-complement}, hence also over the smaller
  open $U\subseteq W$; since $U$ is quasi-compact and scheme-theoretically
  dense in $X$, the inclusion $U\hookrightarrow W$ is scheme-theoretically
  dominant, and transporting it across the isomorphism over $W$ and composing
  with the dense open immersion $\pi^{-1}(W)\hookrightarrow
  \operatorname{Bl}_{\mathcal I}(X)$ shows that $\pi^{-1}(U)$ is
  scheme-theoretically dense.  For a finite composite, the inverse image of
  $U$ at each stage is again a quasi-compact scheme-theoretically dense open,
  so induction on the length of the composite applies the one-step lemma at
  each stage.
\end{proof}

\begin{definition}[Homogeneous polynomial map to the Rees algebra]
  \label{def:nagata-rees-polynomial-map}
  \uses{def:nagata-rees-algebra}
  \lean{ReesAlgebra.gradedPolynomialMap}

  Let $A$ be a ring, let $I\subseteq A$ be an ideal, and fix generators
  $I=(f_0,\ldots,f_r)$.  Define the $A$-algebra map
  \[
    \rho_I\colon A[T_0,\ldots,T_r]\longrightarrow\mathcal R_A(I)
  \]
  by $\rho_I(T_i)=f_it$, with each $T_i$ placed in degree one.
\end{definition}

\begin{lemma}[Properties of the homogeneous polynomial map]
  \label{lemma:nagata-rees-polynomial-surjection}
  \uses{def:nagata-rees-polynomial-map}
  \lean{ReesAlgebra.gradedPolynomial_surjective}

  The constructed map $\rho_I$ is graded and surjective.
\end{lemma}

\begin{proof}
  \uses{lemma:nagata-rees-degree-one-element,
    prop:nagata-rees-finite-generation}
  The assignment is homogeneous because $f_it$ lies in degree one by
  \Cref{lemma:nagata-rees-degree-one-element}.  A monomial
  $T_{i_1}\cdots T_{i_n}$ maps to
  $(f_{i_1}\cdots f_{i_n})t^n$.  Such products generate $I^n$, so every
  degree component is hit.  Every Rees element is a finite sum of homogeneous
  components, proving surjectivity.
\end{proof}

\begin{definition}[Projective-space map of an affine blow-up]
  \label{def:nagata-affine-blowup-projective-embedding}
  \uses{def:nagata-affine-blowup,def:nagata-rees-polynomial-map,
    def:nagata-relative-projective-space,def:nagata-relative-proj,
    prop:nagata-relative-proj-functoriality}
  \lean{AlgebraicGeometry.AffineBlowup.toProjectiveSpace}

  Under the data of \Cref{def:nagata-rees-polynomial-map}, define
  \[
    e_I\colon\operatorname{Bl}_I(\operatorname{Spec}A)
      \longrightarrow\mathbf P^r_A
  \]
  to be the contravariant relative-Proj morphism supplied by
  \Cref{prop:nagata-relative-proj-functoriality}, induced by $\rho_I$.
\end{definition}

\begin{lemma}[Closed immersion into relative projective space]
  \label{lemma:nagata-affine-blowup-projective-embedding}
  \uses{def:nagata-affine-blowup-projective-embedding}
  \lean{AlgebraicGeometry.affineBlowup_to_projectiveSpace_isClosedImmersion}

  The constructed morphism $e_I$ is a closed immersion over
  $\operatorname{Spec}(A)$.
\end{lemma}

\begin{proof}
  \uses{lemma:nagata-rees-polynomial-surjection,
    prop:nagata-relative-proj-functoriality}
  A surjective graded algebra map induces a closed immersion in the opposite
  direction on Proj.  The source of $\rho_I$ has Proj $\mathbf P^r_A$, and
  its degree-zero map is the identity of $A$, so the closed immersion is over
  $\operatorname{Spec}(A)$.
\end{proof}


\begin{proposition}[Projectivity of the affine blow-up]
  \label{prop:nagata-affine-blowup-projective}
  \uses{def:nagata-projective-morphism}
  \lean{AlgebraicGeometry.affineBlowup_isProjective}

  If \(I\) is finitely generated, the affine blow-up morphism
  \[
    \operatorname{Bl}_I(\operatorname{Spec}(A))
      \longrightarrow\operatorname{Spec}(A)
  \]
  is projective.
\end{proposition}

\begin{proof}
  \uses{lemma:nagata-affine-blowup-projective-embedding}
  Choose finitely many generators of \(I\).  By
  \Cref{lemma:nagata-affine-blowup-projective-embedding}, the blow-up is a
  closed subscheme of a finite-dimensional projective space over
  \(\operatorname{Spec}(A)\).  The composite of a closed immersion with
  \(\mathbf P^r_A\to\operatorname{Spec}(A)\) is projective.
\end{proof}

\begin{lemma}[Properness of the affine blow-up]
  \label{lemma:nagata-affine-blowup-proper}
  \lean{AlgebraicGeometry.AffineBlowup.isProper_toSpec}

  If \(I\) is finitely generated, the affine blow-up morphism is proper.
\end{lemma}

\begin{proof}
  \uses{prop:nagata-affine-blowup-projective,
    lemma:nagata-projective-is-proper}
  It is projective by \Cref{prop:nagata-affine-blowup-projective}, and is
  therefore proper by \Cref{lemma:nagata-projective-is-proper}.
\end{proof}

\begin{lemma}[Local properness of the global blow-up]
  \label{lemma:nagata-global-blowup-proper-local}
  \lean{Scheme.IdealSheafData.blowupMap_restrict_isProper}

  If $\mathcal I$ is of finite type, the restriction of
  $b_{X,\mathcal I}$ over every affine open of $X$ is proper.
\end{lemma}

\begin{proof}
  \uses{prop:nagata-global-blowup-affine-restrict,lemma:nagata-affine-blowup-proper}
  Finite type of $\mathcal I$ gives finite generation of the corresponding
  ideal on the affine open.  Transfer
  \Cref{lemma:nagata-affine-blowup-proper} through
  \Cref{prop:nagata-global-blowup-affine-restrict}.
\end{proof}

\begin{definition}[Finite global free generation]
  \label{def:nagata-global-finite-free-generation}
  \lean{Scheme.IdealSheafData.HasFiniteGlobalFreePresentation}

  A quasi-coherent ideal sheaf \(\mathcal I\subseteq\mathcal O_X\) has a
  \emph{finite global free presentation} if there are \(r\in\mathbb N\) and
  a surjective \(\mathcal O_X\)-linear morphism
  \(q\colon\mathcal O_X^{r+1}\twoheadrightarrow\mathcal I\).
  Equivalently, \(\mathcal I\) is generated by finitely many global sections.
\end{definition}

\begin{lemma}[Relative projective-space construction from a global presentation]
  \label{lemma:nagata-global-blowup-relative-projective-embedding}
  \uses{def:nagata-global-finite-free-generation,
    def:nagata-relative-projective-space,
    prop:nagata-relative-proj-functoriality,
    prop:nagata-global-blowup-affine-restrict,
    lemma:nagata-relative-proj-exists,
    lemma:nagata-affine-blowup-projective-embedding,
    def:nagata-rees-polynomial-map,
    lemma:nagata-rees-polynomial-surjection}
  \lean{Scheme.IdealSheafData.blowupMap_projectiveEmbedding_of_globalGenerators}

  Given a finite global free presentation
  \(q\colon\mathcal O_X^{r+1}\twoheadrightarrow\mathcal I\), the affine
  Rees-algebra maps glue to a closed immersion over \(X\)
  \[
    \operatorname{Bl}_{\mathcal I}(X)
      \lhook\joinrel\longrightarrow
    \mathbf P^r_X
      =\operatorname{Proj}_X
        \bigl(\operatorname{Sym}_{\mathcal O_X}(\mathcal O_X^{r+1})\bigr),
  \]
  whose composite with the structural morphism is \(b_{X,\mathcal I}\).
\end{lemma}

\begin{proof}
  Write \(\mathcal R_X(\mathcal I)=\bigoplus_{n\geq0}\mathcal I^n\), with
  multiplication induced by ideal multiplication.  The presentation \(q\)
  induces a surjection of graded algebras
  \[
    \operatorname{Sym}_{\mathcal O_X}(\mathcal O_X^{r+1})
      \twoheadrightarrow
    \mathcal R_X(\mathcal I)=\bigoplus_{n\geq0}\mathcal I^n,
  \]
  by sending the degree-one generators to their images in \(\mathcal I\) and
  then multiplying.  Surjectivity in degree \(n\) follows because
  \(\mathcal I^n\) is generated by products of \(n\) local sections of
  \(\mathcal I\).  Contravariant functoriality of relative Proj therefore
  gives a closed immersion
  \(\operatorname{Proj}_X\mathcal R_X(\mathcal I)\hookrightarrow\mathbf P^r_X\).
  On an affine open, this is the closed immersion induced by the
  polynomial-to-Rees map from
  \Cref{def:nagata-rees-polynomial-map}, whose surjectivity is
  \Cref{lemma:nagata-rees-polynomial-surjection}; hence it agrees with the affine blow-up
  embedding of \Cref{lemma:nagata-affine-blowup-projective-embedding}, after
  the restriction identification of \Cref{prop:nagata-global-blowup-affine-restrict}.
  The affine
  comparison isomorphisms and the relative-Proj restriction isomorphisms
  satisfy the cocycle condition, so these local maps glue to the displayed
  closed immersion of the global blow-up.  Its composite with the projection
  to \(X\) is the blow-up morphism by construction.
\end{proof}

\begin{theorem}[Projectivity of the global blow-up]
  \label{thm:nagata-global-blowup-projective}
  \uses{def:nagata-projective-morphism,
    def:nagata-global-finite-free-generation,
    lemma:nagata-global-blowup-relative-projective-embedding}
  \lean{Scheme.IdealSheafData.blowupMap_isProjective}

  If \(\mathcal I\) is of finite type, then
  \(b_{X,\mathcal I}\colon\operatorname{Bl}_{\mathcal I}(X)\to X\) is
  projective provided that \(\mathcal I\) admits a finite global free
  presentation, i.e. there are \(r\in\mathbb N\) and a surjection
  \(q\colon\mathcal O_X^{r+1}\twoheadrightarrow\mathcal I\).
\end{theorem}

\begin{proof}
  \uses{def:nagata-global-finite-free-generation,
    lemma:nagata-global-blowup-relative-projective-embedding,
    def:nagata-projective-morphism}
  The preceding lemma supplies a closed immersion into the single finite
  projective space \(\mathbf P^r_X\) required by
  \Cref{def:nagata-projective-morphism}; its structural composite is
  \(b_{X,\mathcal I}\), so the morphism is projective.  Without the displayed
  global-generation hypothesis, finite type alone yields the properness
  statement below but does not imply projectivity for this finite-projective-
  space definition.
\end{proof}

\begin{theorem}[Properness of the global blow-up]
  \label{thm:nagata-global-blowup-proper}
  \lean{AlgebraicGeometry.Scheme.IdealSheafData.isProper_blowupMap}

  If \(\mathcal I\) is of finite type, then
  \(b_{X,\mathcal I}\) is proper.
\end{theorem}

\begin{proof}
  \uses{lemma:nagata-global-blowup-proper-local,
    lemma:nagata-projective-target-local}
  Apply \Cref{lemma:nagata-global-blowup-proper-local} to the affine open cover
  of $X$.  Each restricted morphism is proper, and
  \Cref{lemma:nagata-projective-target-local} is the verified target-local
  properness lemma (despite its historical label).  It follows that
  $b_{X,\mathcal I}$ is proper.  This argument does not use the unresolved
  global projectivity theorem.
\end{proof}

\begin{remark}[Formalization route]
  The formal proof (\texttt{isProper\_blowupMap}) follows this target-local
  reduction and settles the affine case
  (\Cref{lemma:nagata-affine-blowup-proper}) by identifying
  $\operatorname{Bl}_I(A)\to\operatorname{Spec}(A)$ with
  \texttt{Proj.toSpecZero} of the Rees algebra and invoking Mathlib's
  properness of $\operatorname{Proj}$ over its degree-zero part, fed by
  \Cref{prop:nagata-rees-finite-generation}.  The projective-embedding chain
  (\Cref{def:nagata-rees-polynomial-map},
  \Cref{lemma:nagata-rees-polynomial-surjection},
  \Cref{prop:nagata-affine-blowup-projective},
  \Cref{thm:nagata-global-blowup-projective}) is therefore not formalized.
\end{remark}

\begin{lemma}[Principalization on one affine chart]
  \label{lemma:nagata-chart-extended-ideal-principal}
  \uses{def:nagata-blowup-chart-ring}
  \lean{AlgebraicGeometry.AffineBlowup.map_awayAlgebraMap_eq_span}

  In $A[I/a]$, the extended ideal satisfies
  \[
    I\,A[I/a]=(a).
  \]
\end{lemma}

\begin{proof}
  For $x\in I$, the identity $x=a(x/a)$ gives
  $IA[I/a]\subseteq(a)$.  Since $a\in I$ and $1\in A[I/a]$, the element
  $a=a\cdot1$ belongs to $IA[I/a]$, so every multiple of $a$ belongs to that
  ideal and $(a)\subseteq IA[I/a]$.
\end{proof}

\begin{lemma}[Regularity of the chart generator]
  \label{lemma:nagata-chart-generator-regular}
  \uses{def:nagata-blowup-chart-ring}
  \lean{AlgebraicGeometry.AffineBlowup.awayAlgebraMap_mem_nonZeroDivisors}

  Multiplication by $a$ on $A[I/a]$ is injective.
\end{lemma}

\begin{proof}
  The chart algebra is a subring of $A_a$.  The image of $a$ is a unit in
  $A_a$, hence multiplication by it is injective on the subring.
\end{proof}

\begin{lemma}[Transition between local exceptional generators]
  \label{lemma:nagata-chart-generators-associated}
  \lean{ReesAlgebra.chartGenerators_associated_on_overlap}

  On the overlap of the $a$-chart and the $b$-chart, the two generators $a$
  and $b$ of the extended ideal differ by a unit.
\end{lemma}

\begin{proof}
  \uses{lemma:nagata-chart-extended-ideal-principal}
  On the overlap, both $(a)$ and $(b)$ equal the same extended ideal.  More
  explicitly, $b=a(b/a)$ and $a=b(a/b)$, and $(a/b)(b/a)=1$.
\end{proof}

\begin{definition}[Invertible ideal]
  \label{def:nagata-invertible-ideal}
  \lean{Scheme.IdealSheafData.IsInvertible}

  A quasi-coherent ideal $\mathcal I\subseteq\mathcal O_X$ is
  \emph{invertible} if it is locally free of rank one as an
  $\mathcal O_X$-module.  The inclusion $\mathcal I\hookrightarrow\mathcal O_X$
  is part of the ideal data; regularity of a local generator is not included
  in this definition.
\end{definition}

\begin{definition}[Effective Cartier ideal]
  \label{def:nagata-effective-cartier-ideal}
  \lean{AlgebraicGeometry.Scheme.IdealSheafData.IsEffectiveCartier}

  A quasi-coherent ideal $\mathcal I\subseteq\mathcal O_X$ is an
  \emph{effective Cartier ideal} if, on an open cover, it has a generator
  which is a non-zero-divisor in the corresponding ring.  Invertibility is
  a consequence recorded separately below.
\end{definition}

\begin{lemma}[An effective Cartier ideal is invertible]
  \label{lemma:nagata-effective-cartier-is-invertible}
  \uses{def:nagata-effective-cartier-ideal,def:nagata-invertible-ideal}
  \lean{Scheme.IdealSheafData.IsEffectiveCartier.isInvertible}

  Every effective Cartier ideal is invertible as an ideal sheaf.
\end{lemma}

\begin{proof}
  The regular local generator gives an isomorphism from the structure sheaf
  to the ideal on each member of the defining cover.  On overlaps the two
  generators differ by a unit, so these local rank-one bases glue.
\end{proof}

\begin{lemma}[An invertible ideal is effective Cartier]
  \label{lemma:nagata-invertible-ideal-is-effective-cartier}
  \uses{def:nagata-invertible-ideal,def:nagata-effective-cartier-ideal}
  \lean{Scheme.IdealSheafData.IsInvertible.isEffectiveCartier}

  An invertible ideal sheaf, viewed as an ideal by its specified inclusion into
  the structure sheaf, is an effective Cartier ideal.
\end{lemma}

\begin{proof}
  On an open set where the invertible ideal is free with basis $a$, the map
  $\mathcal O_X\to\mathcal I$ sending $1$ to $a$ is an isomorphism.  The
  inclusion $\mathcal I\hookrightarrow\mathcal O_X$ is injective because
  $\mathcal I$ is an ideal submodule, so the composite
  $\mathcal O_X\to\mathcal O_X$ is multiplication by $a$ and is injective.
  Thus $a$ is a non-zero-divisor on every member of this cover, which is
  exactly the effective-Cartier condition.
\end{proof}

\begin{theorem}[Principalization of the center ideal]
  \label{thm:nagata-blowup-principalizes-center}
  \uses{def:nagata-global-blowup,def:nagata-effective-cartier-ideal,
    def:nagata-invertible-ideal}
  \lean{AlgebraicGeometry.Scheme.IdealSheafData.extendedToBlowup_isEffectiveCartier}

  If $\mathcal I$ is a finite type quasi-coherent ideal on $X$, then on
  $\operatorname{Bl}_{\mathcal I}(X)$ the extended ideal
  $\mathcal I\mathcal O_{\operatorname{Bl}_{\mathcal I}(X)}$ is locally
  principal, generated locally by a non-zero-divisor, and hence is an
  invertible ideal defining an effective Cartier divisor.
\end{theorem}

\begin{proof}
  \uses{prop:nagata-affine-blowup-chart-cover,
    lemma:nagata-chart-extended-ideal-principal,lemma:nagata-chart-generator-regular,
    lemma:nagata-chart-generators-associated}
  Work locally on an arbitrary affine open of $X$.  Finite type of
  $\mathcal I$ gives finitely many generators on that affine, and the
  corresponding standard charts cover its blow-up.  Use
  \Cref{lemma:nagata-chart-extended-ideal-principal} and
  \Cref{lemma:nagata-chart-generator-regular}.  The local principal ideals
  glue because their generators differ by units by
  \Cref{lemma:nagata-chart-generators-associated}.
\end{proof}

\begin{remark}[Formalization route]
  The formal statement (\texttt{extendedToBlowup\_isEffectiveCartier}) proves
  the effective-Cartier form only, directly from
  \Cref{lemma:nagata-chart-extended-ideal-principal} and
  \Cref{lemma:nagata-chart-generator-regular}: being effective Cartier is a
  local condition, so the gluing step through
  \Cref{lemma:nagata-chart-generators-associated} and the invertible-sheaf
  half (\Cref{def:nagata-invertible-ideal},
  \Cref{lemma:nagata-effective-cartier-is-invertible},
  \Cref{lemma:nagata-invertible-ideal-is-effective-cartier}) are not needed
  and are not formalized.  The Lean theorem also requires no finite-type
  hypothesis on the center.
\end{remark}

\subsection{Universal property, flat base change, and products of ideals}

\begin{lemma}[Associativity of ideal-power multiplication]
  \label{lemma:nagata-sheaf-ideal-powers-assoc}
  \lean{IdealSheafData.idealPowers_mul_assoc}

  For a quasi-coherent ideal $\mathcal I\subseteq\mathcal O_X$, the canonical
  maps $\mathcal I^m\otimes\mathcal I^n\to\mathcal I^{m+n}$ are associative
  and have $\mathcal I^0=\mathcal O_X$ as a two-sided unit.
\end{lemma}

\begin{proof}
  On an affine open $\operatorname{Spec}(A)$ these maps are induced by
  multiplication of the ideals $I^m$ and $I^n$ in $A$.  Associativity and the
  unit law are therefore the corresponding ring identities.  Localization
  preserves these identities, and equality of morphisms of quasi-coherent
  sheaves can be checked on an affine cover, so the identities glue on $X$.
\end{proof}

\begin{definition}[Sheaf Rees algebra]
  \label{def:nagata-sheaf-rees-algebra}
  \uses{lemma:nagata-sheaf-ideal-powers-assoc}
  \lean{AlgebraicGeometry.ReesAlgebra.sheaf}

  Let $\mathcal I\subseteq\mathcal O_X$ be a quasi-coherent ideal sheaf.  Define
  its sheaf Rees algebra by
  \[
    \mathcal R_X(\mathcal I)=\bigoplus_{n\geq0}\mathcal I^n,
  \]
  with multiplication induced by the canonical multiplication maps
  $\mathcal I^m\otimes_{\mathcal O_X}\mathcal I^n\to
  \mathcal I^{m+n}$ and unit the degree-zero copy of $\mathcal O_X$.
  The direct sum is taken in sheaves of $\mathcal O_X$-modules.  Define
  multiplication on homogeneous summands by the displayed maps and extend it
  by the universal property of the direct sum; the degree-zero inclusion is
  the unit.  Associativity and the unit laws are exactly
  \Cref{lemma:nagata-sheaf-ideal-powers-assoc}, so this is a graded
  quasi-coherent $\mathcal O_X$-algebra.  The multiplication maps need not be
  injective, and no such hypothesis is used.
\end{definition}

\begin{definition}[Extended ideal along a morphism]
  \label{def:nagata-extended-ideal}
  \lean{AlgebraicGeometry.Scheme.IdealSheafData.extendedToBlowup}

  For $f\colon Y\to X$ and a quasi-coherent ideal
  $\mathcal I\subseteq\mathcal O_X$, define
  $\mathcal I\mathcal O_Y$ to be the image ideal of the composite
  \[
    f^*\mathcal I\longrightarrow f^*\mathcal O_X
      \longrightarrow\mathcal O_Y.
  \]
\end{definition}

\begin{lemma}[Extended ideal on an affine blow-up chart]
  \label{lemma:nagata-extended-ideal-affine-chart}
  \uses{def:nagata-extended-ideal,
    lemma:nagata-chart-extended-ideal-principal}
  \lean{AlgebraicGeometry.Scheme.IdealSheafData.ideal_comap_specMap_fromSpec}

  On an affine chart $\operatorname{Spec}(A[I/a])$ of the affine blow-up, the
  extended ideal of $I$ is the ideal generated by the image of $I$ in
  $A[I/a]$.
\end{lemma}

\begin{proof}
  \uses{lemma:nagata-chart-extended-ideal-principal}
  This is the affine calculation for the image defining
  \Cref{def:nagata-extended-ideal}.
\end{proof}

\begin{lemma}[Tensor powers of an invertible ideal]
  \label{lemma:nagata-invertible-ideal-tensor-power}
  \uses{def:nagata-invertible-ideal,
    lemma:nagata-invertible-ideal-is-effective-cartier}
  \lean{Scheme.IdealSheafData.IsInvertible.tensorPowerIsoPow}

  If $\mathcal L\subseteq\mathcal O_Y$ is an invertible ideal, then for every
  $n\geq0$ the multiplication map
  \[
    \mathcal L^{\otimes n}\longrightarrow\mathcal L^n
  \]
  is an isomorphism.
\end{lemma}

\begin{proof}
  By \Cref{lemma:nagata-invertible-ideal-is-effective-cartier}, work on an
  open cover where $\mathcal L=(a)$ with $a$ a non-zero-divisor.  The source
  has basis $a^{\otimes n}$ and the target has basis $a^n$; multiplication
  sends the first basis to the second.  Hence it is an isomorphism on this
  cover, and is therefore an isomorphism of sheaves.
\end{proof}

\begin{lemma}[Relative Proj base change for an arbitrary morphism]
  \label{lemma:nagata-relative-proj-base-change-arbitrary}
  \uses{prop:nagata-relative-proj-functoriality}
  \lean{AlgebraicGeometry.RelativeProj.baseChange}

  For any morphism $g\colon Y\to X$ and graded quasi-coherent algebra
  $\mathcal A$ on $X$, coefficientwise pullback gives a canonical isomorphism
  \[
    \operatorname{Proj}_Y(g^*\mathcal A)
      \cong Y\times_X\operatorname{Proj}_X(\mathcal A).
  \]
\end{lemma}

\begin{proof}
  \uses{lemma:nagata-proj-chart-ring-cocycle}
  Cover $X$ by affine opens $V=\operatorname{Spec}(A)$ and cover each affine
  inverse image in $Y$ by $W=\operatorname{Spec}(B)$, with structure map
  $A\to B$.  Put $C=\Gamma(V,\mathcal A)$.  For every positive-degree
  homogeneous $f\in C$, the two standard charts in the asserted comparison
  have coordinate rings
  \[
    B\otimes_A C_{(f),0}\quad\text{and}\quad
    (B\otimes_A C)_{(f\otimes1),0}.
  \]
  For fixed $f$, the degree-zero homogeneous localization is the filtered
  colimit of the modules $C_{m\deg(f)}$ under multiplication by $f$.
  Since tensor product commutes with filtered colimits, there is a canonical
  isomorphism
  \[
    B\otimes_A C_{(f),0}\;\xrightarrow{\sim}\;
    (B\otimes_A C)_{(f\otimes1),0}.
  \]
  On pure fractions this isomorphism and its inverse are given by
  \[
    (b\otimes x/f^m)\longmapsto (b\otimes x)/(f\otimes1)^m,
    \qquad
    (b\otimes x)/(f\otimes1)^m\longmapsto b\otimes x/f^m
  \]
  and extend additively.  The filtered-colimit description, rather than
  injectivity after tensoring, proves that these formulas are well-defined and
  mutually inverse for an arbitrary $A\to B$.  They are compatible with
  localization at a second homogeneous element, so the affine chart
  isomorphisms agree on all overlaps.  The standard opens cover both Proj
  schemes, and \texttt{AlgebraicGeometry.Scheme.GlueData.glued} therefore
  produces an
  isomorphism
  \[
    \operatorname{Proj}_Y(g^*\mathcal A)\cong
    Y\times_X\operatorname{Proj}_X(\mathcal A).
  \]
  The construction is natural in $g$ and in maps of graded algebras because
  every displayed map is induced by tensor product and localization; hence it
  is the required canonical arbitrary-base-change comparison.
\end{proof}

\begin{definition}[Graded map induced by an invertible quotient]
  \label{def:nagata-blowup-universal-graded-map}
  \uses{def:nagata-sheaf-rees-algebra,def:nagata-extended-ideal,
    def:nagata-invertible-ideal}
  \lean{Scheme.IdealSheafData.reesToTensorPowersConstruction}

  Let $f\colon Y\to X$, let $\mathcal I\subseteq\mathcal O_X$ be a
  quasi-coherent ideal, and let
  $\mathcal L\hookrightarrow\mathcal O_Y$ be an invertible ideal together with
  a surjection $q\colon f^*\mathcal I\twoheadrightarrow\mathcal L$ such that
  the composite $\mathcal L\hookrightarrow\mathcal O_Y$ after $q$ is the
  canonical map whose image is the extended ideal
  $\mathcal I\mathcal O_Y$ of
  \Cref{def:nagata-extended-ideal}.  Define $\alpha_f$ degreewise by the
  tensor powers of this quotient:
  \[
    \alpha_f\colon
    f^*\Bigl(\bigoplus_{n\geq0}\mathcal I^n\Bigr)
      \longrightarrow\bigoplus_{n\geq0}\mathcal L^{\otimes n}.
  \]
\end{definition}

\begin{lemma}[The quotient construction is a graded algebra map]
  \label{lemma:nagata-blowup-universal-graded-map}
  \uses{def:nagata-blowup-universal-graded-map,
    lemma:nagata-invertible-ideal-tensor-power}
  \lean{Scheme.IdealSheafData.reesToTensorPowers}

  The degreewise map $\alpha_f$ respects the unit and multiplication and is a
  graded algebra morphism.
\end{lemma}

\begin{proof}
  In degree $n$, tensor $q$ $n$ times.  Let
  $\mu_{\mathcal I}\colon(f^*\mathcal I)^{\otimes n}\to
  f^*(\mathcal I^n)$ be multiplication.  The composite of $q^{\otimes n}$
  with
  multiplication isomorphism
  $\mathcal L^{\otimes n}\cong\mathcal L^n$ from
  \Cref{lemma:nagata-invertible-ideal-tensor-power} and the injection
  $\mathcal L^n\hookrightarrow\mathcal O_Y$ is ordinary $n$-fold
  multiplication.  It therefore kills $\ker(\mu_{\mathcal I})$; injectivity
  of $\mathcal L^n\hookrightarrow\mathcal O_Y$ shows that $q^{\otimes n}$
  itself kills this kernel.  Hence it factors uniquely through
  $f^*(\mathcal I^n)$.  These factorizations commute with multiplication by
  associativity and with the degree-zero unit, giving the graded algebra map.
  Associativity of tensor products and the multiplication maps for ideal
  powers show that these degreewise maps preserve the unit and multiplication.
\end{proof}

\begin{lemma}[Affine chart calculation for an invertible tensor algebra]
  \label{lemma:nagata-relative-proj-invertible-tensor-powers-affine}
  \uses{def:nagata-relative-proj,def:nagata-invertible-ideal,
    prop:nagata-proj-standard-open}

  Let $V=\operatorname{Spec}(A)$ be an affine open on which the invertible
  ideal $\mathcal L$ is generated by one element $a$.  Then the graded
  algebra $\bigoplus_{n\geq0}\mathcal L(V)^{\otimes n}$ is generated by the
  single degree-one element $a$, and its degree-zero homogeneous localization
  at $a$ is canonically $A$.  Consequently the standard open $D_+(a)$ is all
  of $\operatorname{Proj}_A(\bigoplus_n\mathcal L(V)^{\otimes n})$, and the
  resulting affine Proj is canonically $V$ over $V$.
\end{lemma}

\begin{proof}
  Choose the basis $a$ of the rank-one $A$-module $\mathcal L(V)$.  Tensoring
  this basis identifies $\mathcal L(V)^{\otimes n}$ with $A\,a^{\otimes n}$,
  and multiplication sends $a^{\otimes m}\otimes a^{\otimes n}$ to
  $a^{\otimes(m+n)}$.  Thus every homogeneous element is uniquely $r a^n$
  and the algebra is the polynomial algebra $A[a]$ with $a$ in degree one.
  A degree-zero fraction has the form $(r a^n)/a^n$ and is sent to $r$;
  this gives mutually inverse ring maps between the homogeneous localization
  and $A$.  Since $a$ generates the irrelevant ideal, the standard-open cover
  lemma makes $D_+(a)$ the whole Proj, proving the asserted affine
  identification and its compatibility with the degree-zero structure map.
\end{proof}

\begin{lemma}[Proj of the tensor powers of an invertible ideal]
  \label{lemma:nagata-relative-proj-invertible-tensor-powers}
  \uses{def:nagata-relative-proj,def:nagata-invertible-ideal,
    lemma:nagata-relative-proj-exists,
    lemma:nagata-relative-proj-invertible-tensor-powers-affine}
  \lean{AlgebraicGeometry.RelativeProj.proj_tensorPowers_invertible}

  If $\mathcal L$ is an invertible ideal on $Y$, the relative Proj of
  $\bigoplus_{n\geq0}\mathcal L^{\otimes n}$ is canonically isomorphic to $Y$.
  Under this isomorphism $\mathcal O(1)$ identifies with $\mathcal L$, and
  the universal degree-one map $\mathcal L\to\mathcal O(1)$ is the identity
  of $\mathcal L$.  The inclusion
  $\mathcal L\hookrightarrow\mathcal O_Y$ is a separate effective-Cartier
  inclusion and is not the tautological quotient.
\end{lemma}

\begin{proof}
  Cover $Y$ by affine opens $V$ on which $\mathcal L$ has a generator $a$.
  By \Cref{lemma:nagata-relative-proj-invertible-tensor-powers-affine}, the
  relative Proj restricts over each such $V$ to $V$, with the identification
  induced by the degree-zero localization map.  On an overlap $V\cap W$, if
  $b=ua$ for a unit $u$, the two chart identifications send a fraction
  $(r a^n)/a^n$ (respectively $(r u^{-n}b^n)/b^n$) to the same coefficient $r$.
  Hence they agree on overlaps and glue, by the relative-Proj gluing
  compatibility, to a canonical isomorphism with $Y$.  The degree-one
  localization identifies the tautological section with the chosen generator;
  these local identifications glue to $\mathcal O(1)\cong\mathcal L$, and the
  resulting universal degree-one map is the identity on $\mathcal L$.  The
  inclusion $\mathcal L\hookrightarrow\mathcal O_Y$ is independent of this
  tautological quotient and remains the specified effective-Cartier inclusion.
\end{proof}

\begin{definition}[Universal lift construction]
  \label{def:nagata-blowup-universal-lift}
  \uses{def:nagata-global-blowup,def:nagata-sheaf-rees-algebra,
    def:nagata-blowup-universal-graded-map,
    def:nagata-relative-proj,
    lemma:nagata-relative-proj-invertible-tensor-powers,
    lemma:nagata-relative-proj-base-change-arbitrary}
  \lean{Scheme.IdealSheafData.liftToBlowupConstruction}

  Under the data of
  \Cref{def:nagata-blowup-universal-graded-map}, use the inverse of the canonical isomorphism in
  \Cref{lemma:nagata-relative-proj-invertible-tensor-powers}, apply the
  contravariant relative-Proj map induced by $\alpha_f$.  Then use the
  arbitrary base-change comparison in
  \Cref{lemma:nagata-relative-proj-base-change-arbitrary} to identify the
  resulting target with $Y\times_X\operatorname{Bl}_{\mathcal I}(X)$.  Here
  $\operatorname{Proj}_X(\mathcal R_X(\mathcal I))$ is identified with
  \Cref{def:nagata-global-blowup}: both are obtained by gluing the same affine
  Rees-algebra Proj schemes, and the comparison respects the restriction
  cocycles.  Compose
  the resulting map with the second projection
  $Y\times_X\operatorname{Bl}_{\mathcal I}(X)\to
  \operatorname{Bl}_{\mathcal I}(X)$, and
  define
  \[
    \widetilde f\colon Y\longrightarrow\operatorname{Bl}_{\mathcal I}(X)
  \]
  to be the resulting composite.  The affine maps and their gluing data are
  part of the relative-Proj interface; no identification of $Y$ with a Proj
  is left implicit.
\end{definition}

\begin{lemma}[The constructed lift lies over $X$]
  \label{lemma:nagata-blowup-universal-exists}
  \uses{def:nagata-blowup-universal-lift}
  \lean{AlgebraicGeometry.Scheme.IdealSheafData.exists_lift_blowup}

  Under the data of \Cref{def:nagata-blowup-universal-lift}, the constructed
  map $\widetilde f$ is an $X$-morphism: its composite with
  $b_{X,\mathcal I}$ is $f$.
\end{lemma}

\begin{proof}
  \uses{lemma:nagata-blowup-universal-graded-map,
    lemma:nagata-global-blowup-glue-data-valid,
    lemma:nagata-global-blowup-local-maps-agree}
  On each standard chart, the degree-zero component of
  \Cref{lemma:nagata-blowup-universal-graded-map} is the canonical map to
  $\mathcal O_Y$, hence the composite with the blow-up structure map is $f$.
  The chart maps agree on overlaps by functoriality of localization, and the
  glued equality is therefore the specified equality over $X$.
\end{proof}

\begin{lemma}[Uniqueness of the universal lift]
  \label{lemma:nagata-blowup-universal-unique}
  \uses{def:nagata-blowup-universal-lift,
    lemma:nagata-blowup-universal-graded-map,
    lemma:nagata-relative-proj-invertible-tensor-powers,
    lemma:nagata-proj-chart-ring-cocycle,
    prop:nagata-proj-standard-open}
  \lean{AlgebraicGeometry.Scheme.IdealSheafData.hom_ext_of_isEffectiveCartier_comap}

  Under the data of \Cref{def:nagata-blowup-universal-lift}, every $X$-morphism
  $h\colon Y\to\operatorname{Bl}_{\mathcal I}(X)$ inducing the chosen
  invertible quotient equals the constructed map $\widetilde f$.
\end{lemma}

\begin{proof}
  On each affine open of $X$, let $D_+(f)$ range over the standard opens
  indexed by the degree-one elements $f\in\mathcal I$; these opens cover the
  blow-up Proj.  The inverse image of each chart
  under $h$ is determined by the degree-zero map obtained from the given
  quotient $f^*\mathcal I\twoheadrightarrow\mathcal L$: on a representative
  $x/f_i^n$ it is the coefficient of the corresponding quotient fraction in
  $\mathcal O_Y$.  The constructed lift uses exactly the same degreewise maps,
  by definition of $\alpha_f$ and the tensor-power identification, so $h$ and
  $\widetilde f$ agree on each inverse-image chart.  These equalities are
  compatible on chart overlaps because both maps send a localized fraction to
  its further localization; the localization cocycle gives literal equality.
  The standard charts cover the target, hence their inverse images cover $Y$,
  and equality of morphisms on this cover implies $h=\widetilde f$.
\end{proof}

\begin{theorem}[{\cite[\href{https://stacks.math.columbia.edu/tag/0806}{Tag 0806}]{stacks-project}}]
  \label{thm:nagata-blowup-universal-property}
  \uses{def:nagata-blowup-universal-lift}
  \lean{AlgebraicGeometry.Scheme.IdealSheafData.exists_lift_blowup,
    AlgebraicGeometry.Scheme.IdealSheafData.hom_ext_of_isInvertible_comap}

  Under the data of \Cref{def:nagata-blowup-universal-lift}, the constructed
  map $\widetilde f$ is the unique $X$-morphism from $Y$ to
  $\operatorname{Bl}_{\mathcal I}(X)$ inducing the chosen invertible
  quotient.  Thus the blow-up is final among such $X$-schemes.
\end{theorem}

\begin{proof}
  \uses{lemma:nagata-blowup-universal-exists,
    lemma:nagata-blowup-universal-unique,
    lemma:nagata-relative-proj-invertible-tensor-powers}
  Existence is \Cref{lemma:nagata-blowup-universal-exists}, and uniqueness
  among $X$-morphisms inducing the specified quotient is
  \Cref{lemma:nagata-blowup-universal-unique}.  Together these are exactly the
  asserted final-object property.
\end{proof}

\begin{remark}[Formalization route]
  The formal proof (\texttt{exists\_lift\_blowup} together with
  \texttt{hom\_ext\_of\_isInvertible\_comap}) takes the
  Stacks-style hypothesis that the inverse image ideal sheaf of the center is
  effective Cartier, and is entirely chart-local: existence is glued from the
  universal property of the chart ring $A[\mathcal I/a]$
  (\texttt{AffineBlowup.chartRingLift}) on affine probes, and uniqueness
  follows from chart-ring extensionality
  (\texttt{chartRingHom\_ext}) together with the chart-overlap localization.
  The graded-map and tensor-power scaffolding
  (\Cref{def:nagata-blowup-universal-graded-map},
  \Cref{lemma:nagata-invertible-ideal-tensor-power},
  \Cref{lemma:nagata-relative-proj-invertible-tensor-powers},
  \Cref{def:nagata-blowup-universal-lift}) is not formalized.
\end{remark}

\begin{corollary}[Blow-up of an effective Cartier center]
  \label{cor:nagata-blowup-cartier-iso}
  \lean{Scheme.IdealSheafData.blowupMap_isIso_of_isInvertible}

  Let $X$ be a scheme.  If $\mathcal I$ is an invertible ideal on $X$, then
  $b_{X,\mathcal I}$ is an isomorphism.
\end{corollary}

\begin{proof}
  \uses{def:nagata-invertible-ideal,
    thm:nagata-blowup-universal-property}
  The identity of $X$ lifts to the blow-up by the universal property.  Compose
  this lift with the blow-up map in both orders and use uniqueness.
\end{proof}

\begin{lemma}[Flat pullback commutes with ideal powers]
  \label{lemma:nagata-flat-pullback-ideal-pow}
  \uses{def:nagata-extended-ideal}
  \lean{IdealSheafData.flat_pullback_pow}

  Let $\mathcal I\subseteq\mathcal O_X$ be a quasi-coherent ideal.  If
  $g\colon Y\to X$ is flat, then the image of $g^*(\mathcal I^n)$ in
  $\mathcal O_Y$ is $(\mathcal I\mathcal O_Y)^n$ for every $n$.
\end{lemma}

\begin{proof}
  Work on an affine chart, with a flat ring map $A\to B$ and ideal $I\subseteq
  A$.  Tensoring $0\to I^n\to A$ with $B$ gives an injection
  $I^n\otimes_AB\to B$.  Its image is generated by the elements $xb$ with
  $x\in I^n$ and $b\in B$, which is exactly $(IB)^n$: expand products of
  generators for one inclusion and write elements of $I^n$ as finite sums of
  products for the other.  These equalities commute with localization, so
  they glue to the asserted equality of ideal sheaves.
\end{proof}

\begin{definition}[Flat pullback map for Rees algebras]
  \label{def:nagata-flat-pullback-rees-map}
  \uses{def:nagata-sheaf-rees-algebra,def:nagata-extended-ideal}
  \lean{IdealSheafData.flatPullbackReesMap}

  For a quasi-coherent ideal $\mathcal I\subseteq\mathcal O_X$ and flat
  $g\colon Y\to X$, define $\beta_g$ degreewise by the canonical
  maps
  \[
    g^*(\mathcal I^n)\longrightarrow
      (\mathcal I\mathcal O_Y)^n.
  \]
  Thus
  \[
    \beta_g\colon g^*\mathcal R_X(\mathcal I)\longrightarrow
      \mathcal R_Y(\mathcal I\mathcal O_Y).
  \]
\end{definition}

\begin{lemma}[Flat pullback of the Rees algebra]
  \label{lemma:nagata-flat-pullback-rees}
  \uses{def:nagata-flat-pullback-rees-map}
  \lean{IdealSheafData.flat_pullback_rees_equiv}

  The map $\beta_g$ is a graded algebra equivalence.
\end{lemma}

\begin{proof}
  \uses{lemma:nagata-flat-pullback-ideal-pow}
  Apply \Cref{lemma:nagata-flat-pullback-ideal-pow} in every degree and check
  compatibility with multiplication.  The degreewise isomorphisms assemble
  to the claimed graded equivalence.
\end{proof}

\begin{definition}[Flat base-change comparison for blow-ups]
  \label{def:nagata-blowup-flat-base-change-map}
  \uses{def:nagata-global-blowup,
    def:nagata-flat-pullback-rees-map,
    def:nagata-extended-ideal,
    lemma:nagata-flat-pullback-rees,
    prop:nagata-relative-proj-functoriality}
  \lean{Scheme.IdealSheafData.blowup_flatBaseChangeMap}

  Let $g\colon Y\to X$ be flat and let
  $\mathcal I\subseteq\mathcal O_X$ be a quasi-coherent ideal.  With
  $\mathcal I\mathcal O_Y$ defined in
  \Cref{def:nagata-extended-ideal}, define
  \[
    \operatorname{bc}_g\colon
    \operatorname{Bl}_{\mathcal I\mathcal O_Y}(Y)
      \longrightarrow
    Y\times_X\operatorname{Bl}_{\mathcal I}(X)
  \]
  by applying relative Proj to $\beta_g$ on affine opens and gluing.  The
  degreewise graded equivalence used in this construction is
  \Cref{lemma:nagata-flat-pullback-rees}.
\end{definition}

\begin{theorem}[Flat base change of blow-ups]
  \label{thm:nagata-blowup-flat-base-change}
  \uses{def:nagata-blowup-flat-base-change-map,
    prop:nagata-relative-proj-functoriality}
  \lean{Scheme.IdealSheafData.blowup_flatBaseChangeIso}

  The constructed comparison $\operatorname{bc}_g$ is an isomorphism over
  $Y$.
\end{theorem}

\begin{proof}
  \uses{lemma:nagata-flat-pullback-rees,
    lemma:nagata-global-blowup-overlap-cocycle}
  Relative Proj sends the graded equivalence of
  \Cref{lemma:nagata-flat-pullback-rees} to an affine-local isomorphism.
  Compatibility with degree zero identifies the structure maps, and the
  localization cocycle glues the local isomorphisms.
\end{proof}

\begin{lemma}[Product of invertible ideals is invertible]
  \label{lemma:nagata-invertible-ideal-mul}
  \uses{def:nagata-invertible-ideal,
    lemma:nagata-invertible-ideal-is-effective-cartier}
  \lean{AlgebraicGeometry.Scheme.IdealSheafData.IsEffectiveCartier.mul}

  The product of two invertible ideal sheaves is invertible.
\end{lemma}

\begin{proof}
  By \Cref{lemma:nagata-invertible-ideal-is-effective-cartier}, local bases of
  the two ideals are regular elements $a$ and $b$.  Their product is generated
  by $ab$, which is regular because multiplication by $a$ and by $b$ are
  injective.  These local generators give the required rank-one bases.
\end{proof}

\begin{lemma}[Factors of a finite-type effective Cartier product]
  \label{lemma:nagata-ideal-invertible-product-factor}
  \uses{def:nagata-effective-cartier-ideal}
  \lean{Ideal.isInvertible_of_mul_isInvertible}

  Let $R$ be a ring and let $I,J\subseteq R$ be ideals.  If $IJ$
  is an effective Cartier ideal, then $I$ and $J$ are effective Cartier ideals.
\end{lemma}

\begin{proof}
  The assertion is local.  Suppose $IJ=(f)$ with $f$ regular.  Write
  $f=\sum_{i=1}^n x_i y_i$ with $x_i\in I$ and $y_i\in J$.  Since each
  $x_iy_i$ belongs to $(f)$, choose $a_i$ with $x_iy_i=a_if$.  Regularity of
  $f$ and $f=(\sum_i a_i)f$ give $\sum_i a_i=1$, so the opens $D(a_i)$
  cover.  On $D(a_i)$ one has
  $f=a_i^{-1}x_i y_i$.  Since this product is regular, both $x_i$ and $y_i$
  are regular there.  For $x\in I$, write $xy_i=cf$; cancellation of $y_i$
  gives $x=ca_i^{-1}x_i$, and hence $I=(x_i)$ on $D(a_i)$.  The symmetric
  argument gives $J=(y_i)$.  Thus both ideals are locally generated by one
  regular element and are effective Cartier ideals.
\end{proof}

\begin{lemma}[{\cite[\href{https://stacks.math.columbia.edu/tag/07ZV}{Tag 07ZV}]{stacks-project}}]
  \label{lemma:nagata-invertible-product-factor}
  \uses{lemma:nagata-ideal-invertible-product-factor,
    def:nagata-effective-cartier-ideal}
  \lean{AlgebraicGeometry.Scheme.IdealSheafData.IsEffectiveCartier.of_mul_left}

  Let $X$ be a scheme and let $\mathcal I,\mathcal J$ be quasi-coherent
  ideals on $X$.  If the product
  $\mathcal I\mathcal J$ defines an effective Cartier divisor, then each of
  $\mathcal I$ and $\mathcal J$ defines an effective Cartier divisor.
\end{lemma}

\begin{proof}
  Work on affine opens and apply
  \Cref{lemma:nagata-ideal-invertible-product-factor} to the corresponding
  ideals.  The resulting local invertible-module structures and regular
  generators are compatible under localization, so both ideal sheaves define
  effective Cartier divisors.
\end{proof}

\begin{lemma}[Total transform of a Cartier ideal under a blow-up]
  \label{lemma:nagata-cartier-total-transform-blowup}
  \uses{def:nagata-effective-cartier-ideal,
    def:nagata-blowup-chart-ring}
  \lean{AlgebraicGeometry.Scheme.IdealSheafData.IsEffectiveCartier.comap_blowupMap}

  Let $\mathcal L\subseteq\mathcal O_T$ be an effective Cartier ideal and let
  $T'\to T$ be a blow-up.  Then the extended ideal
  $\mathcal L\mathcal O_{T'}$ is again effective Cartier, and in particular
  is invertible.
\end{lemma}

\begin{proof}
  Work over an affine open $\operatorname{Spec}(A)\subseteq T$ on which
  $\mathcal L=(x)$ with $x$ a non-zero-divisor.  A standard chart of the
  blow-up has ring $A[I/a]\subseteq A_a$.  If
  $x(y/a^n)=0$ in this chart, then $a^mxy=0$ in $A$ for some $m$; regularity
  of $x$ gives $a^my=0$, hence $y/a^n=0$.  Thus the image of $x$ is regular
  on every blow-up chart.  It generates the extended ideal there, so these
  local equations make $\mathcal L\mathcal O_{T'}$ an effective Cartier
  ideal, hence an invertible one.
\end{proof}

\begin{definition}[Comparison maps for the product-ideal blow-up]
  \label{def:nagata-blowup-product-comparison}
  \uses{def:nagata-blowup-universal-lift,def:nagata-global-blowup,
    def:nagata-extended-ideal,thm:nagata-blowup-principalizes-center,
    lemma:nagata-cartier-total-transform-blowup,
    lemma:nagata-invertible-ideal-mul,
    lemma:nagata-invertible-product-factor,
    lemma:nagata-effective-cartier-is-invertible}
  \lean{Scheme.IdealSheafData.iteratedBlowupComparison}

  Let $X$ be a scheme and let $\mathcal I,\mathcal J$ be finite-type
  quasi-coherent ideals on $X$.  Let $X_{\mathcal I,\mathcal J}$ be the
  scheme obtained by first blowing up $\mathcal I$ and then the extended
  ideal $\mathcal J\mathcal O_{\operatorname{Bl}_{\mathcal I}(X)}$ from
  \Cref{def:nagata-extended-ideal}.
  The first exceptional ideal is Cartier after the first blow-up.  Applying
  \Cref{lemma:nagata-cartier-total-transform-blowup} on the second blow-up
  (chartwise, where the pulled-back local equation remains regular) shows
  that its total transform is invertible.  This is the special total-transform
  input; no assertion about arbitrary images of invertible ideals is used.
  On this iterated blow-up, the total transform
  of the first exceptional ideal is invertible, and the second exceptional
  ideal is invertible by \Cref{thm:nagata-blowup-principalizes-center}.
  Hence their product, which is the extension of
  $\mathcal I\mathcal J$, is invertible by
  \Cref{lemma:nagata-invertible-ideal-mul}; define
  \[
    \Phi_{\mathcal I,\mathcal J}\colon
      X_{\mathcal I,\mathcal J}
      \longrightarrow\operatorname{Bl}_{\mathcal I\mathcal J}(X)
  \]
  by the universal-lift construction for their product.  On
  $\operatorname{Bl}_{\mathcal I\mathcal J}(X)$, factor the invertible
  product first through the $\mathcal I$-blow-up.  The extension of
  $\mathcal I\mathcal J$ is an effective Cartier ideal, so
  \Cref{lemma:nagata-invertible-product-factor} and
  \Cref{lemma:nagata-effective-cartier-is-invertible} make the extensions
  of both factors invertible.  The universal lift for $\mathcal I$, followed
  by the universal lift for the extended $\mathcal J$, defines the resulting
  map
  \[
    \Psi_{\mathcal I,\mathcal J}\colon
      \operatorname{Bl}_{\mathcal I\mathcal J}(X)
      \longrightarrow X_{\mathcal I,\mathcal J}.
  \]
\end{definition}

\begin{theorem}[{\cite[\href{https://stacks.math.columbia.edu/tag/080A}{Tag 080A}]{stacks-project}}]
  \label{thm:nagata-blowup-product-ideal}
  \uses{def:nagata-blowup-product-comparison}
  \lean{AlgebraicGeometry.Scheme.IdealSheafData.exists_blowup_mul_iso}

  The constructed maps $\Phi_{\mathcal I,\mathcal J}$ and
  $\Psi_{\mathcal I,\mathcal J}$ are inverse over $X$.  Thus the iterated
  blow-up is identified with the blow-up of $\mathcal I\mathcal J$.
\end{theorem}

\begin{proof}
  \uses{thm:nagata-blowup-universal-property,
    lemma:nagata-invertible-ideal-mul,
    lemma:nagata-invertible-product-factor,
    lemma:nagata-effective-cartier-is-invertible}
  Write $p\colon X_{\mathcal I,\mathcal J}\to X$ for the composite and
  $q\colon\operatorname{Bl}_{\mathcal I\mathcal J}(X)\to X$ for the
  product blow-up.  The construction of $\Phi_{\mathcal I,\mathcal J}$ is
  the unique $X$-morphism whose pullback of the tautological invertible
  quotient on the product blow-up is the product of the two tautological
  quotients on $X_{\mathcal I,\mathcal J}$.  Conversely, the construction of
  $\Psi_{\mathcal I,\mathcal J}$ first uses the invertible quotient
  $q^*\mathcal I$ and then the invertible quotient of the pullback of
  $\mathcal J$; their product is the tautological quotient of
  $q^*(\mathcal I\mathcal J)$.  Thus both composites are $X$-morphisms
  pulling back the relevant tautological quotient to itself.  Apply uniqueness
  in \Cref{thm:nagata-blowup-universal-property} first for the product
  blow-up to obtain $\Phi\circ\Psi=\mathrm{id}$.  For the other composite,
  first apply uniqueness for the $\mathcal I$-blow-up to see that its
  projection to $\operatorname{Bl}_{\mathcal I}(X)$ is the original first
  projection.  It is therefore a morphism over
  $\operatorname{Bl}_{\mathcal I}(X)$, and uniqueness for the second blow-up
  gives $\Psi\circ\Phi=\mathrm{id}$.
\end{proof}

\begin{definition}[Domination maps from the product blow-up]
  \label{def:nagata-blowup-product-domination-maps}
  \uses{def:nagata-blowup-product-comparison}
  \lean{Scheme.IdealSheafData.blowupMulDominationMaps}

  Define
  \[
    d_{\mathcal I}\colon\operatorname{Bl}_{\mathcal I\mathcal J}(X)
      \longrightarrow\operatorname{Bl}_{\mathcal I}(X)
  \]
  as $\Psi_{\mathcal I,\mathcal J}$ followed by the first blow-up
  projection.  Define $d_{\mathcal J}$ symmetrically after interchanging
  $\mathcal I$ and $\mathcal J$.
\end{definition}

\begin{corollary}[Common refinement by the product ideal]
  \label{cor:nagata-blowup-product-common-refinement}
  \uses{def:nagata-blowup-product-domination-maps}
  \lean{Scheme.IdealSheafData.blowupMulToLeft}

  The constructed $d_{\mathcal I}$ and $d_{\mathcal J}$ are morphisms over
  $X$.  If both ideals restrict to the unit ideal on an open $U$, the product
  blow-up and both maps restrict to the identity over $U$.
\end{corollary}

\begin{proof}
  \uses{thm:nagata-blowup-product-ideal,
    thm:nagata-global-blowup-complement}
  The maps are composites over $X$ by definition.  On $U$, all three ideals
  are the unit ideal, so all three blow-ups are identified with $U$ by the
  complement theorem.  The universal maps are unique over $U$, hence become
  the identity under these identifications.
\end{proof}


\begin{lemma}[Existence of a finite-type extension of an ideal]
  \label{lemma:nagata-finite-type-ideal-extension-exists}
  \lean{AlgebraicGeometry.Scheme.IdealSheafData.exists_finiteType_le_comap_eq}

  Let $X$ be quasi-compact and quasi-separated, let $W\subseteq X$ be a
  quasi-compact open, and let
  $\mathcal J\subseteq\mathcal O_W$ be a finite type quasi-coherent ideal.
  There exists a finite type quasi-coherent ideal
  $\mathcal I\subseteq\mathcal O_X$ with
  $\mathcal I|_W=\mathcal J$.
  More generally, if $\mathcal G\subseteq\mathcal F|_W$ is a finite type
  quasi-coherent submodule of the restriction of a quasi-coherent module
  $\mathcal F$ on $X$, then $\mathcal G$ is the restriction of a finite type
  quasi-coherent submodule of $\mathcal F$.
\end{lemma}

\begin{proof}
  We prove the more general assertion.  First take the inverse image
  $\mathcal H$ of $j_*\mathcal G$ under
  $\mathcal F\to j_*j^*\mathcal F$, where $j\colon W\hookrightarrow X$.
  Since $j$ is quasi-compact and $X$ is quasi-separated, $j_*\mathcal G$ is
  quasi-coherent; kernels of maps of quasi-coherent modules are
  quasi-coherent.  Thus $\mathcal H\subseteq\mathcal F$ is quasi-coherent and
  $\mathcal H|_W=\mathcal G$.

  Choose affine opens $V_1,\ldots,V_n$ with
  $X=W\cup V_1\cup\cdots\cup V_n$.  Inductively suppose that a finite type
  submodule $\mathcal G_r\subseteq\mathcal H$ has been constructed on
  $W_r=W\cup V_1\cup\cdots\cup V_r$ and restricts to $\mathcal G$ on $W$.
  Put $V_{r+1}=\operatorname{Spec}(R)$.  On $V_{r+1}$ take the inverse image
  $\mathcal K$ of the direct image of
  $\mathcal G_r|_{W_r\cap V_{r+1}}$ under the restriction map from
  $\mathcal H|_{V_{r+1}}$.  The same quasi-compact-open direct-image argument
  shows that $\mathcal K$ is quasi-coherent, and by construction its
  restriction to $W_r\cap V_{r+1}$ is $\mathcal G_r$.  Write
  $\mathcal K=\widetilde N$.  The intersection is quasi-compact.  Cover it by
  finitely many basic opens $D(f_i)$ on which $\mathcal G_r$ is generated by
  finitely many sections.  Because $\mathcal G_r\subseteq\widetilde N$ there,
  after shrinking the basic opens if necessary and clearing denominators, those
  generators are represented by finitely many elements $x_{ij}\in N$.  Let
  $N'\subseteq N$ be generated by all $x_{ij}$.  Enlarging the denominator
  powers when clearing them ensures
  $\widetilde{N'}|_{D(f_i)}=\mathcal G_r|_{D(f_i)}$: one inclusion holds by
  construction, and the other is checked on the chosen finite generators.
  Hence $\widetilde{N'}$ and $\mathcal G_r$ agree on the whole intersection and
  glue to a finite type submodule on $W_{r+1}$.  Finite induction gives the
  required submodule on $X$.  Taking $\mathcal F=\mathcal O_X$ and
  $\mathcal G=\mathcal J$ gives the asserted ideal; compatibility with the
  inclusions into $\mathcal O_X$ makes the glued submodule an ideal.
\end{proof}

\subsection{An arbitrary composition of two blow-ups}

The product-ideal theorem treats a second center pulled back from the original
scheme.  Stability of admissible blow-ups under composition needs the more
general statement in which the second center is an arbitrary finite-presentation
closed subscheme of the first blow-up.

\begin{lemma}[Existence of finite-degree equations for a closed center in Proj]
  \label{lemma:nagata-proj-closed-center-finite-degree-exists}
  \uses{def:nagata-relative-proj}
  \lean{AlgebraicGeometry.Scheme.IdealSheafData.exists_finiteDegree_center_le_blowup}

  Let $\mathcal I$ be a finite type quasi-coherent ideal on a qcqs scheme
  $X$, put $\mathcal A=\bigoplus_{n\geq0}\mathcal I^n$, and let
  $b\colon X'=\operatorname{Proj}_X(\mathcal A)\to X$ be the structure map.
  For every finite-presentation
  closed immersion $Z'\hookrightarrow X'$ there exist $d>0$ and a finite type
  submodule $\mathcal F\subseteq\mathcal I^d$ such that
  \[
    Z'=\operatorname{Proj}_X(\mathcal A/\mathcal F\mathcal A)
  \]
  and the support of $\mathcal I^d/\mathcal F$ is contained in the union of
  $V(\mathcal I)$ and the underlying topological image $b(|Z'|)$ in $X$.
\end{lemma}

\begin{proof}
  \uses{lemma:nagata-relative-proj-exists,
    prop:nagata-affine-blowup-chart-cover,
    thm:nagata-global-blowup-proper,
    thm:nagata-global-blowup-complement,
    lemma:nagata-finite-type-ideal-extension-exists}
  Identify $X'$ with the global blow-up obtained by gluing the same affine
  Rees-algebra Proj charts; under this identification $b$ is
  $b_{X,\mathcal I}$.  Thus the properness and complement results cited below
  apply to $b$.
  Put $Z=V(\mathcal I)$, $S=b(|Z'|)$, and
  $U=X\setminus(Z\cup S)$.  Properness of $b$ makes $S$ closed.  Moreover
  $X\setminus Z$ is quasi-compact because $\mathcal I$ is of finite type,
  and $b^{-1}(X\setminus Z)\to X\setminus Z$ is an isomorphism.  Under this
  isomorphism, $U$ corresponds to the complement of the finite-presentation
  closed subscheme $Z'$ in $b^{-1}(X\setminus Z)$; that complement is
  quasi-compact.  Thus $U$ is a quasi-compact open disjoint from $b(|Z'|)$.

  Let $\mathcal H\subseteq\mathcal O_{X'}$ be the finite-type
  quasi-coherent ideal defining $Z'$.  On an affine chart
  $D_+(a)\simeq\operatorname{Spec}(A[\mathcal I/a])$, choose finitely many
  generators of $\mathcal H$ and clear their powers of $a$.  After increasing
  the common exponent, every generator is represented by a homogeneous
  numerator in one degree.  Taking a finite affine cover of $X$ and of the
  finitely many standard charts, and then a common degree $d>0$, gives a
  finite collection of homogeneous elements in $\mathcal I^d$.

  More invariantly, let $M_d$ be the quasi-coherent submodule of
  $\mathcal I^d$ consisting of sections whose induced section of
  $\mathcal O_{X'}(d)$ lies in $\mathcal H(d)$.  The preceding chart
  calculation says that, for $d$ sufficiently large, the image of $M_d$ in
  $\mathcal H(d)$ generates $\mathcal H(d)$; hence the homogeneous ideal
  generated by $M_d$ has Proj equal to $Z'$.  Because $U$ is quasi-compact,
  choose finitely many affine opens covering it and finitely many local
  generators of $\mathcal I^d|_U=M_d|_U$.  Apply the general finite-type
  submodule-extension result of
  \Cref{lemma:nagata-finite-type-ideal-extension-exists} to the finite-type
  submodule of $M_d|_U$ generated by these sections, and enlarge the resulting
  finite-type submodule $\mathcal F\subseteq M_d$ by the finitely many cleared
  homogeneous equations above.  Consequently
  $\mathcal F\mathcal A$ cuts out $Z'$ on every standard chart, while
  $\mathcal F|_U=\mathcal I^d|_U$.  (On overlaps the descriptions agree after
  localizing at the chart generator; the localization cocycle makes these
  equalities descend.)  The latter equality gives
  $\operatorname{Supp}(\mathcal I^d/\mathcal F)\subseteq
  V(\mathcal I)\cup b(|Z'|)$ and completes the construction.
\end{proof}

\begin{remark}[Formalization route]
  The formal statement
  (\texttt{exists\_finiteDegree\_center\_le\_blowup}) recasts Tag 0803 purely
  in ideal-sheaf terms: $Z'$ is replaced by a finite-type ideal sheaf
  $\mathcal J$ on the blow-up and the conclusion is the inverse-image
  equality $\mathcal L\cdot\mathcal O_{X'}=\mathcal I_E^d\cdot\mathcal J$ for
  a finite-type $\mathcal L\subseteq\mathcal I^d$; no relative-Proj quotient
  is formed, and the support bound is recovered downstream
  (\texttt{coe\_support\_eq\_of\_comap\_pow\_mul}).  The choice-packaged
  \Cref{def:nagata-proj-closed-center-finite-degree} is replaced by the bare
  existential.
\end{remark}

\begin{definition}[Finite-degree data attached to a closed center in Proj]
  \label{def:nagata-proj-closed-center-finite-degree}
  \uses{lemma:nagata-proj-closed-center-finite-degree-exists,
    def:nagata-relative-proj}
  \lean{AlgebraicGeometry.Proj.finiteDegreeCenterData}

  Let $\mathcal I$ be a finite type quasi-coherent ideal on a qcqs scheme
  $X$, let $\mathcal A=\bigoplus_{n\geq0}\mathcal I^n$, let
  $X'=\operatorname{Proj}_X(\mathcal A)$, and let
  $Z'\hookrightarrow X'$ be of finite presentation.  Select, by classical
  choice from the nonempty witness type supplied by
  \Cref{lemma:nagata-proj-closed-center-finite-degree-exists}, a pair
  $(d,\mathcal F)$.  No property of the selected pair is built into this
  definition.  Denote it by
  $\mathsf{FD}(Z')=(d(Z'),\mathcal F(Z'))$.
\end{definition}

\begin{lemma}[Properties of the finite-degree center data]
  \label{lemma:nagata-proj-closed-center-finite-degree}
  \uses{def:nagata-proj-closed-center-finite-degree}
  \lean{AlgebraicGeometry.Proj.exists_finiteDegree_submodule_defining_closed}

  The selected integer $d(Z')$ is positive, $\mathcal F(Z')$ is a finite
  type submodule of $\mathcal I^{d(Z')}$, and
  \[
    Z'=\operatorname{Proj}_X
      \bigl(\mathcal A/\mathcal F(Z')\mathcal A\bigr).
  \]
  Moreover the support of
  $\mathcal I^{d(Z')}/\mathcal F(Z')$ lies over the underlying topological
  image $b(|Z'|)$ together with the first center.
\end{lemma}

\begin{proof}
  \uses{lemma:nagata-proj-closed-center-finite-degree-exists}
  These are the properties required of the witness selected in
  \Cref{def:nagata-proj-closed-center-finite-degree}; existence of such a
  witness is
  \Cref{lemma:nagata-proj-closed-center-finite-degree-exists}.
\end{proof}

\begin{definition}[Center ideal representing a second blow-up]
  \label{def:nagata-composition-center-ideal}
  \uses{def:nagata-proj-closed-center-finite-degree}
  \lean{AlgebraicGeometry.Blowup.compositionCenterIdealConstruction}

  Let $X$ be qcqs, let $\mathcal I\subseteq\mathcal O_X$ be a finite-type
  quasi-coherent ideal, put
  $X'=\operatorname{Bl}_{\mathcal I}(X)$ with structure map
  $b\colon X'\to X$, and let $Z'\hookrightarrow X'$ be a
  finite-presentation closed immersion.  Write
  $Z=V(\mathcal I)\hookrightarrow X$.  With
  $\mathsf{FD}(Z')=(d,\mathcal F)$, define
  \[
    \mathcal J(Z')=\mathcal F\subseteq\mathcal I^d\subseteq\mathcal O_X,
    \qquad
    \mathcal K(Z')=\mathcal I\mathcal J(Z'),
    \qquad
    Y(Z')=V(\mathcal K(Z')).
  \]
\end{definition}

\begin{lemma}[Support of the constructed composition center]
  \label{lemma:nagata-composition-center-ideal}
  \uses{def:nagata-composition-center-ideal}
  \lean{AlgebraicGeometry.Blowup.compositionCenterIdeal}

  Under the hypotheses of
  \Cref{def:nagata-composition-center-ideal}, the ideal $\mathcal K(Z')$ is
  finite type and
  \[
    |Y(Z')|=|V(\mathcal I)|\cup b(|Z'|).
  \]
\end{lemma}

\begin{proof}
  \uses{lemma:nagata-proj-closed-center-finite-degree,
    thm:nagata-global-blowup-complement,
    thm:nagata-global-blowup-proper}
  Write $S=b(|Z'|)$ and $\mathcal J=\mathcal F\subseteq\mathcal I^d$ as in
  the definition.  The module $\mathcal F$ is finite type, and the product of
  the finite-type ideals $\mathcal I$ and $\mathcal J$ is finite type.  If
  $x\notin V(\mathcal I)\cup S$, then $\mathcal I_x=\mathcal O_{X,x}$ and
  the first blow-up is an isomorphism near $x$ by
  \Cref{thm:nagata-global-blowup-complement}; the defining equation
  $\operatorname{Proj}(\mathcal A/\mathcal F\mathcal A)=Z'$ is empty there.
  On the degree-one chart this says $\mathcal F_x=\mathcal I_x^d=\mathcal O_{X,x}$,
  hence $\mathcal K_x=\mathcal O_{X,x}$.  Thus
  $V(\mathcal K)\subseteq V(\mathcal I)\cup S$.

  Conversely, $\mathcal K\subseteq\mathcal I$ gives
  $V(\mathcal I)\subseteq V(\mathcal K)$.  If $x\in S\setminus V(\mathcal I)$,
  then $b$ is an isomorphism near $x$, and the standard chart calculation for
  $\operatorname{Proj}(\mathcal A/\mathcal F\mathcal A)$ identifies
  $x\in S$ with $\mathcal F_x\ne\mathcal O_{X,x}$.  Since
  $\mathcal I_x^d=\mathcal O_{X,x}$, this gives
  $\mathcal K_x=\mathcal F_x\ne\mathcal O_{X,x}$, so
  $x\in V(\mathcal K)$.  Both inclusions follow, and therefore
  $|Y(Z')|=|V(\mathcal I)|\cup b(|Z'|)$.
\end{proof}


\begin{lemma}[Comparison of the two second-center ideals]
  \label{lemma:nagata-composition-extended-ideal-equality}
  \uses{def:nagata-composition-center-ideal}
  \lean{AlgebraicGeometry.Blowup.extended_compositionCenterIdeal}

  Under the hypotheses of
  \Cref{def:nagata-composition-center-ideal}, let $E$ be the exceptional
  divisor of the first blow-up, write $\mathcal I_E$ for its invertible ideal,
  and let $\mathcal I_{Z'}$ define the second center.  On $X'$ one has
  \[
    b^{-1}\mathcal J\,\mathcal O_{X'}
      =\mathcal I_E^d\mathcal I_{Z'}.
  \]
\end{lemma}

\begin{proof}
  \uses{lemma:nagata-composition-center-ideal,
    lemma:nagata-chart-extended-ideal-principal}
  Let $D_+(at)=\operatorname{Spec}(R_a)$ be a standard chart of the first
  blow-up, where $R_a=A[\mathcal I/a]$ and $a\in\mathcal I$.  The extended
  ideal of $\mathcal J$ is, by definition, the image of $\mathcal J$ in
  $R_a$.  The equation
  $Z'=\operatorname{Proj}(\mathcal A/\mathcal F\mathcal A)$ says that its
  ideal on this chart is generated by all fractions $f/a^d$ with
  $f\in\mathcal F$ (and their $R_a$-multiples).  The chart calculation for
  the exceptional ideal gives $\mathcal I_E R_a=(a)$.  Therefore, for every
  generator $f\in\mathcal F$,
  \[
    f=a^d\,(f/a^d)\quad\text{in }R_a,
  \]
  and hence
  $(\mathcal J\mathcal O_{X'})R_a=(a)^d\,\mathcal I_{Z'}R_a$.
  This proves the asserted equality on every standard chart.  On the overlap
  of the $a$- and $c$-charts one has
  $f/a^d=(c/a)^d(f/c^d)$ and $(a)^d=(a/c)^d(c)^d$; the two products are
  literally the same ideal after localization.  The standard charts cover
  $X'$, so equality of the quasi-coherent ideals glues to
  $b^{-1}\mathcal J\,\mathcal O_{X'}=\mathcal I_E^d\mathcal I_{Z'}$.
\end{proof}

\begin{lemma}[Blow-up is unchanged after multiplying by an invertible factor]
  \label{lemma:nagata-blowup-invertible-factor}
  \uses{def:nagata-invertible-ideal,
    thm:nagata-blowup-universal-property}
  \lean{AlgebraicGeometry.Scheme.IdealSheafData.exists_blowup_mul_isEffectiveCartier_iso}

  Let $X$ be a scheme, let $\mathcal L$ be an invertible ideal on $X$, and
  let $\mathcal J$ be a quasi-coherent ideal (finite-typeness is not needed).
  Then the blow-ups of $\mathcal L\mathcal J$ and of $\mathcal J$ are
  canonically isomorphic over $X$.
\end{lemma}

\begin{proof}
  Tensor the tautological quotient for $\mathcal J$ with the pullback of the
  invertible ideal $\mathcal L$ to obtain an invertible quotient of
  $\mathcal L\mathcal J$; this gives a lift from
  $\operatorname{Bl}_{\mathcal J}(X)$ to
  $\operatorname{Bl}_{\mathcal L\mathcal J}(X)$.  Conversely, tensor the
  tautological quotient of $\mathcal L\mathcal J$ with the dual of the
  pulled-back invertible module $\mathcal L$.  The evaluation isomorphism
  gives an invertible quotient of $\mathcal J$ and hence a reverse lift.
  Locally this is division by a basis of $\mathcal L$, so the quotient maps
  are compatible with the specified ideal inclusions and glue independently
  of the chosen bases.  Both composites induce the original tautological
  quotient, so uniqueness in
  \Cref{thm:nagata-blowup-universal-property} makes them identities.
\end{proof}

\begin{definition}[Comparison for a composite of arbitrary blow-ups]
  \label{def:nagata-composition-of-blowups-comparison}
  \uses{def:nagata-composition-center-ideal,
    def:nagata-blowup-product-comparison,
    def:nagata-global-blowup,
    lemma:nagata-blowup-invertible-factor}
  \lean{AlgebraicGeometry.Blowup.compComparisonBlowup}

  Let $X$, $\mathcal I$, $X'$, $b$, $Z'$, $d$, and $\mathcal J(Z')$ be as in
  \Cref{def:nagata-composition-center-ideal}, and let the second blow-up be
  centred at $Z'\hookrightarrow X'$.  Let $E$ be the exceptional divisor of
  $b$ and let $\mathcal I_E$ be its invertible ideal; write
  $\mathcal I_{Z'}$ for the ideal defining $Z'$.  Form the explicit ideal
  $\mathcal K(Z')=\mathcal I\mathcal J(Z')$ on $X$.  The extended ideal on
  $X'$ is $\mathcal I_E^d\mathcal I_{Z'}$.  Define
  \[
    \Omega_{Z,Z'}\colon
      \operatorname{Bl}_{\mathcal I_{Z'}}(X')
      \longrightarrow
      \operatorname{Bl}_{\mathcal K(Z')}(X)
  \]
  as the product-blow-up comparison followed by the comparison which removes
  the invertible factor $\mathcal I_E^d$.  Define the reverse map by the
  corresponding universal lifts.
\end{definition}

\begin{theorem}[{\cite[\href{https://stacks.math.columbia.edu/tag/080B}{Tag 080B}]{stacks-project}}]
  \label{thm:nagata-composition-of-blowups}
  \uses{def:nagata-composition-of-blowups-comparison}
  \lean{AlgebraicGeometry.Scheme.IdealSheafData.exists_blowup_comp}

  With $Z=V(\mathcal I)$ as above, the two constructed comparison maps are
  inverse over $X$.  Hence the
  composite is the blow-up of the finite-presentation center
  $Y(Z')=V(\mathcal K(Z'))\hookrightarrow X$, whose support satisfies
  \[
    |Y(Z')|=|Z|\cup b(|Z'|).
  \]
\end{theorem}

\begin{proof}
  \uses{lemma:nagata-composition-center-ideal,
    lemma:nagata-composition-extended-ideal-equality,
    thm:nagata-blowup-product-ideal,
    cor:nagata-blowup-cartier-iso,
    lemma:nagata-blowup-invertible-factor}
  Put $X_1=\operatorname{Bl}_{\mathcal I}(X)$ and
  $\mathcal J'=\mathcal J\mathcal O_{X_1}$.  By the extended-ideal equality,
  $\mathcal J'=\mathcal I_E^d\mathcal I_{Z'}$.  The product-ideal comparison
  \Cref{thm:nagata-blowup-product-ideal} therefore gives an isomorphism over
  $X$
  \[
    \operatorname{Bl}_{\mathcal J'}(X_1)
      \xrightarrow{\sim}\operatorname{Bl}_{\mathcal I\mathcal J}(X).
  \]
  Since $\mathcal I_E^d$ is invertible, the invertible-factor comparison
  \Cref{lemma:nagata-blowup-invertible-factor} gives a canonical isomorphism
  $\operatorname{Bl}_{\mathcal I_{Z'}}(X_1)\cong
    \operatorname{Bl}_{\mathcal J'}(X_1)$.  Their composite is the map
  $\Omega_{Z,Z'}$ in the definition; the reverse composite of the inverse
  comparisons is the stated reverse map.

  For completeness, these maps are the universal lifts for the same quotient:
  on the iterated blow-up the product of the tautological invertible quotients
  is the tautological quotient of $\mathcal I\mathcal J$, while on the
  $\mathcal I\mathcal J$-blow-up its two factor quotients give the tautological
  quotients for $\mathcal I$ and for $\mathcal J'$; multiplying by the inverse
  of the local generator of $\mathcal I_E^d$ gives the quotient for
  $\mathcal I_{Z'}$.  Thus each composite is an $X$-morphism inducing the
  original tautological quotient.  Uniqueness in the blow-up universal
  property, first over $X$ and then over $X_1$, makes both composites identities.
  The center ideal $\mathcal K=\mathcal I\mathcal J$ is finite type, hence
  defines a finite-presentation closed immersion, and
  \Cref{lemma:nagata-composition-center-ideal} gives its support as
  $|V(\mathcal I)|\cup b(|Z'|)$.
\end{proof}


\section{Admissible blow-ups and strict transforms}

\begin{lemma}[An admissible blow-up is unchanged over the preserved open]
  \label{lemma:nagata-admissible-blowup-over-open}
  \uses{def:nagata-admissible-blowup,
    thm:nagata-global-blowup-complement}
  \lean{AlgebraicGeometry.AdmissibleBlowup.isIso_π_restrict}

  If $\pi\colon X'\to X$ is $U$-admissible, then its center ideal restricts to
  the unit ideal on $U$, and the induced morphism
  $X'\times_XU\to U$ is an isomorphism.  The inverse and the equality over
  $U$ are part of the resulting specified isomorphism.
\end{lemma}

\begin{proof}
  Choose the ideal and blow-up identification from
  \Cref{def:nagata-admissible-blowup}.  The unit-ideal restriction is the
  disjointness condition on the center.  Apply
  \Cref{thm:nagata-global-blowup-complement} to identify the blow-up over $U$
  with $U$, then transport that isomorphism across the chosen identification
  with $X'$.
\end{proof}

\begin{definition}[Strict transform]
  \label{def:nagata-strict-transform}
  \uses{def:nagata-scheme-theoretic-closure,
    def:nagata-global-blowup}
  \lean{AlgebraicGeometry.StrictTransform}

  Let $Z\to X$ be a morphism and let $X'\to X$ be the blow-up of a
  finite-presentation center $C\subseteq X$.  Put $U=X\setminus C$; then
  $U\hookrightarrow X$ is quasi-compact.  The \emph{strict transform} of $Z$
  is the scheme-theoretic closure of
  \[
    (Z\times_XU)\times_U(X'\times_XU)
  \]
  in $Z\times_XX'$, in the sense of
  \Cref{def:nagata-scheme-theoretic-closure}.  It is not enough to use only
  the topological closure.
\end{definition}

\begin{remark}[Formalization route]
  The formal development (\texttt{AlgebraicGeometry.StrictTransform},
  specialized to blow-ups as
  \texttt{Scheme.IdealSheafData.strictTransform}) works with this
  scheme-theoretic-closure definition directly, and proves the identification
  over the preserved open (\Cref{lemma:nagata-strict-transform-over-complement})
  by scheme-image flat base change for a quasi-compact flat monomorphism
  (\texttt{Nagata.ker\_comap\_self}), not through the torsion-quotient chart
  model.  The chart-by-chart infrastructure below
  (\Cref{def:nagata-strict-transform-chart-ring} through
  \Cref{thm:nagata-strict-transform-is-blowup}) is accordingly not
  formalized; of the power-torsion material only the affine algebraic layer
  (\texttt{Submodule.powTorsion} and its localization/exactness lemmas) is.
\end{remark}

\subsection{Power torsion and strict-transform charts}

The affine power-torsion quotient and its localization formula are proved
below.  The only remaining strict-transform infrastructure is the transport
of these affine identifications through the global blow-up restriction
cocycle; uses of that cocycle retain an explicit dependency on the global
blow-up interface.

The strict transform is the correction term for the failure of blow-ups to
commute with non-flat base change.  The following algebraic lemmas make that
correction explicit before it is sheafified and glued.

\begin{definition}[Power-torsion submodule]
  \label{def:nagata-power-torsion}
  \lean{Submodule.powTorsion}

  Let $R$ be a ring, let $M$ be an $R$-module, and let $a\in R$.  Define
  \[
    \Gamma_{(a)}(M)
      =\{m\in M\mid \text{there is }n\geq0\text{ with }a^nm=0\}.
  \]
\end{definition}

\begin{lemma}[Power torsion is a submodule]
  \label{lemma:nagata-power-torsion-submodule}
  \uses{def:nagata-power-torsion}
  \lean{Submodule.powTorsion_isSubmodule}

  The subset $\Gamma_{(a)}(M)$ contains zero and is closed under addition and
  scalar multiplication.
\end{lemma}

\begin{proof}
  Zero is killed by $a^0$.  If $a^mm_1=0$ and $a^nm_2=0$, then
  $a^{m+n}(m_1+m_2)=0$.  If $r\in R$, then
  $a^m(rm_1)=r(a^mm_1)=0$.
\end{proof}

\begin{lemma}[Membership in power torsion]
  \label{lemma:nagata-mem-power-torsion}
  \uses{def:nagata-power-torsion}
  \lean{Submodule.mem_powTorsion_iff}

  One has
  \[
    m\in\Gamma_{(a)}(M)
      \Longleftrightarrow
    \exists n,\ a^nm=0.
  \]
  The right-hand side should be the simp-normal form for membership.
\end{lemma}

\begin{proof}
  This is the defining membership condition of
  \Cref{def:nagata-power-torsion}; it should be registered as the simp theorem
  for that definition.
\end{proof}

\begin{lemma}[Power torsion as a localization kernel]
  \label{lemma:nagata-power-torsion-kernel-localization}
  \uses{def:nagata-power-torsion}
  \lean{Submodule.powTorsion_eq_ker_localization}

  The power-torsion submodule is the kernel of the localization map
  \[
    M\longrightarrow M_a.
  \]
\end{lemma}

\begin{proof}
  An element maps to zero in localization away from the powers of $a$ exactly
  when it is killed by one element of that multiplicative system, namely by a
  power of $a$.
\end{proof}

\begin{lemma}[Localization of power torsion]
  \label{lemma:nagata-power-torsion-localization}
  \lean{Submodule.localized'_powTorsion}

  Let $s\in R$.  Under localization at $s$, the image of
  $\Gamma_{(a)}(M)$ is
  $\Gamma_{(a/1)}(M_s)$.
\end{lemma}

\begin{proof}
  \uses{lemma:nagata-power-torsion-kernel-localization}
  Localize the kernel description in
  \Cref{lemma:nagata-power-torsion-kernel-localization}.  Localization is flat,
  so it preserves this kernel.  The iterated localization $(M_s)_{a/1}$ is
  canonically $M_{sa}$.
\end{proof}

\begin{lemma}[The power-torsion quotient has no exceptional torsion]
  \label{lemma:nagata-power-torsion-quotient-torsionfree}
  \uses{def:nagata-power-torsion}
  \lean{Submodule.powTorsion_quotient_torsionfree}

  Multiplication by $a$ on $M/\Gamma_{(a)}(M)$ is injective.  More generally,
  no nonzero class in the quotient is killed by a power of $a$.
\end{lemma}

\begin{proof}
  If $a^n[m]=0$, then $a^nm$ belongs to the power-torsion submodule, so some
  further power of $a$ kills $m$.  Thus $m$ already represents zero in the
  quotient.
\end{proof}

\begin{definition}[Affine base-change chart ring]
  \label{def:nagata-strict-transform-basechange-chart-ring}
  \uses{def:nagata-blowup-chart-ring}
  \lean{AlgebraicGeometry.strictTransformBaseChangeChartRing}

  Let $A\to B$ be a ring map, let $I\subseteq A$, and let $a\in I$.
  Define
  \[
    C_a(A,B,I)=B\otimes_A A[I/a].
  \]
  We regard $a$ as an element of $C_a(A,B,I)$ through the structural map
  $A\to C_a(A,B,I)$.
\end{definition}

\begin{lemma}[Ordinary base-change chart]
  \label{lemma:nagata-strict-transform-basechange-chart}
  \uses{def:nagata-strict-transform-basechange-chart-ring}
  \lean{AlgebraicGeometry.strictTransform_baseChangeChart}

  Over the $a$-chart of
  $\operatorname{Bl}_I(\operatorname{Spec}A)$, the ordinary base change
  \[
    \operatorname{Spec}(B)\times_{\operatorname{Spec}A}
    \operatorname{Bl}_I(\operatorname{Spec}A)
  \]
  is canonically $\operatorname{Spec}(C_a(A,B,I))$.
\end{lemma}

\begin{proof}
  \uses{prop:nagata-affine-blowup-chart}
  Apply the affine fibre-product formula to the constructed chart
  identification of \Cref{prop:nagata-affine-blowup-chart}.  The tensor
  product appearing in that formula is precisely
  \Cref{def:nagata-strict-transform-basechange-chart-ring}.
\end{proof}

\begin{definition}[Affine strict-transform chart ring]
  \label{def:nagata-strict-transform-chart-ring}
  \uses{def:nagata-power-torsion,
    def:nagata-strict-transform-basechange-chart-ring}
  \lean{AlgebraicGeometry.strictTransformChartRing}

  Define
  \[
    C_a^{\mathrm{st}}(A,B,I)
      =C_a(A,B,I)/\Gamma_{(a)}(C_a(A,B,I)).
  \]
\end{definition}

\begin{lemma}[Closedness of the strict-transform chart]
  \label{lemma:nagata-strict-transform-chart-closed}
  \uses{def:nagata-strict-transform-chart-ring}
  \lean{AlgebraicGeometry.strictTransformChart_closedImmersion}

  The quotient map
  $C_a(A,B,I)\to C_a^{\mathrm{st}}(A,B,I)$ defines a closed immersion of
  the strict-transform chart into the ordinary base-change chart.
\end{lemma}

\begin{proof}
  \uses{lemma:nagata-power-torsion-submodule}
  The power-torsion subset is an ideal by
  \Cref{lemma:nagata-power-torsion-submodule}, applied to the regular module
  $C_a(A,B,I)$.  The spectrum of a quotient ring is a closed subscheme.
\end{proof}

\begin{definition}[Map to the blow-up chart of the extended ideal]
  \label{def:nagata-strict-chart-to-extended-blowup-chart}
  \uses{def:nagata-strict-transform-basechange-chart-ring,
    def:nagata-blowup-chart-ring}
  \lean{AlgebraicGeometry.strictTransformChartToExtendedBlowupChart}

  Let $J=IB\subseteq B$, let $b$ be the image of $a$ in $B$, and write
  $B[J/b]$ for the corresponding chart ring.  Define
  \[
    \sigma_a\colon C_a(A,B,I)\longrightarrow B[J/b]
  \]
  on pure tensors and chart generators by
  \[
    \sigma_a\bigl(y\otimes(x/a)\bigr)=y\,\overline x/b.
  \]
\end{definition}

\begin{lemma}[Well-definedness of the extended-ideal chart map]
  \label{lemma:nagata-strict-chart-to-extended-blowup-chart}
  \uses{def:nagata-strict-chart-to-extended-blowup-chart}
  \lean{AlgebraicGeometry.strictTransformChartToExtendedBlowupChart_apply}

  The assignment defining $\sigma_a$ respects addition, multiplication, the
  localization relations in $A[I/a]$, and the two $A$-algebra structures.
  Consequently $\sigma_a$ is a well-defined $B$-algebra homomorphism.
\end{lemma}

\begin{proof}
  For $x,x'\in I$, the relations
  \[
    \frac{x+x'}a=\frac xa+\frac{x'}a,
    \qquad
    \frac xa\frac{x'}a=\frac{xx'}{a^2}
  \]
  have the same images in $B_b$.  The image of $A$ through the left tensor
  factor equals its image through the chart factor.  The universal properties
  of $A[I/a]$ and of $B\otimes_AA[I/a]$ therefore give the stated map, and
  its displayed formula determines it uniquely.
\end{proof}

\begin{lemma}[Surjectivity onto the extended-ideal chart]
  \label{lemma:nagata-strict-chart-map-surjective}
  \uses{def:nagata-strict-chart-to-extended-blowup-chart}
  \lean{AlgebraicGeometry.strictTransformChartToExtended_surjective}

  The constructed homomorphism $\sigma_a$ is surjective.
\end{lemma}

\begin{proof}
  \uses{lemma:nagata-strict-chart-to-extended-blowup-chart,
    lemma:nagata-chart-ring-normal-form}
  By \Cref{lemma:nagata-chart-ring-normal-form}, the target is generated as a
  $B$-algebra by the fractions $z/b$ with $z\in J$.  Write
  $z=\sum_k y_k\overline{x_k}$ with $x_k\in I$.  Then
  \[
    z/b=\sum_k\sigma_a\bigl(y_k\otimes(x_k/a)\bigr),
  \]
  proving surjectivity.
\end{proof}

\begin{lemma}[Kernel of the map to the extended-ideal chart]
  \label{lemma:nagata-strict-chart-map-kernel}
  \uses{def:nagata-power-torsion,def:nagata-strict-chart-to-extended-blowup-chart}
  \lean{AlgebraicGeometry.ker_strictTransformChartToExtended}

  The kernel of $\sigma_a$ is exactly
  $\Gamma_{(a)}(C_a(A,B,I))$.
\end{lemma}

\begin{proof}
  \uses{lemma:nagata-strict-chart-map-surjective,
    lemma:nagata-power-torsion-kernel-localization}
  Both source and target map to $B_b$.  The target $B[J/b]$ is the subalgebra
  generated there by the fractions from $J$, and $\sigma_a$ has exactly that
  image by \Cref{lemma:nagata-strict-chart-map-surjective}.  Thus its kernel
  is the kernel of localization at $b$.  Under the $A$-action this is the
  kernel of localization at $a$, which equals the $a$-power torsion by
  \Cref{lemma:nagata-power-torsion-kernel-localization}.
\end{proof}

\begin{definition}[Strict-transform chart comparison]
  \label{def:nagata-strict-transform-chart-comparison}
  \uses{def:nagata-strict-transform-chart-ring,
    def:nagata-strict-chart-to-extended-blowup-chart}
  \lean{AlgebraicGeometry.strictTransformChartComparison}

  Define
  \[
    \overline\sigma_a\colon
      C_a^{\mathrm{st}}(A,B,I)\longrightarrow B[IB/b]
  \]
  to be the map induced by $\sigma_a$ on the power-torsion quotient.
\end{definition}

\begin{proposition}[Affine strict-transform chart formula]
  \label{prop:nagata-strict-transform-chart-iso}
  \uses{def:nagata-strict-transform-chart-comparison}
  \lean{AlgebraicGeometry.strictTransformChartIsoExtendedBlowupChart}

  The constructed comparison $\overline\sigma_a$ is a ring equivalence.
\end{proposition}

\begin{proof}
  \uses{lemma:nagata-strict-chart-map-surjective,
    lemma:nagata-strict-chart-map-kernel}
  The map $\sigma_a$ is surjective and its kernel is precisely the ideal by
  which the source was quotiented.  The first isomorphism theorem therefore
  makes $\overline\sigma_a$ bijective and identifies its inverse ring
  structure, so it is a ring equivalence.
\end{proof}

\begin{definition}[Localized strict-transform chart comparison]
  \label{def:nagata-strict-transform-chart-comparison-localized}
  \uses{def:nagata-strict-transform-chart-comparison}
  \lean{AlgebraicGeometry.strictTransformChartComparisonLocalization}

  After localizing the base ring or a chart ring, define
  $(\overline\sigma_a)_s$ by localizing the source and target of
  $\overline\sigma_a$ and transporting along the canonical tensor-product
  and iterated-localization associativity maps.
\end{definition}

\begin{lemma}[Strict-transform chart formula localizes]
  \label{lemma:nagata-strict-transform-chart-localization}
  \uses{def:nagata-strict-transform-chart-comparison-localized}
  \lean{AlgebraicGeometry.strictTransformChartIso_localization}

  The constructed localized comparison $(\overline\sigma_a)_s$ equals the
  strict-transform chart comparison constructed directly from the localized
  data.
\end{lemma}

\begin{proof}
  \uses{lemma:nagata-power-torsion-localization,
    prop:nagata-strict-transform-chart-iso}
  Localization commutes with the tensor product and with the power-torsion
  quotient by \Cref{lemma:nagata-power-torsion-localization}.  Both maps send
  every pure tensor and every generator $x/a$ to the same fraction.  Ring-map
  extensionality gives the equality, and
  \Cref{prop:nagata-strict-transform-chart-iso} supplies the equivalence
  property on both sides.
\end{proof}

\begin{lemma}[Localization data for strict-transform chart quotients]
  \label{lemma:nagata-strict-transform-chart-localization-data}
  \uses{def:nagata-strict-transform-chart-ring,
    lemma:nagata-power-torsion-localization}
  \lean{AlgebraicGeometry.strictTransformChartLocalizationData}

  For $s\in B$, the canonical map
  \[
    (C_a(A,B,I)/\Gamma_{(a)}C_a(A,B,I))_s
      \longrightarrow
    C_a(A,B_s,I)/\Gamma_{(a/1)}C_a(A,B_s,I)
  \]
  is an isomorphism.  These isomorphisms are compatible with further
  localization: the maps obtained by localizing successively at $s$ and $t$
  agree with localization at $st$.
\end{lemma}

\begin{proof}
  Localization commutes with tensor products, so
  $(B\otimes_A A[I/a])_s\cong B_s\otimes_A A[I/a]$.  By
  \Cref{lemma:nagata-power-torsion-localization}, the image of the
  $a$-power-torsion submodule is exactly the $a/1$-power-torsion submodule.
  The universal property of quotients gives the displayed isomorphism.
  Compatibility for $s,t$ follows because both composites send pure tensors
  and chart fractions to the same elements of the $st$-localization.
\end{proof}

\begin{lemma}[Cocycle for strict-transform charts]
  \label{lemma:nagata-strict-transform-chart-cocycle}
  \uses{def:nagata-strict-transform-chart-comparison-localized}
  \lean{AlgebraicGeometry.strictTransformChartIso_cocycle}

  The localized strict-transform chart comparisons agree on double
  localizations and satisfy the cocycle equality on triple chart overlaps.
\end{lemma}

\begin{proof}
  \uses{lemma:nagata-strict-transform-chart-localization,
    lemma:nagata-strict-transform-chart-localization-data,
    lemma:nagata-affine-blowup-restrict-cocycle}
  On an intersection of the $a$- and $a'$-charts, choose a basic affine
  refinement on which both chart restrictions are defined.  The ordinary
  blow-up transition is the specified localization isomorphism; by
  \Cref{lemma:nagata-affine-blowup-restrict-cocycle} its two composites on a
  triple refinement agree.  Localizing the quotient map
  $C_a\twoheadrightarrow C_a^{\mathrm{st}}$ gives the quotient by the
  localized power-torsion ideal, by
  \Cref{lemma:nagata-strict-transform-chart-localization-data}.  Thus the
  ordinary transition descends uniquely to a transition of strict-transform
  chart rings.  The descended maps agree with the comparisons constructed
  after localization by
  \Cref{lemma:nagata-strict-transform-chart-localization}.  On a triple
  refinement their two composites become the same ordinary transition before
  quotienting, hence are equal after quotienting.  This proves both the
  double-overlap compatibility and the cocycle.
\end{proof}

\begin{definition}[Strict transform of a quasi-coherent module]
  \label{def:nagata-module-strict-transform}
  \uses{def:nagata-power-torsion,def:nagata-global-blowup}
  \lean{QuasiCoherent.Sheaf.strictTransform}

  For a morphism $f\colon Y\to X$, a quasi-coherent
  $\mathcal O_Y$-module $\mathcal F$, a finite-type quasi-coherent ideal
  $\mathcal I\subseteq\mathcal O_X$, and its blow-up
  $X'=\operatorname{Bl}_{\mathcal I}(X)$, define the pullback center ideal on
  $Y\times_XX'$ to be the image of
  $\mathcal I\mathcal O_{X'}$ in $\mathcal O_{Y\times_XX'}$.  The strict
  transform of $\mathcal F$ is its pullback to $Y\times_XX'$ modulo the
  subsheaf of sections killed locally by a power of this named ideal.
\end{definition}

\begin{lemma}[Affine formula for a module strict transform]
  \label{lemma:nagata-module-strict-transform-chart}
  \uses{def:nagata-module-strict-transform,def:nagata-power-torsion}
  \lean{QuasiCoherent.Sheaf.strictTransform_restrict_chart}

  Let $A\to B$ be a ring map, let $I\subseteq A$ be an ideal, let $a\in I$,
  and let $\mathcal F$ on $\operatorname{Spec}(B)$ correspond to the
  $B$-module $M$.  For the induced morphism
  $\operatorname{Spec}(B)\to\operatorname{Spec}(A)$ and the blow-up of $I$,
  on the $a$-chart the
  strict transform corresponds to
  \[
    \frac{M\otimes_AA[I/a]}
         {a\text{-power torsion}}.
  \]
\end{lemma}

\begin{proof}
  Pullback of a quasi-coherent module on an affine morphism is tensor product.
  Sections supported on the exceptional closed subset are exactly the elements
  killed by a power of its local equation $a$.
\end{proof}

\begin{lemma}[Localization of module power torsion and its quotient]
  \label{lemma:nagata-module-power-torsion-localization}
  \uses{def:nagata-power-torsion,
    lemma:nagata-power-torsion-kernel-localization}
  \lean{Submodule.localized'_powTorsion}

  Let $R$ be a ring, $N$ an $R$-module, and $a,s\in R$.  Under the canonical
  localization map one has
  \[
    \operatorname{im}\!\left(\Gamma_{(a)}(N)\to N_s\right)
      =\Gamma_{(a/1)}(N_s),
  \]
  and consequently a canonical isomorphism
  \[
    (N/\Gamma_{(a)}(N))_s
      \;\xrightarrow{\sim}\;
    N_s/\Gamma_{(a/1)}(N_s).
  \]
  These isomorphisms commute with successive localization at $s$ and $t$.
\end{lemma}

\begin{proof}
  The power-torsion submodule is the kernel of the map
  $N\to N_a$ by
  \Cref{lemma:nagata-power-torsion-kernel-localization}.  Localization at
  $s$ is flat, hence carries this kernel to the kernel of
  $N_s\to (N_s)_{a/1}$.  The latter kernel is
  $\Gamma_{(a/1)}(N_s)$, proving the first equality.  The universal property
  of the quotient then gives the displayed map; it is an isomorphism because
  localization is exact and the localized kernel is exactly the target
  torsion submodule.  For $s,t$, both composites are induced by the same map
  $N\to N_{st}$ and therefore agree by quotient and localization
  extensionality.
\end{proof}

\begin{lemma}[Quasi-coherence of the module strict transform]
  \label{lemma:nagata-module-strict-transform-qc}
  \uses{def:nagata-module-strict-transform}
  \lean{QuasiCoherent.Sheaf.strictTransform_isQuasiCoherent}

  The module strict transform is quasi-coherent.
\end{lemma}

\begin{proof}
  \uses{lemma:nagata-module-strict-transform-chart,
    lemma:nagata-module-power-torsion-localization}
  On affine charts it is associated to the quotient module in
  \Cref{lemma:nagata-module-strict-transform-chart}.  These modules localize
  correctly by
  \Cref{lemma:nagata-module-power-torsion-localization}.  On a basic overlap,
  the transition map is the quotient of the ordinary localization map; the
  final assertion of that lemma gives equality for two successive
  localizations, hence the cocycle on triple overlaps.  The associated affine
  sheaves therefore glue to a quasi-coherent module.
\end{proof}

\begin{definition}[Closed chart model of the global strict transform]
  \label{def:nagata-strict-transform-closed-chart-model}
  \uses{def:nagata-strict-transform-chart-ring,
    def:nagata-global-blowup,
    lemma:nagata-strict-transform-chart-localization-data}
  \lean{AlgebraicGeometry.Hom.strictTransformClosedChartModel}

  Let $f\colon Y\to X$, let $\mathcal I$ be a finite-type quasi-coherent
  ideal, and put $X'=\operatorname{Bl}_{\mathcal I}(X)$.  Choose an affine
  cover $V_i=\operatorname{Spec}(A_i)$ of $X$, finite generating sets for
  $\mathcal I|_{V_i}$, and affine covers
  $W_{ij}=\operatorname{Spec}(B_{ij})$ of $f^{-1}(V_i)$.  Define
  $Y^{\mathrm{st}}_{\mathrm{chart}}$ by gluing the affine schemes
  $\operatorname{Spec}(C_{i,j,a}^{\mathrm{st}})$ and their quotient closed
  immersions into the corresponding affine pieces of $Y\times_XX'$.
  On an overlap, refine to basic opens in the chosen $W_{ij}$ and in the
  blow-up chart.  The transition maps are the localizations of
  \Cref{lemma:nagata-strict-transform-chart-localization-data}; they preserve
  the quotient ideals and satisfy the double- and triple-overlap identities.
  These maps, together with the affine closed immersions, are the specified
  scheme-glue datum.  No affineness of an unchosen fibre product is used.
\end{definition}

\begin{proposition}[Global strict transform is a closed subscheme]
  \label{prop:nagata-strict-transform-closed-fiber-product}
  \uses{def:nagata-strict-transform-closed-chart-model,
    def:nagata-strict-transform}
  \lean{AlgebraicGeometry.Hom.strictTransform_closedImmersion}

  Under the data of
  \Cref{def:nagata-strict-transform-closed-chart-model}, the constructed chart
  model $Y^{\mathrm{st}}_{\mathrm{chart}}$ is a closed
  subscheme of $Y\times_XX'$ and agrees with the strict transform of
  \Cref{def:nagata-strict-transform}.
\end{proposition}

\begin{proof}
  \uses{lemma:nagata-strict-transform-chart-closed,
    lemma:nagata-strict-transform-chart-cocycle,
    lemma:nagata-strict-transform-chart-localization-data}
  The chartwise closed immersions agree on overlaps by
  \Cref{lemma:nagata-strict-transform-chart-cocycle}; therefore their ideals
  glue to a global quasi-coherent ideal, with localization compatibility
  supplied by \Cref{lemma:nagata-strict-transform-chart-localization-data}.
  On each chart the quotient removes
  exactly the sections supported on the exceptional divisor, so the resulting
  closed subscheme is the scheme-theoretic closure of the preserved open and
  hence is the strict transform.
\end{proof}

\begin{definition}[Comparison from the strict transform to the extended-ideal blow-up]
  \label{def:nagata-strict-transform-blowup-comparison}
  \uses{def:nagata-strict-transform-closed-chart-model,
    def:nagata-strict-transform-chart-comparison,
    def:nagata-global-blowup}
  \lean{AlgebraicGeometry.Hom.strictTransformToBlowup}

  Let $f\colon Y\to X$, let $\mathcal I$ be a finite-type quasi-coherent
  ideal on $X$, and put $X'=\operatorname{Bl}_{\mathcal I}(X)$.  Define
  \[
    \Theta_{f,\mathcal I}\colon
      Y^{\mathrm{st}}\longrightarrow
      \operatorname{Bl}_{\mathcal I\mathcal O_Y}(Y)
  \]
  by gluing the affine morphisms
  $\operatorname{Spec}(\overline\sigma_a^{-1})$ on the standard blow-up
  charts.
\end{definition}

\begin{theorem}[{\cite[\href{https://stacks.math.columbia.edu/tag/080E}{Tag 080E}]{stacks-project}}]
  \label{thm:nagata-strict-transform-is-blowup}
  \uses{def:nagata-strict-transform-blowup-comparison}
  \lean{AlgebraicGeometry.Hom.strictTransformIsoBlowup}

  Under the preceding hypotheses, the constructed morphism
  $\Theta_{f,\mathcal I}$ is an isomorphism over
  $Y$ and identifies the strict transform with the blow-up of $Y$ in the
  extended ideal $\mathcal I\mathcal O_Y$.
\end{theorem}

\begin{proof}
  \uses{prop:nagata-strict-transform-chart-iso,
    prop:nagata-strict-transform-closed-fiber-product,
    prop:nagata-global-blowup-affine-restrict,
    lemma:nagata-strict-transform-chart-cocycle,
    lemma:nagata-strict-transform-chart-localization-data}
  On every standard affine chart, $\Theta_{f,\mathcal I}$ is induced by the
  ring equivalence $\overline\sigma_a$ of
  \Cref{prop:nagata-strict-transform-chart-iso}; hence it is an isomorphism
  there.  The charts cover the target by
  \Cref{prop:nagata-global-blowup-affine-restrict}, and their comparison maps
  agree on overlaps by \Cref{lemma:nagata-strict-transform-chart-cocycle}.
  More explicitly, on every basic refinement the inverse is the localization
  of $\overline\sigma_a^{-1}$; the equality of the two possible localized
  inverses follows from
  \Cref{lemma:nagata-strict-transform-chart-localization-data}.  The
  morphism-gluing theorem therefore gives a morphism in the reverse direction.
  Its composites with $\Theta_{f,\mathcal I}$ restrict to the identity on
  every member of the affine cover, and so are identities globally.  The
  chart maps commute with the maps to $Y$, since all are induced by the
  structural map $B\to B[IB/b]$; this equality also glues.  Hence the
  resulting isomorphism is over $Y$.
\end{proof}


\begin{lemma}[Scheme-theoretic closure characterization]
  \label{lemma:nagata-strict-transform-closure-characterization}
  \lean{AlgebraicGeometry.Hom.strictTransform_eq_schemeTheoreticClosure}

  Let $f\colon Y\to X$, let $\mathcal I$ be a finite-type quasi-coherent
  ideal on $X$, put $X'=\operatorname{Bl}_{\mathcal I}(X)$, and let
  $U=X\setminus V(\mathcal I)$.  The strict transform is the smallest
  closed subscheme of $Y\times_XX'$ through which
  \[
    Y\times_XU\longrightarrow Y\times_XX'
  \]
  factors.  In particular it is the scheme-theoretic closure of this open
  part.
\end{lemma}

\begin{proof}
  \uses{lemma:nagata-power-torsion-quotient-torsionfree,
    lemma:nagata-power-torsion-kernel-localization,
    prop:nagata-strict-transform-closed-fiber-product}
  On an $a$-chart put $C=C_a(A,B,I)$.  A closed subscheme containing the
  preserved open is given by a quotient $C\twoheadrightarrow C/K$ for which
  $K_a=0$.  By
  \Cref{lemma:nagata-power-torsion-kernel-localization}, this says exactly
  that $K\subseteq\Gamma_{(a)}(C)$.  Therefore the quotient
  $C\twoheadrightarrow C/\Gamma_{(a)}(C)$ factors through every such
  quotient, which is precisely the containment of the corresponding closed
  subschemes in the required direction.  It becomes an isomorphism after
  inverting $a$, again by the kernel description.  The quotient has no
  residual $a$-power torsion by
  \Cref{lemma:nagata-power-torsion-quotient-torsionfree}, so this local
  characterization is stable under further restriction.  The chart
  descriptions glue by
  \Cref{prop:nagata-strict-transform-closed-fiber-product}, proving the
  asserted minimality globally.
\end{proof}

\begin{lemma}[Flatness kills the strict-transform torsion]
  \label{lemma:nagata-flat-strict-transform-no-torsion}
  \uses{def:nagata-power-torsion}
  \lean{AlgebraicGeometry.strictTransform_powTorsion_eq_bot_of_flat}

  Let $A\to B$ be flat, let $I\subseteq A$ be an ideal, and let $a\in I$.
  Then the $a$-power-torsion in
  $B\otimes_AA[I/a]$ is zero.
\end{lemma}

\begin{proof}
  \uses{lemma:nagata-chart-generator-regular}
  Multiplication by $a$ is injective on $A[I/a]$ by
  \Cref{lemma:nagata-chart-generator-regular}.  Tensor this injection with the
  flat $A$-module $B$; multiplication by $a$ remains injective, and therefore
  no nonzero element is killed by a power of $a$.
\end{proof}

\begin{corollary}[Flat strict transform equals ordinary base change]
  \label{cor:nagata-flat-strict-transform-base-change}
  \lean{AlgebraicGeometry.Hom.strictTransformIsoBaseChange_of_flat}

  Let $f\colon Y\to X$ and let $\mathcal I$ be a finite-type quasi-coherent
  ideal on $X$.  If $f$ is flat, the strict transform in
  $\operatorname{Bl}_{\mathcal I}(X)$ is all of
  $Y\times_X\operatorname{Bl}_{\mathcal I}(X)$.
\end{corollary}

\begin{proof}
  \uses{lemma:nagata-flat-strict-transform-no-torsion,
    prop:nagata-strict-transform-closed-fiber-product}
  On every affine blow-up chart the quotient is by the zero ideal, so the
  closed immersion in
  \Cref{prop:nagata-strict-transform-closed-fiber-product} is an isomorphism.
\end{proof}

\begin{lemma}[Strict transform on the preserved open]
  \label{lemma:nagata-strict-transform-over-complement}
  \lean{AlgebraicGeometry.isPullback_strictTransformOpen}

  Let $f\colon Y\to X$, let $\mathcal I$ be a finite-type quasi-coherent
  ideal, and put $X'=\operatorname{Bl}_{\mathcal I}(X)$.  Over
  $U=X\setminus V(\mathcal I)$, the strict transform is canonically
  $Y\times_XU$, and its map to $Y$ is the ordinary pullback isomorphism.
\end{lemma}

\begin{proof}
  \uses{thm:nagata-global-blowup-complement,
    lemma:nagata-strict-transform-closure-characterization,
    prop:nagata-strict-transform-closed-fiber-product}
  The blow-up is the identity over $U$ by
  \Cref{thm:nagata-global-blowup-complement}.  Under this specified
  isomorphism the open whose closure defines the strict transform is all of
  $Y\times_XU$.  It is already a closed subscheme of itself, so the minimal
  closed-subscheme characterization of
  \Cref{lemma:nagata-strict-transform-closure-characterization} identifies
  the restriction with $Y\times_XU$.  On charts this is the quotient by the
  zero torsion ideal after $a$ is inverted; these identifications agree on
  overlaps by the closed chart model in
  \Cref{prop:nagata-strict-transform-closed-fiber-product}.  The projection
  equality is the one induced by the fibre-product square.
\end{proof}

\begin{lemma}[Strict transform under a dominating blow-up]
  \label{lemma:nagata-strict-transform-dominating}
  \lean{AlgebraicGeometry.Hom.strictTransform_comp_dominatingBlowup}

  Let $f\colon Y\to X$.  Suppose $X''\to X'\to X$ are blow-ups which are
  isomorphisms over the same open $U\subseteq X$.  The strict transform of
  $Y$ in $X''$ is canonically the strict
  transform, along $X''\to X'$, of the strict transform of $Y$ in $X'$.
\end{lemma}

\begin{proof}
  \uses{lemma:nagata-strict-transform-closure-characterization,
    prop:nagata-strict-transform-closed-fiber-product}
  Use the canonical associativity isomorphism
  $(Y\times_XX')\times_{X'}X''\cong Y\times_XX''$.  Under it, the preserved
  open used in the second strict transform is identified with
  $Y\times_XU$, because both blow-ups restrict to the identity over $U$.
  The second strict transform is closed in the left-hand fibre product by
  \Cref{prop:nagata-strict-transform-closed-fiber-product}; hence, after this
  identification, it is a closed subscheme of $Y\times_XX''$ containing that
  same open.  The fibre product of the first strict transform with $X''$ is a
  closed subscheme of $Y\times_XX''$ containing that open.  Taking the
  scheme-theoretic closure inside this closed subscheme therefore gives the
  same result as taking it in $Y\times_XX''$: the ambient closure is already
  contained in every closed subscheme containing the open.  The closure
  characterization of
  \Cref{lemma:nagata-strict-transform-closure-characterization} now gives the
  canonical identification.
\end{proof}

\begin{lemma}[Disjoint strict transforms stay disjoint]
  \label{lemma:nagata-strict-transform-disjoint-dominating}
  \lean{AlgebraicGeometry.StrictTransform.disjoint_of_dominates}

  Let $Z_1,Z_2\to X$ be morphisms and let $X''\to X'\to X$ be blow-ups
  which are isomorphisms over the same open $U\subseteq X$.  If the strict
  transforms of $Z_1$ and $Z_2$ in $X'$ are disjoint, their strict transforms
  after the further dominating blow-up $X''\to X'$ are disjoint.
\end{lemma}

\begin{proof}
  \uses{lemma:nagata-strict-transform-dominating}
  Each new strict transform is contained in the inverse image of the
  corresponding old strict transform.  Inverse images of disjoint closed
  subschemes are disjoint.
\end{proof}

\begin{lemma}[Composition of admissible blow-ups, {\cite[\href{https://stacks.math.columbia.edu/tag/080L}{Tag 080L}]{stacks-project}}]
  \label{lemma:nagata-admissible-blowup-comp}
  \uses{def:nagata-admissible-blowup}
  \lean{AlgebraicGeometry.AdmissibleBlowup.exists_comp}

  Let $X$ be quasi-compact and quasi-separated and let $U\subseteq X$ be
  quasi-compact.  A composition of two $U$-admissible blow-ups is a
  $U$-admissible blow-up.  The chosen isomorphism over $U$ is the composite of
  the chosen pullback isomorphisms.
\end{lemma}

\begin{proof}
  \uses{thm:nagata-composition-of-blowups,
    lemma:nagata-admissible-blowup-over-open}
  Apply \Cref{thm:nagata-composition-of-blowups}.  It writes the composite as
  the blow-up of a finite-presentation center $Y\hookrightarrow X$ whose
  support is the union of the first center and the image of the second center.
  Both sets are disjoint from $U$, because the first blow-up is the identity
  over $U$ by \Cref{lemma:nagata-admissible-blowup-over-open}, and the second
  center is disjoint from its inverse image.  Thus $Y\cap U=\varnothing$,
  equivalently its defining ideal restricts to the unit ideal on $U$, so the
  resulting blow-up is $U$-admissible.  The isomorphism supplied by the
  composition theorem is over $X$.  Pulling it back to $U$ identifies the
  composite with the composite of the two specified pullback isomorphisms;
  the equality follows after composing with the structure maps, where it is
  the associativity equality for fibre products.
\end{proof}

\begin{lemma}[Properness of a blow-up]
  \label{lemma:nagata-blowup-proper}
  \uses{def:nagata-global-blowup,
    thm:nagata-global-blowup-proper}
  \lean{AlgebraicGeometry.IsBlowup.isProper}

  Let $X$ be a scheme and let $\mathcal I$ be a finite-type quasi-coherent
  ideal on $X$.  Its blow-up is proper over $X$.  In particular,
  every admissible blow-up used in this chapter is proper.
\end{lemma}

\begin{proof}
  \uses{thm:nagata-global-blowup-proper}
  A finite type quasi-coherent ideal defines the global blow-up of
  \Cref{def:nagata-global-blowup}.  Its structure map is proper by
  \Cref{thm:nagata-global-blowup-proper}.  An admissible blow-up is a blow-up
  of precisely this form.
\end{proof}

\begin{lemma}[Strict transform over the complement of the center]
  \label{lemma:nagata-strict-transform-restrict}
  \uses{def:nagata-strict-transform}
  \lean{AlgebraicGeometry.StrictTransform.restrict_isIso_locus}

  Let $X'\to X$ be the blow-up of a center $C\hookrightarrow X$.  Over
  $X\setminus C$, the strict transform is canonically the ordinary pullback.
  This identification is compatible with the projection to the original
  scheme.
\end{lemma}

\begin{proof}
  \uses{lemma:nagata-strict-transform-over-complement}
  This is exactly
  \Cref{lemma:nagata-strict-transform-over-complement}, applied to the ideal
  defining $C$.  The compatibility with the projection is part of that
  canonical identification.
\end{proof}

\begin{lemma}[{\cite[\href{https://stacks.math.columbia.edu/tag/080E}{Tag 080E}]{stacks-project}}]
  \label{lemma:nagata-strict-transform-ambient-closed}
  \uses{def:nagata-strict-transform}
  \lean{AlgebraicGeometry.StrictTransform.of_closedImmersion}

  Let $Z\hookrightarrow X$ be a closed immersion and let $X'\to X$ be a
  blow-up.  Blowing up $Z$ in the pulled-back center gives the strict transform
  of $Z$ in $X'$.  Strict transforms computed inside $Z$ and inside $X$ agree
  under the induced closed immersion.
\end{lemma}

\begin{proof}
  \uses{thm:nagata-strict-transform-is-blowup,
    prop:nagata-strict-transform-chart-iso,
    prop:nagata-strict-transform-closed-fiber-product}
  By \Cref{thm:nagata-strict-transform-is-blowup}, the strict transform is the
  blow-up of $Z$ in the extended center ideal.  On an affine chart
  $X=\operatorname{Spec}(A)$ with $Z=\operatorname{Spec}(B)$, the induced
  map to $Z\times_XX'$ is the quotient map
  $B\otimes_AA[I/a]\twoheadrightarrow
  B\otimes_AA[I/a]/\Gamma_{(a)}$.  The target chart of the blow-up of $Z$ is
  identified with this quotient by
  \Cref{prop:nagata-strict-transform-chart-iso}.  It is therefore precisely
  the closed immersion of the closed chart model.  These identifications
  commute with localization and agree on overlaps by
  \Cref{prop:nagata-strict-transform-closed-fiber-product}; gluing proves
  equality of the two global closed subschemes, as well as the claimed
  compatibility with the induced closed immersion.
\end{proof}

\begin{lemma}[Gluing compatible finite-type ideals on two qc opens]
  \label{lemma:nagata-glue-finite-type-ideals}
  \lean{Scheme.QuasiCoherent.Ideal.glue_two_qcOpen}

  Let $U,W\subseteq X$ be quasi-compact opens of a qcqs scheme.  If
  $\mathcal I_U\subseteq\mathcal O_U$ and
  $\mathcal I_W\subseteq\mathcal O_W$ are finite type quasi-coherent ideals
  equipped with an equality of their restrictions to $U\cap W$, then there is
  a finite type quasi-coherent ideal
  $\mathcal I_{U\cup W}\subseteq\mathcal O_{U\cup W}$ restricting to the two
  given ideals.
\end{lemma}

\begin{proof}
  The specified equality is an isomorphism of ideal sheaves compatible with
  the inclusions into the structure sheaf.  The sheaf gluing theorem for the
  open cover $U\cup W=U\cup W$ therefore glues the ideals and inclusions to a
  subsheaf of $\mathcal O_{U\cup W}$.  Quasi-coherence, the ideal property,
  and finite type are local on this cover, so they hold for the glued ideal.
\end{proof}

\begin{definition}[Finite-type extension of an ideal from a quasi-compact open]
  \label{def:nagata-finite-type-ideal-extension}
  \uses{lemma:nagata-finite-type-ideal-extension-exists}
  \lean{Scheme.QuasiCoherent.finiteTypeIdealExtension}

  Let $X$ be quasi-compact and quasi-separated, let $W\subseteq X$ be a
  quasi-compact open, and let
  $\mathcal J\subseteq\mathcal O_W$ be a finite type quasi-coherent ideal.
  By classical choice from
  \Cref{lemma:nagata-finite-type-ideal-extension-exists}, select a finite type
  quasi-coherent ideal on $X$ restricting to $\mathcal J$, and define
  \[
    \operatorname{Ext}_{X,W}(\mathcal J)\subseteq\mathcal O_X
  \]
  to be this selected ideal.
\end{definition}

\begin{lemma}[Finite-type ideal extension,
  {\cite[\href{https://stacks.math.columbia.edu/tag/01PF}{Tag 01PF}]{stacks-project}}]
  \label{lemma:nagata-extend-finite-type-ideal}
  \uses{def:nagata-finite-type-ideal-extension}
  \lean{AlgebraicGeometry.Scheme.exists_finiteType_ideal_extension}

  The sheaf $\operatorname{Ext}_{X,W}(\mathcal J)$ is a finite type
  quasi-coherent ideal on $X$, and its restriction to $W$ is
  $\mathcal J$.
\end{lemma}

\begin{proof}
  \uses{lemma:nagata-finite-type-ideal-extension-exists}
  The selected ideal is required to have these two properties in
  \Cref{def:nagata-finite-type-ideal-extension}; the witness type is nonempty
  by \Cref{lemma:nagata-finite-type-ideal-extension-exists}.
\end{proof}

\begin{definition}[Extension construction for an admissible blow-up]
  \label{def:nagata-extended-admissible-blowup}
  \uses{def:nagata-admissible-blowup,
    def:nagata-finite-type-ideal-extension,
    def:nagata-global-blowup,
    lemma:nagata-glue-finite-type-ideals}
  \lean{AlgebraicGeometry.AdmissibleBlowup.extensionConstruction}

  Let $X$ be quasi-compact and quasi-separated and let $U,W\subseteq X$ be
  quasi-compact opens.  Write a given $(U\cap W)$-admissible blow-up of $W$
  as $\operatorname{Bl}_{\mathcal J}(W)\to W$.  Its center ideal restricts
  to the unit ideal on $U\cap W$, so glue $\mathcal J$ on $W$ with the unit
  ideal on $U$ to an ideal $\mathcal J^+$ on $U\cup W$ using
  \Cref{lemma:nagata-glue-finite-type-ideals}, and set
  \[
    \mathcal I_{X}(\mathcal J)
      =\operatorname{Ext}_{X,U\cup W}(\mathcal J^+),
  \]
  and define
  \[
    \mathsf{ExtBl}_{X,W,U}(\mathcal J)
      =\operatorname{Bl}_{\mathcal I_X(\mathcal J)}(X)\longrightarrow X.
  \]
\end{definition}

\begin{lemma}[Property of the extended admissible blow-up,
  {\cite[\href{https://stacks.math.columbia.edu/tag/080M}{Tag 080M}]{stacks-project}}]
  \label{lemma:nagata-extend-admissible-blowup}
  \uses{def:nagata-extended-admissible-blowup}
  \lean{AlgebraicGeometry.AdmissibleBlowup.extend_from_open}

  The constructed morphism
  $\mathsf{ExtBl}_{X,W,U}(\mathcal J)\to X$ is $U$-admissible, and its
  restriction to $W$ is canonically the given blow-up
  $\operatorname{Bl}_{\mathcal J}(W)\to W$.
\end{lemma}

\begin{proof}
  \uses{lemma:nagata-extend-finite-type-ideal,
    thm:nagata-blowup-flat-base-change}
  By \Cref{lemma:nagata-extend-finite-type-ideal}, the ideal
  $\mathcal I_X(\mathcal J)$ is finite type, restricts to $\mathcal J$ on
  $W$, and restricts to the unit ideal on $U$.  Its blow-up is therefore
  $U$-admissible.  The open immersion $W\to X$ is flat, so
  \Cref{thm:nagata-blowup-flat-base-change} identifies the pullback of the
  constructed blow-up with $\operatorname{Bl}_{\mathcal J}(W)$.  On an
  affine chart, that comparison is induced degreewise by the canonical map
  from the pullback of the Rees algebra to the Rees algebra of the restricted
  ideal; it sends every standard chart generator to itself.  Hence the local
  comparisons commute with restriction to basic overlaps and glue to the
  stated canonical isomorphism.  Over $U\cap W$ both ideals are the unit
  ideal, so the same formula is the selected identity of the two trivial
  blow-ups.  This also verifies compatibility with the chosen admissibility
  isomorphism.
\end{proof}

\begin{lemma}[{\cite[\href{https://stacks.math.columbia.edu/tag/080A}{Tag 080A}]{stacks-project}}]
  \label{lemma:nagata-product-ideal-common-blowup}
  \lean{AlgebraicGeometry.Blowup.productIdeal_dominates}

  Let $X$ be qcqs, let $U\subseteq X$ be a quasi-compact open, and let
  $\mathcal I_1,\mathcal I_2$ be finite-type quasi-coherent ideals on $X$ with
  $\mathcal I_1|_U=\mathcal I_2|_U=\mathcal O_U$.  The blow-up of
  $\mathcal I_1\mathcal I_2$ canonically dominates the blow-ups of
  $\mathcal I_1$ and of $\mathcal I_2$, and all three blow-ups are
  $U$-admissible.
\end{lemma}

\begin{proof}
  \uses{cor:nagata-blowup-product-common-refinement,
    thm:nagata-global-blowup-complement}
  Let $\pi_i\colon\operatorname{Bl}_{\mathcal I_1\mathcal I_2}(X)
  \to\operatorname{Bl}_{\mathcal I_i}(X)$ be the two maps supplied by
  \Cref{cor:nagata-blowup-product-common-refinement}.  The product ideal is
  also the unit ideal on $U$, so each of the three blow-ups restricts to $U$
  by \Cref{thm:nagata-global-blowup-complement}; hence all are
  $U$-admissible and the displayed maps are the required domination maps.
\end{proof}

\begin{lemma}[Existence of an ideal defining the complement of a qc open]
  \label{lemma:nagata-qc-open-complement-ideal-exists}
  \uses{lemma:nagata-finite-type-ideal-extension-exists}
  \lean{AlgebraicGeometry.Scheme.exists_finiteType_ideal_coe_support_eq_compl}

  Let $W\subseteq Y$ be a quasi-compact open in a qcqs scheme.  There exists
  a finite type quasi-coherent ideal $\mathcal I\subseteq\mathcal O_Y$ such
  that $V(\mathcal I)=Y\setminus W$.
\end{lemma}

\begin{proof}
  Let $\mathcal K\subseteq\mathcal O_Y$ be the quasi-coherent ideal defining
  the reduced induced closed subscheme on $Y\setminus W$.  Its restriction to
  $W$ is the unit ideal.  Apply the general submodule form of
  \Cref{lemma:nagata-finite-type-ideal-extension-exists} with ambient module
  $\mathcal K$ and submodule $\mathcal O_W=\mathcal K|_W$.  It gives a finite
  type quasi-coherent submodule $\mathcal I\subseteq\mathcal K$ with
  $\mathcal I|_W=\mathcal O_W$.  As an $\mathcal O_Y$-submodule of the ideal
  $\mathcal K$, the sheaf $\mathcal I$ is itself an ideal.
  The inclusion $\mathcal I\subseteq\mathcal K$ gives
  $Y\setminus W=V(\mathcal K)\subseteq V(\mathcal I)$, while
  $\mathcal I|_W=\mathcal O_W$ gives the reverse containment.  Hence
  $V(\mathcal I)=Y\setminus W$.
\end{proof}

\begin{lemma}[Preparation of a Cartier boundary,
  {\cite[\href{https://stacks.math.columbia.edu/tag/080Y}{Tag 080Y}]{stacks-project}}]
  \label{lemma:nagata-cartier-boundary-preparation}
  \uses{def:nagata-admissible-blowup,
    def:nagata-effective-cartier-ideal,
    lemma:nagata-qc-open-complement-ideal-exists}
  \lean{AlgebraicGeometry.exists_admissibleBlowup_cartierBoundary}

  Let $X$ be qcqs and let $U\subseteq X$ be a quasi-compact open.  There is a
  $U$-admissible blow-up $X_1\to X$ for which the complement of the inverse
  image of $U$ is the support of an effective Cartier ideal.
\end{lemma}

\begin{proof}
  \uses{lemma:nagata-qc-open-complement-ideal-exists,
    thm:nagata-blowup-principalizes-center,
    thm:nagata-global-blowup-complement}
  Choose the finite type quasi-coherent ideal $\mathcal I$ with
  $V(\mathcal I)=X\setminus U$ supplied by
  \Cref{lemma:nagata-qc-open-complement-ideal-exists}.  Its restriction to
  $U$ is the unit ideal, so its blow-up is $U$-admissible.  Moreover,
  \Cref{thm:nagata-blowup-principalizes-center} makes the extended ideal an
  effective Cartier ideal.  Its vanishing locus is the inverse image of
  $V(\mathcal I)$, which is the complement of the inverse image of $U$;
  the latter inverse image is the preserved open by
  \Cref{thm:nagata-global-blowup-complement}.  Hence this Cartier ideal has
  precisely the required boundary support.
\end{proof}

\begin{lemma}[Scheme-theoretic density survives flat pullback]
  \label{lemma:nagata-dense-flat-pullback}
  \uses{def:nagata-scheme-theoretic-image}
  \lean{AlgebraicGeometry.IsSchemeTheoreticallyDominant.of_isPullback}

  Let \(h\colon U\to X\) be quasi-compact and quasi-separated with
  scheme-theoretic image \(X\), and let
  \(q\colon P\to X\) be flat.  Then
  \[
    U\times_XP\longrightarrow P
  \]
  has scheme-theoretic image \(P\).
\end{lemma}

\begin{proof}
  \uses{lemma:nagata-scheme-image-flat-base-change}
  By \Cref{lemma:nagata-scheme-image-flat-base-change}, the scheme-theoretic
  image after pullback is canonically the pullback of the scheme-theoretic
  image of $h$.  By hypothesis that image is $X$, and its pullback along
  $P\to X$ is the identity closed subscheme of $P$.  Thus the base-changed
  morphism is scheme-theoretically dense in $P$.
\end{proof}

\begin{lemma}[Strict transform preserves universally exact sequences,
  {\cite[\href{https://stacks.math.columbia.edu/tag/080W}{Tag 080W}]{stacks-project}}]
  \label{lemma:nagata-module-strict-transform-exact}
  \uses{def:nagata-module-strict-transform,
    lemma:nagata-module-strict-transform-chart,
    lemma:nagata-power-torsion-localization}
  \lean{QuasiCoherent.Sheaf.strictTransform_exact_of_universallyExact}

  Let $V\to S$ be a morphism and let
  $0\to\mathcal M_1\to\mathcal M_2\to\mathcal M_3\to0$ be a sequence of
  quasi-coherent $\mathcal O_V$-modules which remains exact after every base
  change $T\to S$.  For a specified blow-up $S'\to S$, the sequence of
  module strict transforms
  on $V\times_SS'$,
  $0\to\mathcal M'_1\to\mathcal M'_2\to\mathcal M'_3\to0$, is exact.
\end{lemma}

\begin{proof}
  The claim is local on source and base.  On an affine blow-up chart, let $a$
  be the local equation of the exceptional divisor and let
  $0\to M_1\to M_2\to M_3\to0$ be the resulting universally exact sequence.
  Surjectivity after quotienting by $a$-power torsion follows by taking a
  preimage in $M_2$ and then passing to the quotient; injectivity at $M_1$
  follows because an element of $M_1$ killed in the quotient is already
  $a$-power torsion.  If $x\in M_2$ maps to $a$-power torsion in $M_3$,
  then $a^nx$ is the image of some $y\in M_1$.  The class of $y$ maps to zero
  in $M_2/a^nM_2$; universal injectivity makes
  $M_1/a^nM_1\to M_2/a^nM_2$ injective, so $y=a^nz$ for some $z\in M_1$.
  Hence $a^n(x-z)=0$ in $M_2$.  Thus the class of $x$ comes from the class of
  $z$, proving exactness in the middle.  The affine chart descriptions glue
  by localization compatibility of module strict transform.
\end{proof}

\begin{remark}[Formalization route: the flattening chain is not on the
  formalized path]
  The revised formalization replaces the module-flattening route below
  (\Cref{lemma:nagata-finite-data-flattening-package},
  \Cref{thm:nagata-module-flattening-by-blowup},
  \Cref{lemma:nagata-flattening-after-cartier-boundary},
  \Cref{lemma:nagata-flattening-ideal-exists},
  \Cref{lemma:nagata-dense-open-dimension-package}; Stacks Tags
  0815/081R/081P/081L) by the boundary-separation route: the two-open
  compactification theorem uses the separation blow-up of
  \Cref{lemma:nagata-separate-closed-subschemes} (Tag 0F3V), common
  modifications, and the dense-source valuative criterion
  (\Cref{lemma:nagata-dense-source-valuative}).  In the present document the
  flattening chain survives only as the input to
  \Cref{lemma:nagata-elimination-indeterminacy} (Tag 081S), whose formal
  proof route is still open; that lemma is currently stated with
  \texttt{sorry} in Lean.
\end{remark}

\begin{lemma}[Finite-data flattening induction package]
  \label{lemma:nagata-finite-data-flattening-package}
  \uses{def:nagata-admissible-blowup,
    def:nagata-module-strict-transform,
    lemma:nagata-module-strict-transform-exact}
  \lean{QuasiCoherent.Sheaf.finiteDataFlatteningPackage}

  Let $S$ be qcqs, $V\to S$ quasi-compact and locally of finite
  presentation, and let $\mathcal M$ be finite type.  Suppose that on a
  quasi-compact open $U$ the module is finitely presented and flat.  There
  is a finite sequence of $U$-admissible blow-ups whose composite has strict
  transform finitely presented and flat over $S$.
\end{lemma}

\begin{proof}
  This is the finite-data form of the flattening theorem in Stacks Project,
  Tag~0815.  Instantiate that theorem with the displayed qcqs base $S$, the
  quasi-compact locally finitely presented morphism $V\to S$, the
  quasi-compact open $U\subseteq S$, and the finite-type module
  $\mathcal M$.  Its hypothesis that $\mathcal M|_U$ is finitely presented and
  flat is exactly the hypothesis in this block.  The theorem supplies a
  $U$-admissible blow-up $S'\to S$ and says that the strict transform on
  $V\times_SS'$ is finitely presented and flat.  Taking this single blow-up as
  a sequence of length one gives the stated finite-data package.  The affine
  chart and overlap descriptions required by the blueprint are the
  presentations used in the cited theorem and are identified with the
  strict-transform charts by
  \Cref{lemma:nagata-module-strict-transform-exact}.
\end{proof}

\begin{theorem}[Flattening a finite-type module by an admissible blow-up,
  {\cite[\href{https://stacks.math.columbia.edu/tag/0815}{Tag 0815}]{stacks-project}}]
  \label{thm:nagata-module-flattening-by-blowup}
  \uses{def:nagata-admissible-blowup,
    def:nagata-module-strict-transform,
    lemma:nagata-module-strict-transform-exact,
    lemma:nagata-admissible-blowup-comp,
    lemma:nagata-finite-data-flattening-package}
  \lean{QuasiCoherent.Sheaf.exists_admissibleBlowup_flat_strictTransform}

  Let $S$ be qcqs, let $V\to S$ be quasi-compact and locally of finite
  presentation, let $U\subseteq S$ be a quasi-compact open, and let
  $\mathcal M$ be a finite type quasi-coherent $\mathcal O_V$-module whose
  restriction over $U$ is finitely presented and flat over $U$.  There is a
  $U$-admissible blow-up $S'\to S$ such that the module strict transform
  $\mathcal M'$ on $V\times_SS'$ is finitely presented and flat over $S'$.
\end{theorem}

\begin{proof}
  Apply \Cref{lemma:nagata-finite-data-flattening-package} to the finite affine
  presentations of $\mathcal M$.  Its finite composite of admissible blow-ups
  is a single $U$-admissible blow-up by
  \Cref{lemma:nagata-admissible-blowup-comp}.  The package proves on each
  chart that the strict transform is finitely presented and flat, and the
  overlap equations make these local transforms glue.
\end{proof}

\begin{lemma}[Module and scheme strict transforms agree for a closed immersion]
  \label{lemma:nagata-strict-transform-closed-module-comparison}
  \uses{lemma:nagata-module-strict-transform-chart,
    prop:nagata-strict-transform-closed-fiber-product,
    lemma:nagata-strict-transform-closure-characterization}
  \lean{AlgebraicGeometry.StrictTransform.closedImmersion_moduleComparison}

  Let $i\colon Z\hookrightarrow V$ be a finite-presentation closed immersion
  and let $V'\to V$ be a blow-up with center ideal $\mathcal I$.  On every
  affine blow-up chart, the coordinate algebra of the scheme-theoretic strict
  transform of $Z$ is the strict-transform quotient of the pulled-back
  coordinate module of $Z$.  These identifications are compatible with
  localization and hence identify the resulting quasi-coherent module with
  the structure sheaf of the strict transform.
\end{lemma}

\begin{proof}
  \uses{lemma:nagata-module-power-torsion-localization}
  Work on an affine chart $V=\operatorname{Spec}(A)$, write
  $Z=\operatorname{Spec}(B)$ with $B=A/J$, and let
  $R=A[I/a]$ be a blow-up chart.  The ordinary base-change chart has ring
  $B\otimes_A R$, while the scheme strict-transform chart has ring
  $(B\otimes_A R)/\Gamma_{(a)}(B\otimes_A R)$ by
  \Cref{prop:nagata-strict-transform-closed-fiber-product} and
  \Cref{lemma:nagata-strict-transform-closure-characterization}.  The
  module strict-transform chart for the $A$-module $B$ is the same quotient
  by \Cref{lemma:nagata-module-strict-transform-chart}; the $B$-algebra
  structure is the one induced from the ordinary base-change chart.  On a
  basic overlap, both descriptions localize to the corresponding quotient,
  and the quotient-localization cocycle is
  \Cref{lemma:nagata-module-power-torsion-localization}.  Thus the affine
  identifications glue and identify the module's associated sheaf with the
  structure sheaf of the scheme strict transform.
\end{proof}

\begin{lemma}[Flattening after preparation of the boundary,
  {\cite[\href{https://stacks.math.columbia.edu/tag/081R}{Tag 081R}]{stacks-project}}]
  \label{lemma:nagata-flattening-after-cartier-boundary}
  \uses{thm:nagata-module-flattening-by-blowup,
    lemma:nagata-admissible-blowup-comp,
    lemma:nagata-strict-transform-closed-module-comparison,
    lemma:nagata-strict-transform-dominating,
    lemma:nagata-dense-flat-pullback}
  \lean{AlgebraicGeometry.exists_admissibleBlowup_flatStrictTransform}

  Let $q\colon Z\to X$ be separated and of finite type, with $X$ qcqs, and
  let $U\subseteq X$ be a quasi-compact open over which $q$ is an
  isomorphism.  If $X\setminus U$ supports an effective
  Cartier ideal, there is a $U$-admissible blow-up $X_2\to X$ such that the
  strict transform $Z_2\to X_2$ is flat and of finite presentation.
\end{lemma}

\begin{proof}
  Choose a finite affine cover $Z=\bigcup_j Z_j$.  After restricting each
  $Z_j$ over a suitable affine open of $X$, its finite-type separated map to
  $X$ factors as a closed immersion
  $i_j\colon Z_j\hookrightarrow V_j$ followed by a quasi-compact open
  $V_j\subseteq\mathbb A_X^{n_j}$.  The closed immersion need not be of finite
  presentation; what is needed is that $(i_j)_*\mathcal O_{Z_j}$ is a finite
  type $\mathcal O_{V_j}$-module.  Over $U$ it is finitely presented and flat
  over $U$, because $Z_U\cong U$ and the corresponding section of the
  separated, locally finitely presented morphism $V_{j,U}\to U$ is a
  finite-presentation closed immersion.

  Apply \Cref{thm:nagata-module-flattening-by-blowup} to these finitely many
  modules.  Successively pull back the resulting admissible blow-ups and use
  \Cref{lemma:nagata-admissible-blowup-comp} to express the composite as one
  $U$-admissible blow-up.  The comparison
  \Cref{lemma:nagata-strict-transform-closed-module-comparison} identifies
  each transformed module with the pushforward of the structure sheaf of the
  corresponding scheme strict transform.  It follows that every transformed
  $Z_j$ is flat and locally of finite presentation over the common blow-up.
  The strict transform is quasi-compact and quasi-separated and is covered by
  these finitely many pieces, so its morphism is of finite presentation.
  Finally \Cref{lemma:nagata-strict-transform-dominating} identifies this
  iterated transform with the direct strict transform for the composite.
\end{proof}

\begin{lemma}[Dense-open dimension and finite-presentation package]
  \label{lemma:nagata-dense-open-dimension-package}
  \uses{def:nagata-scheme-theoretically-dense-open}
  \lean{Scheme.Hom.flat_denseOpen_dimensionPackage}

  Let $q\colon Z\to X$ be flat and locally of finite type, and let $U\subseteq X$
  be retrocompact and scheme-theoretically dense.  If $q_U$ is locally of
  finite presentation and has fibre dimension at most $e$, then $q$ is locally
  of finite presentation and every fibre of $q$ has dimension at most $e$.
\end{lemma}

\begin{proof}
  \uses{def:nagata-scheme-theoretically-dense-open}
  Apply Stacks Project, Tag~081P (Lemma~38.11.2), to $q$ and the
  retrocompact scheme-theoretically dense open $U$.  The hypotheses are the
  flatness and locally-finite-type assumptions in the statement together with
  the assumed local finite presentation of $q_U$; the conclusion is that $q$
  is locally of finite presentation.  With this conclusion established,
  apply Tag~081L (Lemma~38.11.3) to the same $q$, $U$, and the integer $e$.
  Its relative-dimension hypothesis is precisely the stated bound on
  $q_U$, and it concludes that every fibre of $q$ has dimension at most $e$.
  Both conclusions are local on source and target, so the affine instances
  supplied by the two cited lemmas glue to the assertions above.
\end{proof}

\begin{lemma}[Relative dimension extends from a retrocompact dense open,
  {\cite[\href{https://stacks.math.columbia.edu/tag/081L}{Tag 081L}]{stacks-project}}]
  \label{lemma:nagata-flat-dense-open-relative-dimension}
  \uses{lemma:nagata-dense-open-dimension-package}
  \lean{Scheme.Hom.relativeDimension_le_of_retrocompact_dense_open}

  Let $q\colon Z\to X$ be flat and locally of finite type, and let
  $U\subseteq X$ be a retrocompact dense open.  If $q^{-1}(U)\to U$ has
  relative dimension at most $e$, then every fibre of $q$ has dimension at
  most $e$.
\end{lemma}

\begin{proof}
  Apply \Cref{lemma:nagata-dense-open-dimension-package} on affine charts of
  $X$.  Its finite-presentation upgrade and generic-point dimension argument
  give the asserted bound, and the conclusions glue since they are local on
  source and target.
\end{proof}

\begin{lemma}[Standard-etale finite-part package]
  \label{lemma:nagata-standard-etale-finite-part-package}
  \lean{Scheme.Hom.standardEtale_finitePartPackage}

  If $q\colon Z\to X$ is separated and quasi-finite and $x\in X$, there are
  finitely many points $z_i\in Z_x$ and an elementary etale neighbourhood
  $(T,t)\to(X,x)$ on which pairwise disjoint standard-etale neighbourhoods of
  the $z_i$ glue to an open-and-closed finite part $V\subseteq Z_T$; the
  complement has empty fibre over $t$.
\end{lemma}

\begin{proof}
  Quasi-finiteness makes $Z_x$ finite.  For each $z_i$, choose an affine
  neighbourhood and an element whose image in the fibre cuts out only $z_i$.
  The corresponding standard-etale algebra is finite over a neighbourhood of
  $x$ and etale at $z_i$.  Shrinking the base and taking the fibre product of
  these finitely many neighbourhoods makes them pairwise disjoint.  Their
  union is open-and-closed and finite over the resulting elementary etale
  neighbourhood; separatedness identifies its complement, which misses $t$.
\end{proof}

\begin{lemma}[Etale finite part of a quasi-finite separated morphism,
  {\cite[\href{https://stacks.math.columbia.edu/tag/02LP}{Tag 02LP}]{stacks-project}}]
  \label{lemma:nagata-quasifinite-etale-finite-part}
  \uses{lemma:nagata-standard-etale-finite-part-package}
  \lean{Scheme.Hom.exists_etaleNeighborhood_finitePart}

  Let $q\colon Z\to X$ be quasi-finite and separated and let $x\in X$.
  There is an elementary etale neighbourhood $(T,t)\to(X,x)$ and a
  decomposition $Z\times_XT=V\amalg W$ into open-and-closed subschemes such
  that $V\to T$ is finite and $W_t=\varnothing$.
\end{lemma}

\begin{proof}
  Apply \Cref{lemma:nagata-standard-etale-finite-part-package}; its $T,t,V$
  and open-and-closed complement $W$ give the required decomposition.
\end{proof}

\begin{lemma}[Affine morphisms descend fpqc-locally]
  \label{lemma:nagata-affine-fpqc-descent-package}
  \lean{Scheme.Hom.affine_of_fpqc_baseChange_affine}

  Let $p\colon X'\to X$ be faithfully flat and quasi-compact.  If
  $f\colon Y\to X$ is quasi-compact and $Y\times_XX'\to X'$ is affine, then
  $f$ is affine.
\end{lemma}

\begin{proof}
  Work over an affine $X_0=\operatorname{Spec}(A)$ and choose a finite affine
  cover $X'_0=\bigcup_i\operatorname{Spec}(A_i)$.  Write the affine pullbacks
  as $\operatorname{Spec}(B_i)$.  On overlaps the $B_i$ carry the canonical
  descent identifications.  Set
  $B=\ker(\prod_iB_i\rightrightarrows\prod_{i,j}B_i\otimes_AA_{ij})$.
  Faithful flatness makes tensoring this equalizer with each $A_i$ exact and
  identifies $B\otimes_AA_i$ with $B_i$.  Hence $\operatorname{Spec}(B)$
  pulls back to the given affine cover and glues to $Y\times_XX'_0$; the
  comparison with $Y$ is an isomorphism after the faithfully flat cover and
  therefore an isomorphism.  Thus $Y$ is affine over $X_0$.
\end{proof}

\begin{lemma}[Finiteness descends from an etale neighbourhood,
  {\cite[\href{https://stacks.math.columbia.edu/tag/02LA}{Tag 02LA}]{stacks-project}}]
  \label{lemma:nagata-finite-etale-local-on-target}
  \uses{lemma:nagata-affine-fpqc-descent-package}
  \lean{Scheme.Hom.finite_etaleLocalOnTarget}

  Finiteness of a morphism of schemes is local on the target for the etale
  topology.  In particular, if a morphism becomes finite after a surjective
  etale base change over a neighbourhood of a point, it is finite after
  shrinking that neighbourhood.
\end{lemma}

\begin{proof}
  After shrinking, take a quasi-compact surjective etale neighbourhood; it is
  faithfully flat and quasi-compact.  By
  \Cref{lemma:nagata-affine-fpqc-descent-package}, the original morphism is
  affine.  On an affine target write it as $\operatorname{Spec}(B)$ over
  $A$.  The base-change algebra is finite over $A'$, and faithful-flat descent
  of finite generators gives that $B$ is finite over $A$.  Hence the morphism
  is finite on the chosen neighbourhood.
\end{proof}

\begin{lemma}[Finite flat degree one is an isomorphism]
  \label{lemma:nagata-finite-flat-degree-one-iso}
  \lean{Scheme.Hom.isIso_of_finiteFlat_degree_one}

  A finite flat morphism of finite presentation whose finite locally free
  degree is one is an isomorphism.
\end{lemma}

\begin{proof}
  Work affine locally, with corresponding finite locally free algebra $B$ over
  $A$.  On every residue field $\kappa(\mathfrak p)$ the fibre
  $B\otimes_A\kappa(\mathfrak p)$ has dimension one.  Its unit is nonzero
  (otherwise the fibre algebra would be zero), hence it is a basis.  Nakayama's
  lemma therefore shows that $1\in B$ generates $B$ as an $A$-module in a
  neighbourhood of every prime.  The map $A\to B$, $a\mapsto a\,1$, is
  consequently a surjective map between locally free rank-one modules and is
  an isomorphism.  These local algebra isomorphisms are compatible on
  overlaps, so the original finite flat morphism is an isomorphism.
\end{proof}

\begin{lemma}[A flat strict transform of a partial isomorphism is open,
  {\cite[\href{https://stacks.math.columbia.edu/tag/081M}{Tag 081M}]{stacks-project}}]
  \label{lemma:nagata-flat-strict-transform-open}
  \uses{lemma:nagata-flat-dense-open-relative-dimension,
    lemma:nagata-quasifinite-etale-finite-part,
    lemma:nagata-finite-etale-local-on-target,
    lemma:nagata-finite-flat-degree-one-iso,
    lemma:nagata-dense-flat-pullback}
  \lean{AlgebraicGeometry.flatStrictTransform_isOpenImmersion}

  Let $q\colon Z\to X$ be separated, flat, and locally of finite type.  If
  $U\subseteq X$ is a retrocompact (equivalently, $U\to X$ is
  quasi-compact) scheme-theoretically dense open and
  $q^{-1}(U)\to U$ is an isomorphism, then $q$ is an open immersion.
\end{lemma}

\begin{proof}
  The image of $q$ is open because $q$ is flat and locally of finite
  presentation (the latter follows from the retrocompact dense-open extension
  input in \Cref{lemma:nagata-flat-dense-open-relative-dimension}); replace
  $X$ by this image.  The assertion is local on an affine open
  $X_0\subseteq X$.  Include the quasi-compact open
  $q^{-1}(U\cap X_0)\cong U\cap X_0$ and finitely many affine neighbourhoods
  of points over $X_0$; their union is a quasi-compact open of $Z$ whose image
  contains $X_0$, so we may reduce to a quasi-compact source without claiming
  that the full inverse image of $X_0$ is quasi-compact.  The
  relative-dimension lemma with
  $e=0$ gives zero-dimensional fibres, hence $q$ is quasi-finite.  At a point
  $x\in X_0$, apply the etale finite-part lemma.  The complementary summand
  $W$ has empty fibre at the chosen etale point; after shrinking, flatness and
  scheme-theoretic density of $U$ (preserved by flat pullback) force $W=\varnothing$.
  Thus the base-changed morphism is finite.  Finiteness descends etale-locally,
  so $q$ is finite near $x$.  Flatness makes it finite locally free.  Its
  rank is one on the dense open $U$; rank is locally constant for a finite
  locally free module, and a component disjoint from $U$ would contradict
  density.  Hence the rank is one everywhere near $x$, and
  \Cref{lemma:nagata-finite-flat-degree-one-iso} makes $q$ an isomorphism on
  that neighbourhood.  These neighbourhoods give an open immersion.
\end{proof}

\begin{lemma}[Flattening ideal with an open strict transform,
  {\cite[\href{https://stacks.math.columbia.edu/tag/081S}{Tag 081S}]{stacks-project}}]
  \label{lemma:nagata-flattening-ideal-exists}
  \uses{def:nagata-strict-transform,
    lemma:nagata-cartier-boundary-preparation,
    lemma:nagata-flattening-after-cartier-boundary,
    lemma:nagata-flat-strict-transform-open,
    lemma:nagata-admissible-blowup-preserves-density}
  \lean{AlgebraicGeometry.exists_flatteningIdeal_openStrictTransform}

  Let $q\colon Z\to X$ be separated and of finite type, let $X$ be qcqs,
  and let $U\subseteq X$ be a quasi-compact open.  Fix an isomorphism
  $e_U\colon q^{-1}(U)\xrightarrow{\sim}U$ over $U$.  There exists a finite type
  quasi-coherent ideal $\mathcal F\subseteq\mathcal O_X$ satisfying
  $\mathcal F|_U=\mathcal O_U$ such that, after blowing up $\mathcal F$, the
  strict-transform projection
  \[
    \operatorname{ST}_{\operatorname{Bl}_{\mathcal F}(X)/X}(Z)
      \longrightarrow \operatorname{Bl}_{\mathcal F}(X)
  \]
  is an open immersion and restricts over $U$ to $e_U$.
\end{lemma}

\begin{proof}
  \uses{lemma:nagata-admissible-blowup-comp,
    lemma:nagata-strict-transform-dominating}
  Apply \Cref{lemma:nagata-cartier-boundary-preparation} to obtain a
  $U$-admissible blow-up $X_1\to X$ with Cartier boundary.  The transported
  isomorphism over $U$ identifies the strict transform $Z_1\to X_1$ with the
  original partial isomorphism.  Apply
  \Cref{lemma:nagata-flattening-after-cartier-boundary} to get a second
  $U$-admissible blow-up $X_2\to X_1$ for which the strict transform
  $Z_2\to X_2$ is flat and finitely presented.  The preserved open is dense
  after both steps by
  \Cref{lemma:nagata-admissible-blowup-preserves-density}, and strict
  transforms identify with the direct strict transform by
  \Cref{lemma:nagata-strict-transform-dominating}.  The flat-partial-isomorphism
  lemma makes $Z_2\to X_2$ an open immersion.  Finally,
  \Cref{lemma:nagata-admissible-blowup-comp} identifies the composite with the
  blow-up of a single finite-type ideal $\mathcal F$ that is the unit on $U$;
  transporting the open immersion and its restriction gives the required
  witness.
\end{proof}

\begin{definition}[Selected flattening ideal]
  \label{def:nagata-flattening-ideal}
  \uses{lemma:nagata-flattening-ideal-exists}
  \lean{AlgebraicGeometry.flatteningIdeal}

  Under the hypotheses of
  \Cref{lemma:nagata-flattening-ideal-exists}, choose one witness and define
  $\mathcal F_U(q)$ to be its ideal component.  The accompanying finite-type,
  unit-on-$U$, and strict-transform fields are not part of this definition;
  they are recorded in the specification lemma below.
\end{definition}

\begin{lemma}[Properties of the selected flattening ideal]
  \label{lemma:nagata-flattening-ideal}
  \uses{def:nagata-flattening-ideal}
  \lean{AlgebraicGeometry.flatteningIdeal_spec}

  The ideal $\mathcal F_U(q)$ is finite type and restricts to the unit ideal
  on $U$; the strict-transform projection after blowing it up is an open
  immersion and is the original isomorphism over $U$.
\end{lemma}

\begin{proof}
  \uses{lemma:nagata-flattening-ideal-exists}
  The witness type is nonempty by
  \Cref{lemma:nagata-flattening-ideal-exists}.  Projecting its specification
  fields along the choice made in
  \Cref{def:nagata-flattening-ideal} gives finite type, unit restriction, the
  open strict-transform projection, and the specified identification over $U$.
\end{proof}

\begin{definition}[Elimination-of-indeterminacy construction]
  \label{def:nagata-elimination-indeterminacy}
  \uses{def:nagata-admissible-blowup,def:nagata-strict-transform,
    def:nagata-flattening-ideal}
  \lean{AlgebraicGeometry.eliminationIndeterminacyConstruction}

  Let $q\colon Z\to X$ be separated and of finite type, let $X$ be qcqs, and
  let $U\subseteq X$ be quasi-compact.  Fix an isomorphism
  $e_U\colon q^{-1}(U)\xrightarrow{\sim}U$ over $U$.
  Let $\mathcal F_U(q)$ be the selected flattening ideal of
  \Cref{def:nagata-flattening-ideal}.  Define
  \[
    X^{\mathrm{el}}(q,U)=\operatorname{Bl}_{\mathcal F_U(q)}(X),
    \qquad
    Z^{\mathrm{el}}(q,U)=
      \operatorname{ST}_{X^{\mathrm{el}}/X}(Z),
  \]
  together with the strict-transform projection
  $q^{\mathrm{el}}\colon Z^{\mathrm{el}}(q,U)\to X^{\mathrm{el}}(q,U)$ and
  the transported open identification induced by $e_U$.
  Thus $\mathcal F_U(q)$ is a labelled selected object rather than an
  unspecified existential variable.
\end{definition}

\begin{lemma}[Elimination of indeterminacy,
  {\cite[\href{https://stacks.math.columbia.edu/tag/081S}{Tag 081S}]{stacks-project}}]
  \label{lemma:nagata-elimination-indeterminacy}
  \uses{def:nagata-elimination-indeterminacy}
  \lean{AlgebraicGeometry.exists_admissibleBlowup_strictTransform_open}

  The morphism $X^{\mathrm{el}}(q,U)\to X$ is $U$-admissible, and the
  constructed strict-transform projection
  $q^{\mathrm{el}}\colon Z^{\mathrm{el}}(q,U)\to X^{\mathrm{el}}(q,U)$ is
  an open immersion.  Its restriction over $U$ is the original chosen
  isomorphism $q^{-1}(U)\cong U$.
\end{lemma}

\begin{proof}
  \uses{lemma:nagata-flattening-ideal}
  By \Cref{lemma:nagata-flattening-ideal}, the selected ideal is finite type,
  is the unit ideal on $U$, and has open strict-transform projection.  The
  first two properties make its blow-up $U$-admissible; the last property and
  the selected comparison over $U$ identifies the restriction of
  $q^{\mathrm{el}}$ with $e_U$, so the named construction has exactly the
  open-immersion and prescribed-over-$U$ fields asserted above.
\end{proof}

\begin{lemma}[A proper dense open immersion is an isomorphism]
  \label{lemma:nagata-proper-open-dense-iso}
  \lean{Scheme.Hom.isIso_of_isProper_isOpenImmersion_dense}

  Let $q\colon Z\to X$ be both proper and an open immersion.  If its image
  contains a scheme-theoretically dense open subscheme of $X$, then $q$ is an
  isomorphism.
\end{lemma}

\begin{proof}
  A proper immersion is a closed immersion.  Thus the image is both open and
  closed.  Its closed complement has structure sheaf supported away from the
  dense open; scheme-theoretic density forces that complement to be empty.
\end{proof}

\section{Separating boundary pieces}

\begin{lemma}[Finite-presentation closed immersions and finite type ideals]
  \label{lemma:nagata-closed-fp-ideal-finite-type}
  \lean{Scheme.Hom.closedImmersion_finitePresentation_iff_ideal_finiteType}

  A closed immersion $Z\hookrightarrow X$ is of finite presentation if and
  only if its defining quasi-coherent ideal sheaf is of finite type.
\end{lemma}

\begin{proof}
  The assertion is local on $X$.  On $\operatorname{Spec}(A)$ the closed
  immersion is $\operatorname{Spec}(A/I)\to\operatorname{Spec}(A)$.  The
  quotient $A/I$ is a finitely presented $A$-algebra exactly when $I$ is
  finitely generated as an ideal.  Translate this equivalence back to the
  associated quasi-coherent ideal sheaf.
\end{proof}

\begin{definition}[Finite-type ideal attached to a quasi-compact open]
  \label{def:nagata-qc-open-complement-ideal}
  \uses{lemma:nagata-qc-open-complement-ideal-exists}
  \lean{Scheme.qcOpenComplementIdeal}

  Let $W\subseteq Y$ be a quasi-compact open in a qcqs scheme.  Choose one
  witness supplied by
  \Cref{lemma:nagata-qc-open-complement-ideal-exists} and define
  $\mathcal I_Y(W)$ to be its ideal component.  Finiteness, quasi-coherence,
  and the vanishing-set identity are specification fields, not additional
  assertions in this definition.
\end{definition}

\begin{lemma}[Vanishing set of the open-complement ideal]
  \label{lemma:nagata-qc-open-complement-ideal}
  \uses{def:nagata-qc-open-complement-ideal}
  \lean{Scheme.zeroLocus_qcOpenComplementIdeal}

  The ideal $\mathcal I_Y(W)$ is finite type and quasi-coherent, and
  \[
    V(\mathcal I_Y(W))=Y\setminus W.
  \]
\end{lemma}

\begin{proof}
  \uses{lemma:nagata-qc-open-complement-ideal-exists}
  The witness supplied by
  \Cref{lemma:nagata-qc-open-complement-ideal-exists} has exactly these
  fields, and the definition selects its ideal component.  Projecting those
  fields proves the statement.
\end{proof}

\begin{lemma}[A finite type nil ideal sheaf is uniformly nilpotent]
  \label{lemma:nagata-finite-type-nil-ideal-nilpotent}
  \lean{Scheme.IdealSheafData.isNilpotent_of_zeroLocus_eq_top}

  Let $X$ be quasi-compact and let $\mathcal I$ be a finite type
  quasi-coherent ideal with $V(\mathcal I)=X$.  Then there exists $n>0$ with
  $\mathcal I^n=0$.
\end{lemma}

\begin{proof}
  On an affine open, let $x_1,\ldots,x_r$ generate the ideal.  Every $x_i$
  lies in every prime and is nilpotent; choose $n_i>0$ with $x_i^{n_i}=0$.
  Set $N=1+\sum_i(n_i-1)$.  Every degree-$N$ monomial in the generators
  contains some $x_i^{n_i}$, so the $N$th power of the ideal is zero.  Choose
  a finite affine cover of $X$ and then a common exponent for the finitely many chart
  exponents.  Equality of quasi-coherent ideals is local, so that exponent
  annihilates $\mathcal I$ globally.
\end{proof}

\begin{lemma}[Preparation data for separating two closed boundaries,
  {\cite[\href{https://stacks.math.columbia.edu/tag/080P}{Tag 080P}]{stacks-project}}]
  \label{lemma:nagata-separate-closed-data-exists}
  \uses{def:nagata-admissible-blowup,def:nagata-strict-transform,
    lemma:nagata-chart-generator-regular}
  \lean{AlgebraicGeometry.exists_closedBoundarySeparationData}

  Let $X$ be qcqs, let $U\subseteq X$ be quasi-compact open, and let
  $Z_i=V(\mathcal I_i)\hookrightarrow X$ be finite-presentation closed
  immersions with $Z_1\cap Z_2\cap U=\varnothing$.  Put
  $Z=V(\mathcal I_1\mathcal I_2)$, $U_Z=Z\times_XU$, and
  $P_i=Z_i\times_XU$.  The two closed subschemes $P_i$ are the disjoint
  open-and-closed components of $U_Z=P_1\amalg P_2$.  There is a finite
  presentation closed center $C\hookrightarrow Z$, disjoint from $U_Z$, such
  that the blow-up $Z' = \operatorname{Bl}_C(Z)\to Z$ is $U_Z$-admissible
  and the strict transforms of $Z_1$ and $Z_2$ in $Z'$ are disjoint.
\end{lemma}

\begin{proof}
  \uses{def:nagata-qc-open-complement-ideal,
    lemma:nagata-finite-type-nil-ideal-nilpotent}
  Put $K_i=\mathcal I_Z(P_i)$, the finite-type ideal whose vanishing set is
  $Z\setminus P_i$ (\Cref{def:nagata-qc-open-complement-ideal}).  Since
  $P_1$ and $P_2$ are disjoint, $V(K_1K_2)=Z$; finite type and
  quasi-compactness therefore give $N>0$ with $(K_1K_2)^N=0$ by
  \Cref{lemma:nagata-finite-type-nil-ideal-nilpotent}.  The restrictions of
  both $K_i$ to the quasi-compact closed boundary $Z\setminus U_Z$ are finite
  type nil ideals, so enlarge $N$ until those restrictions are also killed by
  the $N$th power.  Replace $K_i$ by $K_i^N$; their vanishing sets are
  unchanged, now $K_1K_2=0$, and both ideals vanish on $Z\setminus U_Z$.  Let
  $J=K_1+K_2$ and $C=V(J)$.  On $U_Z=P_1\amalg P_2$, $K_1$ is the unit ideal
  on $P_1$ and zero on $P_2$, while $K_2$ has the opposite behaviour.  Hence
  $J|_{U_Z}=\mathcal O_{U_Z}$ and $J$ vanishes on $Z\setminus U_Z$.  Thus
  $Z\setminus C=U_Z$ and $C\cap U_Z=\varnothing$.  The closed
  immersion $C\hookrightarrow Z$ is finite presentation because $J$ is finite
  type.

  Work on an affine open $\operatorname{Spec}(A)\subseteq Z$, with generators
  $(a_\alpha)$ of $K_1$ and $(b_\beta)$ of $K_2$.  The blow-up of $J$ is
  covered by $A[J/a_\alpha]$ and $A[J/b_\beta]$.  In
  $A[J/a_\alpha]$, the element $a_\alpha$ is a nonzerodivisor by
  \Cref{lemma:nagata-chart-generator-regular}; since
  $a_\alpha K_2=0$, it follows that $K_2A[J/a_\alpha]=0$.  Interchanging
  the roles of $K_1$ and $K_2$ gives $K_1A[J/b_\beta]=0$.  An
  $a_\alpha$-chart and a $b_\beta$-chart have empty
  intersection: on their overlap $a_\alpha b_\beta$ is both zero and a
  nonzerodivisor, so the overlap ring is zero.  Consequently the blow-up is a
  disjoint union of the open-and-closed union $Z'_1$ of the $a_\alpha$-charts
  and the open-and-closed union $Z'_2$ of the $b_\beta$-charts; the construction
  is intrinsic in the extended ideals $K_i\mathcal O_{Z'}$, so these local
  decompositions glue.

  Over $U_Z$, the inverse image of $P_i$ lies in $Z'_i$.  The strict transform
  of $Z_i\to Z$ is the scheme-theoretic closure of that inverse image.  Since
  $Z'_i$ is closed as well as open, this closure remains in $Z'_i$.  The two
  strict transforms are therefore disjoint.
\end{proof}

\begin{definition}[Closed-boundary separation blow-up]
  \label{def:nagata-separate-closed-subschemes}
  \uses{lemma:nagata-separate-closed-data-exists,
    def:nagata-global-blowup,def:nagata-strict-transform}
  \lean{AlgebraicGeometry.separationBlowup}

  Under the hypotheses of
  \Cref{lemma:nagata-separate-closed-data-exists}, choose its center
  $C\hookrightarrow Z$.  Compose it with the finite-presentation closed
  immersion $Z\hookrightarrow X$, and define
  $\mathsf{SepBl}_U(Z_1,Z_2)\to X$ to be the blow-up of $X$ along the
  resulting center $C\hookrightarrow Z\hookrightarrow X$.
\end{definition}

\begin{lemma}[Properties of the closed-boundary separation blow-up,
  {\cite[\href{https://stacks.math.columbia.edu/tag/0F3V}{Tag 0F3V}]{stacks-project}}]
  \label{lemma:nagata-separate-closed-subschemes}
  \uses{def:nagata-separate-closed-subschemes}
  \lean{AlgebraicGeometry.disjoint_range_strictTransformMap_separationBlowup}

  The morphism $\mathsf{SepBl}_U(Z_1,Z_2)\to X$ is $U$-admissible, and the
  strict transforms of $Z_1$ and $Z_2$ in it are disjoint.
\end{lemma}

\begin{proof}
  \uses{lemma:nagata-separate-closed-data-exists,
    lemma:nagata-strict-transform-ambient-closed}
  The selected center is disjoint from $U_Z$, hence its image in $X$ is
  disjoint from $U$.  A composite of finite-presentation closed immersions is
  finite presentation, so the ambient center is finite presentation and its
  blow-up is $U$-admissible.  By
  \Cref{lemma:nagata-strict-transform-ambient-closed}, the strict transforms
  of $Z_1$ and $Z_2$ in this ambient blow-up agree with their strict
  transforms in $\operatorname{Bl}_C(Z)$, which are disjoint by the selected
  preparation witness.
\end{proof}

\begin{remark}[Formalization route]
  The formal proof
  (\texttt{separationBlowup\_disjoint\_strictTransforms}) blows up the sum
  $\mathcal I_1+\mathcal I_2$ directly on $X$
  (\texttt{separationBlowup}) and proves disjointness of the strict
  transforms chart-locally, using the chart inclusion
  $D_+(b_1+b_2)\subseteq D_+(b_1)\cup D_+(b_2)$ for degree-one sections
  $b_i$ coming from $\mathcal I_i$.  The Tag 080P preparation inside
  $Z=V(\mathcal I_1\mathcal I_2)$
  (\Cref{lemma:nagata-separate-closed-data-exists}) and
  \Cref{lemma:nagata-finite-type-nil-ideal-nilpotent} are not formalized.
\end{remark}

\begin{definition}[Ideal whose support is a constructible closed subset]
  \label{def:nagata-constructible-closed-support}
  \uses{def:nagata-qc-open-complement-ideal}
  \lean{Scheme.constructibleClosedSupportIdeal}

  Let $T\subseteq U$ be constructible and closed, with $U$ qcqs.  Its
  complement $W=U\setminus T$ is a quasi-compact open.  Define
  \[
    \mathcal J_U(T)=\mathcal I_U(W).
  \]
\end{definition}

\begin{lemma}[Property of the constructible-support ideal,
  {\cite[\href{https://stacks.math.columbia.edu/tag/01PH}{Tag 01PH}]{stacks-project}}]
  \label{lemma:nagata-constructible-closed-support}
  \uses{def:nagata-constructible-closed-support}
  \lean{Scheme.exists_finiteType_ideal_zeroLocus_eq}

  The ideal $\mathcal J_U(T)$ is finite type and quasi-coherent, and its
  vanishing set is exactly $T$.
\end{lemma}

\begin{proof}
  \uses{lemma:nagata-qc-open-complement-ideal}
  An open constructible subset of a spectral space is quasi-compact.  Apply
  \Cref{lemma:nagata-qc-open-complement-ideal} to $W=U\setminus T$; its
  vanishing set is $U\setminus W=T$.
\end{proof}

\begin{definition}[Constructible-boundary separation blow-up]
  \label{def:nagata-separate-constructible-closed}
  \uses{def:nagata-constructible-closed-support,
    def:nagata-finite-type-ideal-extension,
    def:nagata-separate-closed-subschemes}
  \lean{AlgebraicGeometry.exists_admissibleBlowup_disjoint_closures}

  Let $X$ be qcqs, let $U\subseteq X$ be a quasi-compact open, and let
  $T_1,T_2\subseteq U$ be disjoint constructible subsets closed in $U$.
  Define
  \[
    \mathcal J_i=\mathcal J_U(T_i),\qquad
    \mathcal I_i=\operatorname{Ext}_{X,U}(\mathcal J_i),\qquad
    Z_i=V(\mathcal I_i),
  \]
  and set
  \[
    \mathsf{CSepBl}_U(T_1,T_2)
      =\mathsf{SepBl}_U(Z_1,Z_2).
  \]
\end{definition}

\begin{lemma}[Properties of constructible-boundary separation]
  \label{lemma:nagata-separate-constructible-closed}
  \uses{def:nagata-separate-constructible-closed}
  \lean{AlgebraicGeometry.exists_admissibleBlowup_disjoint_closures}

  The constructed blow-up $\mathsf{CSepBl}_U(T_1,T_2)\to X$ is
  $U$-admissible, and the closures of the inverse images of $T_1$ and $T_2$
  are disjoint.
\end{lemma}

\begin{proof}
  \uses{lemma:nagata-constructible-closed-support,
    lemma:nagata-extend-finite-type-ideal,
    lemma:nagata-separate-closed-subschemes}
  By \Cref{lemma:nagata-constructible-closed-support},
  $V(\mathcal J_i)=T_i$ in $U$.  The restriction equality from
  \Cref{lemma:nagata-extend-finite-type-ideal} consequently gives
  $|Z_i\cap U|=T_i$, so $Z_1\cap Z_2\cap U=\varnothing$.  Apply
  \Cref{lemma:nagata-separate-closed-subschemes}.  In the resulting
  $U$-admissible blow-up $\pi\colon X'\to X$, the closure of
  $\pi^{-1}(T_i)$ is contained in the underlying set of the strict transform
  of $Z_i$: the strict transform is closed in $X'$, and it contains every
  point lying both over the preserved open and over $Z_i$ — in particular
  every point of $\pi^{-1}(T_i)$.  The two strict transforms are disjoint,
  so these two closures are disjoint as claimed.
\end{proof}

\begin{lemma}[Existence of constructible closed separators in a spectral space]
  \label{lemma:nagata-spectral-constructible-separation-exists}
  \lean{AlgebraicGeometry.Scheme.exists_disjoint_constructibleClosed_separation}

  If $T,T'$ are disjoint closed subsets of a spectral space $V$, there exist
  quasi-compact opens $W,W'\subseteq V$ such that
  \[
    T'\subseteq W,\quad W\cap T=\varnothing,\quad
    V\setminus W\subseteq W',\quad W'\cap T'=\varnothing.
  \]
\end{lemma}

\begin{proof}
  Quasi-compact opens form a basis of a spectral space.  For each $x\in T'$,
  choose a quasi-compact open neighbourhood contained in $V\setminus T$.
  The closed subset $T'$ is quasi-compact, so finitely many suffice; their
  union is a quasi-compact open $W$ containing $T'$ and disjoint from $T$.
  The closed subset $V\setminus W$ is quasi-compact and lies in the open
  $V\setminus T'$.  Cover it by finitely many quasi-compact opens contained
  in $V\setminus T'$ and let $W'$ be their union.
\end{proof}

\begin{definition}[Constructible separation data in a spectral space]
  \label{def:nagata-spectral-constructible-separation}
  \uses{lemma:nagata-spectral-constructible-separation-exists}
  \lean{TopologicalSpace.Spectral.constructibleClosedSeparation}

  Let $T,T'\subseteq V$ be disjoint closed subsets of a spectral space.
  Select quasi-compact opens $(W,W')$ from
  \Cref{lemma:nagata-spectral-constructible-separation-exists} and define
  \[
    T_c=V\setminus W,
    \qquad
    T'_c=V\setminus W'.
  \]
\end{definition}

\begin{lemma}[Properties of the constructible separation data]
  \label{lemma:nagata-spectral-constructible-separation}
  \uses{def:nagata-spectral-constructible-separation}
  \lean{TopologicalSpace.Spectral.exists_disjoint_constructibleClosed_supersets}

  The constructed subsets $T_c,T'_c$ are constructible and closed in $V$,
  and they satisfy
  \[
    T\subseteq T_c,\qquad T'\subseteq T'_c,\qquad
    T_c\cap T'_c=\varnothing.
  \]
\end{lemma}

\begin{proof}
  \uses{lemma:nagata-spectral-constructible-separation-exists}
  Complements of quasi-compact opens in a spectral space are constructible
  closed subsets.  The choice of $W$ gives $T\subseteq T_c$, the choice of
  $W'$ gives $T'\subseteq T'_c$, and $V\setminus W\subseteq W'$ gives
  $T_c\cap T'_c=\varnothing$.
\end{proof}

\begin{definition}[Maximal isomorphism open]
  \label{def:nagata-maximal-isomorphism-open}
  \lean{AlgebraicGeometry.Scheme.Hom.maximalIsomorphismOpen}

  For a morphism $q\colon Z\to X$, define $\operatorname{IsoOpen}(q)$ to be
  the union of all open subschemes $W\subseteq X$ for which
  $q^{-1}(W)\to W$ is an isomorphism.
\end{definition}

\begin{lemma}[Properties of the maximal isomorphism open]
  \label{lemma:nagata-maximal-isomorphism-open}
  \uses{def:nagata-maximal-isomorphism-open}
  \lean{AlgebraicGeometry.Scheme.Hom.isIso_restrict_maximalIsomorphismOpen}

  The subset $\operatorname{IsoOpen}(q)$ is open,
  $q^{-1}(\operatorname{IsoOpen}(q))\to\operatorname{IsoOpen}(q)$ is an
  isomorphism, and every open over which $q$ is an isomorphism is contained in
  it.
\end{lemma}

\begin{proof}
  A union of opens is open.  On pairwise intersections, the inverses of the
  restricted morphisms agree because an inverse to a fixed morphism is unique;
  hence the local inverses glue to an inverse over the union.  If an open
  $W$ has the property that $q^{-1}(W)\to W$ is an isomorphism, then $W$ is
  one of the opens in the defining union, and therefore $W\subseteq
  \operatorname{IsoOpen}(q)$; this is the maximality assertion.
\end{proof}

\begin{definition}[Modification making a proper map an isomorphism near a closure]
  \label{def:nagata-isomorphism-near-closure}
  \uses{def:nagata-spectral-constructible-separation,
    def:nagata-separate-constructible-closed,
    def:nagata-elimination-indeterminacy,
    def:nagata-extended-admissible-blowup,
    def:nagata-maximal-isomorphism-open}
  \lean{AlgebraicGeometry.exists_admissibleBlowup_strictTransform_isIso_near_closure}

  Let $X$ be qcqs, let $q\colon Z\to X$ be proper, let $V\subseteq X$ be a
  quasi-compact open, set
  $U=q^{-1}(V)$, and let $T\subseteq V$ be closed.  Assume $U\to V$ is an
  isomorphism over an open neighbourhood of $T$.  Let $E\subseteq V$ be
  the complement of $\operatorname{IsoOpen}(U\to V)$ from
  \Cref{def:nagata-maximal-isomorphism-open}.
  Apply \Cref{def:nagata-spectral-constructible-separation} to $(T,E)$ and
  denote the output by $(T_c,E_c)$.  Perform the explicitly constructed
  $V$-admissible separation blow-up
  \[
    X_1=\mathsf{CSepBl}_V(T_c,E_c).
  \]
  Let $\overline E_c$ denote the topological closure of the underlying subset
  of $E_c$ in $X_1$, and let
  $X_0=X_1\setminus\overline E_c$ and let $q_0\colon Z_0\to X_0$ be the
  restriction of the strict transform of $q$.  Form
  $X_0^{\mathrm{el}}(q_0,V\setminus E_c)$ and extend this admissible blow-up
  to $X_1$ by \Cref{def:nagata-extended-admissible-blowup}.  Define
  $\mathsf{IsoBl}_{V,T}(q)\to X$ to be the composite of these two
  $V$-admissible blow-ups, represented as one blow-up by the composition
  construction.
\end{definition}

\begin{lemma}[Affine scheme-image saturation]
  \label{lemma:nagata-scheme-image-affine-saturation}
  \uses{def:nagata-power-torsion,
    lemma:nagata-power-torsion-kernel-localization,
    lemma:nagata-strict-transform-closure-characterization}
  Let $C\to B$ be a ring map, let $K=\ker(C\to B)$, and let $a\in C$ map
  to a unit of $B$.  Put $C_a=C[1/a]$.  If $C/K\to B$ is injective, then
  \[
    \ker(C\longrightarrow C_a\longrightarrow B)
      =\{c\in C\mid \exists n,\ a^nc\in K\},
  \]
  and this kernel is the inverse image of
  $\Gamma_{(a)}(C/K)$ under $C\to C/K$.
  Consequently the scheme-theoretic image of the induced map from $B$ on an
  $a$-chart is the quotient by the $a$-power torsion, i.e. the strict-transform
  chart of the scheme-theoretic image of $\operatorname{Spec}(B)\to
  \operatorname{Spec}(C)$.
\end{lemma}

\begin{proof}
  If $c$ maps to zero in $B$, then its image in $C_a/KC_a$ is zero because
  $C/K\to B$ is injective.  The localization criterion gives
  $a^nc\in K$ for some $n$.  Conversely, such an equality maps to zero in $B$
  because $a$ is a unit there.  This proves the displayed kernel equality.
  On a standard blow-up chart replace $C$ by $C_a(A,B,I)$ and $a$ by the
  chart equation.  The quotient by the displayed kernel is exactly the
  quotient by $\Gamma_{(a)}$; the affine strict-transform formula and the
  closure characterization identify it with the strict-transform chart.
\end{proof}

\begin{lemma}[Scheme-image transport to a strict transform]
  \label{lemma:nagata-scheme-image-transport-strict-transform}
  \uses{lemma:nagata-scheme-image-restrict-open,
    lemma:nagata-scheme-image-affine-saturation,
    prop:nagata-strict-transform-closed-fiber-product,
    lemma:nagata-strict-transform-closure-characterization,
    thm:nagata-global-blowup-complement}
  Let $h\colon T\to X\times_SY$ be quasi-compact and quasi-separated, with
  first projection factoring through a quasi-compact open $U\subseteq X$.
  Let $Z=\operatorname{SIm}(h)$ and let $\pi\colon X'\to X$ be a
  $U$-admissible blow-up.  Via the canonical identity
  $X'\times_XU\cong U$, let $h'$ be the lifted map to $X'\times_SY$.  Then
  \[
    \operatorname{SIm}(h')
      \;=\;\operatorname{ST}_{X'/X}(Z)
  \]
  as closed subschemes of $X'\times_SY$, compatibly with their maps to $X'$.
\end{lemma}

\begin{proof}
  Restriction of scheme-theoretic image to $U\times_SY$ identifies the image
  there with the image of $h|_U$ by
  \Cref{lemma:nagata-scheme-image-restrict-open}.  Work on an affine chart of
  $X'\times_SY$ on which the pulled-back center has principal equation $a$.
  Write its coordinate ring as $C$ and refine the source by affine opens
  $\operatorname{Spec}(B_j)$.  The image ideal on this chart is
  $K=\bigcap_j\ker(C\to B_j)$, so $C/K\to\prod_jB_j$ is injective; the image
  of $a$ is a unit in every $B_j$ because $h$ factors through $U$.  Applying
  \Cref{lemma:nagata-scheme-image-affine-saturation} to the product ring
  identifies the kernel on this chart with the $a$-saturation of the pulled
  back image ideal.  By
  \Cref{prop:nagata-strict-transform-closed-fiber-product} and
  \Cref{lemma:nagata-strict-transform-closure-characterization}, that
  quotient is precisely the strict-transform chart.  The identifications are
  unchanged after localization, so they agree on overlaps and glue to the
  claimed global equality.
\end{proof}

\begin{lemma}[Property of the isomorphism-near-closure construction,
  {\cite[\href{https://stacks.math.columbia.edu/tag/0F3W}{Tag 0F3W}]{stacks-project}}]
  \label{lemma:nagata-isomorphism-near-closure}
  \uses{def:nagata-isomorphism-near-closure}
  \lean{AlgebraicGeometry.exists_admissibleBlowup_strictTransform_isIso_near_closure}

  If $U\to V$ is an isomorphism over an open neighbourhood of $T$, then
  $\mathsf{IsoBl}_{V,T}(q)\to X$ is $V$-admissible and the strict transform
  of $q$ is an isomorphism over an open neighbourhood of the closure of $T$
  in $\mathsf{IsoBl}_{V,T}(q)$.
\end{lemma}

\begin{proof}
  \uses{lemma:nagata-spectral-constructible-separation,
    lemma:nagata-maximal-isomorphism-open,
    lemma:nagata-separate-constructible-closed,
    lemma:nagata-elimination-indeterminacy,
    lemma:nagata-proper-open-dense-iso,
    lemma:nagata-extend-admissible-blowup,
    lemma:nagata-admissible-blowup-comp,
    lemma:nagata-admissible-blowup-preserves-density,
    prop:nagata-strict-transform-closed-fiber-product,
    lemma:nagata-strict-transform-over-complement}
  By \Cref{lemma:nagata-maximal-isomorphism-open}, the assumed isomorphism
  neighbourhood is contained in $\operatorname{IsoOpen}(U\to V)$; hence the
  bad locus $E$ is disjoint from $T$.  The two separation steps make
  $\overline T_c$ and $\overline E_c$
  disjoint in $X_1$, so $X_0$ contains an open neighbourhood of
  $\overline T_c$, hence of $\overline T$.  On
  $V\setminus E_c$ the restricted proper map is already an isomorphism.

  By \Cref{lemma:nagata-elimination-indeterminacy}, the strict transform after
  the constructed blow-up of $X_0$ is an open immersion.  The strict transform
  is proper because it is a closed subscheme of the base change of the proper
  morphism $q$ (use the closed-chart model in
  \Cref{prop:nagata-strict-transform-closed-fiber-product}); hence
  \Cref{lemma:nagata-proper-open-dense-iso} applies once density is checked.
  The open $V\setminus E_c$ is scheme-theoretically dense in $X_0$ by the
  density-preservation clause for admissible blow-ups, and its pullback is the
  preserved isomorphism open by
  \Cref{lemma:nagata-strict-transform-over-complement}.  Thus the strict
  transform over $X_0$ is an isomorphism.  The extension construction has
  unit center on $X_1\setminus X_0$ and restricts canonically to the preceding
  blow-up on $X_0$ by
  \Cref{lemma:nagata-extend-admissible-blowup}; the restriction and inverse
  maps therefore transport across the canonical strict-transform comparison.
  Since $\overline T\subset X_0$, this transported isomorphism holds on an
  open neighbourhood of $\overline T$.  Finally
  \Cref{lemma:nagata-admissible-blowup-comp} identifies the composite with one
  $V$-admissible blow-up, giving the asserted construction.
\end{proof}


\section{Common modifications and graph extension}

\begin{lemma}[Scheme-theoretic image of a partial section]
  \label{lemma:nagata-scheme-image-partial-section}
  \uses{def:nagata-scheme-theoretic-image,
    lemma:nagata-scheme-image-of-closed-immersion}
  \lean{AlgebraicGeometry.Nagata.ker_comap_partialSectionOpenι}

  Let \(U\hookrightarrow X\) be a quasi-compact open immersion, let
  \(Y\to S\) be separated, and let
  \(U\to X\times_SY\) be a morphism whose first projection is the given open
  immersion.  If \(Z\) is its scheme-theoretic image, then the projection
  \(Z\to X\) has inverse image \(U\) over the open \(U\subseteq X\), and this
  inverse image maps isomorphically to \(U\).
\end{lemma}

\begin{proof}
  \uses{lemma:nagata-scheme-image-restrict-open,
    lemma:nagata-scheme-image-of-closed-immersion}
  By \Cref{lemma:nagata-scheme-image-restrict-open}, restriction to
  \(U\times_SY\) commutes with taking the scheme-theoretic image.  On this
  open target the original map is a section of
  \(U\times_SY\to U\).  That projection is separated, so its section is a
  closed immersion.  By
  \Cref{lemma:nagata-scheme-image-of-closed-immersion}, its
  scheme-theoretic image is canonically its source, compatibly with the
  projection to $U$.
\end{proof}

\begin{lemma}[{\cite[\href{https://stacks.math.columbia.edu/tag/01W6}{Tag 01W6}]{stacks-project}}]
  \label{lemma:nagata-proper-to-separated-target}
  \lean{AlgebraicGeometry.IsProper.of_comp}

  In a commutative triangle $T\to Y\to B$, if $T\to B$ is proper and
  $Y\to B$ is separated, then $T\to Y$ is proper.
\end{lemma}

\begin{proof}
  Factor $T\to Y$ through the graph
  $T\to T\times_BY$.  The graph is closed because $Y\to B$ is separated, and
  the projection $T\times_BY\to Y$ is the base change of the proper morphism
  $T\to B$.
\end{proof}

\begin{lemma}[Finite type over a separated target]
  \label{lemma:nagata-finite-type-to-separated-target}
  \lean{AlgebraicGeometry.Nagata.locallyOfFiniteType_of_comp_of_isSeparated}

  In a commutative triangle \(T\to Y\to B\), if \(T\to B\) is of finite
  type and \(Y\to B\) is separated, then \(T\to Y\) is of finite type.
\end{lemma}

\begin{proof}
  Factor \(T\to Y\) as its graph
  \[
    T\longrightarrow T\times_BY\longrightarrow Y.
  \]
  The graph is a closed immersion because \(Y\to B\) is separated, hence is
  of finite type.  The second map is the base change of the finite-type
  morphism \(T\to B\).  Their composite is therefore of finite type.
\end{proof}

\begin{lemma}[Quasi-compactness of an open map to a separated target]
  \label{lemma:nagata-open-immersion-qc-over-separated-target}
  \lean{AlgebraicGeometry.Nagata.quasiCompact_of_comp_of_isSeparated}

  Let \(T\to Y\to B\) be a commutative triangle.  Assume that \(T\to B\) is
  of finite type, that \(Y\to B\) is separated, and that \(T\to Y\) is an
  open immersion.  Then \(T\to Y\) is a quasi-compact open immersion.
\end{lemma}

\begin{proof}
  \uses{lemma:nagata-finite-type-to-separated-target}
  By \Cref{lemma:nagata-finite-type-to-separated-target}, the open immersion
  is of finite type.  An open immersion is locally of finite type, and it is
  of finite type exactly when it is quasi-compact.
\end{proof}

\begin{definition}[Common-modification construction]
  \label{def:nagata-common-modification}
  \uses{def:nagata-scheme-theoretic-image,
    def:nagata-elimination-indeterminacy,
    def:nagata-global-blowup,
    def:nagata-extended-admissible-blowup}
  \lean{AlgebraicGeometry.CommonModification}

  Let $U\hookrightarrow X_i$ be open immersions over $S$, with
  $U,X_1,X_2$ separated and of finite type over the qcqs scheme $S$.
  Construct the following objects in order.
  \begin{enumerate}
    \item Let
      \[
        X_{12}=\operatorname{SIm}
        \bigl(U\longrightarrow X_1\times_SX_2\bigr)
      \]
      and write $p_i\colon X_{12}\to X_i$ for the projections.
    \item Apply \Cref{def:nagata-elimination-indeterminacy} to $(p_i,U)$.
      Denote the resulting blow-ups by
      $X_i^i= X^{\mathrm{el}}(p_i,U)$, their center ideals by
      $\mathcal I_i$, and the open strict transforms by
      $X_{12}^i\subseteq X_i^i$.
    \item On $X_{12}$ define
      \[
        \mathcal K_{12}
          =p_1^{-1}\mathcal I_1\mathcal O_{X_{12}}\,
           p_2^{-1}\mathcal I_2\mathcal O_{X_{12}},
        \qquad
        X'_{12}=\operatorname{Bl}_{\mathcal K_{12}}(X_{12}).
      \]
    \item The product-center blow-up maps to each $X_{12}^i$.  Extend
      $X'_{12}\to X_{12}^i$ across the open
      $X_{12}^i\subseteq X_i^i$ by
      \Cref{def:nagata-extended-admissible-blowup}; call the resulting
      ambient schemes $X'_i$.
    \item Define
      \[
        \mathsf{CM}(X_1,X_2;U)=X'_1\amalg_{X'_{12}}X'_2
      \]
      using the two constructed open immersions of $X'_{12}$.
  \end{enumerate}
  The construction includes the maps $X'_i\to X_i$, the two open maps
  $X'_i\to\mathsf{CM}(X_1,X_2;U)$, and the identified copy of $U$.
\end{definition}

\begin{lemma}[Properties of the common modification,
  {\cite[\href{https://stacks.math.columbia.edu/tag/0F3X}{Tag 0F3X}]{stacks-project}}]
  \label{lemma:nagata-common-modification}
  \uses{def:nagata-common-modification}
  \lean{AlgebraicGeometry.exists_commonModification}

  Each constructed map $X'_i\to X_i$ is a $U$-admissible blow-up.  Each
  $X'_i\to\mathsf{CM}(X_1,X_2;U)$ is an open immersion, their two copies of
  $U$ agree, and $\mathsf{CM}(X_1,X_2;U)$ is separated and of finite type
  over $S$.
\end{lemma}

\begin{proof}
  \uses{lemma:nagata-scheme-image-partial-section,
    lemma:nagata-elimination-indeterminacy,
    lemma:nagata-product-ideal-common-blowup,
    lemma:nagata-extend-admissible-blowup,
    lemma:nagata-proper-to-separated-target,
    lemma:nagata-separated-two-open-cover,
    lemma:nagata-admissible-blowup-comp}
  Apply Stacks Project, Tag~0F3X (Lemma~38.33.4), with the qcqs base $S$,
  the common quasi-compact open $U$, and the separated finite-type
  morphisms $X_i\to S$.  The two input maps are open immersions and are
  identified over $U$, so the hypotheses of that theorem are exactly the
  hypotheses of the present construction.  Its scheme-theoretic image
  $X_{12}\subseteq X_1\times_SX_2$ is the object named in
  \Cref{def:nagata-common-modification}; the two projections are the partial
  sections used by
  \Cref{lemma:nagata-scheme-image-partial-section}.  The theorem first
  principalizes the two projection ideals and then takes their product
  center, which is the map $X'_{12}\to X_{12}$ supplied by
  \Cref{lemma:nagata-product-ideal-common-blowup}.  Its strict-transform
  comparison gives $X'_{12}\subseteq X'_i$ as an open subscheme, with both
  restrictions identified with $U$; this is the extension step recorded by
  \Cref{lemma:nagata-extend-admissible-blowup}.  Consequently the two open
  immersions glue along exactly $X'_{12}$, not merely along their underlying
  topological intersection.  The graph of this overlap is a closed immersion
  because $X'_2\to S$ is separated.  The composite
  $X'_{12}\to X_1\times_SX_2$ is proper, and the product map from
  $X'_1\times_SX'_2$ to $X_1\times_SX_2$ is separated; hence the graph map
  is proper by
  \Cref{lemma:nagata-proper-to-separated-target}.  It is therefore closed,
  so the two-open diagonal criterion
  \Cref{lemma:nagata-separated-two-open-cover} proves that the glued scheme
  is separated.  The same theorem gives finite type, and the composition
  statement
  \Cref{lemma:nagata-admissible-blowup-comp} shows that each projection is a
  $U$-admissible blow-up.
\end{proof}

\begin{lemma}[{\cite[\href{https://stacks.math.columbia.edu/tag/0F3Y}{Tag 0F3Y}]{stacks-project}}]
  \label{lemma:nagata-scheme-image-strict-transform}
  \uses{def:nagata-scheme-theoretic-image,def:nagata-strict-transform,
    def:nagata-admissible-blowup,thm:nagata-global-blowup-complement}
  \lean{AlgebraicGeometry.Scheme.Hom.schemeTheoreticImage_after_admissibleBlowup}

  Let $V\to X\times_SY$ be quasi-compact and quasi-separated, and assume its
  projection to $X$ factors through an open $U\subseteq X$.  Let $Z$ be its
  scheme-theoretic image.  If $X'\to X$ is a $U$-admissible blow-up, then the scheme-theoretic
  image of the lifted map $V\to X'\times_SY$ is the strict transform of $Z$.
  More precisely, use the canonical isomorphism
  $\varepsilon\colon X'\times_XU\cong U$ from
  \Cref{thm:nagata-global-blowup-complement}; the lift is the map whose
  $X'$-component is obtained from $V\to U$ by
  $\varepsilon^{-1}\colon U\to X'\times_XU\to X'$ and whose $Y$-component is
  the given projection $V\to Y$.
\end{lemma}

\begin{proof}
  \uses{lemma:nagata-scheme-image-transport-strict-transform}
  Apply \Cref{lemma:nagata-scheme-image-transport-strict-transform} to the
  given map $V\to X\times_SY$.  Its lifted map is exactly the map specified in
  the statement through $\varepsilon^{-1}$, so the resulting image equality is
  the claimed identification with the strict transform.
\end{proof}

\begin{definition}[Graph-extension construction]
  \label{def:nagata-graph-extension}
  \uses{def:nagata-scheme-theoretic-image,
    def:nagata-elimination-indeterminacy}
  \lean{AlgebraicGeometry.Compactification.properExtensionConstruction}

  Let $U\to S$ be separated and of finite type, let $V\subseteq U$ be
  quasi-compact, and fix a compactification
  $V\hookrightarrow Y\to S$.  Define
  \[
    Z_V=\operatorname{SIm}
      \bigl(V\xrightarrow{(j,\,\mathrm{incl})}Y\times_SU\bigr).
  \]
  Apply \Cref{def:nagata-elimination-indeterminacy} to the first projection
  $Z_V\to Y$ over $V$.  Denote the resulting blow-up by
  $Y^{\mathrm{gr}}\to Y$ and the open strict transform by
  $Z_V^{\mathrm{st}}\hookrightarrow Y^{\mathrm{gr}}$.  Define
  \[
    V^{\mathrm{gr}}\subseteq Y^{\mathrm{gr}}
      \quad\text{to be the represented open subscheme }Z_V^{\mathrm{st}},
    \qquad
    \psi^{\mathrm{gr}}\colon V^{\mathrm{gr}}\longrightarrow U
  \]
  by transporting the second projection
  $Z_V^{\mathrm{st}}\to U$ across the open-immersion identification.
\end{definition}

\begin{lemma}[Properties of the graph-extension construction,
  {\cite[\href{https://stacks.math.columbia.edu/tag/0F3Z}{Tag 0F3Z}]{stacks-project}}]
  \label{lemma:nagata-graph-extension}
  \uses{def:nagata-graph-extension}
  \lean{AlgebraicGeometry.Compactification.exists_properExtension}

  The map $Y^{\mathrm{gr}}\to Y$ is $V$-admissible,
  $V\subseteq V^{\mathrm{gr}}\subseteq Y^{\mathrm{gr}}$ is open, and
  $\psi^{\mathrm{gr}}\colon V^{\mathrm{gr}}\to U$ is proper.  Its
  restriction to $V$ is the original inclusion, its inverse image of $V$ is
  exactly $V$, and this restriction is an isomorphism.
\end{lemma}

\begin{proof}
  \uses{lemma:nagata-scheme-image-partial-section,
    lemma:nagata-scheme-image-strict-transform,
    lemma:nagata-elimination-indeterminacy,
    lemma:nagata-strict-transform-ambient-closed,
    thm:nagata-global-blowup-proper,
    lemma:nagata-strict-transform-over-complement}
  Apply Stacks Project, Tag~0F3Z (Lemma~38.33.6), to the separated
  finite-type morphism $U\to S$, the quasi-compact open $V\subseteq U$, and
  the chosen compactification $V\hookrightarrow Y\to S$.  The theorem forms
  the scheme-theoretic image of the graph
  $V\to Y\times_SU$, exactly the object $Z_V$ in
  \Cref{def:nagata-graph-extension}, and after a $V$-admissible blow-up it
  makes the strict transform an open subscheme of
  $Y^{\mathrm{gr}}\times_SU$.  The open-immersion assertion is the
  elimination statement in
  \Cref{lemma:nagata-elimination-indeterminacy}; the closed-subscheme
  assertion is
  \Cref{lemma:nagata-strict-transform-ambient-closed}.  Since
  $Y^{\mathrm{gr}}\to S$ is proper by
  \Cref{thm:nagata-global-blowup-proper}, the projection to $U$ is proper,
  and its restriction is the transported map $\psi^{\mathrm{gr}}$.
  Finally,
  \Cref{lemma:nagata-scheme-image-partial-section} identifies the graph
  image over the open $V\subseteq Y$ with $V$, while
  \Cref{lemma:nagata-strict-transform-over-complement} identifies the
  admissible blow-up with the identity over that open.  Therefore
  \[
    V^{\mathrm{gr}}\cap
    (Y^{\mathrm{gr}}\times_YV)=V
  \]
  as schemes, the second projection agrees with the original inclusion on
  this open, and its inverse image of $V\subseteq U$ is exactly $V$.
\end{proof}


\section{The dense-source valuative criterion}

\begin{lemma}[Factorization through a closed subscheme over a valuation ring]
  \label{lemma:nagata-valuation-factor-closed}
  \lean{AlgebraicGeometry.IsClosedImmersion.exists_lift_of_isFractionRing}

  Let $A$ be a valuation ring with fraction field $K$, let $Z\hookrightarrow Y$
  be a closed immersion, and let $g\colon\operatorname{Spec}(A)\to Y$.  If the
  restriction of $g$ to $\operatorname{Spec}(K)$ factors through $Z$, then
  $g$ factors through $Z$.
\end{lemma}

\begin{proof}
  Pull the defining ideal of $Z$ back to $A$.  It becomes zero after applying
  the injective map $A\hookrightarrow K$, so it was already zero in $A$.
\end{proof}

\begin{lemma}[Existence of a valuation probe through a dense source]
  \label{lemma:nagata-dense-valuation-probe-exists}
  \uses{def:nagata-scheme-theoretic-image}
  \lean{AlgebraicGeometry.nonempty_valuationProbe}

  Let $h\colon T\to P$ be quasi-compact, assume its image
  inclusion $\operatorname{SIm}(h)\hookrightarrow P$ is an isomorphism, assume
  $P$ is affine, and let $z\in P$.  There
  exist a valuation ring $A'$ with fraction field $K'$, a morphism
  $\eta\colon\operatorname{Spec}(K')\to T$, and a morphism
  $\alpha\colon\operatorname{Spec}(A')\to P$ such that the generic square
  commutes and the closed point of $\operatorname{Spec}(A')$ maps to $z$.
\end{lemma}

\begin{proof}
  Let $Z=\operatorname{SIm}(h)$.  The assumed image isomorphism identifies
  $z$ with a point of $Z$.  Stacks Project, Tag~02JQ (Lemma~29.6.5),
  applied to the quasi-compact morphism $h$ and the point $z\in Z$, supplies
  a valuation ring $A'$, its fraction field $K'$, a map
  $\operatorname{Spec}(K')\to T$, and a map
  $\operatorname{Spec}(A')\to Z$ whose generic restriction is the supplied
  map and whose closed point is $z$.  Composing the latter map with the image
  closed immersion $Z\hookrightarrow P$ gives $\alpha$, and the generic
  commutativity supplied by the theorem is exactly the displayed square.
\end{proof}

\begin{definition}[Valuation probe through a schematically dense source]
  \label{def:nagata-dense-valuation-probe}
  \uses{lemma:nagata-dense-valuation-probe-exists}
  \lean{AlgebraicGeometry.ValuationProbe}

  Let $h\colon T\to P$ be quasi-compact, assume its image
  inclusion $\operatorname{SIm}(h)\hookrightarrow P$ is an isomorphism, assume
  $P$ is affine, and let $z\in P$.
  Select a witness from
  \Cref{lemma:nagata-dense-valuation-probe-exists} and denote it by
  \[
    \mathsf{VP}(h,z)=
      (A'_{h,z},K'_{h,z},\eta_{h,z},\alpha_{h,z}),
  \]
  where $A'_{h,z}$ is the selected valuation ring, $K'_{h,z}$ its fraction
  field, $\eta_{h,z}\colon\operatorname{Spec}(K'_{h,z})\to T$ the selected
  generic lift, and
  $\alpha_{h,z}\colon\operatorname{Spec}(A'_{h,z})\to P$ the selected centred
  valuation map.
\end{definition}

\begin{lemma}[Properties of the valuation probe]
  \label{lemma:nagata-dense-valuation-probe}
  \uses{def:nagata-dense-valuation-probe}
  \lean{AlgebraicGeometry.ValuationProbe.commSq,
    AlgebraicGeometry.ValuationProbe.closedPoint_eq}

  The maps in $\mathsf{VP}(h,z)$ form a commutative diagram
  \[
    \begin{array}{ccc}
      \operatorname{Spec}(K'_{h,z})&\xrightarrow{\eta_{h,z}}&T\\
      \big\downarrow&&\big\downarrow h\\
      \operatorname{Spec}(A'_{h,z})&\xrightarrow{\alpha_{h,z}}&P,
    \end{array}
  \]
  and the closed point of $\operatorname{Spec}(A'_{h,z})$ maps to $z$.
\end{lemma}

\begin{proof}
  \uses{lemma:nagata-dense-valuation-probe-exists}
  These are exactly the fields of the witness selected in
  \Cref{def:nagata-dense-valuation-probe}; nonemptiness is
  \Cref{lemma:nagata-dense-valuation-probe-exists}.
\end{proof}


\begin{lemma}[Composite of two valuation rings]
  \label{lemma:nagata-composite-valuation-ring}
  \lean{ValuationSubring.comapResidue}

  Let \(A'\) be a valuation ring with maximal ideal \(\mathfrak m'\), and let
  \(A''\) be a valuation subring of the residue field
  \(k'=A'/\mathfrak m'\).  The inverse image
  \[
    C=\{a\in A'\mid a\bmod\mathfrak m'\in A''\}
  \]
  is a valuation ring.  Its fraction field is the fraction field of \(A'\),
  \(\mathfrak m'\subseteq C\), and \(C/\mathfrak m'\cong A''\).
\end{lemma}

\begin{proof}
  Let \(x\) be a nonzero element of \(\operatorname{Frac}(A')\).  Since
  \(A'\) is a valuation ring, either \(x\in A'\) or \(x^{-1}\in A'\).  If
  exactly one lies in \(A'\), that element lies in \(C\) whenever its residue
  is zero.  If both are units of \(A'\), compare their nonzero residue classes
  in \(k'\); the valuation-ring property of \(A''\) puts one of the two
  residues in \(A''\), and hence puts \(x\) or \(x^{-1}\) in \(C\).  Thus
  \(C\) is a valuation ring.  If $\mathfrak m'=0$, then $A'=k'$ is a field,
  $C=A''$, and $\operatorname{Frac}(C)=k'=A'$.  If
  $\mathfrak m'\ne0$, choose $0\ne t\in\mathfrak m'$.  For every $a\in A'$,
  both $at$ and $t$ lie in $\mathfrak m'\subseteq C$, so
  $a=(at)/t\in\operatorname{Frac}(C)$.  Hence in this case also
  $\operatorname{Frac}(C)=\operatorname{Frac}(A')$.  The restricted residue
  map $C\to A''$ is surjective by the inverse-image definition, and its
  kernel is $\mathfrak m'$, giving $C/\mathfrak m'\cong A''$.
\end{proof}

\begin{lemma}[Intersection with a dominating valuation ring]
  \label{lemma:nagata-valuation-intersection}
  \lean{ValuationSubring.comap_eq_of_le_comap}

  Let \(A\) be a valuation ring with fraction field \(K\), let
  \(K\hookrightarrow L\) be a field extension, and let \(B\subseteq L\) be a
  valuation ring dominating the image of \(A\).  Then
  \[
    K\cap B=A
  \]
  inside \(L\).
\end{lemma}

\begin{proof}
  The inclusion \(A\subseteq K\cap B\) follows because $B$ contains the image
  of $A$.  Conversely, let $x\in K\cap B$.  If $x\notin A$, the valuation-ring
  property of $A$ gives $x^{-1}\in A$, and in fact
  $x^{-1}\in\mathfrak m_A$: otherwise $x^{-1}$ would be a unit of $A$ and
  $x$ would belong to $A$.  Since $B$ dominates $A$, the image of
  $\mathfrak m_A$ lies in $\mathfrak m_B$.  But $x\in B$ and
  $x^{-1}\in\mathfrak m_B$ would imply
  $1=xx^{-1}\in\mathfrak m_B$, a contradiction.  Thus $x\in A$, proving
  $K\cap B\subseteq A$ and hence the asserted equality.
\end{proof}

\begin{lemma}[Dense image into a reduced scheme is schematically dense]
  \label{lemma:nagata-reduced-dense-image-schematic}
  \uses{def:nagata-scheme-theoretic-image}
  \lean{AlgebraicGeometry.IsSchemeTheoreticallyDominant.of_isDominant}

  Let $X$ be reduced and let $h\colon U\to X$ be quasi-compact and
  quasi-separated with dense topological image.  Then the image inclusion
  $\operatorname{SIm}(h)\hookrightarrow X$ is an isomorphism.
\end{lemma}

\begin{proof}
  On an affine open $\operatorname{Spec}(A)\subseteq X$, every element of the
  kernel ideal defining the scheme-theoretic image vanishes at every point of
  a dense subset.  It therefore lies in every prime ideal of $A$, hence in
  the nilradical.  Since $A$ is reduced, the element is zero.  The kernel
  vanishes on an affine cover, so the image inclusion is an isomorphism.
\end{proof}

\begin{lemma}[A field-valued point has a flat polynomial presentation]
  \label{lemma:nagata-field-point-polynomial-presentation}

  For every ring homomorphism $B\to K$ with $K$ a field, there is a
  polynomial $B$-algebra $P$ and a surjective $B$-algebra map $P\to K$.
  The structural map $B\to P$ is flat.
\end{lemma}

\begin{proof}
  Take the polynomial algebra over $B$ on variables indexed by the underlying
  set of $K$ and evaluate the variable indexed by $k$ at $k$.  Evaluation is
  surjective.  This polynomial algebra is a filtered colimit of free
  polynomial algebras on finite variable sets, and is therefore flat over
  $B$.
\end{proof}

\begin{lemma}[Valuation domination of a local domain]
  \label{lemma:nagata-local-domain-valuation-dominates}
  \lean{ValuationRing.exists_dominating}

  If $B$ is a local domain with fraction field $L$, there is a valuation ring
  $C\subseteq L$ whose maximal ideal contracts to the maximal ideal of $B$.
\end{lemma}

\begin{proof}
  Consider the set of subrings of $L$ containing $B$ and whose units avoid the
  maximal ideal of $B$, ordered by inclusion.  The union of a chain has the
  same property, so Zorn's lemma gives a maximal such ring $C$.  For any
  $x\in L^\times$, if neither $x$ nor $x^{-1}$ belonged to $C$, adjoining one
  of them and localizing at the complement of the maximal ideal would produce
  a strictly larger ring with the same property, contradicting maximality.
  Thus $C$ is a valuation ring.  Its unique maximal ideal contracts to the
  maximal ideal of $B$ because $B$ is local and every element outside that
  ideal is a unit in $C$.
\end{proof}

\begin{lemma}[Extension of a valuation ring to a field extension,
  {\cite[\href{https://stacks.math.columbia.edu/tag/00IA}{Lemma 10.50.2}]{stacks-project}}]
  \label{lemma:nagata-valuation-ring-extension}
  \lean{ValuationSubring.exists_comap_eq}

  Let $A$ be a valuation ring with fraction field $K$, and let
  $K\hookrightarrow L$ be a field extension.  There exists a valuation ring
  $B\subseteq L$ with fraction field $L$ which dominates $A$.  In Lean the
  conclusion is recorded in the combined form $K\cap B=A$ of
  \Cref{lemma:nagata-valuation-intersection}, from which the domination is
  recovered.
\end{lemma}

\begin{proof}
  Apply Stacks Project, Tag~00IA (Lemma~10.50.2), to the inclusion
  $K\hookrightarrow L$.  Its proof considers the partially ordered set of
  local subrings of $L$ containing $A$ and dominating $A$; a chain union is
  again such a local subring, so Zorn gives a maximal one.  If its fraction
  field is smaller than $L$, choose $t$ outside that fraction field.  In the
  transcendental case, localizing $A[t]$ at the prime generated by $t$ and
  $\mathfrak m_A$ gives a strictly larger dominating local subring.  In the
  algebraic case, choose $0\ne a\in A$ with $at$ integral over $A$, choose a
  prime over $\mathfrak m_A$ in the finite algebra $A[at]$, and localize
  there; this again contradicts maximality.  Hence the maximal subring has
  fraction field $L$.  The same maximality argument applied to $x$ and
  $x^{-1}$ for every $0\ne x\in L$ proves that it is a valuation ring, and
  its maximal ideal contracts to $\mathfrak m_A$, so it dominates $A$.
\end{proof}

\begin{lemma}[{\cite[\href{https://stacks.math.columbia.edu/tag/0894}{Tag 0894}]{stacks-project}}]
  \label{lemma:nagata-dense-source-valuative}
  \lean{AlgebraicGeometry.ValuativeCriterion.universallyClosed_of_isDominant,
    AlgebraicGeometry.ValuativeCriterion.isSeparated_of_isDominant}

  Let \(f\colon X\to S\) be quasi-compact and let
  \(h\colon U\to X\) be quasi-compact and quasi-separated,
  and assume \(h(U)\) is dense in \(X\).  Suppose that for every valuation
  ring \(A\) with fraction field \(K\), every commutative diagram
  \[
    \operatorname{Spec}(K)\longrightarrow U\longrightarrow X,
    \qquad
    \operatorname{Spec}(A)\longrightarrow S
  \]
  has a unique lift \(\operatorname{Spec}(A)\to X\).  Then \(f\) is
  universally closed.  If \(f\) is quasi-separated, it is also separated.
\end{lemma}

\begin{remark}[Formalization route]
  In Lean the two conclusions are split into
  \texttt{ValuativeCriterion.universallyClosed\_of\_isDominant} and
  \texttt{ValuativeCriterion.isSeparated\_of\_isDominant}; both are fully
  proved.  The ingredients — the valuation probe
  (Tag 02JQ, \texttt{nonempty\_valuationProbe}), the
  composite valuation ring (\texttt{ValuationSubring.comapResidue}), the
  intersection lemma (Tag 052K,
  \texttt{ValuationSubring.comap\_eq\_of\_le\_comap}) and the extension lemma
  (Tag 00IA, \texttt{ValuationSubring.exists\_comap\_eq}, via
  Mathlib's Chevalley extension theorem rather than the Zorn argument) — live
  in the \texttt{WeilConjectures.Mathlib} helper files.  The formalized proof
  deviates from the Stacks Project text in two reductions: instead of
  replacing $U,X,S$ by their reductions, the morphism $h$ is factored through
  its scheme-theoretic image — the factor map \texttt{Scheme.Hom.toImage} is
  scheme-theoretically dominant with no reducedness hypothesis, and the
  valuative data factors through the image because $\operatorname{Spec}$ of a
  valuation ring or of a field is reduced
  (\texttt{Scheme.IdealSheafData.le\_ker\_of\_range\_subset\_support}); and the
  polynomial presentation of
  \Cref{lemma:nagata-field-point-polynomial-presentation} is inlined via
  \texttt{MvPolynomial.eval₂Hom}.  The flat-pullback and reduced-density
  inputs (\Cref{lemma:nagata-dense-flat-pullback},
  \Cref{lemma:nagata-reduced-dense-image-schematic}) are provided by Mathlib's
  \texttt{IsSchemeTheoreticallyDominant} API.
\end{remark}

\begin{proof}
  \uses{lemma:nagata-dense-flat-pullback,lemma:nagata-dense-valuation-probe,
    lemma:nagata-composite-valuation-ring,lemma:nagata-valuation-intersection,
    lemma:nagata-reduced-dense-image-schematic,
    lemma:nagata-field-point-polynomial-presentation,
    lemma:nagata-valuation-ring-extension}
  We first verify the existence part of the valuative criterion for \(f\).
  Start with an arbitrary square
  \[
    \operatorname{Spec}(K)\longrightarrow X,
    \qquad
    \operatorname{Spec}(A)\longrightarrow S.
  \]
  Valuation rings and fields are reduced, so replacing \(U,X,S\) by their
  reductions does not change this lifting problem.  After this replacement,
  topological density of \(h(U)\) says that the scheme-theoretic image of
  \(h\) is all of \(X\) by
  \Cref{lemma:nagata-reduced-dense-image-schematic}.  We may also replace
  \(S\) by an affine open through
  which \(\operatorname{Spec}(A)\to S\) factors; write it as
  \(\operatorname{Spec}(R)\).

  Choose an affine open \(\operatorname{Spec}(B)\subseteq X\) through which
  the generic map factors.  Choose the polynomial $B$-algebra $P$ and
  surjection $P\twoheadrightarrow K$ from
  \Cref{lemma:nagata-field-point-polynomial-presentation}.  The map
  \(\operatorname{Spec}(P)\to X\) is flat.  Therefore
  \[
    U\times_X\operatorname{Spec}(P)\longrightarrow\operatorname{Spec}(P)
  \]
  is scheme-theoretically dense by
  \Cref{lemma:nagata-dense-flat-pullback}.

  Apply \Cref{lemma:nagata-dense-valuation-probe} to the point
  \(\operatorname{Spec}(K)\subseteq\operatorname{Spec}(P)\).  We obtain a
  valuation ring \(A'\), with fraction field \(K'\), whose generic point maps
  to \(U\times_X\operatorname{Spec}(P)\) and whose closed point maps to the
  chosen \(K\)-point of \(\operatorname{Spec}(P)\).  If
  \(k'=A'/\mathfrak m_{A'}\), this gives a field embedding
  \(K\to k'\).  Choose a valuation ring \(A''\subseteq k'\) dominating the
  image of \(A\), with fraction field \(k'\), by
  \Cref{lemma:nagata-valuation-ring-extension}.

  Form the inverse-image ring
  \[
    C=\{a\in A'\mid a\bmod\mathfrak m_{A'}\in A''\}.
  \]
  By \Cref{lemma:nagata-composite-valuation-ring}, \(C\) is a valuation ring
  with fraction field \(K'\), and \(C/\mathfrak m_{A'}=A''\).  Since
  \(A''\) contains the image of \(A\), the map \(R\to A'\) factors through
  \(C\).  Thus the generic map
  \(\operatorname{Spec}(K')\to U\) and the map
  \(\operatorname{Spec}(C)\to S\) form a diagram covered by the hypothesis.
  Let \(g_C\colon\operatorname{Spec}(C)\to X\) be its unique lift.

  There are now two maps \(\operatorname{Spec}(A')\to X\): one is obtained
  from \(\operatorname{Spec}(A')\to\operatorname{Spec}(P)\to X\), and the
  other by composing \(\operatorname{Spec}(A')\to\operatorname{Spec}(C)\)
  with \(g_C\).  They have the same generic map through \(U\) and the same
  composite to \(S\).  Applying the assumed uniqueness to the valuation ring
  \(A'\) shows that the two maps are equal.

  Restrict $g_C$ along the closed immersion
  $\operatorname{Spec}(A'')\to\operatorname{Spec}(C)$ induced by
  $C/\mathfrak m_{A'}\cong A''$.  Its generic point is
  $\operatorname{Spec}(k')$, the residue field of the prime
  $\mathfrak m_{A'}\subset C$.  The preceding equality on
  $\operatorname{Spec}(A')$ identifies this generic map with the closed-point
  residue map of $\operatorname{Spec}(A')\to X$ and hence with the base
  change of the original $\operatorname{Spec}(K)\to X$.  Let \(x\in X\) be the image
  of the closed point.  The image of the induced local homomorphism
  \(\mathcal O_{X,x}\to A''\) lies in the subfield \(K\subseteq k'\).
  By \Cref{lemma:nagata-valuation-intersection},
  \(K\cap A''=A\).  Hence the local homomorphism factors through \(A\), and
  therefore defines the required lift \(\operatorname{Spec}(A)\to X\).
  The existence part of the valuative criterion now shows that \(f\) is
  universally closed.

  Finally suppose \(f\) is quasi-separated.  Then its diagonal
  \(\Delta_f\colon X\to X\times_SX\) is quasi-compact.  A map
  \(\operatorname{Spec}(A)\to X\times_SX\) is a pair of maps to \(X\).
  If their common generic restriction factors through \(U\), the uniqueness
  hypothesis for \(f\) forces the pair to be equal, hence to factor uniquely
  through \(\Delta_f\).  Applying the universally-closed part just proved to
  \(\Delta_f\) and \(h\colon U\to X\) makes the diagonal universally closed.
  Since every diagonal is an immersion, a universally closed diagonal is a
  closed immersion.  Thus \(f\) is separated.
\end{proof}

\section{The two-open compactification theorem}

Throughout this section let $S$ be Noetherian, let $U\to S$ be separated and
of finite type, and fix an open cover $U=U_1\cup U_2$ such that
$U_1\cap U_2$ is dense in $U$.  Assume each $U_i\to S$ is compactifiable.

\begin{definition}[Chosen dense compactifications of the two opens]
  \label{def:nagata-two-open-dense-compactifications}
  \uses{def:nagata-compactifiable,def:nagata-dense-refinement}
  \lean{AlgebraicGeometry.exists_compactification_isSchemeTheoreticallyDominant}

  Let $C_i^0$ be the compactification selected by classical choice from the
  nonempty type witnessing compactifiability of $U_i\to S$.  Define
  \[
    C_i=\mathsf{Dense}(C_i^0)
      =\bigl(U_i\hookrightarrow X_i\longrightarrow S\bigr).
  \]
\end{definition}

\begin{lemma}[Properties of the chosen dense compactifications]
  \label{lemma:nagata-two-open-dense-compactifications}
  \uses{def:nagata-two-open-dense-compactifications}
  \lean{AlgebraicGeometry.exists_compactification_isSchemeTheoreticallyDominant}

  For $i=1,2$, the constructed $C_i$ is a scheme-theoretically dense
  compactification of $U_i\to S$; in particular $X_i\to S$ is proper and
  $U_i\hookrightarrow X_i$ is a quasi-compact open immersion.
\end{lemma}

\begin{proof}
  \uses{prop:nagata-dense-refinement}
  This is exactly \Cref{prop:nagata-dense-refinement}, applied to each selected
  compactification $C_i^0$.
\end{proof}

\begin{definition}[Chosen graph extensions]
  \label{def:nagata-two-open-graph-extensions}
  \uses{def:nagata-two-open-dense-compactifications,
    def:nagata-graph-extension}
  \lean{AlgebraicGeometry.CompactifiedProperExtension}

  Apply the graph-extension construction to
  $U_i\subseteq U$ and $C_i$.  Denote its output by
  \[
    X_i^{\mathrm{gr}}\longrightarrow X_i,\qquad
    U_i\subseteq V_i\subseteq X_i^{\mathrm{gr}},\qquad
    \psi_i\colon V_i\longrightarrow U.
  \]
  From now on replace $X_i$ by $X_i^{\mathrm{gr}}$ and retain the notation
  $X_i$ for the resulting proper ambient scheme.
\end{definition}

\begin{lemma}[Properties of the chosen graph extensions]
  \label{lemma:nagata-two-open-graph-extensions}
  \uses{def:nagata-two-open-graph-extensions}
  \lean{AlgebraicGeometry.exists_prepared_twoOpenSide}

  Each replacement $X_i\to X_i^{\mathrm{old}}$ is $U_i$-admissible;
  the resulting structural morphism $X_i\to S$ is proper;
  $\psi_i$ is proper, extends the inclusion $U_i\hookrightarrow U$, satisfies
  $\psi_i^{-1}(U_i)=U_i$, and is an isomorphism over $U_i$.
\end{lemma}

\begin{proof}
  \uses{lemma:nagata-graph-extension}
  For each $i$, \Cref{lemma:nagata-graph-extension} says that the selected
  replacement $X_i\to X_i^{\mathrm{old}}$ is $U_i$-admissible and that its
  associated map $\psi_i$ is proper, extends $U_i\hookrightarrow U$, has
  inverse image $\psi_i^{-1}(U_i)=U_i$, and is an isomorphism there.  The
  old structural map $X_i^{\mathrm{old}}\to S$ is proper by the selected
  compactification, while an admissible blow-up is proper.  Their composite,
  which is the structural map of the replacement, is therefore proper.
\end{proof}

\begin{definition}[The four boundary loci]
  \label{def:nagata-two-open-boundary-loci}
  \uses{def:nagata-two-open-graph-extensions}
  \lean{AlgebraicGeometry.CompactifiedProperExtension.forbiddenBoundary}

  For $\{i,j\}=\{1,2\}$, define
  \[
    Z_i=(U\setminus U_j)_{\mathrm{red}}\hookrightarrow U_i,
    \qquad
    Z_{i,k}=V_i\times_U Z_k\hookrightarrow V_i
    \quad(k=1,2).
  \]
  Thus $Z_i$ is the canonical reduced closed subscheme on the indicated
  complement, and $Z_{i,k}$ carries the scheme-theoretic inverse-image
  structure rather than only its underlying closed set.
\end{definition}


\begin{lemma}[Closedness and disjointness of the four boundary loci]
  \label{lemma:nagata-boundary-loci}
  \uses{def:nagata-two-open-boundary-loci}
  \lean{AlgebraicGeometry.CompactifiedProperExtension.closure_image_preimage_inter}

  The subschemes $Z_i$ are closed in $U$.  Each $Z_{i,k}$ is closed in
  $V_i$, their underlying closed sets are disjoint, and $\psi_i$ is an
  isomorphism over an open neighbourhood of $Z_i$.
\end{lemma}

\begin{proof}
  \uses{lemma:nagata-graph-extension}
  Since the support of $Z_i$ is $U\setminus U_j$, it is closed in $U$ and
  contained in $U_i$.  Base change of this closed immersion along $\psi_i$
  gives each closed immersion $Z_{i,k}\hookrightarrow V_i$.  The supports of
  $Z_i$ and $Z_j$ are disjoint because
  $(U\setminus U_j)\cap(U\setminus U_i)=U\setminus(U_i\cup U_j)=\varnothing$;
  their inverse images in $V_i$ are consequently disjoint.  Finally,
  $Z_i\subseteq U_i$, and \Cref{lemma:nagata-graph-extension} identifies
  $\psi_i^{-1}(U_i)\to U_i$ with an isomorphism.  This gives the required
  isomorphism on the open neighbourhood $U_i$ of $Z_i$.
\end{proof}

\begin{definition}[Prepared ambient schemes with separated boundary closures]
  \label{def:nagata-two-open-prepared-boundaries}
  \uses{def:nagata-two-open-boundary-loci,
    def:nagata-separate-constructible-closed}
  \lean{AlgebraicGeometry.CompactifiedProperExtension.Prepared}

  For $\{i,j\}=\{1,2\}$, define
  \[
    X_i^{\mathrm{prep}}
      =\mathsf{CSepBl}_{V_i}(Z_{i,i},Z_{i,j})
  \]
  and transport $V_i$, $U_i$, and $\psi_i$ along the chosen isomorphism over
  $V_i$.  Replace $X_i$ by $X_i^{\mathrm{prep}}$ and let
  $\overline Z_{i,k}$ denote the closure of the transported $Z_{i,k}$ in the
  new $X_i$.
\end{definition}

\begin{lemma}[Properties of the prepared boundary closures]
  \label{lemma:nagata-prepared-boundaries}
  \uses{def:nagata-two-open-prepared-boundaries}
  \lean{AlgebraicGeometry.exists_prepared_twoOpenSide}

  The constructed replacement $X_i^{\mathrm{prep}}\to X_i^{\mathrm{old}}$
  is $V_i$-admissible, the closures
  $\overline Z_{i,i}$ and $\overline Z_{i,j}$ are disjoint, and the open
  $V_i$ together with the proper map $\psi_i$ is unchanged.  The new
  structural morphism $X_i^{\mathrm{prep}}\to S$ is proper.
\end{lemma}

\begin{proof}
  \uses{lemma:nagata-boundary-loci,
    lemma:nagata-separate-constructible-closed,
    lemma:nagata-blowup-proper}
  By \Cref{lemma:nagata-boundary-loci}, the subsets $Z_{i,i}$ and
  $Z_{i,j}$ are disjoint constructible subsets closed in $V_i$.  Apply the
  property of the explicitly constructed separation blow-up.  It is
  $V_i$-admissible, so the selected isomorphism over $V_i$ transports
  $\psi_i$ without changing it.  The conclusion of
  \Cref{lemma:nagata-separate-constructible-closed} is exactly the
  disjointness of the two closures.  The separation blow-up is proper by
  \Cref{lemma:nagata-blowup-proper}; composing it with the original proper
  morphism $X_i^{\mathrm{old}}\to S$ proves that
  $X_i^{\mathrm{prep}}\to S$ is proper.
\end{proof}


\begin{lemma}[Disjoint boundary closures persist]
  \label{lemma:nagata-boundary-disjoint-persists}
  \lean{AlgebraicGeometry.disjoint_range_strictTransformMap_of_disjoint}

  If $\overline Z_{i,i}$ and $\overline Z_{i,j}$ are disjoint, then their
  strict transforms remain disjoint after every further $V_i$-admissible
  blow-up.
\end{lemma}

\begin{proof}
  \uses{lemma:nagata-strict-transform-restrict}
  Let $b\colon X'_i\to X_i$ be a further $V_i$-admissible blow-up.  By
  \Cref{lemma:nagata-strict-transform-restrict}, the strict transform of
  $\overline Z_{i,k}$ is the closure of its restriction over the open on
  which $b$ is an isomorphism.  It is therefore contained in
  $b^{-1}(\overline Z_{i,k})$.  The inverse images of the disjoint closed
  subschemes $\overline Z_{i,i}$ and $\overline Z_{i,j}$ are disjoint, so
  their contained strict transforms are disjoint as well.
\end{proof}

\begin{definition}[Transport of a side along a $V_i$-admissible blow-up]
  \label{def:nagata-two-open-side-transport}
  \uses{def:nagata-two-open-graph-extensions,
    def:nagata-two-open-boundary-loci,
    def:nagata-admissible-blowup}
  \lean{AlgebraicGeometry.CompactifiedProperExtension.admissibleBlowup}

  Let $b\colon X_i'\to X_i$ be a $V_i$-admissible blow-up.  Transport the
  side $(X_i,V_i,\psi_i,U_i\hookrightarrow V_i)$ along $b$: take $X_i'$ with
  its structural morphism $X_i'\to X_i\to S$, the open $b^{-1}(V_i)$, the
  proper morphism $b^{-1}(V_i)\xrightarrow{\ \sim\ }V_i\xrightarrow{\psi_i}U$,
  and the copy of $U_i$ obtained from the selected isomorphism over $V_i$.
  Its boundary loci are the closures in $X_i'$ of the inverse images
  $b^{-1}(\overline Z_{i,k}\cap V_i)$.
\end{definition}

\begin{lemma}[Properties of the transported side]
  \label{lemma:nagata-two-open-side-transport}
  \uses{def:nagata-two-open-side-transport}
  \lean{AlgebraicGeometry.CompactifiedProperExtension.admissibleBlowup}

  The transported data is again a side: $X_i'\to S$ is proper, the
  transported $\psi_i$ is proper and extends $U_i\hookrightarrow U$, it is an
  isomorphism over $U_i$, and the transported copy of $U_i$ is
  scheme-theoretically dense in $X_i'$.
\end{lemma}

\begin{proof}
  \uses{lemma:nagata-blowup-proper,
    lemma:nagata-admissible-blowup-over-open,
    lemma:nagata-admissible-blowup-preserves-density}
  An admissible blow-up is proper by
  \Cref{lemma:nagata-blowup-proper}, so the composite $X_i'\to S$ is proper,
  and the transported $\psi_i$ differs from $\psi_i$ by the isomorphism over
  $V_i$ of \Cref{lemma:nagata-admissible-blowup-over-open}, hence is proper,
  extends $U_i\hookrightarrow U$ and is an isomorphism over $U_i$.  The copy
  of $U_i$ inside $V_i$ is a scheme-theoretically dense open of $X_i$
  contained in $V_i$; its inverse image is therefore scheme-theoretically
  dense in $X_i'$ by
  \Cref{lemma:nagata-admissible-blowup-preserves-density}, and it is exactly
  the transported copy because $b$ is an isomorphism over $V_i$.
\end{proof}

\begin{lemma}[Disjoint boundary closures persist along a transport]
  \label{lemma:nagata-two-open-side-transport-prepared}
  \uses{def:nagata-two-open-side-transport,
    def:nagata-two-open-prepared-boundaries}
  \lean{AlgebraicGeometry.CompactifiedProperExtension.Prepared.admissibleBlowup}

  If the boundary closures $\overline Z_{i,i}$ and $\overline Z_{i,j}$ of a
  side are disjoint, then the boundary closures of the side transported along
  a $V_i$-admissible blow-up $b$ are disjoint as well.
\end{lemma}

\begin{proof}
  The boundary closures of the transported side are the closures of the
  inverse images under $b$ of the boundary loci of the original side.  A
  continuous map carries the closure of an inverse image into the inverse
  image of the closure, so these are contained in
  $b^{-1}(\overline Z_{i,i})$ and $b^{-1}(\overline Z_{i,j})$, which are
  disjoint because $\overline Z_{i,i}$ and $\overline Z_{i,j}$ are.
\end{proof}

\begin{definition}[The two-open comparison space]
  \label{def:nagata-two-open-comparison-space}
  \uses{def:nagata-two-open-prepared-boundaries,
    def:nagata-scheme-theoretic-image}
  \lean{AlgebraicGeometry.CompactifiedProperExtension.comparisonSpace}

  Define
  \[
    V_{12}=V_1\times_UV_2.
  \]
  The separatedness of $U\to S$ gives a closed immersion
  $V_{12}\to V_1\times_SV_2$, hence an immersion
  $V_{12}\to X_1\times_SX_2$.  Define
  \[
    X_{12}=\operatorname{SIm}
      \bigl(V_{12}\longrightarrow X_1\times_SX_2\bigr)
  \]
  and let $p_i\colon X_{12}\to X_i$ be the restrictions of the two product
  projections.
\end{definition}


\begin{lemma}[Properness of the comparison projections]
  \label{lemma:nagata-comparison-projections-proper}
  \uses{def:nagata-two-open-comparison-space,
    lemma:nagata-prepared-boundaries,
    lemma:nagata-scheme-image-universal}
  \lean{AlgebraicGeometry.CompactifiedProperExtension.isProper_comparisonFst}

  Both projections $p_i\colon X_{12}\to X_i$ are proper.
\end{lemma}

\begin{proof}
  Apply \Cref{lemma:nagata-scheme-image-universal} to
  $V_{12}\to X_1\times_SX_2$.  It gives the canonical closed immersion
  $\iota_{12}\colon X_{12}\hookrightarrow X_1\times_SX_2$ occurring in
  \Cref{def:nagata-two-open-comparison-space}.  For
  $\{i,j\}=\{1,2\}$, the morphism $X_j\to S$ is proper by
  \Cref{lemma:nagata-prepared-boundaries}; hence the projection
  $\pi_i\colon X_1\times_SX_2\to X_i$ is its proper base change.  The
  equality $p_i=\pi_i\circ\iota_{12}$ expresses $p_i$ as a composite of
  proper morphisms, so $p_i$ is proper.
\end{proof}

\begin{definition}[Distinguished comparison boundary]
  \label{def:nagata-two-open-comparison-boundary}
  \uses{def:nagata-two-open-boundary-loci,
    def:nagata-two-open-comparison-space}
  \lean{AlgebraicGeometry.exists_prepared_keyDisjoint_twoOpenSide}

  Define
  \[
    Z_{12,2}=(p_1|_{V_{12}})^{-1}(Z_{1,2}).
  \]
\end{definition}

\begin{lemma}[The two descriptions of the comparison boundary agree]
  \label{lemma:nagata-two-open-comparison-boundary}
  \uses{def:nagata-two-open-comparison-boundary,
    def:nagata-two-open-comparison-space}
  \lean{AlgebraicGeometry.exists_prepared_keyDisjoint_twoOpenSide}

  Inside $V_{12}$ one has
  \[
    Z_{12,2}=(p_2|_{V_{12}})^{-1}(Z_{2,2}).
  \]
\end{lemma}

\begin{proof}
  Let $z=(v_1,v_2)$ be a point of $V_{12}=V_1\times_UV_2$.  The fibre-product
  equality gives $\psi_1(v_1)=\psi_2(v_2)=u\in U$.  Since
  $Z_{1,2}=V_1\times_UZ_2$ and $Z_{2,2}=V_2\times_UZ_2$, membership in either
  inverse image is equivalent to $u\in Z_2$.  Scheme-theoretically, both
  sides are the pullback of the closed subscheme $Z_2\hookrightarrow U$ along
  the equal composites $V_{12}\rightrightarrows U$.  The canonical
  associativity isomorphisms for fibre products therefore identify the two
  closed subschemes.
\end{proof}


\begin{lemma}[Restriction of the comparison image over the proper-extension open]
  \label{lemma:nagata-comparison-restrict-proper-open}
  \uses{def:nagata-two-open-comparison-space,
    lemma:nagata-scheme-image-restrict-open,
    lemma:nagata-proper-to-separated-target,
    lemma:nagata-scheme-image-of-closed-immersion}
  \lean{AlgebraicGeometry.exists_prepared_keyDisjoint_twoOpenSide}

  For \(\{i,j\}=\{1,2\}\), the canonical morphism
  \[
    V_{12}=V_1\times_UV_2\longrightarrow p_i^{-1}(V_i)
  \]
  is an isomorphism.  Under this identification,
  \(p_i^{-1}(V_i)\to V_i\) is the base change of
  \(\psi_j\colon V_j\to U\) along \(\psi_i\colon V_i\to U\).
\end{lemma}

\begin{proof}
  \uses{lemma:nagata-scheme-image-restrict-open,
    lemma:nagata-proper-to-separated-target,
    lemma:nagata-scheme-image-of-closed-immersion}
  Restrict the scheme-theoretic image defining \(X_{12}\) to the open target
  \(V_i\times_SX_j\subseteq X_i\times_SX_j\).  The inverse image of this
  open in \(V_{12}\) is all of \(V_{12}\), so
  \Cref{lemma:nagata-scheme-image-restrict-open} identifies
  \(p_i^{-1}(V_i)\) with the scheme-theoretic image of
  \[
    V_{12}\longrightarrow V_i\times_SX_j.
  \]

  The map \(V_{12}\to V_i\times_SV_j\) is the pullback of the diagonal
  \(U\to U\times_SU\); it is therefore a closed immersion because
  \(U\to S\) is separated.  Composing with the open immersion
  \(V_i\times_SV_j\to V_i\times_SX_j\) shows that the displayed map is an
  immersion.  It is also proper: its composite to \(V_i\) is the base change
  of the proper morphism \(\psi_j\), while
  \(V_i\times_SX_j\to V_i\) is separated; apply
  \Cref{lemma:nagata-proper-to-separated-target}.  A proper immersion is a
  closed immersion, so
  \Cref{lemma:nagata-scheme-image-of-closed-immersion} identifies its
  scheme-theoretic image with \(V_{12}\) itself.
  The first projection of the fibre square is the displayed map to $V_i$,
  and its second component is the map to $V_j$.  The universal property of
  the fibre product gives the inverse identification with the base change of
  $\psi_j$ along $\psi_i$, so the description of the projection is exactly
  the one asserted.
\end{proof}

\begin{definition}[Improved first comparison projection]
  \label{def:nagata-two-open-improved-comparison}
  \uses{def:nagata-two-open-comparison-space,
    def:nagata-two-open-comparison-boundary,
    def:nagata-isomorphism-near-closure,
    def:nagata-strict-transform,
    def:nagata-maximal-isomorphism-open}
  \lean{AlgebraicGeometry.exists_prepared_keyDisjoint_twoOpenSide}

  Let $N_2=\operatorname{IsoOpen}(\psi_2)\subseteq U$.  Apply
  $\mathsf{IsoBl}_{V_1,Z_{1,2}}(p_1)$ to the proper projection
  $p_1\colon X_{12}\to X_1$.  Define $X_1^{\mathrm{cmp}}$ to be this
  constructed blow-up and define $X_{12}^{\mathrm{cmp}}$ to be the strict
  transform of $X_{12}$.  Transport $V_1$, all four boundary loci, and the
  projection maps along the selected isomorphism over $V_1$, and thereafter
  retain the notation $X_1,X_{12},p_i$ for the transported data.
\end{definition}

\begin{lemma}[Properties of the improved first comparison projection]
  \label{lemma:nagata-comparison-isIso-near-boundary}
  \uses{def:nagata-two-open-improved-comparison}
  \lean{AlgebraicGeometry.exists_prepared_keyDisjoint_twoOpenSide}

  The constructed map $X_1^{\mathrm{cmp}}\to X_1^{\mathrm{old}}$ is
  $V_1$-admissible, the transported $p_1$ is an isomorphism over an open
  neighbourhood of $\overline Z_{1,2}$, and all previously established
  boundary-disjointness relations still hold.
\end{lemma}

\begin{proof}
  \uses{lemma:nagata-boundary-loci,
    lemma:nagata-comparison-projections-proper,
    lemma:nagata-comparison-restrict-proper-open,
    lemma:nagata-isomorphism-near-closure,
    lemma:nagata-scheme-image-strict-transform,
    lemma:nagata-boundary-disjoint-persists,
    lemma:nagata-maximal-isomorphism-open}
  By \Cref{lemma:nagata-comparison-restrict-proper-open}, the restriction of
  $p_1$ over $V_1$ is the base change
  $V_1\times_UV_2\to V_1$ of $\psi_2$.  The maximal isomorphism open $N_2$
  contains $Z_2$ by \Cref{lemma:nagata-boundary-loci} and
  \Cref{lemma:nagata-maximal-isomorphism-open}; hence $p_1$ is an
  isomorphism over $\psi_1^{-1}(N_2)$, an open containing $Z_{1,2}$.

  The subset $Z_{1,2}$ is closed in $V_1$, and $p_1$ is proper by
  \Cref{lemma:nagata-comparison-projections-proper}.  Therefore the property
  of $\mathsf{IsoBl}_{V_1,Z_{1,2}}(p_1)$ gives the asserted isomorphism near
  the closure.  The new comparison scheme is the strict transform by
  \Cref{lemma:nagata-scheme-image-strict-transform}.  Finally
  \Cref{lemma:nagata-boundary-disjoint-persists} preserves every earlier
  disjointness relation under this further $V_1$-admissible blow-up.
\end{proof}


\begin{lemma}[Density over the prepared boundary]
  \label{lemma:nagata-comparison-boundary-density}
  \lean{AlgebraicGeometry.exists_prepared_keyDisjoint_twoOpenSide}

  The subset $Z_{12,2}$ is dense in
  $p_1^{-1}(\overline Z_{1,2})$.
\end{lemma}

\begin{proof}
  \uses{lemma:nagata-comparison-isIso-near-boundary}
  Choose an open neighbourhood $N\subseteq X_1$ of
  $\overline Z_{1,2}$ over which $p_1$ is an isomorphism, as supplied by
  \Cref{lemma:nagata-comparison-isIso-near-boundary}.  Then
  $p_1^{-1}(\overline Z_{1,2})$ lies in $p_1^{-1}(N)$, and the restriction
  \[
    p_1^{-1}(\overline Z_{1,2})\longrightarrow\overline Z_{1,2}
  \]
  is an isomorphism.  By the definition of $Z_{12,2}$, its inverse image under
  this isomorphism is $Z_{1,2}$.  Since $Z_{1,2}$ is dense in its closure
  $\overline Z_{1,2}$, the inverse image $Z_{12,2}$ is dense in
  $p_1^{-1}(\overline Z_{1,2})$.
\end{proof}

\begin{lemma}[Key boundary disjointness]
  \label{lemma:nagata-key-boundary-disjointness}
  \lean{AlgebraicGeometry.CompactifiedProperExtension.KeyDisjoint}

  Inside $X_1\times_SX_2$ one has
  \[
    X_{12}\cap
    (\overline Z_{1,2}\times_S\overline Z_{2,1})
    =\varnothing.
  \]
\end{lemma}

\begin{proof}
  \uses{lemma:nagata-prepared-boundaries,
    lemma:nagata-comparison-boundary-density,
    lemma:nagata-two-open-comparison-boundary}
  Put $C=p_1^{-1}(\overline Z_{1,2})$.  By
  \Cref{lemma:nagata-comparison-boundary-density}, $Z_{12,2}$ is dense in
  $C$.  The definition of the comparison boundary and
  \Cref{lemma:nagata-two-open-comparison-boundary} give
  \[
    p_2(Z_{12,2})\subseteq Z_{2,2}\subseteq\overline Z_{2,2}.
  \]
  Equivalently, $Z_{12,2}$ is contained in the closed inverse image
  $p_2^{-1}(\overline Z_{2,2})$.  Because it is dense in $C$, this closed
  inverse image also contains $C$.  Thus
  $p_2(C)\subseteq\overline Z_{2,2}$.
  The two closed subschemes $\overline Z_{2,2}$ and
  $\overline Z_{2,1}$ are disjoint by
  \Cref{lemma:nagata-prepared-boundaries}.  Thus no point of $C$ can map into
  $\overline Z_{2,1}$, which is precisely the asserted empty intersection.
\end{proof}

\begin{definition}[Intermediate two-open schemes]
  \label{def:nagata-two-open-intermediate-schemes}
  \uses{def:nagata-two-open-improved-comparison,
    def:nagata-two-open-prepared-boundaries}
  \lean{AlgebraicGeometry.CompactifiedProperExtension.intermediate}

  For $\{i,j\}=\{1,2\}$, define
  \[
    X_i^\circ=X_i\setminus\overline Z_{i,j}
  \]
  and glue the open subschemes $U_i\subseteq U$ and
  $U_i\subseteq X_i^\circ$ by the identity.  Denote the result by
  \[
    W_i=U\amalg_{U_i}X_i^\circ.
  \]
  The construction includes the two open immersions
  $U\hookrightarrow W_i$ and $X_i^\circ\hookrightarrow W_i$ and the induced
  structural map $W_i\to S$.
\end{definition}


\begin{lemma}[Closed image of the proper extension graph]
  \label{lemma:nagata-proper-extension-graph-closed}
  \uses{def:nagata-two-open-graph-extensions,
    lemma:nagata-proper-to-separated-target}
  \lean{AlgebraicGeometry.CompactifiedProperExtension.isProper_graph}

  The morphism
  \[
    \gamma_i=(\psi_i,\mathrm{incl})\colon
    V_i\longrightarrow U\times_SX_i
  \]
  is proper.  In particular, its image is closed.
\end{lemma}

\begin{proof}
  By \Cref{def:nagata-two-open-graph-extensions}, the morphism
  $\psi_i\colon V_i\to U$ is proper.  Regard
  $\gamma_i\colon V_i\to U\times_SX_i$ as a morphism over $U$.  The target
  projection $U\times_SX_i\to U$ is separated because it is the base change
  of the separated morphism $X_i\to S$.  Apply
  \Cref{lemma:nagata-proper-to-separated-target} with
  $T=V_i$, $Y=U\times_SX_i$, and $B=U$.  It makes $\gamma_i$ proper.
  Proper morphisms are universally closed, so the image of $\gamma_i$ is
  closed.
\end{proof}

\begin{lemma}[Constructibility of the overlap image]
  \label{lemma:nagata-intermediate-overlap-constructible}
  \uses{def:nagata-two-open-intermediate-schemes}
  \lean{AlgebraicGeometry.CompactifiedProperExtension.isClosedImmersion_intermediateGraph}

  The image of
  \[
    U_i\longrightarrow
    U\times_S(X_i\setminus\overline Z_{i,j})
  \]
  is constructible.
\end{lemma}

\begin{proof}
  The source and target are of finite type over the Noetherian scheme $S$.
  The displayed morphism is an immersion and hence of finite type.  Chevalley's
  theorem makes its image constructible.
\end{proof}

\begin{lemma}[A specialization leaving $U_i$ meets the forbidden boundary]
  \label{lemma:nagata-intermediate-specialization-boundary}
  \uses{def:nagata-two-open-intermediate-schemes,
    def:nagata-two-open-prepared-boundaries}
  \lean{AlgebraicGeometry.CompactifiedProperExtension.isClosedImmersion_intermediateGraph}

  Let a point in the image of $U_i$ specialize in $U\times_SX_i$ to a point
  $(u,x)$.  There is a point $v\in V_i$ with
  $\psi_i(v)=u$ and image $x$ in $X_i$.  If $u\notin U_i$, then
  $v\in Z_{i,j}$ and $x\in\overline Z_{i,j}$.
\end{lemma}

\begin{proof}
  \uses{lemma:nagata-proper-extension-graph-closed,lemma:nagata-boundary-loci}
  The image of $U_i$ is contained in the image of
  $\gamma_i\colon V_i\to U\times_SX_i$.  The latter image is closed by
  \Cref{lemma:nagata-proper-extension-graph-closed}; hence the specialization
  $(u,x)$ is $\gamma_i(v)$ for some $v\in V_i$.  Thus
  $\psi_i(v)=u$ and the image of $v$ in $X_i$ is $x$.  If $u\notin U_i$, then
  $u\in Z_j$, so $v\in Z_{i,j}=\psi_i^{-1}(Z_j)$.  Therefore
  $x\in Z_{i,j}\subseteq\overline Z_{i,j}$.
\end{proof}

\begin{lemma}[Specialization stability inside the boundary complement]
  \label{lemma:nagata-intermediate-overlap-specialization-closed}
  \uses{def:nagata-two-open-intermediate-schemes,
    lemma:nagata-intermediate-specialization-boundary,
    lemma:nagata-two-open-graph-extensions}
  \lean{AlgebraicGeometry.CompactifiedProperExtension.isClosedImmersion_intermediateGraph}

  The image of $U_i$ in
  $U\times_S(X_i\setminus\overline Z_{i,j})$ is stable under specialization.
\end{lemma}

\begin{proof}
  \uses{lemma:nagata-intermediate-specialization-boundary}
  Let a point of the image specialize to $(u,x)$ in the displayed target.  If
  $u\notin U_i$, then
  \Cref{lemma:nagata-intermediate-specialization-boundary} gives
  $x\in\overline Z_{i,j}$, contradicting the definition of the target.
  Hence $u\in U_i$.  By
  \Cref{lemma:nagata-two-open-graph-extensions}, the identity
  $\psi_i^{-1}(U_i)=U_i$ forces the point $v$ supplied by
  \Cref{lemma:nagata-intermediate-specialization-boundary} to lie in $U_i$.
  The chosen isomorphism $U_i\cong\psi_i^{-1}(U_i)$ then identifies $x$ with
  the graph point of $u$, so $(u,x)$ is again in the image of $U_i$.
\end{proof}

\begin{lemma}[Closedness of the intermediate overlap]
  \label{lemma:nagata-intermediate-overlap-closed}
  \uses{def:nagata-two-open-intermediate-schemes}
  \lean{AlgebraicGeometry.CompactifiedProperExtension.isClosedImmersion_intermediateGraph}

  The canonical immersion
  \[
    U_i\longrightarrow
    U\times_S(X_i\setminus\overline Z_{i,j})
  \]
  is a closed immersion.
\end{lemma}

\begin{proof}
  \uses{lemma:nagata-intermediate-overlap-constructible,
    lemma:nagata-intermediate-overlap-specialization-closed}
  By \Cref{lemma:nagata-intermediate-overlap-constructible}, the image is
  constructible.  By
  \Cref{lemma:nagata-intermediate-overlap-specialization-closed}, it is stable
  under specialization.  The target is Noetherian, so the image is closed.
  Factor the given immersion as a closed immersion into an open subscheme
  $T'\subseteq T$, where
  $T=U\times_S(X_i\setminus\overline Z_{i,j})$.  Over $T'$ the map is a
  closed immersion by construction.  Over the open complement of its closed
  image in $T$, its inverse image is empty, hence is also a closed immersion.
  These two opens cover $T$: every point outside $T'$ is outside the image.
  Since being a closed immersion is local on the target, the original
  immersion is a closed immersion.
\end{proof}

\begin{lemma}[Separatedness of the intermediate schemes]
  \label{lemma:nagata-intermediate-separated}
  \uses{def:nagata-two-open-intermediate-schemes}
  \lean{AlgebraicGeometry.CompactifiedProperExtension.isSeparated_intermediateHom}

  Each $W_i\to S$ is separated and of finite type.
\end{lemma}

\begin{proof}
  \uses{lemma:nagata-intermediate-overlap-closed,
    lemma:nagata-separated-two-open-cover}
  The opens $U$ and $X_i\setminus\overline Z_{i,j}$ cover $W_i$.  The first
  is separated and of finite type over $S$ by the section hypotheses.  The
  second has those properties because it is open in the proper scheme
  $X_i\to S$.  Their overlap map is the closed immersion supplied by
  \Cref{lemma:nagata-intermediate-overlap-closed}.  Apply
  \Cref{lemma:nagata-separated-two-open-cover} to these two opens to obtain
  separatedness of $W_i\to S$.  Finite type is local on this finite open
  cover, so $W_i\to S$ is also of finite type.
\end{proof}

\begin{lemma}[The pair of valuation extensions factors through $X_{12}$]
  \label{lemma:nagata-valuation-pair-factors-comparison}
  \uses{def:nagata-two-open-comparison-space,
    lemma:nagata-valuation-factor-closed}
  \lean{AlgebraicGeometry.CompactifiedProperExtension.KeyDisjoint.avoids}

  Let $A$ be a valuation ring with fraction field $K$.  Suppose
  $\operatorname{Spec}(K)\to U_1\cap U_2$ is an $S$-morphism and let
  $g_i\colon\operatorname{Spec}(A)\to X_i$ be extensions.  Then the pair
  \[
    (g_1,g_2)\colon\operatorname{Spec}(A)\to X_1\times_SX_2
  \]
  factors through $X_{12}$.
\end{lemma}

\begin{proof}
  The two generic maps have the same image in $U$, so their pair lands in
  $V_1\times_UV_2=V_{12}$.  Composing with the canonical image-factor map
  gives a factorization of the generic pair through $X_{12}$.  The map
  $A\to K$ is injective because a valuation ring is a domain.  Apply
  \Cref{lemma:nagata-valuation-factor-closed} to the closed immersion
  $X_{12}\hookrightarrow X_1\times_SX_2$ and the morphism
  $(g_1,g_2)\colon\operatorname{Spec}(A)\to X_1\times_SX_2$.  It extends the
  generic factorization to all of $\operatorname{Spec}(A)$.
\end{proof}

\begin{lemma}[An open containing the closed point of a valuation spectrum is total]
  \label{lemma:nagata-valuation-open-closed-point}
  \lean{AlgebraicGeometry.CompactifiedProperExtension.KeyDisjoint.avoids}

  Let $A$ be a valuation ring and let $O\subseteq\operatorname{Spec}(A)$ be
  open.  If the closed point $\mathfrak m_A$ belongs to $O$, then
  $O=\operatorname{Spec}(A)$.
\end{lemma}

\begin{proof}
  Choose a basic open $D(a)$ with $\mathfrak m_A\in D(a)\subseteq O$.
  Then $a\notin\mathfrak m_A$, hence $a$ is a unit of the local ring $A$ and
  $D(a)=\operatorname{Spec}(A)$.
\end{proof}

\begin{lemma}[One valuation extension avoids the forbidden boundary]
  \label{lemma:nagata-one-valuation-avoids-boundary}
  \uses{def:nagata-two-open-comparison-space,
    def:nagata-two-open-intermediate-schemes}
  \lean{AlgebraicGeometry.CompactifiedProperExtension.KeyDisjoint.avoids}

  Let $A$ be a valuation ring with fraction field $K$.  Let
  $\operatorname{Spec}(K)\to U_1\cap U_2$ be an $S$-morphism and let
  $g_i\colon\operatorname{Spec}(A)\to X_i$ be extensions.  At least one
  $g_i$ factors through $X_i\setminus\overline Z_{i,j}$, and hence determines
  a map $\operatorname{Spec}(A)\to W_i$ extending the given generic map.
\end{lemma}

\begin{proof}
  \uses{lemma:nagata-key-boundary-disjointness,
    lemma:nagata-valuation-pair-factors-comparison,
    lemma:nagata-valuation-open-closed-point}
  Otherwise the closed point would map into $\overline Z_{1,2}$ under $g_1$
  and into $\overline Z_{2,1}$ under $g_2$.  By
  \Cref{lemma:nagata-valuation-pair-factors-comparison}, the pair factors
  through $X_{12}$, contradicting
  \Cref{lemma:nagata-key-boundary-disjointness}.  If the closed point avoids
  one of the two boundaries, its inverse image of the corresponding open
  complement is an open subset of $\operatorname{Spec}(A)$ containing the
  closed point.  By
  \Cref{lemma:nagata-valuation-open-closed-point} this open is all of
  $\operatorname{Spec}(A)$.  The associated $g_i$ therefore factors
  through $X_i\setminus\overline Z_{i,j}$, and the open immersion
  $X_i\setminus\overline Z_{i,j}\hookrightarrow W_i$ supplies the asserted
  extension.
\end{proof}

\begin{lemma}[Existence part of the proper valuative criterion]
  \label{lemma:nagata-proper-valuative-criterion-existence-api}
  \lean{AlgebraicGeometry.UniversallyClosed.eq_valuativeCriterion}

  Let $f\colon Y\to S$ be proper.  For every valuation ring $A$ with
  fraction field $K$, every $S$-morphism
  $\operatorname{Spec}(K)\to Y$ lying over a morphism
  $\operatorname{Spec}(A)\to S$ extends to an $S$-morphism
  $\operatorname{Spec}(A)\to Y$.
\end{lemma}

\begin{proof}
  \uses{lemma:nagata-local-domain-valuation-dominates,
    lemma:nagata-valuation-intersection}
  Let $G$ be the closure of the graph of the given generic map in
  $Y\times_S\operatorname{Spec}(A)$.  Properness of $Y\to S$ makes
  $G\to\operatorname{Spec}(A)$ proper, and the generic point of $G$ maps to
  the generic point of $\operatorname{Spec}(A)$.  Choose an irreducible
  component $G_0$ dominating $\operatorname{Spec}(A)$ and replace it by its
  reduced induced scheme.  At a point over the closed point of
  $\operatorname{Spec}(A)$, the local ring of $G_0$ is a local domain $B$
  with fraction field $L$ containing $K$.  By
  \Cref{lemma:nagata-local-domain-valuation-dominates}, choose a valuation
  ring $C\subset L$ dominating $B$.  The resulting map
  $\operatorname{Spec}(C)\to G_0$ has generic restriction equal to the
  prescribed map.  Its restriction on $K$ and the map to $A$ agree; the
  intersection identity
  $K\cap C=A$ (applied to $A\subset K\subset L$) therefore makes the local
  homomorphism factor through $A$.  This gives a morphism
  $\operatorname{Spec}(A)\to G\to Y$ extending the generic map.  The
  construction is compatible with the structural maps to $S$ by the
  definition of the graph closure.
\end{proof}

\begin{lemma}[Uniqueness of a valuation lift for a separated morphism]
  \label{lemma:nagata-separated-valuation-lift-unique}
  \lean{AlgebraicGeometry.IsSeparated.valuativeCriterion}

  Let $f\colon Y\to S$ be separated, let $A$ be a valuation ring with
  fraction field $K$, and let $h_1,h_2\colon\operatorname{Spec}(A)\to Y$ be
  $S$-morphisms whose restrictions to $\operatorname{Spec}(K)$ coincide.
  Then $h_1=h_2$.
\end{lemma}

\begin{proof}
  The equalizer of $h_1$ and $h_2$ is the inverse image of the diagonal
  $\Delta_f\colon Y\hookrightarrow Y\times_S Y$.  Since $f$ is separated,
  this is a closed subscheme of $\operatorname{Spec}(A)$.  It contains the
  generic point because the two restrictions to $\operatorname{Spec}(K)$ are
  equal.  A valuation ring is a domain, so its generic point is the prime
  $(0)$, whose closure is all of $\operatorname{Spec}(A)$.  Hence the closed
  equalizer is the whole spectrum, and the two morphisms are equal.
\end{proof}

\begin{lemma}[Valuative extension to a proper separated morphism]
  \label{lemma:nagata-proper-separated-valuation-extension}
  \lean{AlgebraicGeometry.IsProper.eq_valuativeCriterion}

  Let $f\colon Y\to S$ be proper and separated.  Every commutative square
  \[
    \operatorname{Spec}(K)\longrightarrow Y
    \qquad\text{over}\qquad
    \operatorname{Spec}(A)\longrightarrow S
  \]
  with $A$ a valuation ring and $K=\operatorname{Frac}(A)$ admits a unique
  lift $\operatorname{Spec}(A)\to Y$.
\end{lemma}

\begin{proof}
  \uses{lemma:nagata-proper-valuative-criterion-existence-api,
    lemma:nagata-separated-valuation-lift-unique}
  Existence is the output of
  \Cref{lemma:nagata-proper-valuative-criterion-existence-api}.  If $h$ and
  $h'$ are two such lifts, their generic restrictions are the prescribed
  morphism, so
  \Cref{lemma:nagata-separated-valuation-lift-unique} gives $h=h'$.  Thus the
  extension exists and is unique.
\end{proof}

\begin{definition}[Final common modification for the two-open construction]
  \label{def:nagata-two-open-final-modification}
  \uses{def:nagata-two-open-intermediate-schemes,
    def:nagata-common-modification}
  \lean{AlgebraicGeometry.CommonModification}

  Apply the common-modification construction to the two open immersions
  $U\hookrightarrow W_1$ and $U\hookrightarrow W_2$.  Denote its output by
  \[
    W'_i\longrightarrow W_i,
    \qquad
    W=\mathsf{CM}(W_1,W_2;U),
  \]
  with the two constructed open immersions $W'_i\hookrightarrow W$ and their
  common open subscheme $U$.
\end{definition}

\begin{lemma}[Structural properties of the final common modification]
  \label{lemma:nagata-two-open-final-modification}
  \uses{def:nagata-two-open-final-modification}
  \lean{AlgebraicGeometry.exists_commonModification}

  Each $W'_i\to W_i$ is $U$-admissible; the maps $W'_i\hookrightarrow W$
  are open immersions covering $W$; their copies of $U$ agree; and $W\to S$
  is separated and of finite type.  In particular the induced map
  $j_W\colon U\to W$ is an open immersion.
\end{lemma}

\begin{proof}
  \uses{lemma:nagata-intermediate-separated,
    lemma:nagata-common-modification}
  The schemes $W_i$ are separated and of finite type by
  \Cref{lemma:nagata-intermediate-separated}.  Thus all hypotheses of the
  common-modification property apply, and
  \Cref{lemma:nagata-common-modification} gives every stated assertion.
  The map $j_W$ is the composite of $U\hookrightarrow W'_i$ with either open
  immersion $W'_i\hookrightarrow W$; the two composites coincide by the
  gluing construction.
\end{proof}

\begin{definition}[Two-open compactification candidate]
  \label{def:nagata-two-open-compactification-candidate}
  \uses{def:nagata-two-open-final-modification}
  \lean{AlgebraicGeometry.HasCompactification.of_two_open_cover}

  Define $\mathsf C_{12}(U/S)$ to be the constructed factorization
  \[
    U\xrightarrow{j_W}W\longrightarrow S.
  \]
\end{definition}

\begin{lemma}[Density and quasi-compactness under the finite modification diagram]
  \label{lemma:nagata-two-open-density-transport}
  \lean{AlgebraicGeometry.CompactifiedProperExtension.dense_range_intermediateU}

  In the finite diagram defining $W'_1,W'_2$ and $W$, assume that $U$ is
  scheme-theoretically dense in each initial ambient piece and that every
  intervening blow-up is $U$-admissible.  Then the induced map
  $U_1\cap U_2\to W$ is quasi-compact with dense image.
\end{lemma}

\begin{proof}
  \uses{lemma:nagata-admissible-blowup-preserves-density,
    lemma:nagata-admissible-blowup-comp,
    lemma:nagata-two-open-final-modification}
  The initial copies of $U$ are scheme-theoretically dense by hypothesis.
  Every graph, boundary-separation, and final modification in the displayed
  diagram is $U$-admissible, so
  \Cref{lemma:nagata-admissible-blowup-preserves-density} identifies the
  inverse image of $U$ at each stage as a scheme-theoretically dense open.
  The compatibility of the chosen copies of $U$ through a finite composite is
  the conclusion of
  \Cref{lemma:nagata-admissible-blowup-comp}.  Therefore the common copy of
  $U$ is dense in each of the final chart opens $W'_1$ and $W'_2$.

  Let $F\subseteq W$ be closed and contain the image of
  $U_1\cap U_2$.  Its inverse image in each open chart $W'_i$ is closed and
  contains the dense image of $U_1\cap U_2$, hence equals $W'_i$.  Since the
  two charts cover $W$ by
  \Cref{lemma:nagata-two-open-final-modification}, $F=W$, proving density.
  The source $U_1\cap U_2$ is quasi-compact because $S$ is Noetherian,
  $U\to S$ is of finite type, and every open subscheme of the resulting
  Noetherian scheme $U$ is quasi-compact.  Each map to $W'_i$ is a composite
  of a quasi-compact open
  immersion with finitely many blow-ups, which are quasi-compact; the finite
  open cover then proves that $U_1\cap U_2\to W$ is quasi-compact.
\end{proof}

\begin{lemma}[Density after the final common modification]
  \label{lemma:nagata-final-density}
  \uses{def:nagata-two-open-compactification-candidate}
  \lean{AlgebraicGeometry.HasCompactification.of_two_open_cover}

  The constructed map $U_1\cap U_2\to W$ is quasi-compact and has dense
  image.
\end{lemma}

\begin{proof}
  \uses{lemma:nagata-two-open-density-transport}
  Apply \Cref{lemma:nagata-two-open-density-transport} to the construction of
  $W'_1,W'_2,W$.  Its hypotheses hold because the selected compactifications
  are scheme-theoretically dense and every subsequent center is supported
  away from the preserved copy of $U$.  The asserted quasi-compactness and
  density follow.
\end{proof}

\begin{lemma}[Valuative lifting for the two-open candidate]
  \label{lemma:nagata-two-open-valuative-lifting}
  \uses{def:nagata-two-open-compactification-candidate}
  \lean{AlgebraicGeometry.HasCompactification.of_two_open_cover}

  Let $A$ be a valuation ring with fraction field $K$.  Every commutative
  diagram whose generic map
  $\operatorname{Spec}(K)\to W$ factors through $U_1\cap U_2$ has a unique
  lift $\operatorname{Spec}(A)\to W$ over $S$.
\end{lemma}

\begin{proof}
  \uses{lemma:nagata-one-valuation-avoids-boundary,
    lemma:nagata-proper-separated-valuation-extension,
    lemma:nagata-blowup-proper,
    lemma:nagata-separated-valuation-lift-unique,
    lemma:nagata-two-open-final-modification}
  By
  \Cref{lemma:nagata-proper-separated-valuation-extension}, properness and
  separatedness of $X_i\to S$ extend the generic map to maps
  $g_i\colon\operatorname{Spec}(A)\to X_i$.  By
  \Cref{lemma:nagata-one-valuation-avoids-boundary}, one extension factors
  through $X_i\setminus\overline Z_{i,j}$ and hence through $W_i$.  The
  generic map already lifts to the preserved open $U\subseteq W'_i$.
  The admissible blow-up $W'_i\to W_i$ is proper, hence separated, so
  \Cref{lemma:nagata-proper-separated-valuation-extension} supplies a lift
  to $W'_i$; its open map to $W$ gives the desired lift.  Any two branch
  choices have the same generic restriction.  Since $W\to S$ is separated by
  \Cref{lemma:nagata-two-open-final-modification},
  \Cref{lemma:nagata-separated-valuation-lift-unique} makes the two maps
  equal.
\end{proof}

\begin{lemma}[Properness of the two-open candidate]
  \label{lemma:nagata-two-open-candidate-proper}
  \uses{def:nagata-two-open-compactification-candidate}
  \lean{AlgebraicGeometry.HasCompactification.of_two_open_cover}

  The constructed structural morphism $W\to S$ is proper.
\end{lemma}

\begin{proof}
  \uses{lemma:nagata-final-density,
    lemma:nagata-two-open-valuative-lifting,
    lemma:nagata-two-open-final-modification,
    lemma:nagata-dense-source-valuative}
  The morphism $U_1\cap U_2\to W$ is quasi-compact and dense by
  \Cref{lemma:nagata-final-density}, and its valuation diagrams admit unique
  lifts by \Cref{lemma:nagata-two-open-valuative-lifting}.  The dense-source
  valuative criterion makes $W\to S$ universally closed.  It is separated
  and of finite type by
  \Cref{lemma:nagata-two-open-final-modification}; hence it is proper.
\end{proof}

\begin{theorem}[{\cite[\href{https://stacks.math.columbia.edu/tag/0F40}{Tag 0F40}]{stacks-project}}]
  \label{thm:nagata-two-open}
  \uses{def:nagata-compactification,def:nagata-two-open-compactification-candidate}
  \lean{AlgebraicGeometry.HasCompactification.of_two_open_cover}

  The constructed factorization $\mathsf C_{12}(U/S)$ is a compactification
  of $U\to S$.  Consequently $U\to S$ is compactifiable.
\end{theorem}

\begin{proof}
  \uses{lemma:nagata-two-open-final-modification,
    lemma:nagata-two-open-candidate-proper,
    lemma:nagata-open-immersion-qc-over-separated-target}
  By \Cref{lemma:nagata-two-open-final-modification}, the first map
  $j_W\colon U\to W$ is an open immersion and $W\to S$ is separated and of
  finite type.  It is proper by
  \Cref{lemma:nagata-two-open-candidate-proper}.  The open immersion $j_W$ is
  quasi-compact by
  \Cref{lemma:nagata-open-immersion-qc-over-separated-target}, applied to the
  triangle $U\to W\to S$.  The defining gluing maps give the factorization
  equality, so all fields of \Cref{def:nagata-compactification} are present.
\end{proof}


\section{The Noetherian case}

\begin{definition}[Relative affine space]
  \label{def:nagata-relative-affine-space}
  \lean{AlgebraicGeometry.AffineSpace}

  For a scheme $S$ and $n\in\mathbb N$, define
  \[
    \mathbb A^n_S=\operatorname{Spec}_S
      \bigl(\operatorname{Sym}_{\mathcal O_S}(\mathcal O_S^{\oplus n})\bigr).
  \]
  Concretely, on an affine open $V=\operatorname{Spec}(R)$ of $S$ this is
  $\operatorname{Spec}(R[T_1,\ldots,T_n])$; for a principal restriction
  $D(a)\subseteq V$ the transition map is induced by
  $R[T_1,\ldots,T_n]\to R_a[T_1,\ldots,T_n]$.
\end{definition}

\begin{lemma}[Relative affine-space charts and base change]
  \label{lemma:nagata-relative-affine-space-charts}
  \uses{def:nagata-relative-affine-space}
  \lean{AlgebraicGeometry.AffineSpace.chart_spec}

  The affine charts in \Cref{def:nagata-relative-affine-space} glue to a
  scheme over $S$.  Their transition maps satisfy the cocycle condition,
  and for every $T\to S$ there is a canonical isomorphism
  $\mathbb A^n_S\times_ST\cong\mathbb A^n_T$ compatible with the charts.
\end{lemma}

\begin{proof}
  On an affine $V=\operatorname{Spec}(R)$, localization gives
  $R[T_1,\ldots,T_n]_a\cong R_a[T_1,\ldots,T_n]$.  On a triple of principal
  opens both composites send each polynomial coefficient to its image in the
  localization at the product of the denominators, so the cocycle equality is
  literal.  Scheme gluing therefore produces the asserted scheme and map to
  $S$.  After base change to $T$, the same affine charts become
  $\operatorname{Spec}(\Gamma(W,\mathcal O_T)[T_1,\ldots,T_n])$ on affine
  $W\subseteq T$; the chartwise identifications agree on overlaps and hence
  glue to the stated base-change isomorphism.
\end{proof}

\begin{lemma}[The standard affine chart in relative projective space]
  \label{lemma:nagata-relative-affine-space-open-immersion}
  \uses{def:nagata-relative-affine-space,
    def:nagata-relative-projective-space,
    lemma:nagata-relative-proj-exists,
    lemma:nagata-relative-affine-space-charts}
  \lean{AlgebraicGeometry.AffineSpace.openImmersion_into_projectiveSpace}

  For every $n$ and $S$, the map induced by
  $T_i\mapsto T_i/T_0$ identifies $\mathbb A^n_S$ with the standard open
  $D_+(T_0)\subseteq\mathbb P^n_S$.  In particular it is an open immersion
  over $S$ and is compatible with base change.
\end{lemma}

\begin{proof}
  On $V=\operatorname{Spec}(R)$, the degree-zero part of
  $R[T_0,\ldots,T_n]_{T_0}$ is the polynomial ring
  $R[T_1/T_0,\ldots,T_n/T_0]$.  Thus the affine chart is canonically
  $\operatorname{Spec}(R[T_1,\ldots,T_n])$.  These identifications commute
  with localization in $R$, so they glue by
  \Cref{lemma:nagata-relative-affine-space-charts}; the relative-Proj
  restriction isomorphisms from
  \Cref{lemma:nagata-relative-proj-exists} identify the glued target with
  $D_+(T_0)$.  Standard Proj charts are open, proving the assertion.
\end{proof}

\begin{lemma}[Existence of an affine-space immersion for a quasi-affine morphism,
  {\cite[\href{https://stacks.math.columbia.edu/tag/04II}{Tag 04II}]{stacks-project}}]
  \label{lemma:nagata-quasi-affine-immersion-affine-space-exists}
  \uses{def:nagata-relative-affine-space}
  \lean{AlgebraicGeometry.exists_immersion_affineSpace_of_quasiAffine}

  Let $X\to S$ be a quasi-affine morphism of finite type, with $S$
  quasi-compact.  There is an integer $n$ and an immersion
  $X\to\mathbb A^n_S$ over $S$.
\end{lemma}

\begin{proof}
  Choose a finite affine cover $S=\bigcup_{j=1}^mS_j$.  Since $X\to S$ is of
  finite type and quasi-affine, each $X_j=X\times_SS_j$ has a finite affine
  cover by principal opens $D(g_{j\ell})$ in an affine $S_j$-scheme containing
  $X_j$.  On each $D(g_{j\ell})$ choose finitely many functions generating its
  coordinate algebra over $\Gamma(S_j,\mathcal O_S)$.  Adjoin the functions
  $g_{j\ell}$ themselves; the resulting finite tuple gives an immersion of
  $X_j$ into an affine space over $S_j$, because the principal opens are
  exactly the loci on which the chosen denominators are invertible and the
  coordinate maps are closed immersions there.  On $S_j\cap S_k$, replace the
  two tuples by their concatenation.  The induced maps agree after this
  replacement, since both are obtained from the same restriction maps on
  coordinate rings.  Taking the maximum tuple length and padding by zero
  coordinates gives a single integer $n$; the chartwise immersions glue to
  $X\to\mathbb A^n_S$.

  On every affine chart, the induced ring map is a localization of a
  surjection from a polynomial ring, hence gives a closed immersion into the
  corresponding principal open of affine space.  The principal opens cover
  $X_j$, so the chartwise map is an immersion.  On an overlap $S_j\cap S_k$,
  each coordinate is the restriction of the same chosen function on $X$;
  after concatenation and padding, the two polynomial-ring maps are therefore
  literally equal.  The affine-space gluing theorem then produces a global
  morphism, and the local immersion criterion proves that it is an immersion.
  This verifies both the coordinate construction and its overlap
  compatibility.
\end{proof}

\begin{definition}[Affine-space immersion attached to a quasi-affine scheme]
  \label{def:nagata-quasi-affine-immersion-affine-space}
  \uses{def:nagata-relative-affine-space,
    lemma:nagata-quasi-affine-immersion-affine-space-exists}
  \lean{AlgebraicGeometry.quasiAffineImmersionAffineSpace}

  Under the hypotheses of
  \Cref{lemma:nagata-quasi-affine-immersion-affine-space-exists}, choose one
  witness and denote its dimension and morphism by $n$ and
  \[
    \iota_X\colon X\longrightarrow\mathbb A^n_S
  \].
\end{definition}

\begin{lemma}[Property of the quasi-affine affine-space map,
  {\cite[\href{https://stacks.math.columbia.edu/tag/04II}{Tag 04II}]{stacks-project}}]
  \label{lemma:nagata-quasi-affine-immersion-affine-space}
  \uses{def:nagata-quasi-affine-immersion-affine-space,
    lemma:nagata-quasi-affine-immersion-affine-space-exists}
  \lean{AlgebraicGeometry.quasiAffineImmersionAffineSpace_spec}

  The constructed morphism
  $\iota_X\colon X\to\mathbb A^n_S$ is an immersion.
\end{lemma}

\begin{proof}
  The selected witness is supplied by
  \Cref{lemma:nagata-quasi-affine-immersion-affine-space-exists}; its
  specification field is precisely that $\iota_X$ is an immersion over $S$.
\end{proof}

\begin{lemma}[Affine finite-type morphisms are quasi-affine]
  \label{lemma:nagata-affine-finite-type-is-quasi-affine}
  \uses{lemma:nagata-quasi-affine-immersion-affine-space-exists}
  \lean{AlgebraicGeometry.IsQuasiAffine.of_isAffine}

  If $U\to S$ is affine and of finite type, then it is quasi-affine of finite
  type in the sense used in
  \Cref{lemma:nagata-quasi-affine-immersion-affine-space-exists}.
\end{lemma}

\begin{proof}
  An affine morphism is itself an affine scheme over $S$; the identity map of
  $U$ is an open immersion into that affine $S$-scheme.  This is the defining
  quasi-affine factorization, and the finite-type hypothesis is unchanged.
\end{proof}

\begin{definition}[Affine compactification construction]
  \label{def:nagata-affine-compactification}
  \uses{def:nagata-relative-affine-space,
    def:nagata-quasi-affine-immersion-affine-space,
    def:nagata-immersion-compactification,
    def:nagata-relative-projective-space}
  \lean{AlgebraicGeometry.Compactification.ofAffineConstruction}

  Let $U\to S$ be separated and of finite type, with $U$ affine.  Compose
  the constructed map $\iota_U\colon U\to\mathbb A^n_S$ with the standard
  open immersion $\mathbb A^n_S\hookrightarrow\mathbb P^n_S$, and define
  \[
    \mathsf C_{\mathrm{aff}}(U/S)
      =\mathsf C\bigl(U\longrightarrow\mathbb P^n_S\bigr)
  \]
  using \Cref{def:nagata-immersion-compactification}.
\end{definition}

\begin{proposition}[Property of the affine compactification construction]
  \label{prop:nagata-affine-compactification}
  \uses{def:nagata-affine-compactification}
  \lean{AlgebraicGeometry.Compactifiable.of_isAffine}

  The constructed factorization $\mathsf C_{\mathrm{aff}}(U/S)$ is a
  compactification of $U\to S$.
\end{proposition}

\begin{proof}
  \uses{lemma:nagata-affine-finite-type-is-quasi-affine,
    lemma:nagata-quasi-affine-immersion-affine-space,
    lemma:nagata-relative-affine-space-open-immersion,
    lemma:nagata-immersion-into-proper,
    lemma:nagata-projective-is-proper}
  Since $U\to S$ is affine and of finite type, it is quasi-affine by
  \Cref{lemma:nagata-affine-finite-type-is-quasi-affine}.  Apply
  \Cref{lemma:nagata-quasi-affine-immersion-affine-space} to obtain an
  immersion $\iota_U\colon U\to\mathbb A^n_S$.  Its composite with the
  standard open immersion of
  \Cref{lemma:nagata-relative-affine-space-open-immersion} is an immersion
  $U\to\mathbb P^n_S$.  By
  \Cref{lemma:nagata-projective-is-proper}, $\mathbb P^n_S\to S$ is proper.
  The immersion-into-proper construction of
  \Cref{lemma:nagata-immersion-into-proper} therefore supplies a
  compactification, and this is exactly
  $\mathsf C_{\mathrm{aff}}(U/S)$ by definition.
\end{proof}

\begin{lemma}[Existence of an affine neighbourhood through finitely many generic points,
  {\cite[\href{https://stacks.math.columbia.edu/tag/01ZX}{Tag 01ZX}]{stacks-project}}]
  \label{lemma:nagata-affine-neighbourhood-generic-points-exists}
  \lean{Scheme.exists_affineOpen_mem_and_genericPoints}

  Let $X$ be quasi-separated, let $x\in X$, and let
  $\eta_1,\ldots,\eta_r$ be generic points of finitely many distinct
  irreducible components.  There is an affine open of $X$ containing $x$ and
  every $\eta_i$.
\end{lemma}

\begin{proof}
  Reorder the components so that
  $x\in Z_1\cap\cdots\cap Z_m$ and
  $x\notin Z_{m+1},\ldots,Z_r$, where $Z_i$ is the component with generic
  point $\eta_i$.  Choose an affine open $W$ containing $x$ and disjoint from
  $Z_{m+1},\ldots,Z_r$.  Then $W$ contains
  $\eta_1,\ldots,\eta_m$.  By Stacks Project, Tag 01ZV, choose pairwise
  disjoint affine opens $A_i$ containing $\eta_i$ for $i>m$.  Since $X$ is
  quasi-separated, each $W\cap A_i$ is quasi-compact; it does not contain
  $\eta_i$.  Write $A_i=\operatorname{Spec}(R_i)$ and let
  $\mathfrak p_i$ be the minimal prime corresponding to $\eta_i$.  Stacks
  Project, Tag 00EV, applied to $W\cap A_i$ gives
  $f_i\in R_i\setminus\mathfrak p_i$ such that
  $D(f_i)\cap W=\varnothing$.  Replace $A_i$ by $D(f_i)$.  It still contains
  $\eta_i$, remains disjoint from the other $A_j$, and is disjoint from $W$.
  Thus
  \[
    W\amalg A_{m+1}\amalg\cdots\amalg A_r
  \]
  is a finite disjoint union of affine schemes, hence is affine, and it
  contains $x$ and every selected generic point.  This is the construction in
  Stacks Project, Tag 01ZX.
\end{proof}

\begin{definition}[Affine neighbourhood containing selected generic points]
  \label{def:nagata-affine-neighbourhood-generic-points}
  \uses{lemma:nagata-affine-neighbourhood-generic-points-exists}
  \lean{Scheme.affineOpen_mem_and_genericPoints}

  Let $X$ be quasi-separated, let $x\in X$, and let
  $\eta_1,\ldots,\eta_r$ be generic points of distinct irreducible
  components.  Choose one witness from
  \Cref{lemma:nagata-affine-neighbourhood-generic-points-exists} and denote
  its affine open by $A(x;\eta_1,\ldots,\eta_r)$.
\end{definition}

\begin{lemma}[Property of the selected affine neighbourhood,
  {\cite[\href{https://stacks.math.columbia.edu/tag/01ZX}{Tag 01ZX}]{stacks-project}}]
  \label{lemma:nagata-affine-neighbourhood-generic-points}
  \uses{def:nagata-affine-neighbourhood-generic-points,
    lemma:nagata-affine-neighbourhood-generic-points-exists}
  \lean{Scheme.affineOpen_mem_and_genericPoints_spec}

  The constructed open $A(x;\eta_1,\ldots,\eta_r)$ is affine and contains
  $x$ and every $\eta_i$.
\end{lemma}

\begin{proof}
  The selected witness exists by
  \Cref{lemma:nagata-affine-neighbourhood-generic-points-exists}; its
  specification says that the open is affine and contains $x$ and every
  selected generic point.
\end{proof}

\begin{lemma}[Existence of a finite affine cover with dense total intersection]
  \label{lemma:nagata-dense-affine-cover-exists}
  \uses{lemma:nagata-affine-neighbourhood-generic-points}
  \lean{Scheme.exists_finite_affineOpen_cover_dense_iInter}

  Let $X$ be quasi-compact and quasi-separated and assume that it has finitely
  many irreducible components, with generic points
  $\eta_1,\ldots,\eta_r$.  There is a finite affine open cover
  $(U_1,\ldots,U_n)$ whose total intersection is dense in $X$.
\end{lemma}

\begin{proof}
  For every $x\in X$, apply
  \Cref{lemma:nagata-affine-neighbourhood-generic-points} to obtain an affine
  open containing $x$ and every $\eta_i$.  Quasi-compactness gives a finite
  subcover.  Its total intersection contains the generic point of every
  irreducible component, and is therefore dense.
\end{proof}

\begin{definition}[Finite affine cover with selected dense intersection]
  \label{def:nagata-dense-affine-cover}
  \uses{lemma:nagata-dense-affine-cover-exists}
  \lean{Scheme.finiteAffineOpenCoverDenseIInter}

  Under the hypotheses of
  \Cref{lemma:nagata-dense-affine-cover-exists}, choose one witness and denote
  its finite family and its total intersection by
  \[
    \mathcal U_X=(U_1,\ldots,U_n),
    \qquad D_X=\bigcap_{i=1}^nU_i.
  \]
\end{definition}

\begin{lemma}[Properties of the selected dense affine cover]
  \label{lemma:nagata-dense-affine-cover}
  \uses{def:nagata-dense-affine-cover,
    lemma:nagata-dense-affine-cover-exists}
  \lean{Scheme.finiteAffineOpenCoverDenseIInter_spec}

  The $U_i$ form a finite affine open cover of $X$, and $D_X$ is dense in
  $X$.
\end{lemma}

\begin{proof}
  The selected witness is supplied by
  \Cref{lemma:nagata-dense-affine-cover-exists}.  Its specification fields say
  that the family is finite, consists of affine opens, covers $X$, and has
  dense total intersection.
\end{proof}

\begin{definition}[Partial unions of the selected affine cover]
  \label{def:nagata-partial-affine-unions}
  \uses{def:nagata-dense-affine-cover}
  \lean{Scheme.partialAffineOpenUnions}

  For $1\leq r\leq n$, define
  \[
    X_r=U_1\cup\cdots\cup U_r.
  \]
\end{definition}

\begin{lemma}[Density needed for induction on the affine cover]
  \label{lemma:nagata-partial-union-overlap-dense}
  \uses{def:nagata-partial-affine-unions}
  \lean{TopologicalSpace.dense_inter_of_dense_total_inter}

  For $r>1$, the overlap $X_{r-1}\cap U_r$ is dense in $X_r$.
\end{lemma}

\begin{proof}
  \uses{lemma:nagata-dense-affine-cover}
  The overlap contains $D_X$.  Since $X_r$ is open in $X$ and $D_X$ is dense
  in $X$, it is dense in $X_r$; therefore the larger subset
  $X_{r-1}\cap U_r$ is dense in $X_r$.
\end{proof}

\begin{definition}[Recursive Noetherian compactification construction]
  \label{def:nagata-noetherian-compactification}
  \uses{def:nagata-affine-compactification,
    def:nagata-dense-affine-cover,
    def:nagata-partial-affine-unions,
    def:nagata-two-open-compactification-candidate,
    prop:nagata-affine-compactification,thm:nagata-two-open,
    lemma:nagata-partial-union-overlap-dense}
  \lean{AlgebraicGeometry.noetherianCompactificationConstruction}

  Let $S$ be Noetherian and let $X\to S$ be separated and of finite type.
  Then $X$ is Noetherian, hence qcqs, and the finite generic-point argument
  supplies the selected dense affine cover
  $\mathcal U_X=(U_1,\ldots,U_n)$ from
  \Cref{def:nagata-dense-affine-cover}.  If $X$ is empty,
  let $\mathsf C_{\emptyset}(X/S)$ be the explicit factorization
  $\emptyset\hookrightarrow\emptyset\to S$ and set
  $\mathsf C_{\mathrm{Noeth}}(X/S)=\mathsf C_{\emptyset}(X/S)$.  Otherwise
  choose the cover with $n\geq1$ and define pairs
  $(\mathsf C_r,h_r)$ recursively, where $h_r$ is an explicit proof that
  $\mathsf C_r$ is a compactification of $X_r\to S$.  Set
  \[
    (\mathsf C_1,h_1)
      =\bigl(\mathsf C_{\mathrm{aff}}(U_1/S),
        \text{the witness supplied by }
        \Cref{prop:nagata-affine-compactification}\bigr).
  \]
  Given $(\mathsf C_{r-1},h_{r-1})$, apply
  \Cref{thm:nagata-two-open} to the open cover
  $X_r=X_{r-1}\cup U_r$, using $h_{r-1}$, the affine compactification witness
  for $U_r$ from \Cref{prop:nagata-affine-compactification}, and the dense
  overlap from \Cref{lemma:nagata-partial-union-overlap-dense}.  Choose the
  resulting candidate and its compactification proof as
  $(\mathsf C_r,h_r)$.  In the nonempty branch finally define
  \[
    \mathsf C_{\mathrm{Noeth}}(X/S)=\mathsf C_n,
    \qquad h_{\mathrm{Noeth}}=h_n.
  \]
  In the empty branch set $h_{\mathrm{Noeth}}$ to be the explicit
  compactification proof for $\mathsf C_{\emptyset}(X/S)$.
  The recursive datum is therefore a dependent finite sequence of pairs, not
  an unproved sequence of candidates.  At every stage also record that
  $X_r\to S$ is separated and of finite type, as follows by restricting the
  corresponding properties of $X\to S$ to the open $X_r\subseteq X$; these
  are the hypotheses used at the next application of the two-open theorem.
\end{definition}

\begin{theorem}[Nagata compactification over a Noetherian base]
  \label{thm:nagata-noetherian}
  \uses{def:nagata-compactification,def:nagata-noetherian-compactification}
  \lean{AlgebraicGeometry.nagataCompactification_of_isNoetherian}

  The recursively constructed factorization
  $\mathsf C_{\mathrm{Noeth}}(X/S)$ is a compactification of $X\to S$.
  Consequently $X\to S$ is compactifiable.
\end{theorem}

\begin{proof}
  \uses{def:nagata-noetherian-compactification}
  In the empty branch, the explicitly chosen empty factorization has the
  three fields required by \Cref{def:nagata-compactification}.  In the
  nonempty branch, the final component $h_n$ of the dependent recursive datum
  in \Cref{def:nagata-noetherian-compactification} has those fields, and
  $X_n=X$.  Thus the selected branch of
  $\mathsf C_{\mathrm{Noeth}}(X/S)$ is a compactification.  No later induction
  theorem is used to define any earlier $\mathsf C_r$.
\end{proof}


\section{Finite-presentation reduction and approximation}

The following constructions make the finite coefficient choices and all
eventual assertions explicit; each stage is obtained from finitely many
coefficients and relations and all later stages inherit the recorded data.

\subsection{Finite data and absolute Noetherian models}

\begin{definition}[Selected upper bound of a finite set]
  \label{def:nagata-directed-finite-upper-bound}
  \lean{Order.Directed.finsetUpperBound}

  Let $I$ be a nonempty directed preorder.  Fix a base point $i_\varnothing$
  and, for each pair $(i,j)$, fix a selected common upper bound $d(i,j)$
  together with proofs $i\leq d(i,j)$ and $j\leq d(i,j)$.  Define
  $u(F)\in I$ by Finset induction, inserting only an element $i\notin F$:
  \[
    u(\varnothing)=i_\varnothing,
    \qquad
    u(F\cup\{i\})=d(u(F),i)\quad(i\notin F).
  \]
\end{definition}

\begin{lemma}[Property of the selected finite upper bound]
  \label{lemma:nagata-directed-finite-upper-bound}
  \uses{def:nagata-directed-finite-upper-bound}
  \lean{Order.Directed.exists_ge_finset}

  Every $i\in F$ satisfies $i\leq u(F)$.
\end{lemma}

\begin{proof}
  Induct on $F$.  The empty case is vacuous.  At an insertion step, the
  selected $d(u(F),i)$ bounds both $u(F)$ and $i$; transitivity and the
  induction hypothesis prove the assertion for all old and new elements.
\end{proof}

\begin{lemma}[A ring is the filtered colimit of its finite-type
  $\mathbf Z$-subalgebras]
  \label{lemma:nagata-ring-finite-type-subalgebras}
  \lean{Algebra.filteredColimit_fgSubalgebra_int}

  For every ring $A$, its finitely generated $\mathbf Z$-subalgebras form a
  directed system and the canonical map
  \[
    \varinjlim_{A_0\subseteq A,\;A_0/\mathbf Z\ \mathrm{finite\ type}}A_0
      \longrightarrow A
  \]
  is an isomorphism.  In particular every finite set of elements of $A$ is
  contained in one such subalgebra.
\end{lemma}

\begin{proof}
  The subalgebra generated by the union of two finite generating sets contains
  both corresponding subalgebras.  Every element of $A$ lies in the
  subalgebra generated by that element, proving surjectivity.  Since all
  transition maps are inclusions, two representatives with equal image agree
  in a common over-algebra, proving injectivity.  A finite set is contained in
  the subalgebra it generates.
\end{proof}

\begin{lemma}[Finite data and finite relations descend in a filtered colimit]
  \label{lemma:nagata-filtered-colimit-finite-data}
  \lean{Algebra.filteredColimit_exists_stage_finset}

  Let $(A_i)_{i\in I}$ be a directed system of rings with colimit $A$.  Given
  finitely many elements of $A$ and finitely many equalities between ring
  expressions in those elements, there exists one stage $i$ and lifts of all
  selected elements to $A_i$ for which all selected equalities hold.  Once
  they hold at $i$, they hold after every transition $A_i\to A_j$.
\end{lemma}

\begin{proof}
  \uses{lemma:nagata-directed-finite-upper-bound}
  Choose a representative stage for each selected element.  Directedness and
  \Cref{lemma:nagata-directed-finite-upper-bound} place all representatives at
  one common stage.  Equality of two representatives in a filtered colimit is
  witnessed at a later stage.  There are only finitely many equalities, so a
  second finite upper bound makes them all true simultaneously.  Ring
  homomorphisms preserve the equalities at every later stage.
\end{proof}

\begin{definition}[Selected finite affine gluing presentation]
  \label{def:nagata-finite-affine-gluing-presentation}
  \lean{Scheme.finiteAffineGlueData}

  Let $S$ be quasi-compact and quasi-separated.  Select a finite affine cover
  $S=\bigcup_{a=1}^nU_a$.  For every pair $(a,b)$, select a finite affine
  cover $(U_{ab,r})_r$ of $U_a\cap U_b$ and then finite basic-open covers of
  each $U_{ab,r}$ as an open of $U_a$ and of $U_b$.  Record the resulting
  chart rings, inverted elements, two directions of each overlap map, and the
  finite lists of identity, inverse, and triple-cocycle equalities.  Call this
  finite record $\mathsf G_S$.
\end{definition}

\begin{lemma}[Properties of the selected finite affine presentation]
  \label{lemma:nagata-finite-affine-gluing-presentation}
  \uses{def:nagata-finite-affine-gluing-presentation}
  \lean{Scheme.exists_finiteAffineGlueData_of_qcqs}

  Every overlap in $\mathsf G_S$ is quasi-compact and represented by finitely
  many basic affine opens, all gluing identities form finite lists, and
  gluing $\mathsf G_S$ recovers $S$.
\end{lemma}

\begin{proof}
  Quasi-compactness gives the finite affine cover.  Quasi-separatedness makes
  each pairwise intersection quasi-compact, hence finitely affine-covered.
  A quasi-compact open in an affine scheme is a finite union of basic opens.
  There are finitely many pairs and triples, so all recorded equations form a
  finite list.  The recovery assertion is the ordinary gluing theorem applied
  to the restrictions of the original scheme.
\end{proof}

\begin{definition}[Coefficient stage of an affine gluing presentation]
  \label{def:nagata-affine-gluing-coefficient-stage}
  \uses{def:nagata-directed-finite-upper-bound,
    def:nagata-finite-affine-gluing-presentation,
    lemma:nagata-filtered-colimit-finite-data}
  \lean{Scheme.AffineGlueData.coefficientStage}

  Start with the finite lists of basic-open elements in $\mathsf G_S$.
  Close these lists under the finitely many overlap maps and proposed
  inverses, recording the finitely many coefficients occurring when the
  images are expressed in the selected localizations.  The identity and
  cocycle paths in the finite gluing groupoid show that one such finite
  closure suffices.  For each chart, take the finitely generated
  $\mathbf Z$-subalgebra generated by its resulting finite list.  Define
  $\lambda_{\mathrm{coeff}}(S)$ to be the selected finite upper bound of all
  chartwise coefficient indices, and denote the resulting finite-stage raw
  gluing datum by
  $\mathsf G_{S,\lambda_{\mathrm{coeff}}}^{\mathrm{raw}}$.
\end{definition}

\begin{lemma}[All gluing coefficients occur at the selected stage]
  \label{lemma:nagata-affine-gluing-coefficients-descend}
  \uses{def:nagata-affine-gluing-coefficient-stage}
  \lean{Scheme.AffineGlueData.exists_finiteTypeModel}

  Every selected finite coefficient subalgebra, basic-open element, restricted
  overlap map, and proposed inverse is defined at
  $\lambda_{\mathrm{coeff}}$.  Under base change to the original chart rings,
  these objects are the corresponding restrictions of the data in
  $\mathsf G_S$.
\end{lemma}

\begin{proof}
  \uses{lemma:nagata-directed-finite-upper-bound,
    lemma:nagata-ring-finite-type-subalgebras,
    lemma:nagata-finite-affine-gluing-presentation}
  Each finite coefficient list lies in one finitely generated
  $\mathbf Z$-subalgebra by
  \Cref{lemma:nagata-ring-finite-type-subalgebras}.  The selected upper bound
  dominates every such stage by
  \Cref{lemma:nagata-directed-finite-upper-bound}.  Hence all selected
  coefficients and restricted maps are simultaneously defined at the named
  common stage, and base change returns the original finite lists.  The full
  chart rings are recovered only after taking the filtered colimit over all
  finite refinements; no single stage is asserted to contain every coefficient
  of an arbitrary chart ring.
\end{proof}

\begin{definition}[Cocycle-valid stage of an affine gluing presentation]
  \label{def:nagata-affine-gluing-valid-stage}
  \uses{def:nagata-affine-gluing-coefficient-stage,
    def:nagata-directed-finite-upper-bound,
    lemma:nagata-filtered-colimit-finite-data}
  \lean{Scheme.AffineGlueData.validStage}

  For every identity, inverse, open-preservation, and triple-cocycle equation
  in $\mathsf G_S$, select a finite stage at which that equation holds on the
  chosen generators after localization.  Define
  $\lambda_{\mathrm{glue}}(S)$ to be the selected upper bound of
  $\lambda_{\mathrm{coeff}}(S)$ and these finitely many witness stages.  Let
  $\mathsf G_{S,\lambda_{\mathrm{glue}}}$ be the pullback of the raw datum to
  this stage.
\end{definition}

\begin{lemma}[The selected finite-stage datum satisfies the gluing axioms]
  \label{lemma:nagata-affine-gluing-identities-eventually}
  \uses{def:nagata-affine-gluing-valid-stage}
  \lean{Scheme.AffineGlueData.eventually_cocycle}

  The overlap maps in $\mathsf G_{S,\lambda_{\mathrm{glue}}}$ preserve the
  selected opens, are mutually inverse, and satisfy all identity and
  triple-overlap cocycle equations.
\end{lemma}

\begin{proof}
  \uses{lemma:nagata-directed-finite-upper-bound,
    lemma:nagata-affine-gluing-coefficients-descend}
  Every required equality is determined by finitely many generators after
  localization.  Equality in the colimit is witnessed at a finite stage.
  The definition includes all those witness stages in one finite set, and the
  selected upper-bound lemma makes every equation hold simultaneously at
  $\lambda_{\mathrm{glue}}(S)$.
\end{proof}

\begin{definition}[Finite-type Noetherian stage]
  \label{def:nagata-noetherian-stage}
  \uses{def:nagata-affine-gluing-valid-stage}
  \lean{Scheme.AffineGlueData.noetherianStage}

  Define $S_{\lambda_{\mathrm{glue}}}$ to be the scheme obtained by gluing
  $\mathsf G_{S,\lambda_{\mathrm{glue}}}$.  More generally, for every finite
  refinement $\lambda$ of the coefficient system satisfying the same finite
  equations, define $S_\lambda$ by gluing the corresponding datum.
\end{definition}

\begin{lemma}[Valid finite coefficient refinements form a directed set]
  \label{lemma:nagata-valid-refinements-directed}
  \uses{def:nagata-noetherian-stage}
  \lean{Scheme.AffineGlueData.validRefinements_directed}

  The valid finite coefficient refinements occurring in
  \Cref{def:nagata-noetherian-stage} are nonempty and directed under
  inclusion.  If $\lambda\leq\mu$, inclusion of coefficient algebras induces
  compatible affine chart maps $S_\mu\to S_\lambda$; these maps satisfy the
  identity and composition laws.
\end{lemma}

\begin{proof}
  \uses{lemma:nagata-affine-gluing-identities-eventually,
    lemma:nagata-filtered-colimit-finite-data,
    lemma:nagata-directed-finite-upper-bound}
  The selected valid stage is one object.  Given two valid refinements, take
  the union of their finite coefficient lists and then a common upper bound of
  the finitely many stages witnessing their gluing equations.  The finite-data
  lemma makes the resulting refinement valid and above both inputs.  Inclusion
  of coefficient rings induces the affine maps on every chart.  Their
  restrictions agree on overlaps by the gluing equations, so they glue; the
  identity and composition laws hold chartwise and therefore globally.
\end{proof}

\begin{definition}[Selected Noetherian inverse system]
  \label{def:nagata-noetherian-inverse-system}
  \uses{def:nagata-noetherian-stage,
    lemma:nagata-valid-refinements-directed}
  \lean{AlgebraicGeometry.noetherianInverseSystem}

  Order all valid finite coefficient refinements by inclusion.  For
  $\lambda\leq\mu$, define
  $S_\mu\to S_\lambda$ by gluing the affine morphisms induced by the
  inclusions of finite coefficient algebras, after refining the indexing data
  so inverse images of the selected affine charts remain affine.  Denote this
  inverse system by
  \[
    \mathsf N(S)=(I_S,(S_\lambda)_\lambda,
      (t_{\mu\lambda})_{\lambda\leq\mu}).
  \]
\end{definition}

\begin{lemma}[Properties of the finite-type stages]
  \label{lemma:nagata-glue-noetherian-stage}
  \uses{def:nagata-noetherian-inverse-system}
  \lean{Scheme.AffineGlueData.glue_finiteType_over_int}

  Every $S_\lambda$ is of finite type over $\mathbf Z$, and every transition
  morphism $t_{\mu\lambda}$ is affine.
\end{lemma}

\begin{proof}
  \uses{lemma:nagata-affine-gluing-identities-eventually,
    lemma:nagata-valid-refinements-directed}
  The gluing axioms hold at each valid stage, so the schemes are defined.
  Each one is covered by finitely many spectra of finitely generated
  $\mathbf Z$-algebras, proving finite type.  On the chosen charts each
  transition is affine by construction; the indexing refinement makes those
  charts cover every inverse image, proving global affineness.
\end{proof}

\begin{definition}[Comparison from the original scheme to the inverse limit]
  \label{def:nagata-glued-limit-comparison}
  \uses{def:nagata-noetherian-inverse-system}
  \lean{Scheme.AffineGlueData.limitComparison}

  The inclusions of every finite coefficient algebra into the original chart
  algebra give compatible maps $S\to S_\lambda$.  Define
  \[
    \kappa_S\colon S\longrightarrow\varprojlim_{\lambda\in I_S}S_\lambda
  \]
  by the universal property of the inverse limit.
\end{definition}

\begin{lemma}[Recovery of the glued limit]
  \label{lemma:nagata-glued-limit-recovery}
  \uses{def:nagata-glued-limit-comparison}
  \lean{Scheme.AffineGlueData.limit_iso}

  The constructed comparison $\kappa_S$ is an isomorphism.
\end{lemma}

\begin{proof}
  \uses{lemma:nagata-ring-finite-type-subalgebras,
    lemma:nagata-glue-noetherian-stage}
  On every selected affine chart, the inverse limit has coordinate ring the
  filtered colimit of its finite-stage coordinate rings, hence the original
  chart ring.  These chartwise isomorphisms agree on the chosen basic
  overlaps because all transition maps came from the original gluing maps.
  They glue to an inverse for $\kappa_S$.
\end{proof}

\begin{proposition}[{\cite[\href{https://stacks.math.columbia.edu/tag/01ZA}{Tag 01ZA}]{stacks-project}}]
  \label{prop:nagata-noetherian-inverse-system}
  \uses{def:nagata-glued-limit-comparison,def:nagata-noetherian-inverse-system}
  \lean{AlgebraicGeometry.exists_noetherian_inverseSystem}

  The constructed system $\mathsf N(S)$ consists of Noetherian schemes with
  affine transition morphisms, every term is of finite type over
  $\mathbf Z$, and $\kappa_S$ identifies $S$ with its inverse limit.
\end{proposition}

\begin{proof}
  \uses{lemma:nagata-valid-refinements-directed,
    lemma:nagata-glue-noetherian-stage}
  Nonemptiness and directedness of the index type, including the transition
  identities, are supplied by
  \Cref{lemma:nagata-valid-refinements-directed}.  By
  \Cref{lemma:nagata-glue-noetherian-stage}, every chart of every $S_\lambda$
  is the spectrum of a finitely generated $\mathbf Z$-algebra and every
  transition is affine.  Such rings are Noetherian, so the glued schemes are
  Noetherian.  On a selected affine chart $\operatorname{Spec}(A)$ of $S$,
  the chart rings $A_\lambda$ form a filtered system whose colimit is $A$ by
  the finite-subalgebra lemma; the affine inverse-limit calculation therefore
  identifies $\varprojlim\operatorname{Spec}(A_\lambda)$ with
  $\operatorname{Spec}(A)$.  The identifications commute with restriction to
  every recorded basic overlap because all transition maps are induced by the
  same localization maps in $A$.  Since the affine charts and overlaps are
  finite, these local inverses glue and are inverse to the comparison
  $\kappa_S$ of \Cref{def:nagata-glued-limit-comparison}.  Hence $\kappa_S$ is
  an isomorphism and all asserted properties hold.
\end{proof}

\subsection{Embedding a finite-type scheme into a finite-presentation model}

\begin{definition}[Finite-stage data for a morphism]
  \label{def:nagata-relative-stage-morphism-data}
  \uses{def:nagata-noetherian-inverse-system,
    def:nagata-glued-limit-comparison}
  \lean{AlgebraicGeometry.Limits.relativeStageMorphismData}

  Let $X\to S$ be a morphism and choose the inverse systems
  $\mathsf N(X)$ and $\mathsf N(S)$ together with their comparison maps to
  $X$ and $S$.  A finite-stage datum for this morphism consists of indices
  $\lambda_X$ and $\lambda_S$, a morphism
  $\varphi_{\lambda_X,\lambda_S}\colon X_{\lambda_X}\to S_{\lambda_S}$,
  and an equality after composing with the two limit projections: the map
  $X\to S\to S_{\lambda_S}$ equals
  $X\to X_{\lambda_X}\to S_{\lambda_S}$.  The datum also records the finite
  affine charts and generator-and-relation lists used to construct
  $\varphi_{\lambda_X,\lambda_S}$.  Finally it records a finite affine cover
  of $X$, compatible opens in
  $X_{\lambda_X}\times_{S_{\lambda_S}}S$, and chartwise surjections of
  coordinate rings showing that the induced comparison
  $X\to X_{\lambda_X}\times_{S_{\lambda_S}}S$ is an immersion; all overlap
  equalities are included in the datum.
\end{definition}

\begin{lemma}[Existence of finite-stage data for a finite-type morphism]
  \label{lemma:nagata-relative-stage-morphism-data-exists}
  \uses{def:nagata-relative-stage-morphism-data,
    lemma:nagata-filtered-colimit-finite-data,
    lemma:nagata-directed-finite-upper-bound}
  \lean{AlgebraicGeometry.Limits.exists_relativeStageMorphismData}

  If $S$ is qcqs and $X\to S$ is locally of finite type with $X$
  quasi-compact and quasi-separated, then a finite-stage datum exists after enlarging both
  indices.  Its chart maps use only finitely many coefficients and relations,
  and their base change to $S$ is the original morphism.
\end{lemma}

\begin{proof}
  Choose finite affine presentations of source and target charts. Their
  finitely many coefficients, structural maps, and overlap identities occur
  at one common stage by the finite-data lemma; enlarging both indices gives
  the required compatible datum and immersion charts.
\end{proof}

\begin{definition}[Finite-presentation immersion model]
  \label{def:nagata-immersion-into-fp-model}
  \uses{lemma:nagata-relative-stage-morphism-data-exists,
    def:nagata-relative-stage-morphism-data}
  \lean{AlgebraicGeometry.finitePresentationImmersionModel}

  Let $S$ be qcqs and let $X\to S$ be locally of finite type with $X$
  quasi-compact and quasi-separated.  Choose a finite-stage datum from
  \Cref{lemma:nagata-relative-stage-morphism-data-exists}, including indices
  $\lambda_X$ and $\lambda_S$ and the descended morphism
  $\varphi_{\lambda_X,\lambda_S}$.  Define
  \[
    Y^{\mathrm{imm}}(X/S)
      =X_{\lambda_X}\times_{S_{\lambda_S}}S
  \]
  using $\varphi_{\lambda_X,\lambda_S}$ and the comparison map
  $S\to S_{\lambda_S}$.  The compatibility equality in the finite-stage
  datum induces the map $i^{\mathrm{imm}}\colon X\to Y^{\mathrm{imm}}(X/S)$
  over $S$.
\end{definition}

\begin{lemma}[Properties of the finite-presentation immersion model,
  {\cite[\href{https://stacks.math.columbia.edu/tag/01ZE}{Tag 01ZE}]{stacks-project}}]
  \label{lemma:nagata-immersion-into-fp-model}
  \uses{def:nagata-immersion-into-fp-model}
  \lean{AlgebraicGeometry.exists_immersion_finitePresentation_model}

  The morphism $Y^{\mathrm{imm}}(X/S)\to S$ is of finite presentation, and
  the constructed map $i^{\mathrm{imm}}$ is an immersion over $S$.
\end{lemma}

\begin{proof}
  \uses{lemma:nagata-relative-stage-morphism-data-exists,
    prop:nagata-noetherian-inverse-system,
    lemma:nagata-directed-finite-upper-bound}
  The schemes $X_{\lambda_X}$ and $S_{\lambda_S}$ are finite type over
  $\mathbf Z$, and $S_{\lambda_S}$ is Noetherian.  Hence the descended
  morphism $X_{\lambda_X}\to S_{\lambda_S}$ is finite type and therefore
  finite presentation.  Its base change along $S\to S_{\lambda_S}$ is thus
  finite presentation over $S$.  On
  each selected affine chart, the descended generator list gives a surjection
  of coordinate rings, hence a closed immersion into an open of the target.
  The compatibility fields in the finite-stage datum identify these local
  maps on overlaps; once their descent witnesses are supplied, the finitely
  many target opens glue to the asserted immersion.
\end{proof}

\begin{lemma}[Existence of a quasi-compact closed-immersion neighbourhood]
  \label{lemma:nagata-immersion-shrink-to-closed-exists}
  \lean{AlgebraicGeometry.Immersion.exists_qcOpen_closedImmersion}

  Let $i\colon X\to Y$ be an immersion with $X$ quasi-compact and $Y$
  quasi-separated.  There is a quasi-compact open $W\subseteq Y$ containing
  $i(X)$ such that $X\to W$ is a closed immersion.
\end{lemma}

\begin{proof}
  Factor the immersion as a closed immersion $X\hookrightarrow U$ followed by
  an open immersion $U\hookrightarrow Y$.  The image of the quasi-compact
  scheme $X$ is quasi-compact in $U$.  Cover it by affine opens of $Y$
  contained in $U$, and choose a finite subcover $V_1,\ldots,V_n$.  Their
  union $W=\bigcup_rV_r$ is a quasi-compact open of $Y$, contains $i(X)$,
  and is contained in $U$.  Base-changing the closed immersion
  $X\hookrightarrow U$ along $W\hookrightarrow U$ gives a closed immersion
  $X\times_UW\hookrightarrow W$.  Because the image of $X$ lies in $W$, the
  first projection identifies $X\times_UW$ with $X$, so the induced map
  $X\to W$ is a closed immersion.
\end{proof}

\begin{definition}[Quasi-compact open on which an immersion is closed]
  \label{def:nagata-immersion-shrink-to-closed}
  \uses{lemma:nagata-immersion-shrink-to-closed-exists}
  \lean{AlgebraicGeometry.Immersion.qcClosedNeighbourhood}

  Under the hypotheses of
  \Cref{lemma:nagata-immersion-shrink-to-closed-exists}, choose one witness and
  define
  \[
    Y^{\mathrm{cl}}(i)=W.
  \]
\end{definition}

\begin{lemma}[Properties of the closed-immersion neighbourhood,
  {\cite[\href{https://stacks.math.columbia.edu/tag/01ZG}{Tag 01ZG}]{stacks-project}}]
  \label{lemma:nagata-immersion-shrink-to-closed}
  \uses{def:nagata-immersion-shrink-to-closed,
    lemma:nagata-immersion-shrink-to-closed-exists}
  \lean{AlgebraicGeometry.Immersion.qcClosedNeighbourhood_spec}

  The constructed $Y^{\mathrm{cl}}(i)$ is a quasi-compact open containing
  $i(X)$, and $X\to Y^{\mathrm{cl}}(i)$ is a closed immersion.
\end{lemma}

\begin{proof}
  These are the specification fields of the witness supplied by
  \Cref{lemma:nagata-immersion-shrink-to-closed-exists} and selected in
  \Cref{def:nagata-immersion-shrink-to-closed}.
\end{proof}

\begin{definition}[Finite-presentation closed-subscheme system]
  \label{def:nagata-closed-subscheme-fp-system}
  \lean{AlgebraicGeometry.ClosedImmersion.finitePresentationSystem}

  Let $X=V(\mathcal I)\hookrightarrow Y$ be closed with $Y$ qcqs.  Let
  $A_{X/Y}$ be the directed set of finite type quasi-coherent subideals
  $\mathcal I_\alpha\subseteq\mathcal I$, ordered by inclusion.  Define
  \[
    X_\alpha=V(\mathcal I_\alpha)\hookrightarrow Y,
  \]
  and for $\alpha\leq\beta$ define the transition
  $X_\beta\hookrightarrow X_\alpha$ from
  $\mathcal I_\alpha\subseteq\mathcal I_\beta$.
\end{definition}

\begin{lemma}[Properties of the closed-subscheme system,
  {\cite[\href{https://stacks.math.columbia.edu/tag/09ZP}{Tag 09ZP}]{stacks-project}}]
  \label{lemma:nagata-closed-subscheme-fp-system}
  \uses{def:nagata-closed-subscheme-fp-system}
  \lean{AlgebraicGeometry.ClosedImmersion.limit_finitePresentation}

  Every $X_\alpha\hookrightarrow Y$ is of finite presentation, all
  transitions are closed immersions, and the canonical map
  $X\to\varprojlim_\alpha X_\alpha$ is an isomorphism.
\end{lemma}

\begin{proof}
  A quasi-coherent ideal on a qcqs scheme is the filtered union of its finite
  type quasi-coherent subideals.  Finite type of $\mathcal I_\alpha$ is
  equivalent to finite presentation of $X_\alpha\hookrightarrow Y$.  On an
  affine open, the coordinate ring of the limit is the quotient by
  $\varinjlim\mathcal I_\alpha=\mathcal I$, so the limit is $X$; these affine
  identifications glue.
\end{proof}

\begin{definition}[Finite-presentation model system of a finite-type morphism]
  \label{def:nagata-finite-type-fp-model-system}
  \uses{def:nagata-immersion-into-fp-model,
    def:nagata-immersion-shrink-to-closed,
    def:nagata-closed-subscheme-fp-system}
  \lean{AlgebraicGeometry.finiteTypeFinitePresentationModelSystem}

  For $X\to S$ locally of finite type with $X$ qcqs and $S$ qcqs, form
  $i^{\mathrm{imm}}\colon X\to Y^{\mathrm{imm}}$, set
  $Y^\circ=Y^{\mathrm{cl}}(i^{\mathrm{imm}})$, and apply
  \Cref{def:nagata-closed-subscheme-fp-system} to
  $X\hookrightarrow Y^\circ$.  Denote the resulting system by
  \[
    \mathsf{FP}(X/S)=(A_{X/S},(X_\alpha)_\alpha),
  \]
  with structural maps $i_\alpha\colon X\to X_\alpha$.
\end{definition}

\begin{lemma}[Properties of the finite-presentation model system,
  {\cite[\href{https://stacks.math.columbia.edu/tag/09ZQ}{Tag 09ZQ}]{stacks-project}}]
  \label{lemma:nagata-finite-type-fp-model-system}
  \uses{def:nagata-finite-type-fp-model-system}
  \lean{AlgebraicGeometry.finiteType_limit_finitePresentation_models}

  Every $X_\alpha\to S$ is of finite presentation, every transition map is a
  closed immersion, every $i_\alpha$ is a closed immersion, and
  $X\cong\varprojlim_\alpha X_\alpha$ over $S$.
\end{lemma}

\begin{proof}
  \uses{lemma:nagata-immersion-into-fp-model,
    lemma:nagata-immersion-shrink-to-closed,
    lemma:nagata-closed-subscheme-fp-system}
  The immersion model is of finite presentation over $S$.  Since $S$ is
  quasi-separated, its source is quasi-separated; the selected open
  $Y^\circ$ is quasi-compact and remains of finite presentation over $S$.
  The closed-subscheme system then has closed transitions and limit $X$.
  Composition of each finite-presentation closed immersion
  $X_\alpha\hookrightarrow Y^\circ$ with $Y^\circ\to S$ is of finite
  presentation.  The canonical maps from the limit are closed immersions
  because the transition maps are closed immersions.
\end{proof}

\begin{lemma}[Diagonal of a finite-presentation morphism]
  \label{lemma:nagata-diagonal-finite-presentation}
  \lean{Scheme.Hom.diagonal_finitePresentation}

  If $f\colon X\to S$ is of finite presentation, then the diagonal
  $\Delta_{X/S}\colon X\to X\times_S X$ is of finite presentation.
\end{lemma}

\begin{proof}
  On an affine chart $\operatorname{Spec}(A)\subseteq S$, choose a finite
  presentation
  $B=A[x_1,\ldots,x_n]/(r_1,\ldots,r_m)$ for an affine chart
  $\operatorname{Spec}(B)\subseteq X$.  The corresponding chart of
  $X\times_SX$ has ring $B\otimes_A B$, and the diagonal is represented by the
  multiplication map
  \[
    \mu\colon B\otimes_A B\longrightarrow B,\qquad b\otimes b'\longmapsto bb'.
  \]
  This map is surjective, and its kernel is generated by the finite list
  $x_i\otimes1-1\otimes x_i$.  Hence $B$ is finitely presented as a
  $(B\otimes_A B)$-algebra.  Finite presentation is local on source and
  target for immersions; a finite affine cover of the diagonal therefore
  glues these chart calculations to $\Delta_{X/S}$.
\end{proof}

\begin{lemma}[Existence of a separated stage over a fixed base]
  \label{lemma:nagata-fixed-base-separated-stage-exists}
  \uses{lemma:nagata-diagonal-finite-presentation}
  \lean{AlgebraicGeometry.Limits.exists_separatedStage_over_fixedBase}

  Let $X=\varprojlim_\alpha X_\alpha$ over a quasi-separated base $S$, where
  every $X_\alpha\to S$ is of finite presentation, the terms are qcqs, and
  the transition maps are affine.  If $X\to S$ is separated, then there is an
  index $\alpha_0$ such that $X_\alpha\to S$ is separated for every
  $\alpha\geq\alpha_0$.
\end{lemma}

\begin{proof}
  \uses{lemma:nagata-filtered-colimit-finite-data,
    lemma:nagata-directed-finite-upper-bound}
  The diagonal is closed on the limit.  Each stage diagonal is of finite
  presentation by \Cref{lemma:nagata-diagonal-finite-presentation}.  A finite
  affine cover and finite
  presentation provide finitely many target generators and preimages; record
  their relations at one common stage. Surjectivity then persists under every
  later base change, so all later diagonals are closed.
  Separatedness is closedness of the diagonal.  Choose a finite affine cover
  of one stage and finite affine covers of the inverse images of its members
  in the source and its square over $S$.  The limiting diagonal is closed, so
  on each chart its coordinate map is surjective.  Finite presentation gives
  finitely many target generators; lift preimages of those generators and all
  required equalities to finite stages.  By the finite-data lemma and a common
  upper bound, all chartwise diagonal maps are surjective at one stage and at
  every later stage.  Hence the diagonals are closed globally at every such
  stage.
\end{proof}

\begin{definition}[Selected stage for fixed-base separatedness]
  \label{def:nagata-fixed-base-separated-stage}
  \uses{lemma:nagata-fixed-base-separated-stage-exists}
  \lean{AlgebraicGeometry.Limits.separatedStageOverFixedBase}

  Under the hypotheses of
  \Cref{lemma:nagata-fixed-base-separated-stage-exists}, choose one witness
  and define $\alpha_{\mathrm{sep}}$ to be its index component.  Eventual
  separatedness is recorded only in the following specification lemma.
\end{definition}

\begin{lemma}[Eventual separatedness over a fixed base,
  {\cite[\href{https://stacks.math.columbia.edu/tag/01ZH}{Tag 01ZH}]{stacks-project}}]
  \label{lemma:nagata-fixed-base-separated-eventually}
  \uses{def:nagata-fixed-base-separated-stage}
  \lean{AlgebraicGeometry.Limits.separated_eventually_over_fixedBase}

  For every $\alpha\geq\alpha_{\mathrm{sep}}$, the morphism
  $X_\alpha\to S$ is separated.
\end{lemma}

\begin{proof}
  \uses{lemma:nagata-fixed-base-separated-stage-exists}
  The witness supplied by
  \Cref{lemma:nagata-fixed-base-separated-stage-exists} contains the eventual
  separatedness proof.  Project that field along the index selected in
  \Cref{def:nagata-fixed-base-separated-stage}.
\end{proof}

\begin{definition}[Separated finite-presentation envelope]
  \label{def:nagata-finite-type-closed-in-fp}
  \uses{def:nagata-finite-type-fp-model-system,
    def:nagata-fixed-base-separated-stage}
  \lean{AlgebraicGeometry.closedImmersionFinitePresentationEnvelope}

  Let $S$ be qcqs and let $X\to S$ be separated and of finite type.  Form
  \[
    \mathsf{FP}(X/S)=(A_{X/S},(X_\alpha)_\alpha).
  \]
  Apply \Cref{def:nagata-fixed-base-separated-stage} to this system and
  define
  \[
    \alpha^{X/S}_{\mathrm{sep}}
      =\alpha_{\mathrm{sep}},
    \qquad
    X^{\mathrm{fp}}=X_{\alpha^{X/S}_{\mathrm{sep}}},
    \qquad
    i^{\mathrm{fp}}_{X/S}
      =i_{\alpha^{X/S}_{\mathrm{sep}}}\colon X\hookrightarrow X^{\mathrm{fp}}.
  \]
  Denote the resulting explicit pair by
  $\mathsf{Env}_{\mathrm{fp}}(X/S)=(X^{\mathrm{fp}},i^{\mathrm{fp}}_{X/S})$.
\end{definition}

\begin{proposition}[Properties of the separated finite-presentation envelope,
  {\cite[\href{https://stacks.math.columbia.edu/tag/01ZJ}{Tag 01ZJ}]{stacks-project}}]
  \label{prop:nagata-finite-type-closed-in-fp}
  \uses{def:nagata-finite-type-closed-in-fp}
  \lean{AlgebraicGeometry.exists_closedImmersion_finitePresentation_model}

  The constructed morphism $X^{\mathrm{fp}}\to S$ is separated and of
  finite presentation, and the constructed map
  $i^{\mathrm{fp}}_{X/S}\colon X\to X^{\mathrm{fp}}$ is a closed immersion
  over $S$.
\end{proposition}

\begin{proof}
  \uses{lemma:nagata-finite-type-fp-model-system,
    lemma:nagata-fixed-base-separated-eventually}
  A finite-type morphism over a quasi-compact base has quasi-compact source,
  and separatedness makes that source quasi-separated, so the model system is
  available.  Its property lemma makes every term finite presentation and
  every structural map $i_\alpha$ a closed immersion.  At the explicitly
  selected stage $\alpha^{X/S}_{\mathrm{sep}}$, separatedness follows from
  \Cref{lemma:nagata-fixed-base-separated-eventually}.  These are exactly the
  three assertions for the named envelope.
\end{proof}


\subsection{Descent of relative finite-presentation data}

In this subsection fix a directed inverse system
$S=\varprojlim_{\lambda\in I}S_\lambda$ with qcqs terms and affine
transition morphisms.

\begin{definition}[Base-changed system over a chosen stage]
  \label{def:nagata-basechanged-system}
  \lean{AlgebraicGeometry.Limits.baseChangedSystem}

  Fix $\lambda_0\in I$ and a scheme $Y_{\lambda_0}$ equipped with a morphism
  to $S_{\lambda_0}$.  On the cofinal directed subsystem
  $I_{\lambda_0}=\{\lambda\mid\lambda_0\leq\lambda\}$, define
  \[
    Y_\lambda=Y_{\lambda_0}\times_{S_{\lambda_0}}S_\lambda,
    \qquad
    Y=Y_{\lambda_0}\times_{S_{\lambda_0}}S.
  \]
  The transition morphisms $Y_\mu\to Y_\lambda$ for
  $\lambda\leq\mu$ are the canonical base changes of the inverse-system
  transitions $S_\mu\to S_\lambda$, and the maps $Y\to Y_\lambda$ are the
  projections induced by $S\to S_\lambda$.
\end{definition}

\begin{definition}[Finite-stage model of a quasi-compact open]
  \label{def:nagata-qc-open-descent}
  \uses{def:nagata-directed-finite-upper-bound,
    lemma:nagata-filtered-colimit-finite-data,def:nagata-basechanged-system}
  \lean{AlgebraicGeometry.Limits.qcOpenDescentConstruction}

  Fix $\lambda_0$, assume $Y_{\lambda_0}$ is qcqs, and put
  $Y=S\times_{S_{\lambda_0}}Y_{\lambda_0}$, and let $W\subseteq Y$ be a
  quasi-compact open.  Select finitely many affine opens of
  $Y_{\lambda_0}$ meeting the image of $W$.  On their inverse limits, write
  $W$ as a finite union of basic opens $D(f_{ak})$.  Choose finite-stage
  lifts $f_{ak,\lambda_{ak}}$ and define
  \[
    \lambda_W=u(\{\lambda_0\}\cup\{\lambda_{ak}\}_{a,k}),
    \qquad
    W_{\lambda_W}=\bigcup_{a,k}D(f_{ak,\lambda_W})
      \subseteq Y_{\lambda_W}.
  \]
  Quasi-compactness supplies a finite affine cover; on each chart the inverse
  limit identifies the coordinate ring with the filtered colimit, so the open
  has a finite basic-open presentation. Record those functions and lifts in
  the witness before taking the common upper bound.
\end{definition}

\begin{lemma}[Property of the descended quasi-compact open]
  \label{lemma:nagata-qc-open-descends}
  \uses{def:nagata-qc-open-descent}
  \lean{AlgebraicGeometry.Limits.qcOpen_descends}

  The constructed $W_{\lambda_W}$ is quasi-compact, and its inverse image in
  $Y$ is exactly $W$.
\end{lemma}

\begin{proof}
  \uses{lemma:nagata-directed-finite-upper-bound}
  A finite union of basic opens is quasi-compact.  Every lifted defining
  function has reached the common stage $\lambda_W$, and base change sends
  $D(f_{ak,\lambda_W})$ to $D(f_{ak})$.  Taking the finite union therefore
  recovers precisely the selected basic-open presentation of $W$.
\end{proof}

\begin{definition}[Stage of equality for descended quasi-compact opens]
  \label{def:nagata-qc-open-equality-stage}
  \uses{def:nagata-directed-finite-upper-bound,
    lemma:nagata-filtered-colimit-finite-data}
  \lean{AlgebraicGeometry.Limits.qcOpenEqualityStage}

  Let $W_{1,\lambda_0},W_{2,\lambda_0}\subseteq Y_{\lambda_0}$ be
  quasi-compact opens whose inverse images in $Y$ agree.  On a finite affine
  cover, choose finite basic-open presentations for both.  For every required
  radical containment select the finite stage containing a witnessing power
  and its finitely many coefficients.  Define
  $\lambda_{W_1=W_2}$ to be the selected upper bound of those witness stages.
\end{definition}

\begin{lemma}[Eventual equality of descended quasi-compact opens]
  \label{lemma:nagata-qc-open-eq-eventually}
  \uses{def:nagata-qc-open-equality-stage}
  \lean{AlgebraicGeometry.Limits.qcOpen_eq_eventually}

  For every $\lambda\geq\lambda_{W_1=W_2}$, the inverse images of
  $W_{1,\lambda_0}$ and $W_{2,\lambda_0}$ in $Y_\lambda$ are equal.
\end{lemma}

\begin{proof}
  Equality of finite unions of basic opens is equivalent, chartwise, to the
  two finite lists of radical containments.  Each chosen power identity holds
  at its witness stage and at every later stage.  Their common upper bound
  therefore imposes both inclusions on every affine chart, hence equality of
  the global opens.
\end{proof}

\begin{definition}[Finite-stage model of a finitely presented affine algebra]
  \label{def:nagata-affine-fp-descent}
  \uses{def:nagata-directed-finite-upper-bound,
    lemma:nagata-filtered-colimit-finite-data}
  \lean{AlgebraicGeometry.Limits.affineFinitePresentationModel}

  Suppose $S_{\lambda_0}=\operatorname{Spec}(A_{\lambda_0})$,
  $S=\operatorname{Spec}(A)$, and fix a finite presentation
  \[
    B=A[x_1,\ldots,x_n]/(r_1,\ldots,r_m).
  \]
  Select finite-stage lifts of every coefficient of every $r_j$, let
  $\lambda_B$ be their selected common upper bound, let
  $r_{j,\lambda_B}$ be the lifted polynomials, and define
  \[
    B_{\lambda_B}
      =A_{\lambda_B}[x_1,\ldots,x_n]/
        (r_{1,\lambda_B},\ldots,r_{m,\lambda_B}).
  \]
  Define $\beta_B$ to be the comparison from
  $B_{\lambda_B}\otimes_{A_{\lambda_B}}A$ to $B$ induced by the presentation.
\end{definition}

\begin{lemma}[Property of affine finite-presentation descent]
  \label{lemma:nagata-affine-fp-descends}
  \uses{def:nagata-affine-fp-descent}
  \lean{AlgebraicGeometry.Limits.affine_finitePresentation_descends}

  The algebra $B_{\lambda_B}$ is finitely presented over
  $A_{\lambda_B}$, and the constructed map $\beta_B$ is an isomorphism.
\end{lemma}

\begin{proof}
  The displayed quotient is a finite presentation by construction.  Tensoring
  it with $A$ sends every lifted relation to the original $r_j$ and preserves
  the finite polynomial presentation.  Thus the base-changed presentation is
  the chosen presentation of $B$, with $\beta_B$ and its generatorwise inverse
  mutually inverse.
\end{proof}

\begin{definition}[Finite-stage model of a morphism]
  \label{def:nagata-morphism-descent}
  \uses{def:nagata-qc-open-descent,
    def:nagata-affine-fp-descent,
    def:nagata-directed-finite-upper-bound,
    lemma:nagata-filtered-colimit-finite-data}
  \lean{AlgebraicGeometry.Limits.morphismDescentConstruction}

  Let $X_{\lambda_0},Y_{\lambda_0}$ be of finite presentation over
  $S_{\lambda_0}$, write
  $X_{\lambda_0,S}=X_{\lambda_0}\times_{S_{\lambda_0}}S$ and
  $Y_{\lambda_0,S}=Y_{\lambda_0}\times_{S_{\lambda_0}}S$, and let
  $\varphi\colon X_{\lambda_0,S}\to Y_{\lambda_0,S}$ be an $S$-morphism.
  Choose a finite affine target cover, descend its quasi-compact inverse-image
  opens, choose finite affine covers of those opens, and descend each affine
  ring map using its finite generator-and-relation data.  Include the finitely
  many pairwise-overlap equality stages in one selected upper bound
  $\lambda_\varphi$, and define
  \[
    \varphi_{\lambda_\varphi}\colon
      X_{\lambda_\varphi}\longrightarrow Y_{\lambda_\varphi}
  \]
  by gluing the descended affine maps.
  Retain comparison isomorphisms for each descended open, transport every
  affine map through them, and include the finitely many overlap equalities in
  the common upper bound; the transported maps glue at the selected stage.
\end{definition}

\begin{lemma}[Property of the descended morphism,
  {\cite[\href{https://stacks.math.columbia.edu/tag/01ZM}{Tag 01ZM(2)}]{stacks-project}}]
  \label{lemma:nagata-morphism-descends}
  \uses{def:nagata-morphism-descent}
  \lean{AlgebraicGeometry.Limits.hom_extends}

  The base change of the constructed
  $\varphi_{\lambda_\varphi}$ to $S$ is $\varphi$.
\end{lemma}

\begin{proof}
  \uses{lemma:nagata-qc-open-descends,
    lemma:nagata-affine-fp-descends}
  The inverse-image opens pull back to the original opens by the quasi-compact
  open descent property.  On each affine piece, the descended ring map pulls
  back to the original map by the affine finite-presentation property.  The
  selected equality stages make those local maps agree on all overlaps, so
  their gluing pulls back to the gluing of the original local maps, namely
  $\varphi$.
\end{proof}

\begin{definition}[Stage of equality for two descended morphisms]
  \label{def:nagata-morphism-equality-stage}
  \uses{def:nagata-qc-open-equality-stage,
    def:nagata-directed-finite-upper-bound,
    lemma:nagata-filtered-colimit-finite-data}
  \lean{AlgebraicGeometry.Limits.morphismEqualityStage}

  Let $\varphi_{\lambda_0},\psi_{\lambda_0}\colon
  X_{\lambda_0}\to Y_{\lambda_0}$ have equal base change to $S$.  On a
  finite affine target cover, include the stages at which their inverse-image
  opens agree and the stages at which the induced ring maps agree on finite
  generator lists.  Define $\lambda_{\varphi=\psi}$ to be the selected upper
  bound of this finite set.
\end{definition}

\begin{lemma}[Eventual equality of morphisms,
  {\cite[\href{https://stacks.math.columbia.edu/tag/01ZM}{Tag 01ZM(3)}]{stacks-project}}]
  \label{lemma:nagata-morphism-eq-eventually}
  \uses{def:nagata-morphism-equality-stage}
  \lean{AlgebraicGeometry.Limits.hom_eq_eventually}

  For every $\lambda\geq\lambda_{\varphi=\psi}$, the base changes of
  $\varphi_{\lambda_0}$ and $\psi_{\lambda_0}$ to $S_\lambda$ are equal.
\end{lemma}

\begin{proof}
  \uses{lemma:nagata-qc-open-eq-eventually,
    lemma:nagata-directed-finite-upper-bound}
  Above the selected stage, the two maps have equal inverse images of every
  target chart.  On a finite affine cover of each common inverse image, their
  ring maps agree on a finite generating set and hence agree.  The local
  equalities glue to the asserted global equality and persist at later stages.
\end{proof}

\begin{definition}[Selected finite relative affine presentation]
  \label{def:nagata-relative-finite-affine-presentation}
  \lean{AlgebraicGeometry.Limits.relativeFiniteAffinePresentation}

  Let $S$ be qcqs and let $X\to S$ be of finite presentation.  Select a
  finite affine cover of
  $S$, finite affine covers of its inverse images in $X$, finite
  presentations of all chart algebras over the base chart algebras, finite
  basic-open presentations of the overlaps, and the two directions of every
  overlap isomorphism.  Record these data and their finite identity, inverse,
  and cocycle equations as $\mathsf R_X$.
  Quasi-compactness gives finite affine covers and quasi-separatedness makes
  overlaps quasi-compact. Finite presentation supplies finite generators and
  relations; record both overlap directions and the finite identity and
  cocycle equations as $\mathsf R_X$.
\end{definition}

\begin{definition}[Finite-stage raw relative gluing datum]
  \label{def:nagata-relative-gluing-datum-descent}
  \uses{def:nagata-relative-finite-affine-presentation,
    def:nagata-qc-open-descent,
    def:nagata-affine-fp-descent,
    def:nagata-morphism-descent,
    lemma:nagata-affine-fp-descends,
    lemma:nagata-qc-open-descends,
    lemma:nagata-morphism-descends,
    def:nagata-directed-finite-upper-bound}
  \lean{AlgebraicGeometry.Limits.relativeGlueDataConstruction}

  Descend every chart algebra in $\mathsf R_X$ using
  \Cref{lemma:nagata-affine-fp-descends}, every quasi-compact overlap open
  using \Cref{lemma:nagata-qc-open-descends}, and both directions of every
  overlap isomorphism using \Cref{lemma:nagata-morphism-descends}.  Define
  $\lambda_{\mathrm{raw}}(X/S)$ to be the selected
  upper bound of their finitely many stages, and denote the resulting raw
  datum over $S_{\lambda_{\mathrm{raw}}}$ by
  $\mathsf R_{X,\lambda_{\mathrm{raw}}}^{\mathrm{raw}}$.
  Each object is transported by the canonical base-change map to the selected
  upper bound. Functoriality supplies the comparison maps and coherence, so
  the raw datum is well typed.
\end{definition}

\begin{lemma}[Property of the descended finite affine presentation]
  \label{lemma:nagata-relative-gluing-datum-descends}
  \uses{def:nagata-relative-gluing-datum-descent}
  \lean{AlgebraicGeometry.Limits.finiteAffinePresentation.descends}

  Every chart, overlap open, and overlap map in the constructed raw datum
  base changes to its named counterpart in $\mathsf R_X$.
\end{lemma}

\begin{proof}
  \uses{lemma:nagata-qc-open-descends,
    lemma:nagata-affine-fp-descends,
    lemma:nagata-morphism-descends}
  Apply the corresponding property lemma to each finite-presentation chart,
  quasi-compact overlap, and overlap morphism.  There are only finitely many
  objects, and the common upper bound transports all of them to one stage
  without changing their base changes.
\end{proof}

\begin{definition}[Cocycle-valid stage for the relative datum]
  \label{def:nagata-relative-gluing-valid-stage}
  \uses{def:nagata-relative-gluing-datum-descent,
    def:nagata-qc-open-equality-stage,
    def:nagata-morphism-equality-stage,
    def:nagata-directed-finite-upper-bound,
    lemma:nagata-filtered-colimit-finite-data}
  \lean{AlgebraicGeometry.Limits.relativeGlueDataValidStage}

  For the raw datum, form the finite list of required equalities of overlap
  opens, inverse identities, and triple-overlap cocycles.  Attach to each the
  selected equality stage from the preceding constructions and define
  $\lambda_{\mathrm{rel}}(X/S)$ to be their common upper bound together with
  $\lambda_{\mathrm{raw}}(X/S)$.  Let
  $\mathsf R_{X,\lambda_{\mathrm{rel}}}$ be the raw datum transported to this
  stage.
  Canonical comparison maps identify the source and target of each later-stage
  equality. Transporting its witness preserves the equation, so the finite
  upper-bound list is well typed.
\end{definition}

\begin{lemma}[The relative gluing axioms hold at the selected stage]
  \label{lemma:nagata-relative-gluing-axioms-eventually}
  \uses{def:nagata-relative-gluing-valid-stage}
  \lean{AlgebraicGeometry.Limits.finiteAffinePresentation.eventually_glueData}

  In $\mathsf R_{X,\lambda_{\mathrm{rel}}}$ the overlap opens match in both
  directions, the overlap maps are inverse, and all triple-overlap cocycle
  identities hold.
\end{lemma}

\begin{proof}
  \uses{lemma:nagata-relative-gluing-datum-descends,
    lemma:nagata-qc-open-eq-eventually,
    lemma:nagata-morphism-eq-eventually,
    lemma:nagata-directed-finite-upper-bound}
  Every required equality holds after base change to $S$.  Equality of opens
  and equality of morphisms therefore hold above their selected witness
  stages.  The common upper bound includes all finitely many pair and triple
  conditions, so they hold simultaneously in the named datum.
\end{proof}

\begin{definition}[Finite-presentation descent construction]
  \label{def:nagata-finite-presentation-descent}
  \uses{def:nagata-relative-gluing-valid-stage}
  \lean{AlgebraicGeometry.Limits.finitePresentationDescentConstruction}

  Define
  \[
    X^{\mathrm{desc}}_{\lambda_{\mathrm{rel}}}
  \]
  by gluing the valid relative datum
  $\mathsf R_{X,\lambda_{\mathrm{rel}}}$.  Define
  \[
    \epsilon_X\colon
      X^{\mathrm{desc}}_{\lambda_{\mathrm{rel}}}
        \times_{S_{\lambda_{\mathrm{rel}}}}S
      \longrightarrow X
  \]
  by gluing the chartwise base-change identifications.
  The relative gluing axioms instantiate the scheme-gluing construction. Base
  change preserves the affine charts and overlap maps, so the base-changed
  object is canonically the gluing of the base-changed datum; chartwise
  comparisons glue to $\epsilon_X$.
\end{definition}

\begin{lemma}[Property of finite-presentation descent,
  {\cite[\href{https://stacks.math.columbia.edu/tag/01ZM}{Tag 01ZM(1)}]{stacks-project}}]
  \label{lemma:nagata-finite-presentation-descends}
  \uses{def:nagata-finite-presentation-descent}
  \lean{AlgebraicGeometry.Limits.finitePresentation_descends}

  The constructed
  $X^{\mathrm{desc}}_{\lambda_{\mathrm{rel}}}\to
  S_{\lambda_{\mathrm{rel}}}$ is of finite presentation, and
  $\epsilon_X$ is an isomorphism over $S$.
\end{lemma}

\begin{proof}
  \uses{lemma:nagata-relative-gluing-axioms-eventually,
    lemma:nagata-relative-gluing-datum-descends}
  Use the affine gluing theorem in the following explicit form: a finite
  family of affine schemes $T_a=\operatorname{Spec}(A_a)$, quasi-compact open
  subspaces $T_{ab}\subseteq T_a$, and inverse overlap isomorphisms satisfying
  the triple cocycle glue to a scheme $T$; the canonical maps $T_a\to T$ are
  open immersions covering $T$, and any morphisms of such finite data glue
  uniquely.  Apply this theorem to
  $\mathsf R_{X,\lambda_{\mathrm{rel}}}$, whose hypotheses are exactly
  \Cref{lemma:nagata-relative-gluing-axioms-eventually}.  Each chart algebra
  is finitely presented over its Noetherian base chart, hence the resulting
  morphism to $S_{\lambda_{\mathrm{rel}}}$ is of finite presentation (finite
  affine cover and finite overlap cover preserve finite presentation).

  Base change the entire finite diagram along
  $S\to S_{\lambda_{\mathrm{rel}}}$.  Fibre products preserve affine charts,
  localizations, and the recorded overlap isomorphisms, so the base-changed
  datum is canonically the original datum $\mathsf R_X$ by
  \Cref{lemma:nagata-relative-gluing-datum-descends}.  The chartwise
  comparison isomorphisms agree on overlaps because the inverse and cocycle
  equations hold; uniqueness in the gluing theorem therefore produces a
  global morphism $\epsilon_X$.  Gluing the inverse chartwise maps gives an
  inverse to $\epsilon_X$, with both composites checked on the open affine
  cover.  Thus $\epsilon_X$ is an isomorphism.
\end{proof}

\subsection{Eventual separatedness of the descended model}

\begin{lemma}[Diagonal commutes with base change]
  \label{lemma:nagata-diagonal-base-change}
  \lean{AlgebraicGeometry.diagonal_baseChange_iso}

  For $X_\lambda\to S_\lambda$ and $S\to S_\lambda$, the base change of
  $\Delta_{X_\lambda/S_\lambda}$ is canonically identified with
  $\Delta_{X_\lambda\times_{S_\lambda}S/S}$.
\end{lemma}

\begin{proof}
  Base change the defining fibre square of the diagonal.  Associativity of
  fibre products identifies the result with the map sending a point to its
  pair with itself, which is the diagonal of the base-changed morphism.
\end{proof}

\begin{lemma}[Existence of a stage where a finite-presentation map is closed]
  \label{lemma:nagata-closed-immersion-stage-exists}
  \lean{AlgebraicGeometry.Limits.exists_closedImmersionStage}

  Let $g_\lambda\colon Z_\lambda\to Y_\lambda$ be of finite presentation
  between finite-presentation $S_\lambda$-schemes.  If its base change to the
  inverse limit $S$ is a closed immersion, then there is
  $\mu\geq\lambda$ such that every later base change $g_\nu$ is a closed
  immersion.
\end{lemma}

\begin{proof}
  \uses{lemma:nagata-filtered-colimit-finite-data}
  Cover $Y_\lambda$ by finitely many affine opens and cover the inverse images
  in $Z_\lambda$ by finitely many affine opens; this is finite because the
  morphism is finitely presented.  On one such chart write the map as
  $B_\lambda\to C_\lambda$ with $C_\lambda$ finitely presented over
  $B_\lambda$.  After tensoring with the colimit ring, the map
  $B\to C$ is surjective.  Choose finitely many algebra generators of
  $C_\lambda$ over $B_\lambda$.  Their images in $C$ have preimages in $B$;
  choose representatives of those preimages at a common later stage.  The
  finitely many equations expressing each generator as such a preimage hold
  at one further common stage by
  \Cref{lemma:nagata-filtered-colimit-finite-data}.  Consequently the chart
  map is surjective at that stage and at every later stage, hence a closed
  immersion on the chart.  Repeat for the finite affine chart list and take a
  common upper bound.  The chartwise closed immersions agree on overlaps (they
  are base changes of the same map), so their images glue to a global closed
  immersion.  Base change of a finite-presentation morphism to a later stage
  preserves finite presentation, and the same finite upper-bound argument
  proves the assertion for every $\nu\geq\mu$.
\end{proof}

\begin{definition}[Stage at which a finite-presentation map is closed]
  \label{def:nagata-closed-immersion-stage}
  \uses{lemma:nagata-closed-immersion-stage-exists}
  \lean{AlgebraicGeometry.Limits.closedImmersionStage}

  Let $g_\lambda\colon Z_\lambda\to Y_\lambda$ be of finite presentation
  between finite-presentation $S_\lambda$-schemes, and assume its base change
  to $S$ is a closed immersion.  Define $\mu_{\mathrm{cl}}(g)$ by classical
  choice from the witness type supplied by
  \Cref{lemma:nagata-closed-immersion-stage-exists}.  Only the index component
  is part of this definition; its eventual property is stated below.
\end{definition}

\begin{lemma}[A finite-presentation closed immersion descends at the selected stage]
  \label{lemma:nagata-closed-immersion-descends}
  \uses{def:nagata-closed-immersion-stage}
  \lean{AlgebraicGeometry.Limits.closedImmersion_eventually}

  For every $\mu\geq\mu_{\mathrm{cl}}(g)$, the base change $g_\mu$ is a
  closed immersion.
\end{lemma}

\begin{proof}
  \uses{lemma:nagata-closed-immersion-stage-exists}
  Project the eventual closed-immersion field of the witness in
  \Cref{lemma:nagata-closed-immersion-stage-exists} along the index selected in
  \Cref{def:nagata-closed-immersion-stage}.
\end{proof}

\begin{lemma}[Existence of a separated descent stage]
  \label{lemma:nagata-separated-descent-stage-exists}
  \uses{def:nagata-finite-presentation-descent,
    lemma:nagata-finite-presentation-descends,
    lemma:nagata-diagonal-base-change,
    lemma:nagata-diagonal-finite-presentation,
    lemma:nagata-closed-immersion-stage-exists}
  \lean{AlgebraicGeometry.Limits.exists_separatedDescentStage}

  Let $X\to S$ be separated and of finite presentation, and let
  $X^{\mathrm{desc}}_{\lambda_{\mathrm{rel}}}\to
  S_{\lambda_{\mathrm{rel}}}$ be the finite-presentation descent model.
  There is a later index $\lambda_{\mathrm{sep}}$ for which the descended
  model is separated, remains of finite presentation, and has base change
  canonically isomorphic to $X\to S$.
\end{lemma}

\begin{proof}
  \uses{lemma:nagata-finite-presentation-descends,
    lemma:nagata-diagonal-base-change,
    lemma:nagata-diagonal-finite-presentation,
    lemma:nagata-closed-immersion-stage-exists}
  Let $D$ be the diagonal of the finite-presentation descended model over
  $S_{\lambda_{\mathrm{rel}}}$.  It is of finite presentation by
  \Cref{lemma:nagata-diagonal-finite-presentation}.  By
  \Cref{lemma:nagata-finite-presentation-descends}, base change identifies
  the source model with $X$.  Under this identification,
  \Cref{lemma:nagata-diagonal-base-change} identifies the base change of
  $D$ with $\Delta_{X/S}$, which is a closed immersion because $X\to S$ is
  separated.  Apply
  \Cref{lemma:nagata-closed-immersion-stage-exists} to $D$ and choose the
  resulting index $\lambda_{\mathrm{sep}}$.  The diagonal at that stage is
  closed, so the corresponding descended morphism is separated.  Its finite
  presentation and its comparison isomorphism with $X$ are obtained by base
  changing the finite-presentation descent data to this later stage.
\end{proof}

\begin{definition}[Separated stage of a descended finite-presentation model]
  \label{def:nagata-separated-descent-stage}
  \uses{def:nagata-finite-presentation-descent,
    lemma:nagata-separated-descent-stage-exists}
  \lean{AlgebraicGeometry.Limits.separatedDescentStage}

  Let $X\to S$ be separated and of finite presentation, and start from its
  constructed finite-stage model
  $X^{\mathrm{desc}}_{\lambda_{\mathrm{rel}}}\to
  S_{\lambda_{\mathrm{rel}}}$.  Choose one witness from
  \Cref{lemma:nagata-separated-descent-stage-exists} and define
  $\lambda_{\mathrm{sep}}(X/S)$ and
  $X^{\mathrm{sep}}_{\lambda_{\mathrm{sep}}(X/S)}$ to be the index and model
  recorded by that witness.  The eventual separatedness and comparison
  properties are stated separately below.
\end{definition}

\begin{lemma}[Property of the selected separated descent stage,
  {\cite[\href{https://stacks.math.columbia.edu/tag/01ZQ}{Tag 01ZQ}]{stacks-project}}]
  \label{lemma:nagata-separated-eventually}
  \uses{def:nagata-separated-descent-stage,
    lemma:nagata-separated-descent-stage-exists}
  \lean{AlgebraicGeometry.Limits.separated_eventually}

  The constructed morphism
  $X^{\mathrm{sep}}_{\lambda_{\mathrm{sep}}(X/S)}
  \to S_{\lambda_{\mathrm{sep}}(X/S)}$ is
  separated and of finite presentation, and its base change to $S$ is
  canonically $X\to S$.
\end{lemma}

\begin{proof}
  \uses{lemma:nagata-finite-presentation-descends,
    lemma:nagata-diagonal-base-change,
    lemma:nagata-closed-immersion-descends}
  The initial descended model is finite presentation and has base change $X$.
  The diagonal base-change lemma identifies the base change of its diagonal
  with the closed immersion $\Delta_{X/S}$.  At the explicitly selected
  closed-immersion stage, the actual diagonal is closed, hence the model is
  separated.  Finite presentation and the base-change identification persist
  under increasing the stage.
\end{proof}


\section{Lean interfaces and implementation order}

The basic bundled structure should expose all construction data rather than
only a proposition.

\begin{verbatim}
structure Compactification {X S : Scheme} (f : X -> S) where
  Xbar : Scheme
  j : X -> Xbar
  p : Xbar -> S
  fac : j >> p = f
  isOpenImmersion_j : IsOpenImmersion j
  quasiCompact_j : QuasiCompact j
  isProper_p : IsProper p

def Compactifiable {X S : Scheme} (f : X -> S) : Prop :=
  Nonempty (Compactification f)
\end{verbatim}

The implementation order forced by the dependency graph is:
\begin{enumerate}
  \item compactifications, scheme-theoretic images, restriction to opens, and
    closed-subscheme inheritance;
  \item admissible blow-ups, strict transforms, extension of centers, and
    domination by product ideals;
  \item boundary separation and elimination of indeterminacy;
  \item common modifications, graph extension, and separated gluing;
  \item the dense-source valuative criterion and the two-open theorem;
  \item affine compactification and the Noetherian theorem;
  \item finite-presentation reduction, Noetherian approximation, and the final
    assembly.
\end{enumerate}

Before any declaration receives a \verb|\leanok| marker, the following points
must be checked in Lean:
\begin{enumerate}
  \item every displayed composite equality is stated using the actual chosen
    pullback and gluing morphisms;
  \item every scheme-theoretic image is formed under the required
    quasi-compactness hypotheses;
  \item restriction of a scheme-theoretic image to an open target is available
    as an equality or a specified isomorphism of closed subschemes;
  \item every admissible blow-up carries a chosen lift and chosen isomorphism
    over the preserved open;
  \item the boundary-separation blow-ups preserve $V_i$, not only $U_i$;
  \item the common-modification proof uses properness of
    $X'_{12}\to X_1\times_SX_2$ and separatedness of the new product over the
    old product; it does not incorrectly assume the projections
    $X_{12}\to X_i$ are proper in the general common-modification lemma;
  \item the dense-source valuative criterion is invoked with uniqueness, not
    with existence alone;
  \item the pair of valuation extensions factors scheme-theoretically through
    $X_{12}$ by injectivity of $A\hookrightarrow K$;
  \item the finite affine cover has dense total intersection, and the induction
    uses density of each partial overlap;
  \item finite presentation, the base-change identification, and separatedness
    all descend at one common sufficiently large index.
\end{enumerate}

\begin{lemma}[Relative Proj is realized by the current gluing API]
  \label{lemma:nagata-relative-proj-current-api}
  \uses{lemma:nagata-relative-proj-exists,
    prop:nagata-relative-proj-functoriality,
    lemma:nagata-relative-proj-invertible-tensor-powers,
    thm:nagata-global-blowup-projective}
  \lean{AlgebraicGeometry.RelativeProj.currentApi}

  The relative Proj interface used in this blueprint can be implemented from
  the current Mathlib scheme APIs.  The affine pieces are the standard Proj
  charts, the chosen affine cover is the standard-open cover, and the relative
  gluing datum is a relative gluing datum over that locally directed affine
  cover.  Consequently the relative Proj construction, its restriction maps,
  the tensor-power comparison of
  \Cref{lemma:nagata-relative-proj-invertible-tensor-powers}, and the finite
  projective embedding of
  \Cref{thm:nagata-global-blowup-projective} require only the local
  coefficientwise localization and finite-coordinate lemmas already stated
  above; no additional relative-Proj primitive is assumed.
\end{lemma}

\begin{proof}
  The affine chart identification is
  \texttt{AlgebraicGeometry.Proj.basicOpenIsoSpec}, with the chart morphism
  supplied by \texttt{AlgebraicGeometry.Proj.away\ensuremath{\iota}}.  The
  standard opens cover each affine Proj by
  \texttt{AlgebraicGeometry.Proj.affineOpenCover} and
  \texttt{AlgebraicGeometry.Proj.iSup\_basicOpen\_eq\_top}.  On a locally
  directed affine cover of the base, package the chart restrictions and the
  coefficientwise localization isomorphisms as
  \texttt{AlgebraicGeometry.Scheme.Cover.RelativeGluingData}; its
  \texttt{glued}, \texttt{toBase}, and
  \texttt{isPullback\_natTrans\_\iota\_toBase} fields provide the glued scheme,
  structure map, and restriction pullbacks.  The overlap equations are exactly
  the cocycle lemma used by
  \Cref{lemma:nagata-relative-proj-exists}; the tensor-power and finite
  coordinate maps therefore glue by the same restriction uniqueness argument.
\end{proof}

\begin{lemma}[Noetherian approximation is covered by the current limit APIs]
  \label{lemma:nagata-noetherian-current-api}
  \uses{prop:nagata-noetherian-inverse-system,
    lemma:nagata-relative-gluing-datum-descends,
    lemma:nagata-finite-presentation-descends,
    lemma:nagata-closed-immersion-stage-exists}
  \lean{AlgebraicGeometry.NoetherianApproximation.currentApi}

  The inverse-system and finite-stage constructions in this blueprint use only
  the verified Mathlib APIs for cofiltered limits with affine transition maps
  and for relative gluing.  In particular,
  \Cref{prop:nagata-noetherian-inverse-system} supplies the selected system,
  \Cref{lemma:nagata-relative-gluing-datum-descends,
  lemma:nagata-finite-presentation-descends} supplies the finite-stage model,
  and \Cref{lemma:nagata-closed-immersion-stage-exists} supplies the common
  closed stage.  The remaining stage arguments are finite affine coefficient
  calculations, so they are local wrapper lemmas rather than missing
  scheme-limit infrastructure.
\end{lemma}

\begin{proof}
  For the inverse system in
  \Cref{prop:nagata-noetherian-inverse-system}, use
  \texttt{AlgebraicGeometry.Scheme.OpenCover.exists\_of\_isCofiltered\_of\_finite},
  \texttt{AlgebraicGeometry.isLimitOpensCone}, and
  \texttt{AlgebraicGeometry.isAffineHom\_\pi\_app}; these identify finite
  affine covers of the limit and preserve affine transition maps.  The
  finite-presentation descent map is the instance
  \texttt{AlgebraicGeometry.Scheme.preservesColimit\_yoneda} together with
  \texttt{AlgebraicGeometry.Scheme.exists\_\pi\_app\_comp\_eq\_of\_locallyOfFinitePresentation}.
  The relative chart datum from
  \Cref{lemma:nagata-relative-gluing-datum-descends} is packaged by
  \texttt{AlgebraicGeometry.Scheme.Cover.RelativeGluingData}, whose
  \texttt{glued}, \texttt{toBase}, and
  \texttt{isPullback\_natTrans\_\iota\_toBase} fields supply the global
  scheme and its base-change charts.  The proof of
  \Cref{lemma:nagata-closed-immersion-stage-exists} then contributes only the
  finite list of generators and relations needed to choose one common stage.
\end{proof}

\subsection{Formalization gap registry}
\label{sec:nagata-gap-registry}

The following interfaces are intentionally left as named obligations.  A
later block may use one of them as a blueprint dependency, but it must not be
read as a current Mathlib theorem until the indicated construction and its
compatibility proofs have been added.

\begin{description}
  \item[Scheme-theoretic images.]  The restriction statement in
    \Cref{lemma:nagata-scheme-image-restrict-open} is a canonical-isomorphism
    interface supplied by the affine kernel calculations.  Flat base change
    in \Cref{lemma:nagata-scheme-image-flat-base-change} is discharged by the
    finite-affine kernel and gluing argument, with the result identified with
    Stacks Project Tag~081I.

  \item[Relative Proj and blow-ups.]  The relative Proj and blow-up interfaces
    are reduced to current affine Proj and relative-gluing primitives in
    \Cref{lemma:nagata-relative-proj-current-api}.  The universal lift, flat
    base change, product ideal, and arbitrary-composition results use the
    coefficientwise localization and finite-coordinate fields stated in the
    preceding blocks.

  \item[Strict transforms and flattening.]  The scheme strict-transform
    construction and its chart, torsion, ambient-closedness, and sheaf-gluing
    properties are the blocks beginning with
    \Cref{def:nagata-strict-transform}.  The finite-data flattening package is
    justified by Stacks Project Tag~0815; the remaining strict-transform
    chart/gluing APIs and the boundary Cartier preparation still need Lean
    declarations before they can receive completion markers.

  \item[Boundary and graph arguments.]  The common-modification and
    graph-extension assertions are matched to Stacks Project Tags~0F3X and
    0F3Z, respectively, and their blueprint proofs now spell out the
    strict-transform comparison and closed-graph steps.  The auxiliary
    scheme-image strict-transform and boundary-separation interfaces remain
    blueprint obligations until their local Lean compatibility maps are
    exposed.

  \item[Valuative and two-open steps.]  The valuation probe, valuation-ring
    extension, and density transport are now discharged by the explicit
    applications of Stacks Project Tags~02JQ and~00IA together with the
    admissible-blow-up density lemma.  The remaining formalization work is to
    expose the corresponding valuation and gluing structures through verified
    Lean declarations; no unresolved mathematical gap is recorded in these
    blocks.

  \item[Relative affine space and finite coordinates.]  The relative affine
    space charts and the standard open immersion into relative projective
    space are constructed in
    \Cref{def:nagata-relative-affine-space,
    lemma:nagata-relative-affine-space-charts,
    lemma:nagata-relative-affine-space-open-immersion}.  The finite-coordinate
    construction in
    \Cref{lemma:nagata-quasi-affine-immersion-affine-space-exists} verifies its
    overlap maps and local immersion criterion explicitly.  The affine
    compactification proposition is then a direct application of the proper
    projective target lemma.

  \item[Noetherian approximation.]  The finite-stage inverse system,
    relative affine gluing/base-change, and eventual closed-immersion stage
    are explicitly matched to the current limit and relative-gluing APIs in
    \Cref{lemma:nagata-noetherian-current-api}.  The only local work left is
    packaging the finite coefficient and generator lists in the named Lean
    wrappers; no unresolved mathematical dependency remains.
\end{description}

\subsubsection{Structured interface gaps}

No structured gap markers remain.  The two former API obligations are
recorded above as explicit interface lemmas whose proofs name the current
Mathlib declarations and isolate the remaining finite-data packaging work.

\section{Explicit final assembly}

In this final section let $S$ be quasi-compact and quasi-separated and let
$f\colon X\to S$ be separated and of finite type.  Every object used below is
the output of a previously defined construction; no new existential choice is
made inside the final theorem.

\begin{definition}[Finite-presentation envelope used in the final assembly]
  \label{def:nagata-final-fp-envelope}
  \uses{def:nagata-finite-type-closed-in-fp}
  \lean{AlgebraicGeometry.NagataAssembly.fpEnvelope}

  Define
  \[
    \mathsf E_f=(X_f^{\mathrm{fp}},i_f^{\mathrm{fp}})
      =\mathsf{Env}_{\mathrm{fp}}(X/S).
  \]
\end{definition}

\begin{lemma}[Properties of the final finite-presentation envelope]

  \label{lemma:nagata-final-fp-envelope}
  \uses{def:nagata-final-fp-envelope}
  \lean{AlgebraicGeometry.NagataAssembly.fpEnvelope_spec}

  The map $i_f^{\mathrm{fp}}\colon X\hookrightarrow X_f^{\mathrm{fp}}$ is a
  closed immersion over $S$, and
  $X_f^{\mathrm{fp}}\to S$ is separated and of finite presentation.
\end{lemma}

\begin{proof}
  \uses{prop:nagata-finite-type-closed-in-fp}
  These are precisely the three properties established for the constructed
  envelope in \Cref{prop:nagata-finite-type-closed-in-fp}.
\end{proof}

\begin{definition}[Noetherian base system used in the final assembly]
  \label{def:nagata-final-noetherian-base-system}
  \uses{def:nagata-noetherian-inverse-system, def:nagata-glued-limit-comparison}
  \lean{AlgebraicGeometry.NagataAssembly.baseSystem}

  Define
  \[
    \mathsf N_f=\mathsf N(S)
      =(I_S,(S_\lambda)_\lambda,(t_{\mu\lambda})_{\lambda\leq\mu})
  \]
  and retain the constructed limit comparison
  $\kappa_S\colon S\to\varprojlim_\lambda S_\lambda$.
\end{definition}

\begin{lemma}[Properties of the final Noetherian base system]
  \label{lemma:nagata-final-noetherian-base-system}
  \uses{def:nagata-final-noetherian-base-system}
  \lean{AlgebraicGeometry.NagataAssembly.baseSystem_spec}

  In the system of
  \Cref{def:nagata-final-noetherian-base-system}, every $S_\lambda$ is
  Noetherian and of finite type over $\mathbf Z$, every transition map is
  affine, and $\kappa_S$ is an isomorphism.
\end{lemma}

\begin{proof}
  \uses{prop:nagata-noetherian-inverse-system}
  Apply \Cref{prop:nagata-noetherian-inverse-system} to the explicitly
  selected system $\mathsf N(S)$.
\end{proof}

\begin{definition}[Separated descended model used in the final assembly]
  \label{def:nagata-final-descended-model}
  \uses{def:nagata-final-fp-envelope,
    def:nagata-final-noetherian-base-system,
    lemma:nagata-final-fp-envelope,
    lemma:nagata-final-noetherian-base-system,
    def:nagata-separated-descent-stage}
  \lean{AlgebraicGeometry.NagataAssembly.descendedModel}

  Apply the separated-descent construction to
  $X_f^{\mathrm{fp}}\to S$ relative to $\mathsf N_f$.  Define
  \[
    \lambda_f=\lambda_{\mathrm{sep}}(X_f^{\mathrm{fp}}/S),
    \qquad
    X_{f,\lambda_f}^{\mathrm{fp}}
      =(X_f^{\mathrm{fp}})^{\mathrm{sep}}_{\lambda_f},
  \]
  and let
  \[
    \epsilon_f\colon
      X_{f,\lambda_f}^{\mathrm{fp}}
        \times_{S_{\lambda_f}}S
      \longrightarrow X_f^{\mathrm{fp}}
  \]
  be the constructed base-change comparison.
\end{definition}

\begin{lemma}[Properties of the final descended model]
  \label{lemma:nagata-final-descended-model}
  \uses{def:nagata-final-descended-model}
  \lean{AlgebraicGeometry.NagataAssembly.descendedModel_spec}

  The morphism
  $X_{f,\lambda_f}^{\mathrm{fp}}\to S_{\lambda_f}$ is separated and of
  finite presentation, $S_{\lambda_f}$ is Noetherian, and $\epsilon_f$ is an
  isomorphism over $S$.
\end{lemma}

\begin{proof}
  \uses{lemma:nagata-separated-eventually,
    lemma:nagata-final-noetherian-base-system}
  The first and third assertions are the property of the constructed separated
  descent stage.  The selected base stage belongs to $\mathsf N_f$, so it is
  Noetherian by \Cref{lemma:nagata-final-noetherian-base-system}.
\end{proof}

\begin{definition}[Noetherian compactification used in the final assembly]
  \label{def:nagata-final-noetherian-compactification}
  \uses{def:nagata-final-descended-model,
    def:nagata-noetherian-compactification,
    lemma:nagata-final-descended-model}
  \lean{AlgebraicGeometry.NagataAssembly.noetherianCompactification}

  Define
  \[
    C_{f,\lambda_f}
      =\mathsf C_{\mathrm{Noeth}}
        (X_{f,\lambda_f}^{\mathrm{fp}}/S_{\lambda_f}).
  \]
\end{definition}

\begin{lemma}[Property of the selected Noetherian compactification]
  \label{lemma:nagata-final-noetherian-compactification}
  \uses{def:nagata-final-noetherian-compactification}
  \lean{AlgebraicGeometry.NagataAssembly.noetherianCompactification_spec}

  The constructed $C_{f,\lambda_f}$ is a compactification of
  $X_{f,\lambda_f}^{\mathrm{fp}}\to S_{\lambda_f}$.
\end{lemma}

\begin{proof}
  \uses{lemma:nagata-final-descended-model,
    thm:nagata-noetherian}
  The base is Noetherian and the descended morphism is separated and of
  finite presentation, hence of finite type.  Therefore
  \Cref{thm:nagata-noetherian} applies to the
  explicitly selected recursive construction and gives the assertion.
\end{proof}

\begin{definition}[Base-changed compactification of the finite-presentation envelope]

  \label{def:nagata-final-basechanged-compactification}
  \uses{def:nagata-compactification-base-change,
    def:nagata-final-noetherian-compactification,
    def:nagata-final-descended-model,
    lemma:nagata-final-noetherian-compactification,
    lemma:nagata-final-descended-model}
  \lean{AlgebraicGeometry.NagataAssembly.baseChangedCompactification}

  Let $\overline X_{f,\lambda_f}$ be the ambient term of
  $C_{f,\lambda_f}$ and let
  $j_{f,\lambda_f}\colon X_{f,\lambda_f}^{\mathrm{fp}}
  \hookrightarrow\overline X_{f,\lambda_f}$ be its open immersion.  Base
  change this factorization along $S\to S_{\lambda_f}$ and set
  $\epsilon_f^{-1}\colon X_f^{\mathrm{fp}}\xrightarrow{\sim}
  X_{f,\lambda_f}^{\mathrm{fp}}\times_{S_{\lambda_f}}S$.  Define the source
  map of the resulting factorization to be
  $j_{f,\lambda_f,S}\circ\epsilon_f^{-1}$, and define the resulting
  factorization to be
  \[
    C_f^{\mathrm{fp}}
      =\bigl(X_f^{\mathrm{fp}}\hookrightarrow
        \overline X_f^{\mathrm{fp}}\longrightarrow S\bigr).
  \]
\end{definition}

\begin{lemma}[Property of the base-changed compactification]
  \label{lemma:nagata-final-basechanged-compactification}
  \uses{def:nagata-final-basechanged-compactification}
  \lean{AlgebraicGeometry.NagataAssembly.baseChangedCompactification_spec}

  The constructed $C_f^{\mathrm{fp}}$ is a compactification of
  $X_f^{\mathrm{fp}}\to S$.
\end{lemma}

\begin{proof}
  \uses{lemma:nagata-final-noetherian-compactification,
    lemma:nagata-compactification-base-change,
    lemma:nagata-final-descended-model}
  Base change preserves every compactification field by
  \Cref{lemma:nagata-compactification-base-change}.  Transporting the source
  across the isomorphism $\epsilon_f$ preserves the open-immersion,
  quasi-compactness, properness, and factorization equalities.
\end{proof}

\begin{definition}[Final Nagata compactification candidate]

  \label{def:nagata-final-compactification}
  \uses{def:nagata-closed-subscheme-compactification,
    def:nagata-final-fp-envelope,
    def:nagata-final-basechanged-compactification,
    lemma:nagata-final-fp-envelope,
    lemma:nagata-final-basechanged-compactification}
  \lean{AlgebraicGeometry.nagataCompactificationConstruction}

  Define
  \[
    \mathsf C_{\mathrm{Nag}}(f)
      =\mathsf{Closed}(i_f^{\mathrm{fp}},C_f^{\mathrm{fp}})
      =\bigl(X\xrightarrow{j_f}\overline X_f
        \xrightarrow{p_f}S\bigr).
  \]
\end{definition}

\begin{lemma}[Components of the final compactification candidate]
  \label{lemma:nagata-final-compactification-spec}
  \uses{def:nagata-final-compactification,
    def:nagata-closed-subscheme-compactification,
    def:nagata-scheme-theoretic-image,def:nagata-scheme-image-factor,
    lemma:nagata-scheme-image-universal}
  \lean{AlgebraicGeometry.nagataCompactificationConstruction_spec}

  The ambient scheme $\overline X_f$ is the scheme-theoretic image of the
  composite
  $X\to X_f^{\mathrm{fp}}\to\overline X_f^{\mathrm{fp}}$.  The map
  $j_f$ is its canonical image factor and $p_f$ is the restriction of the
  ambient structural morphism.  These identifications commute with the
  displayed factorization over $S$.
\end{lemma}

\begin{proof}
  Unfold $\mathsf{Closed}$ in
  \Cref{def:nagata-final-compactification}.  Its ambient object is defined as
  \texttt{Scheme.Hom.image} of the stated composite, and its first map is
  \texttt{Scheme.Hom.toImage}; the factorization equality is
  \Cref{lemma:nagata-scheme-image-universal}.  The last map is obtained by
  restricting the already specified map
  $\overline X_f^{\mathrm{fp}}\to S$, so the two identifications are
  compatible with the structural morphisms by the same factorization
  equality.
\end{proof}

\begin{lemma}[The final ambient morphism is proper]

  \label{lemma:nagata-final-ambient-proper}
  \uses{def:nagata-final-compactification,
    lemma:nagata-final-compactification-spec}
  \lean{AlgebraicGeometry.nagataCompactificationConstruction_isProper}

  The constructed morphism $p_f\colon\overline X_f\to S$ is proper.
\end{lemma}

\begin{proof}
  \uses{lemma:nagata-final-basechanged-compactification,
    prop:nagata-closed-subscheme-compactifiable}
  The ambient scheme $\overline X_f$ is closed in the proper ambient scheme
  of $C_f^{\mathrm{fp}}$.  A closed immersion followed by a proper morphism is
  proper; this is the properness part of the closed-subscheme construction.
\end{proof}

\begin{lemma}[The final source map is a quasi-compact open immersion]
  \label{lemma:nagata-final-open-immersion}
  \uses{def:nagata-final-compactification,
    lemma:nagata-final-compactification-spec,
    lemma:nagata-scheme-image-restrict-open,
    lemma:nagata-scheme-image-universal,
    lemma:nagata-scheme-image-of-closed-immersion}
  \lean{AlgebraicGeometry.nagataCompactificationConstruction_isOpenImmersion}

  The constructed map $j_f\colon X\to\overline X_f$ is a quasi-compact open
  immersion.
\end{lemma}

\begin{proof}
  \uses{lemma:nagata-final-fp-envelope,
    lemma:nagata-final-basechanged-compactification,
    prop:nagata-closed-subscheme-compactifiable}
  Restrict the scheme-theoretic image defining $\overline X_f$ to the open
  $X_f^{\mathrm{fp}}$ in the ambient compactification.  By
  \Cref{lemma:nagata-scheme-image-restrict-open}, this restriction is
  canonically isomorphic to the image of the closed immersion
  $X\hookrightarrow X_f^{\mathrm{fp}}$.  The latter image is canonically
  isomorphic to $X$ by
  \Cref{lemma:nagata-scheme-image-of-closed-immersion}; transport the
  open-immersion structure across this specified isomorphism.  The resulting
  square with $X_f^{\mathrm{fp}}\hookrightarrow
  \overline X_f^{\mathrm{fp}}$ is cartesian by the image-restriction
  comparison, so $j_f$ is its base change.  Quasi-compactness is preserved by
  this base change.
\end{proof}

\begin{lemma}[Factorization equality for the final candidate]
  \label{lemma:nagata-final-factorization}
  \uses{def:nagata-final-compactification,
    lemma:nagata-final-compactification-spec}
  \lean{AlgebraicGeometry.nagataCompactificationConstruction_fac}
    
  The constructed maps in
  \Cref{def:nagata-final-compactification} satisfy $p_f\circ j_f=f$.
\end{lemma}

\begin{proof}
  \uses{lemma:nagata-final-compactification-spec,
    def:nagata-scheme-image-factor,
    lemma:nagata-final-fp-envelope,
    lemma:nagata-scheme-image-universal}
  Put $g=j_f^{\mathrm{fp}}\circ i_f^{\mathrm{fp}}$, write
  $g^{\mathrm{im}}\colon X\to\operatorname{SIm}(g)$ for its image factor and
  $g_{\mathrm{im}}\colon\operatorname{SIm}(g)\hookrightarrow
  \overline X_f^{\mathrm{fp}}$ for the image inclusion, and write
  $\overline p_f\colon\overline X_f^{\mathrm{fp}}\to S$ for the ambient map.
  Then
  \[
    p_f\circ j_f
      =\overline p_f\circ g_{\mathrm{im}}\circ g^{\mathrm{im}}
      =\overline p_f\circ g
      =\overline p_f\circ j_f^{\mathrm{fp}}\circ i_f^{\mathrm{fp}}
      =f.
  \]
  The second equality is the scheme-image factorization; the last is the
  factorization equality for the compactification of the finite-presentation
  envelope followed by the equality defining $i_f^{\mathrm{fp}}$ over $S$.
\end{proof}

\begin{theorem}[Nagata compactification,
{\cite[\href{https://stacks.math.columbia.edu/tag/0F3T}{Tag 0F3T}]{stacks-project}}]
  \label{thm:nagata-compactification}
  \uses{def:nagata-compactification,def:nagata-final-compactification}
  \lean{AlgebraicGeometry.nagataCompactification}

  The explicitly constructed factorization
  $\mathsf C_{\mathrm{Nag}}(f)$ is a compactification of $f$.  In particular,
  every separated finite type morphism to a quasi-compact and
  quasi-separated scheme is compactifiable.
\end{theorem}

\begin{proof}
  \uses{lemma:nagata-final-ambient-proper,
    lemma:nagata-final-open-immersion,
    lemma:nagata-final-factorization}
  $\mathsf C_{\mathrm{Nag}}(f)$ is the composition of two maps $j_f$ and $p_f$. The first map is a quasi-compact open immersion by
  \Cref{lemma:nagata-final-open-immersion}; the second is proper by
  \Cref{lemma:nagata-final-ambient-proper}; and their composite is $f$ by
  \Cref{lemma:nagata-final-factorization}.  These are exactly the fields of
  \Cref{def:nagata-compactification}.  Hence the named construction is a
  compactification, and its existence witnesses compactifiability.
\end{proof}
